\documentclass[11pt]{amsart}

\usepackage[T1]{fontenc}
\usepackage{lmodern}
\usepackage[a4paper,hmargin=1.15in,vmargin=1.1in]{geometry}
\usepackage{amssymb,mathrsfs}
\usepackage{graphicx}
\usepackage{booktabs}
\usepackage{xcolor}
\usepackage{enumitem}
\usepackage{microtype}
\usepackage[colorlinks=true,linkcolor=blue!50!black,citecolor=blue!50!black,urlcolor=blue!50!black,bookmarks=false]{hyperref}
\usepackage{xurl}

\tolerance=2000
\emergencystretch=3em
\hfuzz=1pt
\raggedbottom

\newtheorem{theorem}{Theorem}[section]
\newtheorem{lemma}[theorem]{Lemma}
\newtheorem{proposition}[theorem]{Proposition}
\newtheorem{corollary}[theorem]{Corollary}
\newtheorem{mainthm}{Theorem}
\renewcommand{\themainthm}{\Alph{mainthm}}
\theoremstyle{definition}
\newtheorem{definition}[theorem]{Definition}
\theoremstyle{remark}
\newtheorem{remark}[theorem]{Remark}
\numberwithin{equation}{section}

\newcommand{\R}{\mathbb{R}}
\newcommand{\Z}{\mathbb{Z}}
\newcommand{\C}{\mathbb{C}}
\newcommand{\E}{\mathbb{E}}
\newcommand{\SU}{\mathrm{SU}}
\newcommand{\SO}{\mathrm{SO}}
\newcommand{\U}{\mathrm{U}}
\newcommand{\g}{\mathfrak{g}}
\newcommand{\Tr}{\operatorname{Tr}}
\newcommand{\ReTr}{\operatorname{Re}\operatorname{Tr}}
\newcommand{\Hess}{\operatorname{Hess}}
\newcommand{\Ric}{\operatorname{Ric}}
\newcommand{\Cov}{\operatorname{Cov}}
\newcommand{\Var}{\operatorname{Var}}
\newcommand{\dist}{\operatorname{dist}}
\newcommand{\ip}[2]{\langle #1,#2\rangle}
\newcommand{\BP}{\mathrm{BP}}
\newcommand{\Lip}{\operatorname{Lip}}

\begin{document}
\hypersetup{pdftitle={Yang-Mills Theory and the Mass Gap}}

\title{Yang--Mills Theory and the Mass Gap}

\author[D. Bhattacharjee]{Deep Bhattacharjee$^{*}$}
\address{Formerly, Electro-Gravitational Space Propulsion Laboratory
(EGSPL), Bhubaneswar, Odisha 751030, India; formerly, University of Calcutta (Heramba Chandra College), 23/49 Gariahat
Road, Golpark, Kolkata 700029, West Bengal, India}
\email{itsdeep@live.com, d.bhattacharjee@erl-forschung.de, egspl@proton.me}
\urladdr{\href{https://orcid.org/0000-0003-0466-750X}{orcid.org/0000-0003-0466-750X}}
\thanks{$^{*}$Corresponding author: Deep Bhattacharjee, itsdeep@live.com.}

\author[U. Bhattacharya]{Ushashi Bhattacharya}
\address{Formerly, National Taiwan University, Taipei 106319, Taiwan}
\email{rai.bhattacharya.in@gmail.com}
\urladdr{\href{https://orcid.org/0009-0002-2254-3914}{orcid.org/0009-0002-2254-3914}}

\author[J. M. Lockwood]{James Montgomery Lockwood}
\address{Independent Researcher, San Antonio, USA}
\email{spectralityworkinggroup@gmail.com}

\author[T. Romanov]{Tristan Romanov}
\address{Independent Researcher, San Francisco, CA, USA}
\email{romanov360@gmail.com}

\author[I. J. Sessford]{Ikailo John Sessford}
\address{Ikailo Inc., Whitby, Ontario, Canada}
\email{john@ikailo.com}
\urladdr{\href{https://orcid.org/0009-0002-0150-4027}{orcid.org/0009-0002-0150-4027}}

\author[S. Singha Roy]{Sanjeevan Singha Roy}
\address{Department of Physics, Birla Institute of Technology, Mesra, Ranchi,
Jharkhand 835215, India}
\email{sanjeevan9905@gmail.com}
\urladdr{\href{https://orcid.org/0000-0002-6148-1421}{orcid.org/0000-0002-6148-1421}}

\date{10 November 2026}

\subjclass[2020]{81T25, 82B20, 81T13, 60K35, 22E30. PACS 11.15.Ha, 11.15.Tk, 05.50.+q}
\keywords{Lattice gauge theory, mass gap, Bakry--\'Emery criterion,
Dobrushin uniqueness, continuum limit}

\begin{abstract}
We prove that four-dimensional lattice Yang--Mills theory with any compact
simple gauge group has a mass gap with an explicit lower bound at strong
coupling, and that a spectral gap $\varpi$ for resampling cubes of physical
size $\ell$, uniform in the lattice spacing, gives every nontrivial continuum
limit the mass gap $16\varpi/(625e^2\ell)$. We prove this cube gap at strong
coupling, for $\SU(2)$ up to $\beta_W=\frac7{20}$, and reduce it to a decay of
correlations below $\frac12$ across one slab of width $\ell/2$, which we prove
in two dimensions and which remains open in four.
\end{abstract}

\maketitle
\markboth{D. BHATTACHARJEE ET AL.}{YANG--MILLS THEORY AND THE MASS GAP}


\section{Introduction}\label{sec:intro}

The Yang--Mills existence and mass gap problem asks for a quantum
Yang--Mills theory on $\R^4$ with a compact simple gauge group that satisfies
the axioms of constructive quantum field theory, and for a number $\Delta>0$
such that the spectrum of the Hamiltonian lies in $\{0\}\cup[\Delta,\infty)$
\cite{JW06}. In the Euclidean formulation the theory is a probability measure
on gauge fields whose Schwinger functions satisfy the Osterwalder--Schrader
axioms \cite{OS73,OS75}; the mass gap is then exponential decay of the
truncated two-point functions of gauge-invariant observables at the rate
$\Delta$. The natural regularisation is Wilson's lattice gauge theory
\cite{Wil74}: an element $Q_e$ of $G\subset\U(N)$ sits on every link $e$ of the
lattice $a\Z^4$, and a configuration has the weight
$\exp(N\beta\sum_p\ReTr Q_p)$, where $Q_p$ is the ordered product around the
plaquette~$p$.\footnote{This is the normalisation of~\cite{SZZ23,BB26a}. The
conventional inverse coupling is $\beta_W=N^2\beta$, and for $\SU(2)$,
identified with the unit quaternions, $\beta_W=4\beta$.} Asymptotic freedom
forces the continuum limit $a\to0$ to be taken with $\beta\to\infty$, that is,
at weak coupling.

At small $\beta$ the lattice theory has a unique state, exponential clustering
and a transfer-matrix gap, by convergent expansions \cite{OSei78,Sei82};
Ba{\l}aban's renormalisation group \cite{Bal84a,Bal85a,Bal89b} gives ultraviolet
stability in finite volume at weak coupling (Dimock \cite{Dim13} presents the
method for the scalar $\phi^4$ model in three dimensions), and Magnen, Rivasseau and S\'en\'eor constructed the theory with an
infrared cutoff \cite{MRS93}. No existing result gives a mass gap in physical
units. This paper proves a lattice mass gap with explicit constants for every
compact simple group, reduces the continuum gap to one estimate on one
physical scale, proves that estimate at strong coupling, reduces it further to
a decay of correlations across a slab of physical width $\ell/2$, and
determines what the methods used here can and cannot reach. It also states
precisely what is missing.

\subsection{The lattice gap with explicit constants}

Let $G$ be a compact connected Lie group with semisimple Lie algebra,
identified with its image under a faithful unitary representation in $\U(N)$,
and let $\kappa_G>0$ be its least Ricci curvature for the bi-invariant metric
$\ip XY=-\Tr(XY)$. For a simple group $\kappa_G=h^\vee/(2\ell_\rho)$, with $h^\vee$
the dual Coxeter number and $\ell_\rho$ the Dynkin index of the representation;
for $\SU(N)$, $\kappa_G=N/2$.

\begin{mainthm}\label{thm:A}
Suppose $N|\beta|<\frac2{29}\kappa_G$. For every gauge-invariant local
observable the expectations in boxes with free boundary converge, along every
exhausting sequence of boxes, to a state $\omega$ that is invariant under the
lattice symmetries; truncated correlations of Lipschitz observables decay at
every rate below
\[
m_G(\beta)=\operatorname{arccosh}\Bigl(1+\frac1{12}\Bigl(\frac{\kappa_G}{N|\beta|}
-\frac{29}2\Bigr)\Bigr)>0 ,
\]
and for $\beta\ge0$ the state is reflection positive and its transfer matrix
has spectral gap $m\ge m_G(\beta)$. For $G=\SU(2)$ the same holds for
$|\beta_W|<\frac13$ with $m\ge-\log(3\beta_W)$.
\end{mainthm}

Two ingredients carry the proof. The first is a gauge: a field of straight
lines on $\Z^4$ that contains exactly one link of every plaquette
(Sec.~\ref{sec:linefield}). Setting these links to the identity loses no
gauge-invariant information, and every remaining link then interacts with
exactly twelve others instead of eighteen. The second is a Hessian bound in
this gauge, $|\Hess S|\le\frac{29}2$ for the action $S=\sum_p\ReTr Q_p$, proved
by a divergence identity and a positivity statement for one plaquette
(Sec.~\ref{sec:curvature}); with the Ricci curvature of $G$ it gives a
Bakry--\'Emery bound \cite{BE85,BGL14} with an explicit decay rate. Without
the gauge the sharp Hessian bound is $16$ \cite{BB26a}, so the gauge enlarges
the range of every group by the factor $32/29$ and supplies the rate. The
$\SU(2)$ statement combines the gauge with the sharp Kantorovich modulus
$\frac14$ of the von Mises--Fisher family \cite{BB26a} in Dobrushin's
criterion \cite{Dob68,Dob70}.

\subsection{A continuum mass gap from one cube}

The lattice gap of Theorem~\ref{thm:A} is measured in lattice units and
vanishes in physical units as $a\to0$. The second result isolates a single
estimate, posed on one scale, that controls the gap in physical units. On a
periodic lattice let $\mathcal P_x$ resample the links of the cube of side $n$
at $x$ from their conditional law under the Wilson measure $\mu$. The
condition $\BP(n,\varpi)$ asks that
\[
\varpi\Var_\mu(f)\le\frac1{n^4}\sum_x\mu\bigl((f-\mathcal P_xf)^2\bigr)
\]
for every gauge-invariant $f$: the dynamics that resamples cubes, each link
at rate one, has spectral gap at least $\varpi$ on gauge-invariant functions.
The condition involves no gauge fixing and no supremum over boundary
conditions.

\begin{mainthm}\label{thm:B}
Let $a_k\to0$ and $\beta_k\ge0$ be couplings, and suppose that there are
$\varpi>0$, $\ell>0$ and integers $n_k\ge2$ with $n_ka_k\le\ell$ such that
$\BP(n_k,\varpi)$ holds for the Wilson measure at $\beta_k$ on all large tori.
Then every continuum limit of the corresponding lattice theories has
Hamiltonian $H$ with $\sigma(H)\subset\{0\}\cup[\Delta,\infty)$,
\[
\Delta=\frac{16\varpi}{625e^2\ell},
\]
and $0$ is a simple eigenvalue; if the limit is nontrivial, $H$ has a mass
gap of at least $\Delta$.
\end{mainthm}

The proof (Sec.~\ref{sec:cube}) turns the spectral gap into decay of
covariances through a finite-speed bound for the cube dynamics, at a rate
of order $1/n$ in lattice units, hence of order $1/\ell$ in physical units,
and passes this decay through the Osterwalder--Schrader reconstruction.

\subsection{Where the cube condition holds}

\begin{mainthm}\label{thm:C}
\begin{enumerate}[label=\textup{(\roman*)},leftmargin=*]
\item At $\beta=0$, $\BP(n,1)$ holds for every $n$. For
$21N|\beta|<\kappa_G$ and $n$ large, $\BP(n,1-\alpha_n)$ holds with
$\alpha_n=N|\beta|(18n-3)/((\kappa_G-21N|\beta|)n^2)\to0$; for $\SU(2)$ the
range is $|\beta_W|<\frac5{18}$, and for $\SU(3)$ it contains
$0\le\beta_W\le0.45$.
\item For the free scalar field of lattice mass $\mu$ in $d$ dimensions the
best constant is
$\varpi=\frac{\mu^2}{n^d}\ip{\mathbf 1}{(-\Delta_D+\mu^2)^{-1}\mathbf 1}$, with
$-\Delta_D$ the Dirichlet Laplacian of the cube, and with $s=n\mu$ (the mass
times the physical size of the cube)
\[
\Bigl(\frac56\Bigr)^d\frac{s^2}{s^2+10d}\le\varpi\le\frac{s^2}{s^2+16d/9}.
\]
\item For $G=\SU(2)$ and $|\beta_W|\le\frac7{20}$, $\BP(n,\varpi)$ holds for
every $n\ge2$ with $\varpi=\frac{106}{1933}\cdot517^{-4}$, independent of $n$
and of the torus, and $\BP(n,1-259/n)$ holds for $n\ge260$. For
$0\le\beta_W\le\frac9{25}$ every torus limit is reflection positive and its
transfer matrix has spectral gap at least $7/(625\cdot2957\,e^2)$.
\end{enumerate}
\end{mainthm}

Part (ii) is the model case of the continuum tuning: the constant is uniform
in the lattice spacing at fixed $m\ell$, and it vanishes as $m\to0$. The cube
condition is therefore a quantitative statement of mass generation at the
scale $\ell$. Part (i) is proved in Sec.~\ref{sec:bpstrong} by the gauge of
Theorem~\ref{thm:A} and a weighted Dobrushin criterion for blocks of links.
No argument built on single-link updates in this gauge passes
$\beta_W=\frac13$ for $\SU(2)$ (Proposition~\ref{prop:methodlimit}). Part
(iii) passes that value by resampling \emph{stars}, the eight links at a
vertex of one parity (Sec.~\ref{sec:stars}): the perturbation of one
plaquette at a vertex enters the law of a star only through the relative
variable of two of its links, and no plaquette is internal to a star.

\subsection{Mass generation at one scale}

Statement (D) below asks for the cube condition along the continuum tuning.
Section~\ref{sec:massgen} reduces it to a hypothesis of a different kind,
closer to physics. Split a box of the lattice, all of whose widths lie
between $2d-1$ and $md-1$, into two outer parts of width $d$ and a middle of
width $d-1$ in one coordinate, fix an arbitrary configuration outside the
box, and let $\rho$ be the maximal correlation between the gauge-invariant
observables of the two outer parts under the conditional law. The hypothesis
$(\mathrm M_{d,m}(c))$ asks that $\rho\le c$ for all such boxes, splittings
and boundary conditions. Along the tuning $d\approx\ell/(2a)$, so this is a
decay of correlations across a slab of physical width about $\ell/2$,
uniform in the boundary condition.

Theorem~\ref{thm:D} proves that $(\mathrm M_{d,m}(c))$ with $c<\frac12$ and $m$
large in terms of $c$ gives $\BP(2d-1,(1-2c)^4/32)$, hence, along the tuning,
statement (D) and the continuum gap $(1-2c)^4/(1250\,e^2\ell)$ in every
nontrivial limit. The proof passes through a gap for the dynamics that
resamples overlapping patches of a box, obtained from the decorrelation of
the slabs by a Markov factorisation, and a finite-size criterion for
gauge-invariant functions in the spirit of those of Knabe~\cite{Kna88} and of
Gosset and Mozgunov~\cite{GM16} for quantum spin systems. The hypothesis also
gives exponential decay of gauge-invariant correlations under every boundary
condition, at rate $\log(1/c)/d$ per lattice spacing
(Proposition~\ref{prop:alldecay}). It holds at strong coupling at fixed $d$ and
$m$, and in two dimensions along the weak-coupling tuning
(Theorem~\ref{thm:twod}). In four dimensions it is not proved.

\subsection{The limits of the methods}

Appendix~\ref{app:obstructions} proves that several natural routes to (D)
cannot work, and locates the threshold of Theorem~\ref{thm:D}.
\begin{enumerate}[label=\textup{(\alph*)},leftmargin=*]
\item \emph{The threshold $\frac12$ is sharp.} For a massless Gaussian chain
the maximal correlation of the slab mechanism equals $\frac12$ exactly, and
the gap of the corresponding block dynamics on $M$ blocks is at most
$16/(3M)$; so no decorrelation constant $c\ge\frac12$ can give a gap uniform in
the size (Proposition~\ref{prop:gausschain}).
\item \emph{Two dimensions.} For $\SU(2)$ in two dimensions the middle slab of
a box integrates out, after an axial gauge, to a convolution power of the
heat kernel, and $(\mathrm M_{d,m}(c))$ holds with
$c=20\,e^{-3g^2(da)^2/4}$, which tends to zero as the physical size of the box
grows (Theorem~\ref{thm:twod}). The same condition for strips one link wide
fails along the tuning: a pinned strip is a Brownian bridge on $\SU(2)$, and
the antipodal boundary condition breaks a global symmetry, with correlation
close to one (Proposition~\ref{prop:strip}).
\item \emph{Curvature.} The Bakry--\'Emery constant of the Wilson measure
without gauge fixing is at most $N/2-16N|\beta|$, so the Hessian bound $16$ of
\cite{BB26a} is sharp (Theorem~\ref{thm:fullspace}); in the gauge of
Theorem~\ref{thm:A} the constant $\frac{29}2$ cannot be replaced by anything
below $\frac{38}3$, and no Bakry--\'Emery argument for the gauge-fixed measure
reaches $|\beta_W|=\frac3{19}$ (Theorem~\ref{thm:slice}). Every smooth
plaquette action, such as an effective action at a coarse scale, obeys a bound
of the same kind (Proposition~\ref{prop:coarse}).
\item \emph{Entropy, covariances and fixed blocks.} Block entropy
factorisation implies the cube condition and so cannot be easier
(Proposition~\ref{prop:entbp}); the link covariances of every gauge-invariant
measure are $\frac14I$ and carry no information on the coupling
(Proposition~\ref{prop:linkcov}); and a single-link spectral gap, or any gap
for blocks of fixed lattice size, with constant $C_k$ satisfying $C_ka_k\to0$,
forces every continuum limit to be trivial
(Proposition~\ref{prop:fixedblocks}).
\item \emph{A gap in the limit is not enough.} If the lattice gaps $m_k$ tend
to zero, a choice of lattice spacings proportional to $m_k$ produces a
nontrivial continuum limit with a mass gap; this limit is kinematic and is not
identified with Yang--Mills theory (Theorem~\ref{thm:gaptuning}). The content
of (D) and (E) is that the gap and the nontriviality hold along the tuning of
asymptotic freedom.
\end{enumerate}

\subsection{What remains}

Theorems~\ref{thm:A} and~\ref{thm:C}(i), (iii) are strong-coupling theorems.
The Clay problem needs the opposite end of the coupling range: the continuum
limit lives at $\beta\to\infty$, which the ranges of these theorems do not
reach. Two
statements would complete the proof of a continuum mass gap, and neither is
proved here:
\begin{enumerate}[label=\textup{(\Alph*)},leftmargin=*]
\setcounter{enumi}{3}
\item $\BP(n_k,\varpi)$ with $n_ka_k\le\ell$ along the asymptotic-freedom
tuning $\beta_k\to\infty$, uniformly in $k$;
\item convergence of the renormalised Schwinger functions to a nontrivial
limit satisfying the Osterwalder--Schrader axioms.
\end{enumerate}
Given (D) and (E), Theorem~\ref{thm:B} gives the mass gap $\Delta\ge
16\varpi/(625e^2\ell)$, and by Theorem~\ref{thm:D} statement (D) follows from
$(\mathrm M_{d_k,m}(c))$ along the tuning with $c<\frac12$. Section~\ref{sec:weak} explains why (D) is the heart of
the matter: by Theorem~\ref{thm:C}(ii) it cannot hold without a dynamically
generated mass, it fails at fixed cube size at weak coupling, and it fails in
the Coulomb phase of the abelian theory, so its proof has to use the
non-abelian structure of $G$ at every scale up to $\ell$.

\subsection{Related work}

Strong-coupling clustering goes back to Osterwalder and Seiler \cite{OSei78}
and Seiler \cite{Sei82}; probabilistic treatments include the survey of
Chatterjee \cite{Chat19}, the Langevin approach of Shen, Zhu and Zhu
\cite{SZZ23}, whose Bakry--\'Emery argument we sharpen, the area law in
enlarged regions of Cao, Nissim and Sheffield \cite{CNS25}, and correlation
decay for finite gauge groups at weak coupling by Adhikari and Cao \cite{AC25}.
In two dimensions the continuum measure is known
\cite{GKS89,Dri89,Lev03}. The Kantorovich modulus of the $\SU(2)$ single-link
law and the sharp Hessian bound without gauge fixing are proved in
\cite{BB26a}, and an earlier and much longer account of the one-cube
reduction is \cite{BBLR26}.

One of us (I.~J.~S.) has formalised in Lean~4 with Mathlib~\cite{Ses26}, without
unproved statements and with only the standard axioms, several ingredients of
the strong-coupling analysis on periodic lattices: the Haar probability on
$\SU(N)$, an infinite-volume Gibbs state obtained as a limit of torus states,
site reflection positivity and the transfer matrix, the Holley--Stroock
comparison, Dobrushin's covariance bounds, and the passage from decay of
correlations to a transfer-matrix gap (one direction of
Proposition~\ref{prop:gapdecay}). For the periodic state it proves a
transfer-matrix gap for every $N$ at $\beta_W<5.9\cdot10^{-5}$, and explicit
single-link heat-bath gaps on boxes of side fifteen for $\SU(2)$ at
$\beta_W\le0.166$ and for $\SU(3)$ at $\beta_W\le0.38$ (the coupling there is
$\beta_W/2$). It also proves, by a Poincar\'e inequality and a light-cone
argument different from Lemma~\ref{lem:speed}, that an ungauged one-cube
influence condition gives a lattice gap of order $1/n$; this is a variant of
the lattice step, Theorem~\ref{thm:decay} and Corollary~\ref{cor:latticegap},
at a smaller rate. The continuum step, the gauge
fixing and the curvature bounds are not part of that development, which shares
with this paper the open hypothesis (D).


\section{Lattice gauge theory and the transfer matrix}\label{sec:lattice}

Write $e_0,\dots,e_3$ for the unit vectors of $\Z^4$, $(x,\mu)$ for the link
from $x$ to $x+e_\mu$, and $Q_{(x,\mu)}^{-1}$ for the variable of the reversed
link. A box $\Lambda$ is a product of four finite intervals, with its set
$E(\Lambda)$ of links and its set $P(\Lambda)$ of plaquettes. For a plaquette
at $x$ in the plane $\mu<\nu$ with links $e_1=(x,\mu)$, $e_2=(x+e_\mu,\nu)$,
$e_3=(x+e_\nu,\mu)$, $e_4=(x,\nu)$, put $Q_p=Q_{e_1}Q_{e_2}Q_{e_3}^*Q_{e_4}^*$.
With $dQ$ the Haar probability measure of $G$, the Wilson measure of a box
with free boundary is
\[
W_\Lambda(dQ)=\frac1{Z_\Lambda}\exp\Bigl(N\beta\sum_{p\in P(\Lambda)}\ReTr
Q_p\Bigr)\prod_{e\in E(\Lambda)}dQ_e ,
\]
and the Wilson measure $\mu_L$ of the torus $\Lambda_L=(\Z/L\Z)^4$ is defined
in the same way with all links and plaquettes of $\Lambda_L$. Both are
invariant under the gauge transformations $Q_{(x,y)}\mapsto g_xQ_{(x,y)}
g_y^{-1}$. A function of finitely many links is a local observable, and it
is gauge invariant if it is invariant under these transformations. For
$\beta\ge0$ the measures are reflection positive for reflections in lattice
hyperplanes through sites and in hyperplanes bisecting links, on boxes
symmetric under the reflection and on tori of even side~\cite{OSei78}.

Let $\omega$ be a state on the gauge-invariant local observables of $\Z^4$
that is invariant under translations and reflection positive in this sense.
The Osterwalder--Schrader construction~\cite{OS73,OS75,OSei78} gives a
Hilbert space $\mathcal H$ with a unit vector $\Omega$, a map $F\mapsto\Psi_F$
from the observables at nonnegative times to $\mathcal H$ with dense range,
and a self-adjoint transfer matrix $0\le\mathbb T\le1$ with $\mathbb T\Omega=
\Omega$ and
\[
\ip{\Psi_F}{\mathbb T^t\Psi_H}=\omega\bigl(\vartheta(\bar F)\,\tau_tH\bigr),
\qquad t\in\Z_{\ge0},
\]
where $\vartheta$ is the time reflection and $\tau_t$ the time translation;
$\Psi_F$ is the class of $F$ for the positive semidefinite form
$\omega(\vartheta(\bar F)H)$, and $\|\Psi_F\|^2\le\omega(|F|^2)$. The
\emph{spectral gap} of $\mathbb T$ is $m$ if its spectrum on $\Omega^\perp$
lies in $[0,e^{-m}]$.

\begin{proposition}\label{prop:gapdecay}
Let $m>0$ and let $\mathcal F$ be a class of gauge-invariant local
observables, supported at nonnegative times, that is dense in $L^2(\omega)$
on each bounded time support. If for every $F\in\mathcal F$ there is
$C_F<\infty$ with
\[
\omega(\vartheta(\bar F)\tau_tF)-|\omega(F)|^2\le C_Fe^{-mt}\qquad(t\ge0),
\]
then $\mathbb T$ has spectral gap at least $m$.
\end{proposition}

\begin{proof}
Put $\Psi'_F=\Psi_F-\ip\Omega{\Psi_F}\Omega$, so that the left side equals
$\ip{\Psi'_F}{\mathbb T^t\Psi'_F}=\int\lambda^t\,d\sigma_F(\lambda)$, with
$\sigma_F$ the spectral measure of $\mathbb T$ at $\Psi'_F$. For $\delta>0$
this is at least $\sigma_F((e^{-m+\delta},1])e^{(-m+\delta)t}$, which exceeds
$C_Fe^{-mt}$ for large $t$ unless $\sigma_F((e^{-m+\delta},1])=0$. So the
spectral projection of $(e^{-m},1]$ annihilates every $\Psi'_F$. These vectors
are dense in $\Omega^\perp$, because $\|\Psi_F\|^2\le\omega(|F|^2)$ and the
class is dense, so the projection vanishes on $\Omega^\perp$.
\end{proof}

The decay of a single two-point function does not suffice: it places only
one spectral measure in $[0,e^{-m}]$, and the density of the class is what
gives the gap.


\section{The line field}\label{sec:linefield}

A \emph{direction field} is a map $\sigma:\Z^4\to\{0,1,2,3\}$; its links
are $F=\{(x,\sigma(x)):x\in\Z^4\}$. At each site put $\pi_x(\mu)=
\sigma(x+e_\mu)$. The plaquette at $x$ in the plane $\{\mu,\nu\}$ contains
\begin{equation}\label{eq:count}
[\sigma(x)\in\{\mu,\nu\}]+[\pi_x(\mu)=\nu]+[\pi_x(\nu)=\mu]
\end{equation}
links of $F$.

\begin{lemma}\label{lem:cycle}
Every plaquette contains exactly one link of $F$ if and only if, at every
site $x$, $\pi_x$ fixes $\sigma(x)$ and permutes the other three directions
cyclically.
\end{lemma}

\begin{proof}
Fix $x$, let $s=\sigma(x)$ and let $R$ be the other three directions. For the
plane $\{s,\nu\}$ the count~\eqref{eq:count} is $1+[\pi_x(s)=\nu]+[\pi_x(\nu)
=s]$, which equals one for all $\nu\in R$ exactly when $\pi_x$ fixes $s$ and
maps $R$ into $R$. For a plane $\{\mu,\nu\}\subset R$ the count is
$[\pi_x(\mu)=\nu]+[\pi_x(\nu)=\mu]$. Draw an arrow $\mu\to\pi_x(\mu)$ whenever
$\pi_x(\mu)\in R\setminus\{\mu\}$; each vertex has at most one outgoing arrow,
and the condition asks for exactly one arrow on each of the three pairs. Three
arrows on three vertices with out-degree at most one give out-degree one
everywhere, and such a tournament is a directed three-cycle.
\end{proof}

From now on $\sigma$ is the field of period two of Table~\ref{tab:sigma} in
Appendix~\ref{app:table}; in each row $\pi_x$ fixes $\sigma(x)$ and is a
three-cycle on the other directions, so every plaquette contains exactly one
link of $F$.\footnote{A field with one link of $F$ at every site and one in
every plaquette can exist only in four dimensions: a site carries
$d(d-1)/2$ plaquettes at which it is the base point and a link lies in
$2(d-1)$ plaquettes, and the two counts agree only for $d=4$.} Call the links
outside $F$ \emph{active} and write $A$ for their set.

\begin{lemma}\label{lem:structure}\leavevmode
\begin{enumerate}[label=\textup{(\roman*)},leftmargin=*]
\item $F$ is a disjoint union of bi-infinite straight lines, one through every
site; hence $F\cap E(\Lambda)$ is a forest for every box $\Lambda$.
\item An active link $e=(x,\mu)$ lies in six plaquettes. In four of them the
link of $F$ shares an end point with $e$; in the other two it is opposite to
$e$.
\item Every active link shares a plaquette with exactly twelve other active
links, two in each of its plaquettes, and two links share at most one
plaquette.
\end{enumerate}
\end{lemma}

\begin{proof}
(i) Since $\pi_x$ fixes $\sigma(x)$, $\sigma(x+e_{\sigma(x)})=\sigma(x)$, and since
$\sigma$ has period two, also $\sigma(x-e_{\sigma(x)})=\sigma(x)$. If $(x-e_\nu,\nu)
\in F$, straightness at $x-e_\nu$ gives $\sigma(x)=\nu$, so the incoming link of
$F$ at $x$ is unique and continues the outgoing one. A line meets a box in a
segment.

(ii) A link of $F$ sharing an end point with $e$ is the outgoing or incoming
link of $F$ at $x$, of direction $s=\sigma(x)\ne\mu$, or at $x+e_\mu$, of
direction $t=\pi_x(\mu)\in R\setminus\{\mu\}$. They lie in the two plaquettes of
$e$ in the plane $\{\mu,s\}$ and the two in the plane $\{\mu,t\}$, four distinct
plaquettes as $s\ne t$. In the two plaquettes of the remaining plane
$\{\mu,\rho\}$ the unique link of $F$ touches neither end of $e$, so it is the
opposite link.

(iii) Each plaquette of $e$ contains $e$, one link of $F$ and two other links.
\end{proof}


\section{Gauge fixing and the free-boundary state}\label{sec:state}

For a function $O$ of the links let $O^F$ be $O$ with every link of $F$ set to
the identity.

\begin{lemma}[gauge fixing]\label{lem:gaugefix}
Let $\mu^F_\Lambda$ be the probability measure on $G^{A\cap E(\Lambda)}$ with
density proportional to $\exp(N\beta\sum_{p\in P(\Lambda)}\ReTr Q^F_p)$. For
every gauge-invariant $O$ on $G^{E(\Lambda)}$, $W_\Lambda(O)=\mu^F_\Lambda(O^F)$.
\end{lemma}

\begin{proof}
Each segment of $F\cap E(\Lambda)$ is a path $x_0,\dots,x_k$. Given the
variables on $F$, put $g_{x_0}=1$ and $g_{x_{i+1}}=g_{x_i}Q_{(x_i,x_{i+1})}$
along every segment. The transformed field equals the identity on $F\cap
E(\Lambda)$, and for fixed values on $F$ the map from the active variables to
the transformed ones is a product of translations, which preserve Haar
measure. Gauge invariance of the action and of $O$ gives the claim.
\end{proof}

No holonomy remains because $F\cap E(\Lambda)$ is a forest. The measures
$\mu^F_\Lambda$ are finite-volume versions of the specification $\gamma^F$ on
$G^A$ in which $\gamma^F_e$ resamples an active $Q_e$ from the law with density
proportional to $\exp(N\beta\sum_{p\ni e}\ReTr Q^F_p)$ given the other active
links. Fix a bi-invariant metric $\delta$ on $G$ of diameter $D$ and define the
Kantorovich coefficients
\[
C_{ee'}=\sup\frac{W_1\bigl(\gamma^F_e(\cdot\mid Q),\gamma^F_e(\cdot\mid Q')\bigr)}
{\delta(Q_{e'},Q'_{e'})},
\]
the supremum over configurations that differ only at $e'$; $W_1$ is the
Kantorovich distance~\cite{Vil09}, and $C_{ee'}=0$ unless $e$ and $e'$ share a
plaquette.

\begin{lemma}[comparison]\label{lem:comparison}
Suppose every row and column of $(C_{ee'})$ sums to at most $\alpha<1$, and let
$\lambda>0$ with $\theta=\alpha e^\lambda<1$. Let $S\subset A$ and let $\mu_1,\mu_2$
be probability measures on $G^A$ invariant under $\gamma^F_e$ for $e\notin S$.
Let $T\subset A\setminus S$ be finite, $w_b=e^{-\lambda\dist(b,T)}$ with the
distance of link midpoints, and suppose $\sum_{z\in S}w_z<\infty$. There is a
coupling of $\mu_1$ and $\mu_2$ with $\varrho(b)=\E\,\delta(\sigma_b,\sigma'_b)$
satisfying
\[
\sum_{b\notin S}w_b\varrho(b)\le\frac{\theta D}{1-\theta}\sum_{z\in S}w_z ,
\]
and $\mu_1=\mu_2$ if $S=\emptyset$.
\end{lemma}

\begin{proof}
The couplings form a weakly compact set on which $\Psi(\pi)=\sum_bw_b\E_\pi
\delta(\sigma_b,\sigma'_b)$ is continuous; let $\pi^*$ minimise it. For $e\notin S$
replace $(\sigma_e,\sigma'_e)$ by a sample of an optimal coupling of
$\gamma^F_e(\cdot\mid\sigma)$ and $\gamma^F_e(\cdot\mid\sigma')$, chosen
measurably~\cite[Cor.~5.22]{Vil09}. Both $\mu_i$ are invariant, so this is again
a coupling, and changing the links that enter $\gamma^F_e$ one at a time gives
$\varrho(e)\le\sum_{e'}C_{ee'}\varrho(e')$ by minimality. Multiply by $w_e$ and sum
over $e\notin S$; where $C_{ee'}\ne0$, $w_e\le e^\lambda w_{e'}$, so
$\sum_{e\notin S}w_e\varrho(e)\le\theta[\sum_{e'\notin S}w_{e'}\varrho(e')+D\sum_{z
\in S}w_z]$. If $S=\emptyset$ this forces $\varrho\equiv0$.
\end{proof}

Conditioning on the links of the support of $H$~\cite[Ch.~3]{Kal21} and
applying the lemma to the two conditional measures gives, for the Gibbs
measure $\mu$ of $\gamma^F$, a local $F'$ that is $L$-Lipschitz in each link and
a bounded $H$ with finite support disjoint from that of $F'$,
\begin{equation}\label{eq:decay}
|\Cov_\mu(F',H)|\le\|H\|_\infty L\,\frac{\theta D}{1-\theta}\,|\operatorname{supp}H|\,
e^{-\lambda\dist(\operatorname{supp}F',\operatorname{supp}H)} .
\end{equation}

\begin{proposition}\label{prop:state}
Suppose every row and column of $(C_{ee'})$ sums to at most $\alpha<1$. Then
$\gamma^F$ has a unique Gibbs measure $\mu_F$; for every gauge-invariant local $O$
the limit $\omega(O)=\lim W_\Lambda(O)$ exists along every sequence of boxes
exhausting $\Z^4$ and equals $\mu_F(O^F)$; $\omega$ is invariant under
translations and the hypercubic group; truncated correlations of Lipschitz
gauge-invariant observables decay at every rate below $-\log\alpha$; and for
$\beta\ge0$, $\omega$ is reflection positive and its transfer matrix has
spectral gap at least $-\log\alpha$.
\end{proposition}

\begin{proof}
Gibbs measures exist by compactness and coincide by Lemma~\ref{lem:comparison}
with $S=\emptyset$. Extend $\mu^F_\Lambda$ by Haar measure outside $\Lambda$ and
let $S$ be the set of active links with a plaquette outside $P(\Lambda)$; off $S$
the extended measure is invariant under $\gamma^F_e$. Lemma~\ref{lem:comparison}
with $T$ the support of $O^F$ bounds $|\mu^F_\Lambda(O^F)-\mu_F(O^F)|$ by
$C_TR^3e^{-\lambda R}$, $R$ the distance from $T$ to the complement of $\Lambda$,
for Lipschitz $O$, and uniform approximation handles continuous $O$;
Lemma~\ref{lem:gaugefix} identifies $\mu^F_\Lambda(O^F)$ with $W_\Lambda(O)$. Since
$W_\Lambda(\tau_vO)=W_{\Lambda-v}(O)$ and the limit does not depend on the
exhausting sequence, $\omega$ is invariant; the line field is a device for
computing $W_\Lambda$ and is never moved. Reflection positivity of $W_\Lambda$ on
symmetric boxes passes to the limit. The decay is~\eqref{eq:decay}, since
$\Cov_\omega(O,O')=\Cov_{\mu_F}(O^F,O'^F)$. Lipschitz gauge-invariant functions
are dense in the gauge-invariant part of $L^2(\omega)$ on every support, being
Lipschitz functions on the compact quotient by the gauge group, so
Proposition~\ref{prop:gapdecay} gives the gap.
\end{proof}

\begin{proof}[Proof of Theorem~\ref{thm:A} for $\SU(2)$]
Identify $\SU(2)$ with the unit quaternions and use the chordal metric. For an
active link $e$ and a plaquette $p\ni e$ let $s_p$ be the product of the other
three variables, ordered so that $\operatorname{Re}(Q_es_p)=\frac12\Tr Q_p$, with
links of $F$ equal to one. Since $N\beta\ReTr Q_p=4\beta\ip{Q_e}{\bar s_p}$ in
$\R^4$, the conditional law of $Q_e$ is the von Mises--Fisher law $\nu_b(dx)
\propto e^{\ip xb}dx$ on $S^3$ with $b=4\beta\sum_{p\ni e}\bar s_p$. Moving a
partner $e'$ by $\delta$ moves one staple by an isometric image of that
displacement, hence $b$ by $4|\beta|\delta$, and the sharp
quarter (Theorem~\ref{q-thm:quarter} in Appendix~\ref{app:sharp}), $W_1(\nu_b,\nu_{b'})\le\frac14|b-b'|$, gives
$C_{ee'}\le|\beta|$. By Lemma~\ref{lem:structure}(iii) every row and column has
twelve nonzero entries, so $\alpha\le12|\beta|=3|\beta_W|$, and
Proposition~\ref{prop:state} applies for $|\beta_W|<\frac13$.
\end{proof}

The value $\frac13$ is the end of this argument: at $b=0$ the derivative of
$\nu_b$ has Kantorovich norm exactly $\frac14$, the six staples of an active
link can be arranged with sum zero while one of them moves freely, and so
every row sum of the matrix is exactly $3|\beta_W|$. Without the gauge every
link has eighteen partners and the same argument stops at $\beta_W=\frac29$
\cite[Prop.~13]{BB26a}.


\section{Curvature in the gauge}\label{sec:curvature}

In this section $G\subset\U(N)$ is closed and connected with Lie algebra $\g$
and the inner product $\ip XY=-\Tr(XY)$. A tangent vector at a field $Q$ is
$v=(X_eQ_e)$ with $X_e\in\g$, and $|v|^2=\sum_e|X_e|^2$. For the plaquette $p$
at $x$ in the plane $\mu<\nu$, with links $e_1,\dots,e_4$ as in
Sec.~\ref{sec:lattice}, transport the vectors to $x$,
\[
Y_1=X_{e_1},\quad Y_2=Q_{e_1}X_{e_2}Q_{e_1}^*,\quad
Y_3=Q_{e_4}X_{e_3}Q_{e_4}^*,\quad Y_4=X_{e_4},
\]
and write $W=Q_p$. Along the geodesics $Q_e(t)=e^{tX_e}Q_e$ the plaquette
variable is $Q_p(t)=e^{tY_1}e^{tY_2}We^{-tY_3}e^{-tY_4}$, and differentiating
twice gives the second derivative of $\ReTr Q_p$,
\begin{equation}\label{eq:fp}
f_p=\sum_{i=1}^4\ReTr(Y_i^2W)+2\ReTr(Y_1Y_2W)+2\ReTr(Y_3Y_4W)
-2\!\!\sum_{i\in\{1,2\},\,j\in\{3,4\}}\!\!\ReTr(Y_iWY_j),
\end{equation}
so that $\Hess S_0(v,v)=\sum_pf_p$ for $S_0=\sum_p\ReTr Q_p$. The
\emph{corner form} of $p$ is
\begin{equation}\label{eq:corner}
c_p=2\ip{Y_1}{Y_4}-2\ip{Y_1}{Y_2}+2\ip{Y_2}{WY_3W^*}-2\ip{Y_3}{Y_4};
\end{equation}
it collects the inner products of the vectors of two links that meet at a
right angle at a corner of $p$.

\begin{lemma}[divergence identity]\label{lem:identity}
On $\Z^4$, or on a box with free boundary, for every field $Q$ and every
square-summable $X$,
\[
\sum_y\Bigl[|D_yX|^2+\sum_\mu\bigl|X_{y,\mu}+Q^*_{y-e_\mu,\mu}X_{y-e_\mu,\mu}
Q_{y-e_\mu,\mu}\bigr|^2\Bigr]=4|X|^2+\sum_pc_p ,
\]
where $D_yX=\sum_\mu[X_{y,\mu}-Q^*_{y-e_\mu,\mu}X_{y-e_\mu,\mu}Q_{y-e_\mu,\mu}]$
is the covariant divergence and terms with a link outside the box are
omitted.
\end{lemma}

\begin{proof}
At a site $y$ let $o_\mu$ be the vector of the outgoing link of direction $\mu$
and $i_\mu$ the transported vector of the incoming one, so that $D_yX=\sum_\mu
(o_\mu-i_\mu)$. Every link is outgoing at one end and incoming at the other, so
the squares $|o_\mu|^2+|i_\mu|^2$, which occur once in $|D_yX|^2$ and once in
$|o_\mu+i_\mu|^2$, add up to $4|X|^2$. The collinear cross term
$-2\ip{o_\mu}{i_\mu}$ of $|D_yX|^2$ cancels the term $+2\ip{o_\mu}{i_\mu}$ of
$|o_\mu+i_\mu|^2$. Every remaining cross term of $|D_yX|^2$ pairs two
perpendicular ends at $y$, and such a pair is a corner of exactly one
plaquette, with the sign and transport of~\eqref{eq:corner}; in a box two
perpendicular ends at $y$ span a plaquette of the box.
\end{proof}

\begin{lemma}[one plaquette]\label{lem:local}
Let $k\in\{1,2,3,4\}$, let $k\pm1$ \textup{(}mod $4$\textup{)} be the
positions adjacent to $k$ and $k+2$ the opposite one. For every $N$, every
$W\in\U(N)$ and all $Y_1,\dots,Y_4\in\mathfrak u(N)$ with $Y_k=0$,
\[
\pm f_p+\tfrac9{10}c_p\le\tfrac{21}{10}\bigl(|Y_{k-1}|^2+|Y_{k+1}|^2\bigr)
+\tfrac54|Y_{k+2}|^2 .
\]
\end{lemma}

\begin{proof}
Replacing $W$ by $-W$ changes the sign of $f_p$ and fixes $c_p$, so the sign
$+$ suffices, and conjugation by a unitary matrix preserves both sides, so we
may take $W=\operatorname{diag}(w_a)$. The difference $\mathcal F$ of the two
sides then splits into a block for each diagonal position $a$, in the real
variables $y_i=-i(Y_i)_{aa}$, and a block for each pair $a\ne b$, in the
complex variables $z_i=(Y_i)_{ab}$, since $(Y_i)_{ba}=-\bar z_i$.
Appendix~\ref{app:blocks} shows that every block is positive definite.
\end{proof}

\begin{theorem}[Hessian bound]\label{thm:hess}
For every closed subgroup $G\subset\U(N)$, every box or $\Z^4$, every field
$Q$ and every tangent vector $v$ that vanishes on the links of $F$,
\[
|\Hess S_0(v,v)|\le\tfrac{29}2|v|^2 .
\]
\end{theorem}

\begin{proof}
Every plaquette has exactly one link of $F$, on which $X$ vanishes, and by
Lemma~\ref{lem:structure}(ii) each active link is adjacent to that link in
four of its plaquettes and opposite to it in two. Summing
Lemma~\ref{lem:local} over the plaquettes,
\[
\sum_p\bigl(\pm f_p+\tfrac9{10}c_p\bigr)\le\bigl(4\cdot\tfrac{21}{10}+2\cdot
\tfrac54\bigr)|X|^2=\tfrac{109}{10}|X|^2 .
\]
By Lemma~\ref{lem:identity}, $\frac9{10}\sum_pc_p=\frac9{10}\Sigma-\frac{36}{10}
|X|^2$ with $\Sigma\ge0$ the left side of that lemma. Hence $\pm\Hess S_0(v,v)
\le(\frac{109}{10}+\frac{36}{10})|X|^2=\frac{29}2|X|^2$; links near the boundary
of a box lie in fewer plaquettes and receive less. Both lemmas concern
vectors in $\mathfrak u(N)$, so the bound holds on every closed subgroup, whose
geodesics $e^{tX}Q_e$ with $X\in\g$ are geodesics of $\U(N)$.
\end{proof}

The factor $\frac9{10}$ lets the divergence identity pay for the corners at
the price $\frac{36}{10}$ while the blocks of one plaquette stay positive.
Without the gauge the sharp bound is $16$ \cite{BB26a}.

\begin{lemma}[Ricci curvature]\label{lem:ricci}
Let $G\subset\U(N)$ be closed and connected with simple Lie algebra, and let
$\mathfrak t\subset\g$ be a maximal abelian subalgebra. Write the eigenvalues
of $\operatorname{ad}H$ on $\g\otimes\C$ and of $H$ on $\C^N$, $H\in\mathfrak t$, as
$i\alpha(H)$ and $i\mu(H)$, over the roots $\alpha$ and the weights $\mu$ counted
with multiplicity. Then $\Ric_G=\kappa_G\ip\cdot\cdot$ with
\[
\kappa_G=\frac{\sum_\alpha\alpha(H)^2}{4\sum_\mu\mu(H)^2}=\frac{h^\vee}{2\ell_\rho}
\qquad(0\ne H\in\mathfrak t).
\]
\end{lemma}

\begin{proof}
For a bi-invariant metric $R(X,Y)Z=-\frac14[[X,Y],Z]$ on left-invariant fields
\cite[Sect.~7]{Mil76}, so for an orthonormal basis $(E_i)$,
$\Ric(X,X)=\frac14\sum_i|[X,E_i]|^2=-\frac14\Tr(\operatorname{ad}X)^2$. The
Killing form and $\ip\cdot\cdot$ are invariant and $\g$ is simple, so the
Killing form is $-4\kappa\ip\cdot\cdot$ for a constant $\kappa$, and on
$\mathfrak t$ the two forms are $-\sum_\alpha\alpha(H)^2$ and $\sum_\mu\mu(H)^2$.
For $H$ the coroot of the highest root, $\sum_\alpha\alpha(H)^2=4h^\vee$ and
$\sum_\mu\mu(H)^2=2\ell_\rho$~\cite[Ch.~6]{Kac90}.
\end{proof}

\begin{table}[t]
\caption{The curvature $\kappa_G=h^\vee/(2\ell_\rho)$ in a faithful
representation of least dimension $N$, and the range of Theorem~\ref{thm:A}
for $\beta_W=N^2\beta$.}
\label{tab:groups}
\begin{tabular}{llcccc}
\toprule
Group&$\rho$&$h^\vee$&$\ell_\rho$&$\kappa_G$&$\beta_W<$\\
\midrule
$\SU(N)$, $N\ge2$&$\C^N$&$N$&$1$&$\frac N2$&$\frac{N^2}{29}$\\
$\SO(N)$, $N\ge5$&$\R^N$&$N-2$&$2$&$\frac{N-2}4$&$\frac{N(N-2)}{58}$\\
$\mathrm{Sp}(n)$, $n\ge2$&$\C^{2n}$&$n+1$&$1$&$\frac{n+1}2$&$\frac{2n(n+1)}{29}$\\
$G_2$&$7$&$4$&$2$&$1$&$\frac{14}{29}$\\
$F_4$&$26$&$9$&$6$&$\frac34$&$\frac{39}{29}$\\
$E_6$&$27$&$12$&$6$&$1$&$\frac{54}{29}$\\
$E_7$&$56$&$18$&$12$&$\frac34$&$\frac{84}{29}$\\
$E_8$&$248$&$30$&$60$&$\frac14$&$\frac{124}{29}$\\
\bottomrule
\end{tabular}
\end{table}

For two links $e\ne e'$ of a plaquette $p$ write $\Hess\ReTr Q_p(v,v)=
\sum_{e,e'\in p}H^p_{ee'}(X_e,X_{e'})$ with $H^p_{ee'}(a,b)=H^p_{e'e}(b,a)$. By
\eqref{eq:fp} a mixed block is, after isometric transport, $\pm\ReTr(abW)$ or
$\pm\ReTr(aWb)$, so
\begin{equation}\label{eq:mixed}
|H^p_{ee'}(a,b)|\le|a|\,|b| .
\end{equation}

\subsection{Decay with an explicit rate}

Let $A'$ be a finite set of active links and $P'$ a finite set of plaquettes
whose active links lie in $A'$. Call $S=\sum_{p\in P'}c_p\ReTr Q^F_p$, $0\le
c_p\le1$, \emph{admissible} if $|\Hess S(v,v)|\le\frac{29}2|v|^2$ for $v$
tangent to $G^{A'}$. By Theorem~\ref{thm:hess} the action $S_\Lambda=\sum_{p\in
P(\Lambda)}\ReTr Q^F_p$ of a box is admissible, and so is every convex combination
of admissible actions; by~\eqref{eq:mixed} and Lemma~\ref{lem:structure}(iii)
the mixed blocks of $\Hess S$ have norm at most one and each link has at most
twelve partners. Let $\mu$ have density proportional to $e^{N\beta S}$ on
$G^{A'}$, let $L=\Delta+N\beta\ip{\nabla S}{\nabla\cdot}$ be its Langevin generator,
self-adjoint in $L^2(\mu)$, and $P_t=e^{tL}$. For weights $w>0$ on $A'$ put
$\Gamma^w(f)=\sum_ew_e|\nabla_ef|^2$, with $\nabla_e$ the gradient in the factor of
$e$; $\Gamma$ is the case $w\equiv1$.

\begin{lemma}\label{lem:weighted}
Let $G$ have semisimple Lie algebra, $\lambda\ge0$, $w_e=e^{\lambda d_e}$ with
$|d_e-d_{e'}|\le1$ for partners, and
\[
K_\lambda=\kappa_G-N|\beta|\Bigl(\tfrac{29}2+12\bigl(\cosh\tfrac\lambda2-1\bigr)\Bigr).
\]
Then $\Gamma^w(P_tf)\le e^{-2K_\lambda t}P_t\Gamma^w(f)$ for smooth $f$ and $t\ge0$.
\end{lemma}

\begin{proof}
Put $V=-N\beta S$. On the product $G^{A'}$ the projection to one factor is
parallel and the Ricci tensor is block diagonal, so Bochner's formula
applied factor by factor gives
\[
\tfrac12L\Gamma^w(g)-\Gamma^w(g,Lg)=\sum_ew_e\bigl(|\nabla\nabla_eg|^2+\Ric_G(\nabla_e
g,\nabla_eg)\bigr)+\sum_{e,e'}w_e\Hess V_{ee'}(\nabla_eg,\nabla_{e'}g).
\]
With $u_e=\nabla_eg$, $v_e=\sqrt{w_e}u_e$ and the symmetry of $\Hess V$, the last
sum is
\[
\Hess V(v,v)+\sum_{e\ne e'}\Bigl(\cosh\tfrac{\lambda(d_e-d_{e'})}2-1\Bigr)\Hess
V_{ee'}(v_e,v_{e'}),
\]
since $(w_e+w_{e'})/(2\sqrt{w_ew_{e'}})=\cosh\frac{\lambda(d_e-d_{e'})}2$. The first
term is at least $-\frac{29}2N|\beta||v|^2$; in the second only partners
contribute, each with coefficient at most $\cosh\frac\lambda2-1$ and a block of
norm at most $N|\beta|$, and with twelve partners it is at least
$-12(\cosh\frac\lambda2-1)N|\beta||v|^2$. As $\Ric_G\ge\kappa_G$ and $|v|^2=
\Gamma^w(g)$, the left side is at least $K_\lambda\Gamma^w(g)$. For
$\varphi(s)=P_s\Gamma^w(P_{t-s}f)$ this gives $\varphi'\ge2K_\lambda\varphi$, hence
$\varphi(0)\le e^{-2K_\lambda t}\varphi(t)$.
\end{proof}

\begin{lemma}\label{lem:decayrate}
Suppose $K_\lambda>0$. Let $f,g$ be smooth functions of the links of $A_f$ and
$B_g$, and let $D$ be the graph distance from $A_f$ to $B_g$ in the graph on
$A'$ whose edges join partners. Then
\[
|\Cov_\mu(f,g)|\le\frac{e^{-\lambda D/2}}{K_\lambda}\,\|\nabla f\|_\infty
\sum_{e\in B_g}\|\nabla_eg\|_\infty .
\]
\end{lemma}

\begin{proof}
Lemma~\ref{lem:weighted} with $\lambda=0$ shows that $P_tf\to\mu(f)$ uniformly,
so $\Cov_\mu(f,g)=\int_0^\infty\mu(\Gamma(P_tf,g))\,dt$. Let $d_e$ be the smaller of
$D$ and the graph distance from $e$ to $A_f$; then $w=1$ on $A_f$, $w=e^{\lambda
D}$ on $B_g$, and $\Gamma^w(f)\le\|\nabla f\|_\infty^2$. For $e\in B_g$,
Lemma~\ref{lem:weighted} gives $|\nabla_eP_tf|^2\le e^{-\lambda D}e^{-2K_\lambda t}
\|\nabla f\|_\infty^2$, and integration in $t$ gives the bound.
\end{proof}

\begin{proof}[Proof of Theorem~\ref{thm:A}]
Let $\lambda^*$ solve $K_{\lambda^*}=0$, so that $\lambda^*/2=m_G(\beta)$, and let
$0<\lambda<\lambda^*$. Lemma~\ref{lem:decayrate} with $S=S_\Lambda$ gives
$|\Cov_{\mu^F_\Lambda}(f,g)|\le C_{f,g}e^{-\lambda D/2}$ uniformly in the box. For
boxes $\Lambda\subset\Lambda'$ interpolate $S_s=(1-s)S_\Lambda+sS_{\Lambda'}$, which is
admissible; for smooth $f$ of the active links of $\Lambda$,
\[
\frac d{ds}\mu_s(f)=N\beta\sum_{p\in P(\Lambda')\setminus P(\Lambda)}\Cov_{\mu_s}
\bigl(f,\ReTr Q^F_p\bigr),
\]
with $|\nabla_e\ReTr Q_p|\le\sqrt N$. A plaquette outside $P(\Lambda)$ is at graph
distance at least $R-1$ from the support of $f$, $R$ the distance to the
complement of $\Lambda$, and at most $C(k+1)^4$ plaquettes lie at graph distance
$k$, so $|\mu^F_{\Lambda'}(f)-\mu^F_\Lambda(f)|\le C_fR^4e^{-\lambda R/2}$. The limit
exists along every increasing sequence of boxes, two sequences are compared
through their intersections, and continuous gauge-invariant observables are
uniform limits of smooth ones $e^{\tau\Delta}O$, the heat semigroup commuting
with translations of $G$. With Lemma~\ref{lem:gaugefix} the limit $\omega$ exists;
invariance and reflection positivity follow as in Proposition~\ref{prop:state}.
Since $\Ric_G\ge0$, $e^{\tau\Delta}O$ has the Lipschitz constant of $O$
\cite[Sect.~3.3]{BGL14}, so the decay passes to Lipschitz observables; two links
of a plaquette have times differing by at most one, so observables at time
separation $t$ are at graph distance at least $t-c$, and truncated
correlations decay at every rate below $m_G(\beta)$. Proposition~\ref{prop:gapdecay}
gives the gap.
\end{proof}

The rate behaves like $\log\frac{\kappa_G}{6N|\beta|}$ as $\beta\to0$. The proof
also gives a Poincar\'e inequality with constant $1/(\kappa_G-\frac{29}2N|\beta|)$
for every $\mu^F_\Lambda$, uniformly in the box~\cite[Cor.~5.7.2]{BGL14}.


\section{One cube and the continuum gap}\label{sec:cube}

\emph{Setting.} Let $G\subset\U(N)$ be compact and nontrivial and $\varphi$ a
positive continuous class function on $G$, such as $e^{N\beta\ReTr(\cdot)}$. On
the torus $\Lambda_L=(\Z/L\Z)^4$, $L$ even and $L\ge2n+4$, with links $E$ and
plaquettes $P$, let
\[
\mu(dU)=\frac1Z\prod_{p\in P}\varphi(U_p)\prod_{e\in E}dU_e .
\]
A link $(y,\nu)$ has base point $y$; differences of base points are taken in
$(-L/2,L/2]^4$ and $|\cdot|_\infty$ is the corresponding distance. For
$x\in\Lambda_L$ let $V_x$ be the set of links with base point in $x+[0,n)^4$;
every link lies in exactly $n^4$ of the $L^4$ cubes $V_x$. For $V\subset E$ let
$\partial V$ be the links outside $V$ that share a plaquette with a link of
$V$; the conditional law $\gamma^0_V(\cdot\mid U)$ of $U_V$ given the other links
has density proportional to $\prod_{p\cap V\ne\emptyset}\varphi(U_p)$ and depends
only on $U_{\partial V}$. Let $\mathcal P_xf(U)=\int f\,d\gamma^0_{V_x}(\cdot\mid U)$,
the conditional expectation given the links outside $V_x$, an orthogonal
projection on $L^2(\mu)$, and let $\mathcal A$ be the closed subspace of
gauge-invariant functions. By gauge covariance of $\gamma^0_{V_x}$, $\mathcal P_x$
maps $\mathcal A$ into itself.

\begin{definition}\label{def:bp}
For $n\ge2$ and $\varpi>0$, the condition $\BP(n,\varpi)$ for $\mu$ is
\[
\varpi\Var_\mu(f)\le\frac1{n^4}\sum_{x\in\Lambda_L}\mu\bigl((f-\mathcal P_xf)^2\bigr)
\qquad\text{for every real }f\in\mathcal A .
\]
\end{definition}

Put $q_n=(2+1/n)^4$, so that $16<q_n\le\frac{625}{16}$ for $n\ge2$, and let
$\mathcal K=L^{-4}\sum_x\mathcal P_x$ be the chain that resamples one uniformly
chosen cube; for $Y\subset\Lambda_L$ let $\mathcal K_Y=L^{-4}(\sum_{x\in Y}\mathcal P_x
+(L^4-|Y|)I)$ be the same chain with the moves outside $Y$ suppressed. Write
$y\rightsquigarrow z$ if $\partial V_z\cap V_y\ne\emptyset$. A link of $\partial
V_z$ has its base point in $z+[-1,n]^4$, so $y\rightsquigarrow z$ implies
$|z-y|_\infty\le n$, and for each $z$ at most $(2n+1)^4=q_nn^4$ points $y$ have
$y\rightsquigarrow z$. For finite $S\subset E$ let $Y_1(S)=\{y:V_y\cap S\ne
\emptyset\}$ and $Y_{j+1}(S)=Y_j(S)\cup\{y:y\rightsquigarrow z\text{ for some }
z\in Y_j(S)\}$, $Y_0(S)=\emptyset$.

\begin{lemma}[finite speed]\label{lem:speed}
Let $f$ be a bounded function of the links of a finite set $S$, let $m_0\ge0$
and $Y=Y_{m_0}(S)$. Then $\mathcal K_Y^tf$ depends only on the links of
$S\cup\bigcup_{y\in Y}\partial V_y$, and for every integer $t\ge0$
\[
\bigl\|\mathcal K^tf-\mathcal K_Y^tf\bigr\|_\infty\le\frac{2|S|}{q_n}\,\|f\|_\infty
\sum_{m>m_0}\frac{(q_n\tau)^m}{m!},\qquad\tau=\frac{tn^4}{L^4}.
\]
\end{lemma}

\begin{proof}
If $h$ depends only on the links of $W$, then $\mathcal P_xh=h$ when $V_x\cap W=
\emptyset$, and otherwise $\mathcal P_xh$ depends only on the links of $(W\setminus
V_x)\cup\partial V_x$; induction gives the first claim. For the second run two
chains from the same configuration with the same uniformly chosen cubes
$x_1,x_2,\dots$. If $x_s\notin Y$ the first chain resamples $V_{x_s}$ and the
second does nothing; if $x_s\in Y$ and the two configurations agree on
$\partial V_{x_s}$, both receive one common sample, and otherwise they resample
independently. The marginals are the chains $\mathcal K$ and $\mathcal K_Y$. A
link where the configurations differ at time $t$ lies in $V_{x_{s_m}}$ for some
times $s_1<\dots<s_m\le t$ with $x_{s_1}\notin Y$ and $x_{s_l}\rightsquigarrow
x_{s_{l+1}}$, by induction on $t$. If the set of such links meets $S$, then
$x_{s_m}\in Y_1(S)$ and $x_{s_{m-l}}\in Y_{l+1}(S)$, so $m>m_0$. For fixed $m$
there are at most $t^m/m!$ choices of times, $n^4|S|$ of $x_{s_m}$ and $q_nn^4$
of each earlier cube, each sequence having probability $L^{-4m}$; the
difference of $f$ is at most $2\|f\|_\infty$.
\end{proof}

For finite $S,S'\subset E$ let $d(S,S')$ be the least $|b-b'|_\infty$ over base
points $b$ of links of $S$ and $b'$ of links of $S'$.

\begin{theorem}[decay from one cube]\label{thm:decay}
Suppose $\BP(n,\varpi)$ holds for $\mu$. Then $\varpi\le4$, and for bounded
real $f,g\in\mathcal A$ depending on the links of finite sets $S_f,S_g$ with
$d=d(S_f,S_g)\ge3n-2$,
\[
|\Cov_\mu(f,g)|\le\Bigl(e+\frac{k|S_f||S_g|}{64e^2}\Bigr)\|f\|_\infty\|g\|_\infty\,
e^{-\varpi k/(e^2q_n)},\qquad k=\Bigl\lfloor\frac{d+2}n\Bigr\rfloor-2 .
\]
In particular $|\Cov_\mu(f,g)|\le C(1+d)e^{-\varpi d/(e^2q_nn)}$ for every $d$,
with $C$ depending only on $|S_f|,|S_g|,\|f\|_\infty,\|g\|_\infty$.
\end{theorem}

\begin{proof}
\emph{The bound $\varpi\le4$.} Take $f=\ReTr U_p$ for one plaquette. Each
$I-\mathcal P_x$ is an orthogonal projection annihilating constants, so
$\mu((f-\mathcal P_xf)^2)\le\Var_\mu(f)$, and it vanishes unless $V_x$ contains
a link of $p$, which happens for at most $4n^4$ points $x$.

\emph{Spectral part.} On the functions of $\mathcal A$ with mean zero, $\mathcal K$
is self-adjoint with $0\le\mathcal K\le I$, and $\BP(n,\varpi)$ reads
$\ip h{(I-\mathcal K)h}\ge\epsilon\|h\|^2$ with $\epsilon=\varpi n^4/L^4$. So
$\|\mathcal K^Th\|_2\le e^{-\epsilon T}\|h\|_2$, and from
$f-\mu(f)=\mathcal K^T(f-\mu(f))+\sum_{t<T}\mathcal K^t(I-\mathcal K)f$,
\[
\Cov_\mu(f,g)=\ip{\mathcal K^T(f-\mu f)}{g-\mu g}+\sum_{t=0}^{T-1}
\ip{\mathcal K^tf}{(I-\mathcal K)g},
\]
whose first term is at most $e^{-\epsilon T}\|f\|_\infty\|g\|_\infty$.

\emph{Local part.} Let $Z_g=Y_1(S_g)$, so $|Z_g|\le n^4|S_g|$ and $(I-\mathcal K)g=
L^{-4}\sum_{x\in Z_g}(g-\mathcal P_xg)$; if $h$ does not depend on the links of
$V_x$ then $\ip h{g-\mathcal P_xg}=0$. Put $m_0=k-1$ and $Y=Y_{m_0}(S_f)$. A link
of $V_x$, $x\in Z_g$, lying in $S_f$ or in $\partial V_y$ with $y\in Y$ would join
base points of $S_f$ and $S_g$ at distance at most $n(m_0+3)-3\le n\lfloor
(d+2)/n\rfloor-3<d$, which is impossible. So $\mathcal K_Y^tf$ does not depend on
the links of $V_x$, $x\in Z_g$, and by Lemma~\ref{lem:speed}
\[
\bigl|\ip{\mathcal K^tf}{(I-\mathcal K)g}\bigr|=\bigl|\ip{\mathcal K^tf-\mathcal K_Y^tf}
{(I-\mathcal K)g}\bigr|\le\frac{4|S_f||S_g|n^4}{q_nL^4}\|f\|_\infty\|g\|_\infty
\sum_{m\ge k}\frac{(q_n\tau_t)^m}{m!}.
\]
Summing over $t<T$ with $\tau=Tn^4/L^4$, the second term is at most
\[
\frac{4|S_f||S_g|\tau}{q_n}\,\|f\|_\infty\|g\|_\infty\sum_{m\ge k}
\frac{(q_n\tau)^m}{m!}.
\]

\emph{Choice of $T$.} Take $T=\lfloor kL^4/(e^2q_nn^4)\rfloor$, so that $z=q_n
\tau\le k/e^2$ and $\epsilon T\ge\varpi k/(e^2q_n)-1$, using $\varpi\le4$ and
$L^4\ge4n^4$. Since $k!\ge(k/e)^k$, $\sum_{m\ge k}z^m/m!\le(ez/k)^ke^z\le
e^{-k(1-e^{-2})}$, and $\tau/q_n\le k/(256e^2)$. As $\varpi/(e^2q_n)<1-e^{-2}$,
both terms are bounded as stated. For the second form use $k>d/n-3$, and for
$d<3n-2$ the trivial bound $\|f\|_\infty\|g\|_\infty$.
\end{proof}

A \emph{torus limit} at $\beta$ is a state $\omega$ on the gauge-invariant local
functions of $\Z^4$ with $\omega(O)=\lim_j\mu_{L_j}(O)$ along some $L_j\to\infty$;
by compactness every sequence has a subsequence along which one exists.

\begin{corollary}[lattice gap]\label{cor:latticegap}
Let $\beta\ge0$ and suppose $\BP(n,\varpi)$ holds for $\mu_L$ for all $L$ in an
infinite set of even integers $L\ge2n+4$. Then every torus limit along this
set is invariant under translations, reflection positive, and its transfer
matrix has spectral gap at least
\[
m_n=\frac{\varpi}{e^2q_nn}\ge\frac{16\varpi}{625e^2n}.
\]
\end{corollary}

\begin{proof}
Translation invariance and reflection positivity are identities and
inequalities among finitely many local functions of $\mu_L$~\cite{OSei78}, and
pass to $\omega$. For a continuous gauge-invariant $F$ at positive times the
supports of $\vartheta(\bar F)$ and $\tau_tF$ are at distance $d$ with $|d-t|\le
c_F$, and Theorem~\ref{thm:decay} on $\Lambda_L$ gives $|\Cov_{\mu_L}(\vartheta(\bar
F),\tau_tF)|\le C_F(1+t)e^{-m_nt}$ uniformly in $L$. In the limit
$|\omega(\vartheta(\bar F)\tau_tF)-|\omega(F)|^2|\le C_F(1+t)e^{-m_nt}$, and
Proposition~\ref{prop:gapdecay}, applied with every $m<m_n$, gives the gap.
\end{proof}

\subsection{The continuum limit}

Let $a_k\to0$ be lattice spacings with $a_k=2^{-k}a_0$, and $\beta_k$ couplings.
Let $\omega_k$ be lattice states with transfer matrices $\mathbb T_k$ on spaces
$\mathcal H_k$ with vacua $\Omega_k$. Let $\mathcal T$ be a complex vector space of
test functions on $\R^4$ supported in $\{t>0\}$, and for each $k$ let $f\mapsto
O_k(f)$ be a linear map to gauge-invariant local observables at positive
times; the example is $O_k(f)=Z_k\sum_xa_k^4f(a_kx)\Phi(x)$ with a local field
$\Phi$ and a renormalisation constant $Z_k$.

\begin{definition}\label{def:limit}
A \emph{continuum limit} of $(\omega_k,O_k)$ is a Hilbert space $\mathcal H$ with
a unit vector $\Omega$, a self-adjoint $H\ge0$ with $H\Omega=0$, and a linear
map $\psi:\mathcal T\to\mathcal H$ with dense range such that for all
$f,g\in\mathcal T$ and dyadic $\tau\ge0$
\[
\ip{\psi(f)}{e^{-\tau H}\psi(g)}=\lim_k\omega_k\bigl(\vartheta(\overline{O_k(f)})
\,\tau_{\tau/a_k}O_k(g)\bigr),\qquad
\ip\Omega{\psi(f)}=\lim_k\omega_k(O_k(f)),
\]
over the $k$ with $\tau/a_k\in\Z$. It is \emph{nontrivial} if $\mathcal H\ne\C
\Omega$.
\end{definition}

The Osterwalder--Schrader reconstruction produces such a limit from
convergent renormalised Schwinger functions satisfying the
axioms~\cite{OS73,OS75}; the definition records only what is used below.

\begin{proof}[Proof of Theorem~\ref{thm:B}]
Let $\omega_k$ be torus limits at $\beta_k$ along the tori on which
$\BP(n_k,\varpi)$ holds. By Corollary~\ref{cor:latticegap}, $\mathbb T_k$ has gap
at least $16\varpi/(625e^2n_k)$ on $\Omega_k^\perp$. For $f\in\mathcal T$ put
$O=O_k(f)$ and $\varphi_k=\Psi_O-\omega_k(O)\Omega_k$. For dyadic $\tau$ and $k$
large, $t=\tau/a_k$ is an integer and $16\varpi t/(625e^2n_k)\ge\Delta\tau$ since
$n_ka_k\le\ell$, so
\[
\omega_k\bigl(\vartheta(\bar O)\tau_tO\bigr)-|\omega_k(O)|^2=\ip{\varphi_k}{\mathbb T_k^t
\varphi_k}\le e^{-\Delta\tau}\|\varphi_k\|^2 .
\]
Both sides converge by Definition~\ref{def:limit}: with $\varphi=\psi(f)-\ip\Omega
{\psi(f)}\Omega$, $\ip\varphi{e^{-\tau H}\varphi}\le e^{-\Delta\tau}\|\varphi\|^2$. These
vectors form a subspace dense in $\Omega^\perp$, and the inequality extends to
$\Omega^\perp$ and to all $\tau\ge0$ by continuity. For $\varphi\in\Omega^\perp$ with
spectral measure $\mu_\varphi$ of $H$ this says $\int e^{-\tau E}d\mu_\varphi(E)\le e^{-\Delta
\tau}\mu_\varphi([0,\infty))$ for all $\tau$, which forces $\mu_\varphi([0,\Delta-\eta])=0$
for every $\eta>0$. So $0$ is a simple eigenvalue and the spectrum on
$\Omega^\perp$ lies in $[\Delta,\infty)$, and is not empty if $\Omega^\perp\ne\{0\}$.
\end{proof}

The proof uses the cube condition only through the decay of two-point
functions of gauge-invariant observables, and the constant $\Delta$ depends on
$\varpi$ and $\ell$ alone.

\section{The cube condition at strong coupling}\label{sec:bpstrong}

Keep the setting of Sec.~\ref{sec:cube}, with a bi-invariant metric $\delta$
of diameter $D$ on $G$, invariant under inversion, and $d_E(U,U')=\sum_e
\delta(U_e,U'_e)$. Two distinct links are \emph{partners} if they share a
plaquette; two links share at most one plaquette, and every link has eighteen
partners, three in each of its six plaquettes. The direction field $\sigma$ of
Sec.~\ref{sec:linefield} has period two, so it is defined on $\Lambda_L$. By
Lemma~\ref{lem:structure}(i) the links of $F$ with an end point $v$ are exactly
$(v,\sigma(v))$ and $(v-e_{\sigma(v)},\sigma(v))$, and $\sigma(v-e_{\sigma(v)})=
\sigma(v)$.

\begin{proposition}\label{prop:bpsup}
At $\beta=0$, $\BP(n,1)$ holds for every $n$.
\end{proposition}

\begin{proof}
For $\varphi\equiv1$, $\mu$ is a product of Haar measures. Write $f=\sum_Sf_S$
with $f_S$ a function of the links of $S$ and $\int f_S\,dU_e=0$ for $e\in S$;
the $f_S$ are orthogonal and $\mathcal P_xf=\sum_{S\cap V_x=\emptyset}f_S$, so
$\sum_x\mu((f-\mathcal P_xf)^2)=\sum_{S\ne\emptyset}\|f_S\|^2\#\{x:S\cap V_x\ne
\emptyset\}\ge n^4\Var_\mu(f)$.
\end{proof}

\subsection*{Gauge fixing inside a cube}
For a cube $Q_x=x+[0,n)^4$ write $\hat v=v-x$ and call $v$ \emph{interior} if
$\hat v\in[1,n-1]^4$; let $I_x$ be the set of interior vertices. For $z\in I_x$
put $\tau(z)=(z-e_{\sigma(z)},\sigma(z))$, a link of $F$ with head $z$, and
$T_x=\{\tau(z):z\in I_x\}$. Every link with an end point in $I_x$ has its base
point in $Q_x$, so $T_x\subset V_x$, the map $\tau$ is injective, and no link
outside $V_x$ touches $I_x$. Call the links of $V_x\setminus T_x$ \emph{free}.
For $\psi\in G^{V_x\setminus T_x}$ let $[\psi,U]$ be the configuration equal to
$\psi$ on the free links, to $1$ on $T_x$ and to $U$ outside $V_x$, and let
$\gamma^T_{x,U}$ be the probability measure on $G^{V_x\setminus T_x}$ with
density proportional to $\prod_{p\cap V_x\ne\emptyset}\varphi([\psi,U]_p)$.

\begin{lemma}\label{lem:cubegauge}
For every bounded measurable $f\in\mathcal A$,
$\mathcal P_xf(U)=\int f([\psi,U])\,\gamma^T_{x,U}(d\psi)$.
\end{lemma}

\begin{proof}
Let $\sigma'\in G^{V_x}$ and write $(\sigma',U)$ for the configuration equal to
$\sigma'$ on $V_x$ and to $U$ elsewhere. Order $I_x$ by the sum of the
coordinates of $\hat z$; then $z-e_{\sigma(z)}$ precedes $z$ or is not
interior. Put $g_v=1$ for $v\notin I_x$ and recursively $g_z=g_{z-e_{\sigma(z)}}
\sigma'_{\tau(z)}$. The gauge transformation $g$ sets every link of $T_x$ to $1$
and fixes every link outside $V_x$, so $g\cdot(\sigma',U)=[\psi,U]$ with
$\psi=(g\cdot\sigma')|_{V_x\setminus T_x}$. Since $g$ depends only on
$\sigma'_{T_x}$, for fixed $\sigma'_{T_x}$ the map $\sigma'_{V_x\setminus T_x}
\mapsto\psi$ multiplies each coordinate on both sides by fixed elements, so
$\sigma'\mapsto(\psi,\sigma'_{T_x})$ preserves Haar measure. Plaquette variables
change by conjugation and $\varphi$ is a class function, so the density of
$\gamma^0_{V_x}(\cdot\mid U)$ at $\sigma'$ is proportional to
$\prod_{p\cap V_x\ne\emptyset}\varphi([\psi,U]_p)$, while $f(\sigma',U)=
f([\psi,U])$. Integrating out $\sigma'_{T_x}$ gives the claim.
\end{proof}

A link $(b,\mu)$ of $V_x$ is a \emph{face link} if $\hat b_i=0$ for some
$i\ne\mu$, that is, if both its end points lie in one lower face
$\{\hat v_i=0\}$ of the cube.

\begin{lemma}\label{lem:fifteen}
\begin{enumerate}[label=\textup{(\roman*)},leftmargin=*]
\item If all four links of a plaquette are free, its link of $F$ is a face
link.
\item Every free link has at most fifteen free partners.
\end{enumerate}
\end{lemma}

\begin{proof}
(i) Let $\phi=(y,s)$ be the link of $F$ of the plaquette $p$ and $h=y+e_s$.
If $h$ were interior, then $\sigma(h)=s$ and $\phi=\tau(h)\in T_x$. Since $y\in
Q_x$, $\hat h_i=\hat y_i\in[0,n-1]$ for $i\ne s$ and $\hat h_s\in[1,n]$, so
$\hat h_s=n$ or $\hat y_i=0$ for some $i\ne s$. If $\hat h_s=n$, the link of
$p$ parallel to its second direction $\nu$ through $h$ has a base point with
$s$-coordinate $n$ and is not in $V_x$. Hence $\phi$ is a face link.

(ii) Let $k=(b,\mu)$ be free and $m(p)$ the number of free links of a
plaquette $p\ni k$; $k$ has $\sum_{p\ni k}(m(p)-1)$ free partners. If $k$ is a
face link, $\hat b_i=0$ with $i\ne\mu$, the plaquette of $k$ with base point
$b-e_i$ in the plane $\{\mu,i\}$ has its other three links based at
$i$-coordinate $-1$, outside $V_x$; the other five plaquettes give at most $3$
each, $15$ in all. Otherwise $k$ has at most $12+M$ free partners, where $M$
counts the plaquettes $p\ni k$ with $m(p)=4$, whose link $\phi$ of $F$ is a face
link by (i), so $\phi\ne k$. If $k\in F$ then $M=0$; let $k$ be active and
$\phi$ lie in the face $\{\hat v_i=0\}$. If $\phi$ shares an end point $v$ with
$k$, then $v=b$ (for $v=b+e_\mu$ either $k$ is a face link or $\hat b_i=-1$),
$\hat b_i=0$ and $i=\mu$; there are two links of $F$ at $b$, each in at most
one plaquette of $k$. If $\phi$ is opposite to $k$, it is $(b\pm e_\nu,\mu)$
in a plaquette of the plane $\{\mu,\nu\}$; as $k$ is not a face link, $\nu=i$,
the sign is $-$ and $\hat b_i=1$, and by Lemma~\ref{lem:structure}(ii) the
plane $\{\mu,\nu\}$ is the one plane in which the link of $F$ is opposite to
$k$, so this happens for at most one plaquette. Hence $M\le3$.
\end{proof}

For a free link $e$ the conditional law $\gamma^T_e(\cdot\mid\psi,U)$ of
$\psi_e$ under $\gamma^T_{x,U}$ has density proportional to $\prod_{p\ni e}
\varphi([\psi,U]_p)$. Let $C^T_{ek}$ be the supremum of $W_1(\gamma^T_e(\cdot
\mid\psi,U),\gamma^T_e(\cdot\mid\psi',U'))/\delta([\psi,U]_k,[\psi',U']_k)$ over
pairs for which $[\psi,U]$ and $[\psi',U']$ differ only at a link $k$, free or
outside $V_x$; it vanishes unless $e$ and $k$ are partners. For $j\notin V_x$
let $b(V_x,j)$ be the number of links of $V_x$ that are partners of $j$, and
\[
R_x(j)=\sup\frac{W_1^{d_x}(\gamma^T_{x,U},\gamma^T_{x,U'})}{\delta(U_j,U'_j)},
\qquad d_x(\psi,\psi')=\sum_{e\in V_x\setminus T_x}\delta(\psi_e,\psi'_e),
\]
the supremum over configurations that differ only at $j$.

\begin{lemma}\label{lem:bpblock}
If $C^T_{ek}\le c$ for all free $e$ and all $k$, with $15c<1$, then
$R_x(j)\le c\,b(V_x,j)/(1-15c)$.
\end{lemma}

\begin{proof}
Let $U,U'$ differ only at $j$, let $\pi$ minimise $\int d_x\,d\pi$ over
couplings of $\gamma^T_{x,U}$ and $\gamma^T_{x,U'}$, and let $\varrho(e)=\int
\delta(\psi_e,\psi'_e)\,d\pi$. Replacing $(\psi_e,\psi'_e)$ by a measurably
chosen optimal coupling of the two conditional laws of $\psi_e$
\cite[Cor.~5.22]{Vil09} gives another coupling, so by minimality and the
triangle inequality $\varrho(e)\le c\sum_k\varrho(k)+c[e\text{ a partner of }j]
\delta(U_j,U'_j)$, the sum over the free partners $k$ of $e$. Summing over $e$
and using Lemma~\ref{lem:fifteen}(ii), $(1-15c)\sum_e\varrho(e)\le
c\,b(V_x,j)\delta(U_j,U'_j)$.
\end{proof}

\begin{lemma}\label{lem:bpcount}
For $n\ge2$, $L\ge2n+4$ and a link $j$, $\sum_{x:\,j\notin V_x}b(V_x,j)=
18n^3-3n^2$.
\end{lemma}

\begin{proof}
Let $b$ be the base point of $j$ and $\nu$ its direction. For each $\mu\ne\nu$
the partners in the two plaquettes of the plane $\{\mu,\nu\}$ have base points
$b+\Delta$, $\Delta\in\{e_\nu,e_\mu,0,-e_\mu,-e_\mu,e_\nu-e_\mu\}$. With $z=b-x$
and $Q=[0,n)^4$, a link based at $b+\Delta$ lies in $V_x$ exactly when $z+\Delta\in
Q$, and $j\in V_x$ when $z\in Q$. The sum is $\sum_\Delta(n^4-|Q\cap(Q-\Delta)|)$,
which is $0$ for $\Delta=0$, $n^3$ for $\Delta=\pm e_\mu,e_\nu$ and $2n^3-n^2$ for
$\Delta=e_\nu-e_\mu$, that is $6n^3-n^2$ for each of the three $\mu$.
\end{proof}

\begin{lemma}\label{lem:onelink}
\begin{enumerate}[label=\textup{(\roman*)},leftmargin=*]
\item For $G=\SU(2)$, the chordal metric and $\varphi=e^{2\beta\ReTr(\cdot)}$,
$C^T_{ek}\le|\beta|=|\beta_W|/4$.
\item Let $G$ have semisimple Lie algebra, $\delta$ the geodesic distance of
$\ip XY=-\Tr(XY)$, and $6N|\beta|<\kappa_G$. Then $C^T_{ek}\le N|\beta|/
(\kappa_G-6N|\beta|)$.
\end{enumerate}
\end{lemma}

\begin{proof}
(i) The conditional law of $\psi_e$ is the von Mises--Fisher law $\nu_b$ with
$b=4\beta\sum_{p\ni e}\bar s_p$, where $s_p$ is the product of the other three
variables of $p$, as in Sec.~\ref{sec:state}. A partner $k$ enters one
$s_p$, and moving it by $\delta_k$ moves $\bar s_p$ by $\delta_k$, since left and
right multiplication and inversion are isometries; the sharp quarter,
Theorem~\ref{q-thm:quarter}, gives the bound. Links of $T_x$ do not move.

(ii) The conditional law is $\nu\propto e^{-V}d\psi_e$ with $V=-N\beta
\sum_{p\ni e}\ReTr[\psi,U]_p$; by~\eqref{eq:fp} each term has second
derivative $\ReTr(Y^2W)$ with $|Y|=|X|$, so $\Ric_G+\Hess V\ge K_1=
\kappa_G-6N|\beta|$. Lemma~\ref{lem:weighted} with one factor and the argument
of Lemma~\ref{lem:decayrate} give $|\Cov_\nu(g,h)|\le\|\nabla g\|_\infty
\|\nabla h\|_\infty/K_1$. Move the partner $k$ along a unit-speed geodesic; the
derivative of $\nu_s(g)$ is $-\Cov_{\nu_s}(g,\partial_sV_s)$, and $\partial_sV_s$
comes from the one plaquette containing $e$ and $k$, with gradient in $\psi_e$
of norm at most $N|\beta|$ by~\eqref{eq:mixed}. For smooth $1$-Lipschitz $g$
this gives $|\nu_0(g)-\nu_\delta(g)|\le N|\beta|\delta/K_1$; convolution on $G$
preserves Lipschitz constants for a bi-invariant metric, and
Kantorovich--Rubinstein duality concludes.
\end{proof}

\begin{proposition}\label{prop:interiorbp}
If $\sum_{x:\,i\notin V_x}R_x(i)\le\alpha n^4$ for every link $i$, with
$\alpha<1$, then $\BP(n,1-\alpha)$ holds.
\end{proposition}

\begin{proof}
For a function $h$ let $\mathrm{Lip}(h)$ be the least constant with
$|h(U)-h(U')|\le\mathrm{Lip}(h)\delta(U_i,U'_i)$ whenever $U,U'$ differ only at
one link $i$; then $|h(U)-h(U')|\le\mathrm{Lip}(h)d_E(U,U')$. Let $h\in\mathcal A$
be Lipschitz and $U,U'$ differ only at $i$. If $i\in V_x$,
$\mathcal P_xh(U)=\mathcal P_xh(U')$. If $i\notin V_x$, Lemma~\ref{lem:cubegauge}
and an optimal coupling $\pi$ of $\gamma^T_{x,U}$ and $\gamma^T_{x,U'}$ give
\[
|\mathcal P_xh(U)-\mathcal P_xh(U')|\le\int|h([\psi,U])-h([\psi',U'])|\,d\pi
\le\mathrm{Lip}(h)\bigl(R_x(i)+1\bigr)\delta(U_i,U'_i).
\]
Since $i$ lies in $n^4$ of the cubes, $\mathrm{Lip}(\mathcal Kh)\le(1-\kappa_0)
\mathrm{Lip}(h)$ with $\kappa_0=(1-\alpha)n^4/L^4$, and as $\mathcal K$
preserves $\mathcal A$, $\Var_\mu(\mathcal K^th)\le\frac12(1-\kappa_0)^{2t}
\mathrm{Lip}(h)^2D^2|E|^2$. The restriction of $\mathcal K$ to the mean-zero
part of $\mathcal A$ is self-adjoint with $0\le\mathcal K\le I$, and
$\Var_\mu(\mathcal K^th)=\int\lambda^{2t}d\nu_h$ for its spectral measure $\nu_h$
at $h-\mu(h)$, so $\nu_h$ lives on $[0,1-\kappa_0]$. Lipschitz functions are
dense in $L^2(\mu)$, and averaging over the compact gauge group, the
orthogonal projection onto $\mathcal A$, does not increase $\mathrm{Lip}$,
because gauge transformations act link by link by isometries; so Lipschitz
functions are dense in $\mathcal A$, and $\ip f{(I-\mathcal K)f}\ge\kappa_0
\Var_\mu(f)$ for real $f\in\mathcal A$, which is $\BP(n,1-\alpha)$.
\end{proof}

\begin{theorem}\label{thm:bpstrong}
Let $n\ge2$ and $L\ge2n+4$. If $C^T_{ek}\le c$ for all cubes, free links $e$
and links $k$, with $15c<1$, then $\BP(n,1-\alpha_n)$ holds whenever
\[
\alpha_n=\frac{c\,(18n-3)}{(1-15c)\,n^2}<1 .
\]
In particular:
\begin{enumerate}[label=\textup{(\roman*)},leftmargin=*]
\item for $G=\SU(2)$ and $|\beta_W|<\frac4{15}$, with $\alpha_n=|\beta_W|(18n-3)/
((4-15|\beta_W|)n^2)$;
\item for $G\subset\U(N)$ with semisimple Lie algebra and $21N|\beta|<\kappa_G$,
with $\alpha_n=N|\beta|(18n-3)/((\kappa_G-21N|\beta|)n^2)$.
\end{enumerate}
For fixed $\beta$ in these ranges $\alpha_n$ decreases to $0$, so $\BP(n,\varpi)$
holds with $\varpi\to1$ as $n\to\infty$, uniformly in the torus.
\end{theorem}

\begin{proof}
Lemmas~\ref{lem:bpblock} and~\ref{lem:bpcount} give $\sum_{x:\,i\notin V_x}
R_x(i)\le c(18n^3-3n^2)/(1-15c)=\alpha_nn^4$, and
Proposition~\ref{prop:interiorbp} applies. Lemma~\ref{lem:onelink} gives
$c=|\beta_W|/4$ for $\SU(2)$, and $c=N|\beta|/(\kappa_G-6N|\beta|)$ in
general, for which $15c<1$ reads $21N|\beta|<\kappa_G$ and $c/(1-15c)=
N|\beta|/(\kappa_G-21N|\beta|)$.
\end{proof}

For $\SU(2)$ the condition $\alpha_n<1$ reads $|\beta_W|(15n^2+18n-3)<4n^2$; at
$|\beta_W|=\frac14$ it holds for $n\ge18$, and its threshold increases to
$\frac4{15}$ as $n\to\infty$. At strong coupling the condition is therefore
\emph{easier} for large cubes, because the response of a cube to one boundary
link is spread over $n^4$ links while it touches only $O(n^3)$ of them.

\subsection*{A weighted count}
For a free link $e$ let $N(e)$ be its set of free partners. The proof of
Lemma~\ref{lem:bpblock} gives, for each free $e$,
\begin{equation}\label{eq:rho}
\varrho(e)\le\sum_{k\in N(e)}C^T_{ek}\varrho(k)+C^T_{ej}[e\text{ a partner of }j]
\,\delta(U_j,U'_j).
\end{equation}

\begin{lemma}\label{lem:weightedresp}
Let $w$ be a function on the free links with values in $[w_-,w_+]$, $w_->0$,
with $\sum_{k\in N(e)}w(k)\le\lambda w(e)$ for every free $e$. If $C^T_{ek}\le
c$ for all free $e$ and all $k$, and $c\lambda<1$, then
$R_x(j)\le(w_+/w_-)\,c\,b(V_x,j)/(1-c\lambda)$.
\end{lemma}

\begin{proof}
Multiply~\eqref{eq:rho} by $w(e)$ and sum over $e$. Since $k\in N(e)$ exactly
when $e\in N(k)$, $\sum_ew(e)\varrho(e)\le c\lambda\sum_kw(k)\varrho(k)+
c\,w_+b(V_x,j)\delta(U_j,U'_j)$, and $\sum_e\varrho(e)\le w_-^{-1}\sum_ew(e)
\varrho(e)$.
\end{proof}

Write a link of $V_x$ as $(x+b,\mu)$ with $b\in[0,n-1]^4$, let
$z(e)=\#\{i\ne\mu:b_i=0\}$ be the number of lower faces containing it, and
call it a \emph{foot} if $b_\mu=0$. On free links put $w=45$ if $z=1$ and $e$
is active and not a foot, $w=47$ if $z=1$ and $e\in F$ or $e$ is a foot,
$w=48$ if $z=0$ and $e$ is an active foot, and $w=40$ otherwise.

\begin{proposition}\label{prop:weights}
For every $n\ge2$, every cube and every free link $e$,
$\sum_{k\in N(e)}w(k)\le\frac{72}5w(e)$.
\end{proposition}

\begin{proof}
Work in the coordinates $b$, so that $(y,\nu)\in F$ when $\sigma(x+y)=\nu$.
Let $e=(b,\mu)$. For $\nu\ne\mu$ the other links of the plaquette $p^+_\nu$ of
base point $b$ in the plane $\{\mu,\nu\}$ are $P_1=(b+e_\mu,\nu)$,
$P_2=(b+e_\nu,\mu)$, $P_3=(b,\nu)$, and those of the plaquette $p^-_\nu$ of base
point $b-e_\nu$ are $P_4=(b-e_\nu,\mu)$, $P_5=(b-e_\nu+e_\mu,\nu)$,
$P_6=(b-e_\nu,\nu)$. If $e$ is active, with $s=\sigma(x+b)$,
$t=\sigma(x+b+e_\mu)$ and $\rho$ the fourth direction, the links of $F$ among
them are $P_3,P_6$ for $\nu=s$, $P_1,P_5$ for $\nu=t$ and $P_2,P_4$ for
$\nu=\rho$: the first four by Lemma~\ref{lem:structure}(i) at $b$ and at
$b+e_\mu$, and in the plane $\{\mu,\rho\}$ the link of $F$ is opposite to $e$
by Lemma~\ref{lem:structure}(ii). If $e\in F$, none of them is in $F$.

\emph{Reduction to the orthant.} Call a link with base point in
$O=\Z^4_{\ge0}$ \emph{$O$-free} if it is active or has $z\ge1$, extend $w$ to
such links by the same rule, and count as partners of $e$ the $O$-free links
of the plaquettes of $e$ with base point in $O$. If $b_\nu=0$, the links of
$p^-_\nu$ other than $e$ lie outside $V_x$. A link $(y,\mu')\in V_x$ lies in
$T_x$ exactly when it is in $F$ and its head is interior, that is, $z=0$ and
$y_{\mu'}\le n-2$; so free links are $O$-free except the \emph{exceptional}
ones, links of $F$ with $z=0$ and $y_{\mu'}=n-1$, of weight $40$. An
exceptional link $\phi$ in a plaquette $p$ is its link of $F$, and $p$ has a
link $\ell$ based at $\mu'$-coordinate $n$, which lies outside $V_x$, is
active, has its base point in $O$ (as $z(\phi)=0$), and has $w(\ell)\ge40$;
exchanging $\phi$ for $\ell$ shows that the orthant sum dominates the sum over
$N(e)$ for non-exceptional $e$. An exceptional $e$ has at most two partners in
$V_x$ in each plaquette, each of weight at most $48$, and $12\cdot48=576=
\frac{72}5\cdot40$.

\emph{The orthant count.} Fix $\nu\ne\mu$, put $f=[b_\mu=0]$, $\beta=b_\nu$,
and let $z'$ be the number of zeros among $b_i$, $i\notin\{\mu,\nu\}$. From
the base points, $z(P_1)=z(P_2)=z(P_5)=z'$, $z(P_3)=z(P_6)=z'+f$,
$z(P_4)=z'+[\beta=1]$; $P_2,P_4$ are feet when $f=1$, $P_1,P_3$ when
$\beta=0$, $P_5,P_6$ when $\beta=1$; and $p^-_\nu$ counts only when
$\beta\ge1$. An $O$-free link has weight $40$ if $z\ge2$, $47$ or $45$ if
$z=1$ according as it is in $F$ or a foot or neither, and $48$ or $40$ if
$z=0$ (then it is active) according as it is a foot or not; a link of $F$ with
$z=0$ contributes nothing. Let $C(\text{role},f,\beta,z')$ be the resulting
weight of the $O$-free links among $P_1,\dots,P_6$, the role of $\nu$ being
$s$, $t$, $\rho$, or $\mathrm F$ when $e\in F$. For $z'=2$, $C=120$ if
$\beta=0$ and $C=240$ otherwise; for $z'\le1$ the values are
\begin{center}
\small
\begin{tabular}{cc|rrr|rrr}
\toprule
&&\multicolumn{3}{c|}{$z'=0$, $\beta=$}&\multicolumn{3}{c}{$z'=1$, $\beta=$}\\
role&$f$&$0$&$1$&${\ge}2$&$0$&$1$&${\ge}2$\\
\midrule
$s$&0&88&173&160&139&271&274\\
$s$&1&143&277&270&134&259&264\\
$t$&0&88&173&160&139&271&274\\
$t$&1&95&187&186&134&261&268\\
$\rho$&0&96&223&160&141&271&274\\
$\rho$&1&95&227&170&134&259&264\\
$\mathrm F$&0&136&261&240&139&269&270\\
$\mathrm F$&1&143&275&266&134&259&264\\
\bottomrule
\end{tabular}
\end{center}
For instance, for the role $s$, $f=1$, $\beta=1$, $z'=0$ the weights of
$P_1,\dots,P_6$ are $40,48,47,47,48,47$. The orthant sum is
$\sum_{\nu\ne\mu}C(\text{role}(\nu),f,b_\nu,z'_\nu)$, where $z'_\nu=z(e)-1$ if
$b_\nu=0$ and $z'_\nu=z(e)$ otherwise, and $\frac{72}5w(e)$ is $576$, $648$,
$676.8$ or $691.2$ for $w(e)=40,45,47,48$. For active $e$ with $z=0$ the sum
is at most $173+173+223=569$ ($f=0$, $w=40$) or $277+187+227=691$ ($f=1$,
$w=48$). For active $e$ with $z=1$ it is at most $96+274+274=644$ ($f=0$,
$w=45$), and for $f=1$, $w=47$, at most $143+268+264$, $95+264+264$ or
$95+264+268$ as the zero direction has role $s$, $t$ or $\rho$, all at most
$675$. For active $e$ with $z=2$ it is at most $141+139+240=520$, for $z=3$
it is $360$, with $w=40$. For $e\in F$ (then $z\ge1$) it is at most
$136+270+270=676$ or $143+264+264=671$ when $z=1$, $w=47$; at most $518$ when
$z=2$; and $360$ when $z=3$, with $w=40$.
\end{proof}

\begin{theorem}\label{thm:bpweighted}
Let $n\ge2$ and $L\ge2n+4$. If $C^T_{ek}\le c$ for all cubes, free links $e$
and links $k$, with $\frac{72}5c<1$, then $\BP(n,1-\alpha_n)$ holds whenever
$\alpha_n=6c(18n-3)/((5-72c)n^2)<1$. In particular:
\begin{enumerate}[label=\textup{(\roman*)},leftmargin=*]
\item for $G=\SU(2)$ and $|\beta_W|<\frac5{18}$, with
$\alpha_n=3|\beta_W|(18n-3)/(2(5-18|\beta_W|)n^2)$;
\item for $G=\SU(3)$, the Hilbert--Schmidt distance and $0\le\beta_W\le0.45$.
\end{enumerate}
\end{theorem}

\begin{proof}
Lemma~\ref{lem:weightedresp} with Proposition~\ref{prop:weights},
$\lambda=\frac{72}5$ and $w_+/w_-=\frac65$, and Lemma~\ref{lem:bpcount} give
$\sum_{x:\,i\notin V_x}R_x(i)\le\frac65c(18n^3-3n^2)/(1-\frac{72}5c)=
\alpha_nn^4$, and Proposition~\ref{prop:interiorbp} applies; its proof uses
only that $\delta$ is bi-invariant and invariant under inversion. For (i),
$c=|\beta_W|/4$ by Lemma~\ref{lem:onelink}. For (ii) let $r=2\beta_W$,
$Z_*(r)=\int_{\SU(3)}e^{r\ReTr U}dU$, $\epsilon(r)=1-5e^{-r/5}/(9Z_*(r))$ and
$\gamma(r)=\frac{283}{128}-\frac r2$. The conditional law of a free link is
$\nu_b\propto e^{\ReTr(Ub)}dU$ with $b$ a sum of $N\beta=r/6$ times six
products of elements of $\SU(3)$, the links of $T_x$ being $1$, so the
one-link estimate of the $\SU(3)$ theorem of~\cite{BBLRS26} applies and gives
$c=r\sqrt{\epsilon(r)}/(6\gamma(r))$, increasing in $r$. With the moments
$m_k=\int(\ReTr U)^kdU$, $(m_0,\dots,m_8)=(1,0,\frac12,\frac14,\frac34,
\frac{15}{16},\frac{65}{32},\frac{231}{64},\frac{973}{128})$, and $|m_k|\le
3^{k-8}m_8$ for $k\ge8$,
\[
Z_*(r)\le\sum_{k=0}^8\frac{m_kr^k}{k!}+\frac{m_8}{3^8}\frac{(3r)^9}{9!}
\Bigl(1-\frac{3r}{10}\Bigr)^{-1}.
\]
At $r=0.9$ the sum is at most $1.2599146$ and the tail at most
$3.4\cdot10^{-5}$, so $Z_*(0.9)\le1.2599486$; with $e^{-0.18}\ge1-0.18+0.0162
-0.000972$ this gives $\epsilon(0.9)\le0.63172<0.7949^2$, and with
$\gamma(0.9)=\frac{1127}{640}$, $\frac{72}5c\le2.16\cdot0.7949\cdot640/1127
\le0.9751$.
\end{proof}

For fixed $n$ the threshold of (i) exceeds that of
Theorem~\ref{thm:bpstrong}(i) exactly when $n^2-6n+1>0$, that is for $n\ge6$;
both theorems hold. The same arithmetic at $r=0.88$, where $Z_*\le1.246559$,
$\epsilon\le0.62627$ and $\gamma=\frac{5667}{3200}$, gives $15c\le0.9832$, so
Theorem~\ref{thm:bpstrong} holds for $\SU(3)$ with $0\le\beta_W\le0.44$.

Without the gauge fixing every link has eighteen partners, and the same
argument needs $18c<1$, that is $|\beta_W|<\frac29$ for $\SU(2)$. No choice of
gauge-fixed links and weights goes much further:

\begin{proposition}\label{prop:methodlimit}
Let $T\subset V_x$ have at most $(n-1)^4$ links and $A$ be the partner matrix
of $V_x\setminus T$. Then $\lambda_{\max}(A)\ge6r(4-r)(3-r)/(4-r^2)$ with
$r=(1-1/n)^2$, which tends to $12$. If $c\lambda_{\max}(A)\ge1$, the
inequalities $0\le\varrho\le cA\varrho+s$ used in Lemma~\ref{lem:bpblock}
admit solutions with $\sum_e\varrho(e)$ arbitrarily large, so for $\SU(2)$ the
argument cannot pass $|\beta_W|=\frac13$.
\end{proposition}

\begin{proof}
The $P=6n^2(n-1)^2=6rn^4$ plaquettes with all links in $V_x$ contain at most
$6|T|=aP$ links of $T$ in all, $a\le r$, so their numbers $m(p)$ of links outside
$T$ have mean at least $4-a\ge3$; by convexity $\mathbf1^{\top}A\mathbf1\ge
\sum_pm(p)(m(p)-1)\ge P(4-a)(3-a)$, and the Rayleigh quotient of $\mathbf1$ is
at least $6r(4-a)(3-a)/(4-ar)$, which decreases in $a$ because
$r/(4-ar)\le1/(4-a)$. For a Perron vector $v\ge0$ of $A$, $\varrho=tv$
satisfies $cA\varrho\ge\varrho$ for every $t>0$.
\end{proof}


\section{Stars: the cube condition for \texorpdfstring{$\SU(2)$}{SU(2)} past \texorpdfstring{$\beta_W=\frac13$}{beta\_W=1/3}}\label{sec:stars}

Proposition~\ref{prop:methodlimit} shows that no argument built on single-link
updates and gauge-fixed links passes $\beta_W=\frac13$ for $\SU(2)$. This
section replaces single links by \emph{stars}, the eight links at a vertex of
one parity, and passes that value: the cube condition holds uniformly in the
cube size for $|\beta_W|\le\frac7{20}$
(Theorem~\ref{thm:starsuniform}), and a block condition giving a lattice gap
holds for $|\beta_W|\le\frac9{25}$ (Theorem~\ref{thm:stars}). The gain is
visible at first order: the perturbation of one plaquette at a vertex enters
the law of the star only through the relative variable of two of its links,
so a star responds with the coefficient $c$ where two single links respond
with $2c$, and no plaquette is internal to a star.\footnote{A Dobrushin
certificate bounds the response of a block to a change of its boundary by a
matrix of local influences, and contraction means that this matrix has
spectral radius below one. Grouping links into blocks changes the matrix:
interactions inside a block are treated exactly instead of being bounded.}

\emph{Setting.} Throughout the section $G=\SU(2)$, identified with the unit
quaternions $S^3\subset\R^4$, with the Euclidean inner product
$\ip\cdot\cdot$ of $\R^4$ and the chordal metric $\delta(U,V)=|U-V|$, which is
bi-invariant, invariant under inversion, and of diameter $2$. For unit
quaternions $\ReTr U=2\ip U1$, and for unit $p,q$ and $x,y\in\R^4$,
$\ip{px}y=\ip x{\bar py}$, $\ip{xq}y=\ip x{y\bar q}$ and $|pxq|=|x|$. The
Wilson measure is $\mu(dU)=Z^{-1}\prod_pe^{2\beta\ReTr U_p}\prod_edU_e$ on
$\Lambda_L$, $L$ even, with $\beta_W=4\beta$, and $\mathcal A$, $V_x$,
$\mathcal P_x$ and $\BP(n,\varpi)$ are as in Sec.~\ref{sec:cube}. For a set
$V$ of links, $\mathcal P_V$ is the conditional expectation given the links
outside $V$. For $b\in\R^4$, $\nu_b(dx)\propto e^{\ip xb}dx$ is the von
Mises--Fisher law on $S^3$, and the sharp quarter (Theorem~\ref{q-thm:quarter}
in Appendix~\ref{app:sharp}) states
\begin{equation}\label{eq:Q}
W_1(\nu_b,\nu_{b'})\le\theta\,|b-b'|\quad\text{for all }b,b'\in\R^4,
\qquad\theta=\tfrac14,
\end{equation}
with $W_1$ the Kantorovich distance for $\delta$. This is the only property of
$\nu_b$ used in the section. We put
\[
c=4\theta|\beta|=|\beta|=\tfrac14|\beta_W| .
\]
Every argument below uses only~\eqref{eq:Q}, and with a constant $\theta$ in
place of $\frac14$ the ranges become $4\theta|\beta|\le\frac7{80}$ and
$4\theta|\beta|\le\frac9{100}$.

\subsection{Even stars and two-paths}

Call a vertex \emph{even} if the sum of its coordinates is even. For an even
vertex $v$ the \emph{star} $B_v$ is the set of the eight links with an end
point at $v$. Write the eight directions as $\pm1,\dots,\pm4$, and for the link
of $B_v$ in direction $\sigma$ let $X_\sigma$ be its parallel transporter from
$v$ outwards: $X_{+\mu}=U_{(v,\mu)}$ and $X_{-\mu}=U_{(v-e_\mu,\mu)}^{-1}$. Two
directions are \emph{adjacent} if they are distinct and not opposite; each
direction is adjacent to six others, and there are $24$ adjacent pairs.

\begin{lemma}\label{lem:stars}
Let $L\ge4$ be even.
\begin{enumerate}[label=\textup{(\roman*)},leftmargin=*]
\item Every link has exactly one even end point, so the stars partition $E$.
\item Every plaquette has exactly two even corners $v,v'$, which are opposite
corners, and two of its links lie in $B_v$, two in $B_{v'}$. The plaquettes
with corner $v$ correspond bijectively to the adjacent pairs $\{\sigma,\tau\}$:
the plaquette $p_{\sigma\tau}$ has corners $v$, $v+\sigma$, $v+\sigma+\tau$,
$v+\tau$.
\item For such a plaquette, $\ReTr U_p=2\ip{h_p}{X_\sigma^{-1}X_\tau}$, where
the \emph{two-path} $h_p$ is the transporter from $v+\sigma$ to $v+\tau$
through the opposite even corner $v'=v+\sigma+\tau$, the product of the two
links of $p$ in $B_{v'}$, suitably inverted. $h_p$ is unchanged by a gauge
transformation at $v'$, and changing one of its two links by $\eta$ in
$\delta$ changes $h_p$ by exactly $\eta$.
\item The conditional law of $(X_\sigma)_\sigma$ given the links outside $B_v$ is
\[
\gamma_{v,h}(dX)=\frac1{Z_v}\prod_{\{\sigma,\tau\}}e^{4\beta\ip{h_{\sigma\tau}}
{X_\sigma^{-1}X_\tau}}\prod_\sigma dX_\sigma ,
\]
with one ordering fixed in each adjacent pair. It depends on the
configuration only through the $24$ two-paths $h=(h_{\sigma\tau})$, and it is
invariant under $X_\sigma\mapsto gX_\sigma$, the gauge transformation at $v$.
\item A link $i$ with even end point $v_i$ lies in exactly six plaquettes;
their even corners other than $v_i$ are six distinct even vertices, and for
each of them, $v$, the link $i$ is a factor of exactly one two-path $h_p$ with
$p$ at $v$.
\end{enumerate}
\end{lemma}

\begin{proof}
(i) The end points of a link differ by a unit vector, so their parities
differ. (ii) The corners of a plaquette are $y$, $y+e_\mu$, $y+e_\mu+e_\nu$,
$y+e_\nu$, of parities $\epsilon,\epsilon+1,\epsilon,\epsilon+1$, and each
link joins consecutive corners. A plaquette with even corner $v$ is spanned by
two adjacent directions at $v$, and conversely; since $L\ge4$ these $24$
squares are distinct. (iii) Going around $p$ from $v$, $U_p$ is conjugate to
$X_\sigma h_pX_\tau^{-1}$ or to its inverse, and $\ReTr$ is a class function,
invariant under inversion on $\SU(2)$; so $\ReTr U_p=\ReTr(X_\sigma h_p
X_\tau^{-1})=2\ip{X_\sigma h_p}{X_\tau}=2\ip{h_p}{X_\sigma^{-1}X_\tau}$. A gauge
transformation $g$ at $v'$ multiplies the first factor of $h_p$ on the right
by $g^{-1}$ and the second on the left by $g$, and multiplication by unit
quaternions is an isometry. (iv) The plaquettes containing a link of $B_v$
are those with corner $v$, by (ii), and their weights are given by (iii); the
gauge transformation at $v$ changes no two-path. (v) If $i$ joins $v_i$ to
$w=v_i+\sigma$, its plaquettes are spanned by $\sigma$ and the six directions
$\tau\perp\sigma$, with opposite even corners $v_i+\sigma+\tau$, which are
distinct since $L\ge4$; in the plaquette with corner $v$, the link $i$ is one
of the two links at $v_i$, hence a factor of the two-path.
\end{proof}

For the star at $v$ define the quotient pseudometric
\[
\bar d_v(X,X')=\min_{g\in G}\sum_\sigma\delta(gX_\sigma,X'_\sigma).
\]

\subsection{The seven-spin system}

Fix an adjacent pair $\{a,b\}$, let $S_a$ be the set of the seven directions
other than $a$, $A_7$ its adjacency matrix and $e_b$ the indicator of $b$.
In $S_a$, the direction $b$ is adjacent to $\bar a=-a$ and to the four
directions $C$ of the other two axes; $\bar a$ is adjacent to $b$,
$\bar b=-b$ and $C$; $\bar b$ to $\bar a$ and $C$; and each element of $C$ to
$b$, $\bar a$, $\bar b$ and the two elements of $C$ on the other axis. The
largest degree is $6$. For $0\le c<\frac16$ put
\[
s(c)=\mathbf 1^{\mathsf T}(I-cA_7)^{-1}e_b=\sum_{k\ge0}c^k\,\mathbf 1^{\mathsf T}
A_7^ke_b ,
\]
a nondecreasing function of $c$ with $s(c)\ge1$. The symmetries fixing $a$ and
$b$ permute $C$ transitively, so $x=(I-cA_7)^{-1}e_b$ is constant, equal to
$z$, on $C$, and $(x_b,x_{\bar a},x_{\bar b},z)$ is the unique solution of
\begin{equation}\label{eq:sevensys}
\begin{aligned}
x_b&=1+c(x_{\bar a}+4z),&\quad x_{\bar a}&=c(x_b+x_{\bar b}+4z),\\
x_{\bar b}&=c(x_{\bar a}+4z),&\quad z&=c(x_b+x_{\bar a}+x_{\bar b}+2z),
\end{aligned}
\end{equation}
with $s(c)=x_b+x_{\bar a}+x_{\bar b}+4z$.

\begin{lemma}\label{lem:svalues}
\begin{enumerate}[label=\textup{(\roman*)},leftmargin=*]
\item $s(\frac1{12})=\frac{26}{15}$, with $(x_b,x_{\bar a},x_{\bar b},z)=
\frac1{210}(221,28,11,26)$.
\item $s(\frac7{80})=\frac{3480}{1933}$, with $(x_b,x_{\bar a},x_{\bar b},z)=
\frac1{181702}(192531,26320,10829,24360)$.
\item $s(\frac9{100})=\frac{5450}{2957}$, with $(x_b,x_{\bar a},x_{\bar b},z)=
\frac1{348926}(371363,53100,22437,49050)$.
\end{enumerate}
Hence $6c\,s(c)$ equals $\frac{13}{15}$ at $c=\frac1{12}$, $\frac{1827}{1933}$
at $c=\frac7{80}$ and $\frac{2943}{2957}$ at $c=\frac9{100}$, and
$6c\,s(c)\le\frac{2943}{2957}$ for $0\le c\le\frac9{100}$.
\end{lemma}

\begin{proof}
Since $c\lambda_{\max}(A_7)\le6c<1$, the matrix $I-cA_7$ is invertible and it
suffices to check~\eqref{eq:sevensys}. (i) Multiplied by $12\cdot210$ the four
equations read $12\cdot221-(28+104)=2520$, $12\cdot28-(221+11+104)=0$,
$12\cdot11-132=0$ and $12\cdot26-(221+28+11+52)=0$; the sum is
$364/210=26/15$. (ii) Multiplied by $80\cdot181702=14536160$:
\begin{align*}
80\cdot192531-7(26320+4\cdot24360)&=15402480-866320=14536160,\\
80\cdot26320-7(192531+10829+4\cdot24360)&=2105600-7\cdot300800=0,\\
80\cdot10829-7(26320+4\cdot24360)&=866320-7\cdot123760=0,\\
80\cdot24360-7(192531+26320+10829+2\cdot24360)&=1948800-7\cdot278400=0.
\end{align*}
The sum is $327120/181702=3480/1933$, as $327120=94\cdot3480$ and
$181702=94\cdot1933$. (iii) Multiplied by $100\cdot348926$:
$100\cdot371363-9(53100+196200)=34892600$,
$100\cdot53100-9(371363+22437+196200)=5310000-9\cdot590000=0$,
$100\cdot22437-9(53100+196200)=0$ and
$100\cdot49050-9(371363+53100+22437+98100)=4905000-9\cdot545000=0$. The sum
is $643100/348926=5450/2957$, as $348926=118\cdot2957$ and
$643100=118\cdot5450$. The last claim follows from the monotonicity of $s$.
\end{proof}

\begin{lemma}[one star]\label{lem:onestar}
Assume~\eqref{eq:Q} and $c<\frac16$. Let $h,h'$ be two families of two-paths
at $v$ that differ only at one adjacent pair $\{a,b\}$. Then
\[
W_1^{\bar d_v}(\gamma_{v,h},\gamma_{v,h'})\le c\,s(c)\,|h_{ab}-h'_{ab}| .
\]
\end{lemma}

\begin{proof}
Under $\gamma_{v,h}$ put $Z_\sigma=X_a^{-1}X_\sigma$. The map $X\mapsto(X_a,Z)$
preserves Haar measure, and by Lemma~\ref{lem:stars}(iv) the density depends
only on $Z$, with $Z_a=1$; so $X_a$ is Haar distributed and independent of
$Z_S=(Z_\sigma)_{\sigma\in S_a}$, whose law $\mu^{(a)}_h$ has density
proportional to $\prod e^{4\beta\ip{h_{\sigma\tau}}{Z_\sigma^{-1}Z_\tau}}$. Given
a coupling of $\mu^{(a)}_h$ and $\mu^{(a)}_{h'}$, take $X_a=X'_a$ Haar and
independent, $X_\sigma=X_aZ_\sigma$ and $X'_\sigma=X_aZ'_\sigma$; this couples
$\gamma_{v,h}$ and $\gamma_{v,h'}$, and with $g=1$ in the definition,
$\bar d_v(X,X')\le\sum_{\sigma\in S_a}\delta(Z_\sigma,Z'_\sigma)$. It suffices
to bound the Kantorovich distance of $\mu^{(a)}_h$ and $\mu^{(a)}_{h'}$ for the
cost $\sum_{\sigma\in S_a}\delta$.

For $\sigma\in S_a$ the conditional law of $Z_\sigma$ given the others is
$\nu_{B_\sigma}$ with $B_\sigma=4\beta\sum_{\tau\text{ adj.\ }\sigma}w_{\sigma\tau}$,
where $w_{\sigma\tau}=Z_\tau\bar h_{\sigma\tau}$ or $Z_\tau h_{\tau\sigma}$
according to the ordering of the pair, since $\ip{h_{\sigma\tau}}{Z_\sigma^{-1}
Z_\tau}=\ip{Z_\sigma}{Z_\tau\bar h_{\sigma\tau}}$; changing $Z_\tau$ by $\eta$
changes $w_{\sigma\tau}$ by $\eta$, and changing $h_{ab}$ changes only $B_b$, by
at most $4|\beta||h_{ab}-h'_{ab}|$. Let $\pi$ minimise
$\int\sum_\sigma\delta(Z_\sigma,Z'_\sigma)\,d\pi$ over couplings of
$\mu^{(a)}_h$ and $\mu^{(a)}_{h'}$ (one exists by compactness) and put
$\varrho_\sigma=\int\delta(Z_\sigma,Z'_\sigma)\,d\pi$. Replacing
$(Z_\sigma,Z'_\sigma)$ by a measurably chosen optimal coupling of the two
conditional laws~\cite[Cor.~5.22]{Vil09} gives another coupling that changes
only the $\sigma$-term of the cost, so by minimality and~\eqref{eq:Q}
\[
\varrho_\sigma\le\theta\int|B_\sigma-B'_\sigma|\,d\pi\le c\sum_{\tau\in S_a,\
\tau\text{ adj.\ }\sigma}\varrho_\tau+c\,|h_{ab}-h'_{ab}|\,[\sigma=b],
\]
that is, $\varrho\le cA_7\varrho+c|h_{ab}-h'_{ab}|e_b$ componentwise. Iterating,
$\varrho\le\sum_{k<K}(cA_7)^kc|h_{ab}-h'_{ab}|e_b+(cA_7)^K\varrho$, and
$(cA_7)^K\to0$ since the spectral radius of $cA_7$ is at most $6c<1$.
Summing the components gives $\sum_\sigma\varrho_\sigma\le c\,s(c)|h_{ab}-
h'_{ab}|$.
\end{proof}

\subsection{The star condition and a lattice gap up to \texorpdfstring{$\beta_W=\frac9{25}$}{beta\_W=9/25}}

Let $\mathcal P_v=\mathcal P_{B_v}$ and let $N_s=L^4/2$ be the number of even
vertices. For a function $h$ on $G^E$ let $\Lip(h)$ be the least constant
with $|h(U)-h(U')|\le\Lip(h)\delta(U_i,U'_i)$ whenever $U,U'$ differ only at
one link $i$; changing links one at a time, $|h(U)-h(U')|\le\Lip(h)\sum_e
\delta(U_e,U'_e)$.

\begin{proposition}[star condition]\label{prop:starcond}
Assume~\eqref{eq:Q} and put $\gamma=c\,s(c)$, with $c<\frac16$ and $6\gamma<1$.
Then for every real $f\in\mathcal A$
\[
\sum_{v\text{ even}}\mu\bigl((f-\mathcal P_vf)^2\bigr)\ge(1-6\gamma)\Var_\mu(f).
\]
\end{proposition}

\begin{proof}
Let $h\in\mathcal A$ be Lipschitz, gauge invariant at every configuration, and
let $U,U'$ differ only at $i$. If $i\in B_v$, then $\mathcal P_vh(U)=\mathcal
P_vh(U')$. Otherwise let $\pi$ be a coupling of $\gamma_v(\cdot\mid U)$ and
$\gamma_v(\cdot\mid U')$ minimising the integral of $\bar d_v$, a continuous
pseudometric on a compact space. For $\sigma,\sigma'\in G^{B_v}$ and $g\in G$
the gauge transformation $g$ at $v$ changes only the links of $B_v$, so by
gauge invariance $h(\sigma',U')=h(g\sigma',U')$ and
\[
|h(\sigma,U)-h(\sigma',U')|\le\Lip(h)\Bigl(\sum_{e\in B_v}\delta(\sigma_e,
(g\sigma')_e)+\delta(U_i,U'_i)\Bigr);
\]
the infimum over $g$ gives $\Lip(h)(\bar d_v(\sigma,\sigma')+\delta(U_i,U'_i))$.
By Lemma~\ref{lem:stars}(iv),~(v), the laws $\gamma_v(\cdot\mid U)$ and
$\gamma_v(\cdot\mid U')$ differ only through the two-path of one plaquette,
which moves by $\delta(U_i,U'_i)$, if $v$ is one of the six vertices of
Lemma~\ref{lem:stars}(v), and coincide otherwise. Lemma~\ref{lem:onestar}
gives
\[
|\mathcal P_vh(U)-\mathcal P_vh(U')|\le\Lip(h)\,\bigl(1+\gamma\,[v\text{ one of
the six}]\bigr)\,\delta(U_i,U'_i).
\]
For $\mathcal K_s=N_s^{-1}\sum_v\mathcal P_v$, summing over $v$ (exactly one
star contains $i$) gives $\Lip(\mathcal K_sh)\le(1-\kappa_0)\Lip(h)$ with
$\kappa_0=(1-6\gamma)/N_s$. Each $\mathcal P_v$ maps $\mathcal A$ to itself, by
gauge covariance of $\gamma_v$, and for a Lipschitz $g\in\mathcal A$,
$\Var_\mu(g)\le\frac12\iint(g(U)-g(U'))^2\,\mu(dU)\mu(dU')\le2\Lip(g)^2|E|^2$.
So $\Var_\mu(\mathcal K_s^th)\le2(1-\kappa_0)^{2t}\Lip(h)^2|E|^2$. On the
mean-zero part of $\mathcal A$ the operator $\mathcal K_s$ is self-adjoint with
$0\le\mathcal K_s\le I$, and $\Var_\mu(\mathcal K_s^th)=\int\lambda^{2t}\,
d\nu_h(\lambda)$ for its spectral measure $\nu_h$ at $h-\mu(h)$; hence $\nu_h$
lives on $[0,1-\kappa_0]$. Lipschitz functions are dense in $L^2(\mu)$, and
the orthogonal projection onto $\mathcal A$, averaging over the compact gauge
group, which acts link by link by isometries, does not increase $\Lip$; so
such $h$ are dense in $\mathcal A$, and the spectral projection of $\mathcal K_s$
on $(1-\kappa_0,1]$ vanishes on the mean-zero part of $\mathcal A$. Thus
$\ip f{(I-\mathcal K_s)f}\ge\kappa_0\Var_\mu(f)$, and $\ip f{(I-\mathcal P_v)f}=
\mu((f-\mathcal P_vf)^2)$.
\end{proof}

\begin{theorem}\label{thm:stars}
Let $G=\SU(2)$ and $|\beta_W|\le\frac9{25}$. Then for every even $L\ge4$ and
every real $f\in\mathcal A$,
\[
\sum_{v\text{ even}}\mu\bigl((f-\mathcal P_vf)^2\bigr)\ge\frac{14}{2957}
\Var_\mu(f),
\]
and for every $n\ge2$ and even $L\ge2n+4$, $\BP(n,\varpi_n)$ holds with
\[
\varpi_n=\frac{14}{2957}\cdot\frac{(n-1)^4}{n^4\lceil(n-1)^4/2\rceil},\qquad
\varpi_2=\frac{14}{16\cdot2957}.
\]
For $0\le\beta_W\le\frac9{25}$ every torus limit is translation invariant and
reflection positive, and its transfer matrix has spectral gap at least
$\varpi_2/(2e^2q_2)=7/(625\cdot2957\,e^2)$.
\end{theorem}

\begin{proof}
Here $c=|\beta|\le\frac9{100}$, and by Lemma~\ref{lem:svalues}
$1-6c\,s(c)\ge\frac{14}{2957}$; the first claim is
Proposition~\ref{prop:starcond}. For a cube $V_x$, an even vertex $v$ with
$v-x\in[1,n-1]^4$ has its eight links based in $x+[0,n)^4$, so $B_v\subset
V_x$. The $\sigma$-algebra of the links outside $V_x$ is contained in that of
the links outside $B_v$, so $\mu((f-\mathcal P_xf)^2)=\mu(f^2)-\mu((\mathcal
P_xf)^2)\ge\mu(f^2)-\mu((\mathcal P_vf)^2)=\mu((f-\mathcal P_vf)^2)$. Each cube
contains at most $N=\lceil(n-1)^4/2\rceil$ such vertices, and each even $v$ is
in this position in exactly $(n-1)^4$ cubes. Averaging over those $v$ in each
cube,
\[
\sum_x\mu((f-\mathcal P_xf)^2)\ge\frac{(n-1)^4}{N}\sum_{v\text{ even}}
\mu((f-\mathcal P_vf)^2)\ge\frac{(n-1)^4}N\cdot\frac{14}{2957}\Var_\mu(f),
\]
which is $\BP(n,\varpi_n)$. The last claim is Corollary~\ref{cor:latticegap}
with $n=2$ and $q_2=\frac{625}{16}$.
\end{proof}

The same proof works whenever $6c\,s(c)<1$. The exact threshold is the root
$c_*$ of $6c\,s(c)=1$, and Lemma~\ref{lem:svalues}(iii) shows
$c_*>\frac9{100}$.\footnote{Solving~\eqref{eq:sevensys} numerically gives
$c_*\approx0.09023$, that is $\beta_W\approx0.3609$. No statement of the paper
uses this value.} At $\beta_W=\frac13$ the star condition holds with
$6c\,s(c)=\frac{13}{15}$, while the single-link argument of
Sec.~\ref{sec:bpstrong} is exactly marginal there. The constant $\varpi_n$ of
Theorem~\ref{thm:stars} decays like $\frac{28}{2957}n^{-4}$, so that theorem is
a cube condition for each fixed $n$, enough for the lattice gap but not for
Theorem~\ref{thm:B}. The lower faces of a cube are the reason, and the next
proposition shows that they stop every certificate that treats their links
one at a time.

\begin{proposition}[the face barrier]\label{prop:facebarrier}
For a lower face $\{\hat v_i=0\}$ of the cube $V_x$, $\hat v=v-x$, let $F_i$ be
the set of links $(b,\mu)\in V_x$ with $\hat b_i=0$ and $\mu\ne i$, and let
$A_F$ be the matrix on $F_i$ with $(A_F)_{ek}=1$ when $e\ne k$ lie in a common
plaquette all of whose links are in $F_i$. Then $\lambda_{\max}(A_F)\ge
12(1-\frac1n)^2$. Consequently, if a certificate for the response of a cube
treats the links of $F_i$ as single sites with influence matrix $C\ge cA_F$ on
$F_i\times F_i$, the inequalities $0\le\varrho\le C\varrho+\mathfrak s$ admit
solutions with $\sum\varrho$ arbitrarily large as soon as
$12c(1-\frac1n)^2\ge1$; as $n\to\infty$ such a certificate cannot pass
$c=\frac1{12}$, that is $|\beta_W|=\frac13$.
\end{proposition}

\begin{proof}
$|F_i|=3n^3$. A plaquette in a plane $\{\mu,\nu\}\not\ni i$ with base point
$y$, $\hat y_i=0$, has all four links in $F_i$ exactly when its three base
points lie in $x+[0,n)^4$, that is $\hat y_\mu,\hat y_\nu\le n-2$; there are
$3n(n-1)^2$ of them. Two links share at most one plaquette, so
$\mathbf 1^{\mathsf T}A_F\mathbf 1=12\cdot3n(n-1)^2$, and the Rayleigh quotient
of $\mathbf 1$ is $12(1-\frac1n)^2$. If $c\lambda_{\max}(A_F)\ge1$ and $w\ge0$
is a Perron vector of $A_F$, extended by $0$, then $\varrho=tw$ satisfies
$C\varrho\ge tcA_Fw\ge\varrho$ on $F_i$ for every $t>0$.
\end{proof}

For single-link updates the hypothesis $C\ge cA_F$ holds with $c=|\beta|$: at
$b=0$ the derivative of $b\mapsto\nu_b$ has Kantorovich norm exactly
$\frac14$ (Lemma~\ref{q-lem:rate}(i)), and the staples of a face link can
cancel while one of them moves. The forest of Sec.~\ref{sec:bpstrong} does not
help, since its links $T_x$ have interior heads and the face links are free.
What does help is to let the gauge transformation act at the vertices of the
faces as well, and to pay for it.

\subsection{The cube as a block system}

Fix $x$ and $n\ge2$, and write $\hat v=v-x$, computed in a window
$x+[-2,n+2]^4$ that embeds in $\Z^4$ because $L\ge2n+4$. Let $V=V_x$ and
$Q=\{v:\hat v\in[0,n-1]^4\}$. Call $v$ \emph{lower} if $v\in Q$, with zero set
$Z(v)=\{a:\hat v_a=0\}$ and $z(v)=|Z(v)|$, and \emph{upper} if $\hat v\in
[0,n]^4$ has exactly one coordinate equal to $n$. If $k$ is a link with odd end
point $o$ and direction $\mu$, its six plaquettes are indexed by $(\nu,
\epsilon)$, $\nu\ne\mu$, $\epsilon=\pm1$; the even corner of the plaquette
$(\nu,\epsilon)$ other than the even end point of $k$ is $o+\epsilon e_\nu$,
and there $k$ is a factor of the two-path of this plaquette only
(Lemma~\ref{lem:stars}(v)).

\begin{lemma}\label{lem:cubestars}
Let $n\ge2$.
\begin{enumerate}[label=\textup{(\roman*)},leftmargin=*]
\item If $v$ is lower and even, $B_v\cap V$ consists of $B_v$ minus the
$z(v)$ \emph{outside links} $(v-e_a,a)$, $a\in Z(v)$; we call the directions
$-a$, $a\in Z(v)$, \emph{fixed} and the others \emph{free}. If $v$ is upper
and even with $\hat v_a=n$, then $B_v\cap V=\{(v-e_a,a)\}$. Other even vertices
carry no link of $V$. Hence $V$ is partitioned into the \emph{blocks}
$B_v\cap V$, $v$ even and lower or upper, and the \emph{owner} $w(k)$ of
$k\in V$ is its even end point.
\item Let $k\in V$ have odd end point $o$, direction $\mu$ and owner
$o+\varsigma e_\mu$, $\varsigma=\pm1$, and let $p$ be its plaquette
$(\nu,\epsilon)$, $v=o+\epsilon e_\nu$. The links of $p$ at $v$ are in the
directions $-\epsilon\nu$ and $\varsigma\mu$. If $v$ is lower, both lie in
$V$. If $v$ is upper with $\hat v_a=n$, the links of $p$ at $v$ other than
$(v-e_a,a)$ lie outside $V$.
\item Let $v$ be lower and even and $a\in Z(v)$. The outside link $(v-e_a,a)$
is a factor of no two-path $h_p$ at a lower or upper vertex. If a plaquette
$p$ at $v$ has one of its links at $v$ outside $V$, then both links of $p$ at
the opposite even corner lie outside $V$.
\end{enumerate}
\end{lemma}

\begin{proof}
(i) The outgoing links $(v,b)$ of a lower $v$ have base point in $Q$, and an
incoming link $(v-e_b,b)$ lies in $V$ if and only if $\hat v_b\ge1$. If
$\hat v_a=n$ for exactly one $a$, no outgoing link and no incoming link
$(v-e_b,b)$, $b\ne a$, has its base point in $Q$, while $(v-e_a,a)$ has. If
$\hat v\notin[0,n]^4$ or two coordinates equal $n$, no link at $v$ has its base
point in $Q$.

(ii) The corners of $p$ are $o$, $o+\varsigma e_\mu$, $o+\varsigma e_\mu+
\epsilon e_\nu$ and $v$, so the links at $v$ go to $o$, in direction
$-\epsilon\nu$, and to $v+\varsigma e_\mu$, in direction $\varsigma\mu$. Let $v$
be lower. If $\epsilon=+1$ the first link is $(o,\nu)$, and $\hat o=\hat v-e_\nu$
lies in $Q$, its $\nu$-coordinate being $\hat o_\nu\ge0$; if $\epsilon=-1$ it is
$(v,\nu)$, with base point $v\in Q$. The second link is $(v,\mu)$ if
$\varsigma=+1$; if $\varsigma=-1$ it is $(v-e_\mu,\mu)$, which lies in $V$
because $\hat v_\mu=\hat o_\mu\ge1$, as $k=(o-e_\mu,\mu)\in V$. If $v$ is upper,
(i) gives the claim.

(iii) The plaquettes of $(v-e_a,a)$ have the even corners $v$ and
$v-e_a\pm e_\nu$, whose $a$-coordinate is $-1$, so these are neither lower nor
upper; and a link is a factor of a two-path only at the even corner opposite
to its own even end point. If $p$ at $v$ is spanned by $-a$ and $\tau$, the
opposite corner $v'=v-e_a+\tau$ and both its neighbours in $p$, $v-e_a$ and
$v+\tau$, are joined to it by links based at $a$-coordinate $-1$.
\end{proof}

For a block vertex $v$ the conditional law $\gamma_v(\cdot\mid\sigma,U)$ of
the links of $B_v\cap V$, given the other links of $V$ (values $\sigma$) and
the links off $V$ (values $U$), has density proportional to the product of
the weights $e^{4\beta\ip{h_p}{X_\sigma^{-1}X_\tau}}$ of the plaquettes $p$ at
$v$ with a link at $v$ in $V$. It depends on $(\sigma,U)$ only through these
two-paths and the values of the links at $v$ off $V$.

\subsection{The soft quotient}

At a vertex of a lower face the gauge transformation moves the outside links
at that vertex, and the conditional law of the cube does not see it. The
following pseudometric allows it anyway, at the price of the displacement it
causes on those links. The price is legitimate in the Lipschitz contraction
argument of Proposition~\ref{prop:interiorbp}, because the test function
depends on all links.

\begin{definition}\label{def:dtilde}
For a lower even $v$ and configurations $X,X'$ of $B_v\cap V$, written as
outward transporters, put
\[
\tilde d_v(X,X')=\min_{g\in G}\Bigl[\sum_{\sigma\text{ free}}\delta(gX_\sigma,
X'_\sigma)+z(v)\,\delta(g,1)\Bigr],
\]
and for an upper even $v$ let $\tilde d_v$ be $\delta$ on its single link. For
$\sigma,\sigma'\in G^V$ let $D(\sigma,\sigma')=\sum_v\tilde d_v(\sigma_{B_v},
\sigma'_{B_v})$, the sum over the block vertices, and for $j\notin V$
\[
\tilde R_x(j)=\sup\Bigl\{\frac{W^D(\gamma^0_{V,U},\gamma^0_{V,U'})}
{\delta(U_j,U'_j)}:U,U'\text{ differ only at }j\Bigr\},
\]
where $\gamma^0_{V,U}=\gamma^0_V(\cdot\mid U)$ and $W^D$ is the Kantorovich
distance for the cost $D$.
\end{definition}

The minimum exists by compactness, and $\tilde d_v$ and $D$ are continuous,
symmetric, since $\delta(gX,X')=\delta(X,g^{-1}X')$ and
$\delta(g,1)=\delta(g^{-1},1)$, and satisfy the triangle inequality, by
composing the minimising elements. For $z(v)=0$, $\tilde d_v=\bar d_v$.

\begin{lemma}\label{lem:softlip}
Let $h$ be a function on $G^E$, gauge invariant at every configuration, with
$|h(W)-h(W')|\le\Lip(h)\sum_e\delta(W_e,W'_e)$. If $U,U'$ differ only at
$j\notin V$, then for all $\sigma,\sigma'\in G^V$
\[
|h(\sigma,U)-h(\sigma',U')|\le\Lip(h)\bigl(D(\sigma,\sigma')+\delta(U_j,U'_j)
\bigr).
\]
\end{lemma}

\begin{proof}
For each lower even $v$ choose $g_v$ attaining $\tilde d_v(\sigma_{B_v},
\sigma'_{B_v})$, and let $g$ be the gauge transformation equal to $g_v$ at
these vertices and to $1$ elsewhere; then $h(\sigma,U)=h(g\cdot(\sigma,U))$. A
link of $V$ has exactly one even end point $v$, and $g$ acts on it by $X\mapsto
g_vX$ on its outward transporter at $v$, or trivially if $v$ is upper. A link
$e\notin V$ has at most one end point where $g\ne1$, namely an even lower
$v$, and then $e$ is one of the $z(v)$ outside links at $v$, by
Lemma~\ref{lem:cubestars}(i); there $\delta((g\cdot U)_e,U'_e)\le\delta(g_v,1)+
\delta(U_e,U'_e)$ by bi-invariance. Summing over all links,
$|h(g\cdot(\sigma,U))-h(\sigma',U')|\le\Lip(h)\bigl(\sum_v\tilde d_v+
\delta(U_j,U'_j)\bigr)$.
\end{proof}

\subsection{Block responses}

\begin{lemma}[Dobrushin on $S^3$]\label{lem:dobS3}
Assume~\eqref{eq:Q}. Let $S$ be finite and $\nu,\nu'$ probability measures on
$G^S$ with continuous positive densities with respect to Haar measure, such
that for every $s\in S$ the conditional law of $x_s$ given $x_{S\setminus s}$
is $\nu_{B_s(x)}$ under $\nu$ and $\nu_{B'_s(x)}$ under $\nu'$, where
$|B_s(x)-B_s(y)|\le4|\beta|\sum_ta_{st}\delta(x_t,y_t)$ and
$|B_s(x)-B'_s(x)|\le4|\beta|\eta_s$ for all $x,y$, with $a_{st}\ge0$,
$a_{ss}=0$ and $\eta_s\ge0$. Then there is a coupling $\pi$ of $\nu,\nu'$ whose
$\varrho_s=\int\delta(x_s,x'_s)\,d\pi$ satisfy $\varrho\le cA\varrho+c\eta$
componentwise. Consequently:
\begin{enumerate}[label=\textup{(\alph*)},leftmargin=*]
\item if $\sum_sa_{st}\le d$ for every $t$ and $cd<1$, then
$\sum_s\varrho_s\le c\sum_s\eta_s/(1-cd)$;
\item if the spectral radius of $cA$ is less than $1$ and $y\in[0,\infty)^S$
satisfies $y\ge cAy+c\eta$, then $\varrho\le y$;
\item if $y\ge0$ satisfies $y\ge c_0Ay+\eta$ and $c\le c_0$, then $cy\ge
cA(cy)+c\eta$, so $\varrho\le cy$ when the spectral radius of $cA$ is less
than $1$.
\end{enumerate}
\end{lemma}

\begin{proof}
Let $\pi$ minimise $\int\sum_s\delta(x_s,x'_s)\,d\pi$ over couplings. Fix $s$
and replace under $\pi$ the pair $(x_s,x'_s)$ by a sample of an optimal
coupling of $\nu_{B_s(x)}$ and $\nu_{B'_s(x')}$, chosen measurably in
$(x_{S\setminus s},x'_{S\setminus s})$~\cite[Cor.~5.22]{Vil09}. This is again
a coupling of $\nu,\nu'$, and only the $s$-term of the cost changes, so by
minimality,~\eqref{eq:Q} and the triangle inequality, $\varrho_s\le\theta\int
|B_s(x)-B'_s(x')|\,d\pi\le c\sum_ta_{st}\varrho_t+c\eta_s$. (a) Summing over
$s$, $\sum\varrho\le cd\sum\varrho+c\sum\eta$. (b) The vector $y-\varrho$
satisfies $y-\varrho\ge cA(y-\varrho)$, hence $y-\varrho\ge(cA)^K(y-\varrho)
\to0$. (c) $cy-c^2Ay\ge c(y-c_0Ay)\ge c\eta$; apply (b).
\end{proof}

From now on $c\le c_0=\frac7{80}$, that is $|\beta_W|\le\frac7{20}$, and
$\kappa_0=s(c_0)=\frac{3480}{1933}<1.8004$, so that $6c_0\kappa_0=
\frac{1827}{1933}$. For a direction $a$ and a family $h$ of two-paths let
$\mu^{(a)}_h$ be the law on $G^{S_a}$ with density proportional to
$\prod_{\{\sigma,\tau\}}e^{4\beta\ip{h_{\sigma\tau}}{Z_\sigma^{-1}Z_\tau}}$,
$Z_a=1$, as in the proof of Lemma~\ref{lem:onestar}. By that proof,
Lemma~\ref{lem:svalues}(ii) and Lemma~\ref{lem:dobS3}(c), if $h,h'$ differ only
at the pair $\{a,b\}$ there is a coupling of $\mu^{(a)}_h$ and $\mu^{(a)}_{h'}$
with
\begin{equation}\label{eq:sevenbound}
\int\sum_{\sigma\in S_a}\delta(Z_\sigma,Z'_\sigma)\le\kappa_0c\,
|h_{ab}-h'_{ab}|\qquad(c\le c_0).
\end{equation}

\begin{lemma}\label{lem:partial}
Let $c\le c_0$ and $v$ a block vertex, and write $\gamma^{h,Y}$ for its block
law as a function of the two-paths $h$ at $v$ and of the values $Y$ of the
links at $v$ off $V$. Put $\kappa_1=\kappa_0$, $\kappa_2=\frac94$,
$\kappa_3=\frac{40}{13}$ and $\kappa_4=\frac{160}{59}$.
\begin{enumerate}[label=\textup{(\alph*)},leftmargin=*]
\item Let $v$ be lower, and let $(h,Y)$ and $(h',Y)$ differ only at the
two-path $h_p$ of a plaquette $p$ whose two links at $v$ are free. Then
$W^{\tilde d_v}(\gamma^{h,Y},\gamma^{h',Y})\le\kappa_{z(v)}c\,|h_p-h'_p|$, with
$\kappa_0$ for $z(v)=0$.
\item In particular, for $z(v)=1$ the constant is the full-star constant
$\kappa_0c$.
\item Let $v$ be lower and $p$ a plaquette at $v$ with one link at $v$ free and
one fixed. Then the bound of \textup{(a)} holds with $\kappa_0c$ if $z(v)=1$
and with $\frac{16}9c$ if $z(v)\ge2$.
\item Let $v$ be lower and let $(h,Y)$ and $(h,Y')$ differ only at one
outside link, by $\eta$ in $\delta$. Then $W^{\tilde d_v}(\gamma^{h,Y},
\gamma^{h,Y'})\le\eta$.
\item Let $v$ be upper. If $(h,Y)$ and $(h',Y')$ differ only at one two-path
$h_p$ with $p\ni(v-e_a,a)$, or only at one outside link at $v$, the distance is
at most $c|h_p-h'_p|$, respectively $c\eta$.
\end{enumerate}
\end{lemma}

\begin{proof}
\emph{Representation.} Let $v$ be lower and $F$ its set of fixed directions,
let $\Phi_h(X)=\prod e^{4\beta\ip{h_{\sigma\tau}}{X_\sigma^{-1}X_\tau}}$ over all
$24$ adjacent pairs, so that $\Phi_h(gX)=\Phi_h(X)$, and note that
$\gamma^{h,Y}$ has density proportional to $\Phi_h$ evaluated at $X_F=Y_F$. Fix
a free direction $\sigma_0$. If $F=\emptyset$, the map
$X\mapsto(X_{\sigma_0},(X_{\sigma_0}^{-1}X_\rho)_{\rho\ne\sigma_0})$ preserves
Haar measure and carries $\gamma^h$ to $\mathrm{Haar}\otimes\mu^{(\sigma_0)}_h$,
as in the proof of Lemma~\ref{lem:onestar}. If $F\ne\emptyset$, fix $f_1\in F$
and let $\Psi(X)$ consist of $W_\rho=X_{\sigma_0}^{-1}X_\rho$ for the free
$\rho\ne\sigma_0$ and of $W_m=X_{\sigma_0}^{-1}Y_{f_1}$. Then $\Psi$ is a
bijection of the free coordinates, with inverse $X_{\sigma_0}=Y_{f_1}W_m^{-1}$,
$X_\rho=Y_{f_1}W_m^{-1}W_\rho$, and it preserves Haar measure: $X_{\sigma_0}
\mapsto X_{\sigma_0}^{-1}Y_{f_1}$ does, and for fixed $X_{\sigma_0}$ so does
each $X_\rho\mapsto X_{\sigma_0}^{-1}X_\rho$. Since $\Phi_h(X)=\Phi_h(
X_{\sigma_0}^{-1}X)$ and $X_{\sigma_0}^{-1}X_f=W_mK_f$ with $K_f=Y_{f_1}^{-1}
Y_f$, the image of $\gamma^{h,Y}$ has density proportional to $\Phi_h$ at
$X_{\sigma_0}=1$, $X_\rho=W_\rho$, $X_f=W_mK_f$. If $(W,W')$ couples the images
for $(h,Y)$ and $(h',Y)$ and $X=\Psi^{-1}W$, $X'=\Psi^{-1}W'$, then
$g=Y_{f_1}W'^{-1}_mW_mY_{f_1}^{-1}$ satisfies $gX_{\sigma_0}=X'_{\sigma_0}$,
$\delta(gX_\rho,X'_\rho)=\delta(W_\rho,W'_\rho)$ and $\delta(g,1)=
\delta(W_m,W'_m)$, so
\begin{equation}\label{eq:transfer}
\tilde d_v(X,X')\le\sum_\rho\delta(W_\rho,W'_\rho)+z(v)\,\delta(W_m,W'_m).
\end{equation}

(a), (b) for $z(v)\le1$. Take $\sigma_0$ and $\tau$ to be the two directions of
$p$ at $v$. For $z(v)=1$, $F=\{f_1\}$, the image is exactly
$\mu^{(\sigma_0)}_h$, the coordinate $W_m$ being $Z_{f_1}$, and the change of
$h_p=h_{\sigma_0\tau}$ is a change of the field on $Z_\tau$;
\eqref{eq:sevenbound} and~\eqref{eq:transfer} give $\kappa_0c|h_p-h'_p|$. For
$z(v)=0$ couple with the same $X_{\sigma_0}$ and use $g=1$; this is
Lemma~\ref{lem:onestar} with~\eqref{eq:sevenbound}.

(a) for $z(v)=2$. Let $F=\{-i,-j\}$; the free directions are $P=+i$, $Q=+j$,
$K_\pm=\pm k$ and $L_\pm=\pm l$. Up to the symmetries $i\leftrightarrow j$,
$k\leftrightarrow l$, $k\mapsto-k$ and $l\mapsto-l$ of the star, which fix
$F$, the pair of $p$ is one of $\{P,Q\}$, $\{P,K_+\}$ and $\{K_+,L_+\}$. In the
image the terms of the density are: for adjacent free $\rho,\rho'\ne\sigma_0$,
$\ip h{W_\rho^{-1}W_{\rho'}}$; for $\{\sigma_0,\rho\}$, a field on $W_\rho$; for
$\{f,\rho\}$, $\ip h{(W_mK_f)^{-1}W_\rho}=\ip{K_fh}{W_m^{-1}W_\rho}$, linear in
$W_m$ and in $W_\rho$; for $\{\sigma_0,f\}$, a field on $W_m$; and for
$\{-i,-j\}$, a constant. So Lemma~\ref{lem:dobS3} applies with
$a_{\rho\rho'}=1$ for adjacent free $\rho,\rho'\ne\sigma_0$,
$a_{\rho m}=a_{m\rho}$ equal to the number of $f\in F$ adjacent to $\rho$
(one for $P,Q$, two for $K_\pm,L_\pm$), and $\eta=|h_p-h'_p|e_\tau$. Every row
sum of $A$ is at most $9$, so $c_0A$ has spectral radius at most
$\frac{63}{80}$. By~\eqref{eq:transfer} the cost is at most
$\sum_\rho\varrho_\rho+2\varrho_m$. We give $\hat y\ge c_0A\hat y+e_\tau$ and
use Lemma~\ref{lem:dobS3}(c); each item lists $\hat y$ and the inequalities
multiplied by $80$ and by the denominator of $\hat y$.

\emph{The pair $\{P,Q\}$}, $\sigma_0=P$, $\tau=Q$. The spins are $Q$, $K_\pm$,
$L_\pm$, each of the last four adjacent to $Q$ and to the two of the other
axis, and $m$; $100\hat y=(Q,K_\pm,L_\pm,m)=(108,16,21)$. The inequalities
are $8640\ge7(4\cdot16+21)+8000=8595$, $1280\ge7(108+2\cdot16+2\cdot21)=1274$
and $1680\ge7(108+8\cdot16)=1652$, and the cost is at most
$(108+4\cdot16+2\cdot21)/100=2.14$.

\emph{The pair $\{P,K_+\}$}, $\sigma_0=K_+$, $\tau=P$. The spins are $P$ and
$Q$, adjacent to each other and to $K_-$, $L_\pm$; $K_-$, adjacent to $P,Q,
L_\pm$; $L_\pm$, adjacent to $P,Q,K_-$; and $m$; $200\hat y=(P,Q,K_-,L_\pm,m)=
(216,32,35,32,40)$. The inequalities are $17280\ge7(32+35+64+40)+16000=17197$,
$2560\ge7(216+35+64+40)=2485$, $2800\ge7(216+32+64+80)=2744$,
$2560\ge7(216+32+35+80)=2541$ and $3200\ge7(216+32+70+128)=3122$, and the
cost is at most $(216+32+35+64+80)/200=2.135$.

\emph{The pair $\{K_+,L_+\}$}, $\sigma_0=K_+$, $\tau=L_+$, with the same spins
and adjacencies; $200\hat y=(P,Q,K_-,L_+,L_-,m)=(32,32,36,219,19,54)$. The
inequalities are $2560\ge7(32+36+219+19+54)=2520$ for $P$ and for $Q$,
$2880\ge7(64+219+19+108)=2870$, $17520\ge7(64+36+108)+16000=17456$,
$1520\ge7(64+36+108)=1456$ and $4320\ge7(64+72+2\cdot238)=4284$, and the cost
is at most $(64+36+219+19+108)/200=2.23$.

In all cases the response is at most $2.23\,c|h_p-h'_p|\le\kappa_2c|h_p-h'_p|$.

(a) for $z(v)\ge3$, and (c), (d) for $z(v)\ge2$. Use the coordinates
$X_\sigma$, $\sigma$ free, directly, and $g=1$ in $\tilde d_v$. Each conditional
law is $\nu_B$ with $B=4\beta$ times the sum of the unit vectors of the
plaquettes at $v$ through $\sigma$, those through free neighbours depending
isometrically on them; so Lemma~\ref{lem:dobS3}(a) applies with $A$ the
adjacency matrix of the free directions, whose largest degree is $d=5,4,3$
for $z(v)=2,3,4$ (for $z=2$, $+i$ is adjacent to $+j$, $\pm k$, $\pm l$; for
$z=3$, $+i$ to $+j$, $+k$, $\pm l$; for $z=4$ the free directions are the four
positive ones). A change of $h_p$ with both links free gives $\sum\eta=
2|h_p-h'_p|$, with one link free $\sum\eta=|h_p-h'_p|$, and a change of the
outside link in direction $f$ by $\eta$ gives $\eta$ on each free neighbour of
$f$, at most $d$ of them. With $c\le c_0$: $\frac{2c}{1-4c}\le\frac{40}{13}c$,
$\frac{2c}{1-3c}\le\frac{160}{59}c$, $\frac c{1-cd}\le\frac c{1-5c_0}=
\frac{16}9c$ and $\frac{dc}{1-dc}\le\frac{5c_0}{1-5c_0}=\frac79$.

(c), (d) for $z(v)=1$. Let $F=\{f\}$ and $X_f=Y$. The map $X\mapsto Z=(Y^{-1}
X_\rho)_\rho$ carries $\gamma^{h,Y}$ to $\mu^{(f)}_h$, which does not depend on
$Y$. A change of $h_p$, $p$ spanned by $f$ and $\tau$, is a change of the
field on $Z_\tau$, so~\eqref{eq:sevenbound} gives a coupling with
$\int\sum\delta(YZ_\rho,YZ'_\rho)\le\kappa_0c|h_p-h'_p|$, with $g=1$. For a
change $Y\to Y'$ couple $X=YZ$ and $X'=Y'Z$ and take $g=Y'Y^{-1}$; then
$\tilde d_v(X,X')\le\delta(Y'Y^{-1},1)=\delta(Y,Y')$.

(e) The block law is $\nu_b$ with $b=4\beta$ times the sum of the six vectors
of the plaquettes through $(v-e_a,a)$; such a change moves $b$ by at most
$4|\beta||h_p-h'_p|$, respectively $4|\beta|\eta$, since an outside link lies
in at most one plaquette with $(v-e_a,a)$. Apply~\eqref{eq:Q}.
\end{proof}

Part (b) is the point of the soft quotient: a star with one fixed link is
exactly as good as a full star, and the price $z(v)\delta(g,1)$ never feeds
back into the system, because an outside link of a lower vertex is a factor
of no two-path at a block (Lemma~\ref{lem:cubestars}(iii)). The margin is lost
at the edges of the cube, $z(v)=2$, and the next lemma pays for them with
weights.

\subsection{The column inequality}

For a block vertex $v$ and a plaquette $p$ at $v$ let $\Gamma_v(p)$ be the
constant of Lemma~\ref{lem:partial}(a), (c) or (e), and $\Gamma_v(p)=0$ if no
link of $p$ at $v$ lies in $V$. Put $u_v=1$ for upper $v$ and $u_v=u_{z(v)}$
for lower $v$, with
\[
u_0=u_1=1,\quad u_2=\tfrac56,\quad u_3=\tfrac7{10},\quad u_4=\tfrac35;
\qquad\lambda=\tfrac{97}{100}.
\]

\begin{lemma}\label{lem:column}
Let $n\ge4$ and $c\le c_0$. For every $k\in V$, with plaquettes
$p_{(\nu,\epsilon)}$ and even corners $v_{(\nu,\epsilon)}=o+\epsilon e_\nu$ as
above,
\[
\sum_{(\nu,\epsilon)}u_{v_{(\nu,\epsilon)}}\Gamma_{v_{(\nu,\epsilon)}}
(p_{(\nu,\epsilon)})\le\lambda\,u_{w(k)} .
\]
\end{lemma}

\begin{proof}
Put $f(z)=u_z\kappa_zc$: $f(0)=f(1)=\kappa_0c<1.8004\,c$, $f(2)=\frac{15}8c$,
$f(3)=\frac{28}{13}c<2.1539\,c$ and $f(4)=\frac{96}{59}c<1.6272\,c$, all at
least $c$. By Lemma~\ref{lem:cubestars}(ii) a lower $v_{(\nu,\epsilon)}$ has
both links of $p$ free, so it contributes $f(z(v_{(\nu,\epsilon)}))$; an upper
one contributes at most $c$, and other vertices contribute $0$. Let $o$, $\mu$
and $\varsigma$ be as in Lemma~\ref{lem:cubestars}(ii).

\emph{Case A: $\hat o$ has a coordinate equal to $n$.} Since $k\in V$ is
$(o,\mu)$, which needs $o\in Q$, or $(o-e_\mu,\mu)$, which needs $\hat o-e_\mu
\in[0,n-1]^4$, this coordinate is $\hat o_\mu$, and $w(k)=o-e_\mu$ with
$\hat w_\mu=n-1\ge1$; so $z(w)\le3$ and $u_w\ge\frac7{10}$. Every
$v_{(\nu,\epsilon)}$ has $\mu$-coordinate $n$, so the sum is at most
$6c\le6c_0=0.525<\frac7{10}\cdot0.97$.

\emph{Case B: $o\in Q$.} Let $Z=z(o)$ be the number of zeros of $\hat o$,
$a=\#\{\nu\ne\mu:\hat o_\nu=0\}$ and $b=\#\{\nu\ne\mu:\hat o_\nu=1\}$. As
$n\ge4$, the classes $\{0\}$, $\{1\}$, $[2,n-2]$ and $\{n-1\}$ are disjoint.
For $\nu\ne\mu$: if $\hat o_\nu=0$, $\epsilon=-1$ gives no block and
$\epsilon=+1$ a lower vertex with $z=Z-1$; if $\hat o_\nu=1$, $\epsilon=+1$
gives $z=Z$ and $\epsilon=-1$ gives $z=Z+1$; if $2\le\hat o_\nu\le n-2$ both
give $z=Z$; and if $\hat o_\nu=n-1$, $\epsilon=-1$ gives $z=Z$ and
$\epsilon=+1$ an upper vertex, which contributes at most $c\le f(Z)$. Hence the
sum is at most
\[
a\,f(Z-1)+b\bigl(f(Z)+f(Z+1)\bigr)+2(3-a-b)f(Z).
\]
If $\hat o_\mu=0$, then $\varsigma=+1$, $Z=a+1$ and $z(w)=a$, and the sum is at
most $h(a,b)=af(a)+b(f(a+1)+f(a+2))+2(3-a-b)f(a+1)$, to be compared with
$\lambda u_a$. If $\hat o_\mu\ge1$, then $Z=a$ and $w$ is upper, with $u_w=1$,
or lower with $z(w)\in\{a,a+1\}$, so $u_w\ge u_{a+1}$; the sum is at most
$g(a,b)=af(a-1)+b(f(a)+f(a+1))+2(3-a-b)f(a)$, to be compared with $\lambda
u_{a+1}$. Here $a+b\le3$. Since $f(0)=f(1)\le f(2)\le f(3)$, the quantities
$h(0,b)$, $g(1,b)$, $h(1,b)$ and $g(2,b)$ increase with $b$, and
$g(0,b)=6f(0)$. Every term is $c$ times a constant, so it suffices to take
$c=c_0$; with $\hat\lambda=\lambda/c_0=\frac{388}{35}>11.0857$ and $f$ in units
of $c$, the cases are
\begin{align*}
h(0,3)&=3f(0)+3f(2)<5.4012+5.625=11.0262<\hat\lambda,\\
g(0,b)&=6f(0)<10.8024<\hat\lambda,\\
g(1,2)&=3f(0)+2f(2)<9.1512<\tfrac56\hat\lambda,\quad\tfrac56\hat\lambda>9.2380,\\
h(1,2)&=f(0)+2f(2)+2f(3)<9.8582<\hat\lambda,\\
g(2,1)&=2f(0)+f(2)+f(3)<7.6297<\tfrac7{10}\hat\lambda,\quad
\tfrac7{10}\hat\lambda>7.7599,\\
h(2,0)&=2f(2)+2f(3)<8.0578<\tfrac56\hat\lambda,\\
h(2,1)&=2f(2)+f(3)+f(4)<7.5311<\tfrac56\hat\lambda,\\
g(3,0)&=3f(2)=5.625<\tfrac35\hat\lambda,\quad\tfrac35\hat\lambda>6.6514,\\
h(3,0)&=3f(3)<6.4617<\tfrac7{10}\hat\lambda.
\end{align*}
All inequalities hold.
\end{proof}

\subsection{The cube condition uniformly in the cube size}

\begin{proposition}\label{prop:starresponse}
Let $c\le c_0$, $n\ge4$ and $L\ge2n+4$ even. For every link $j$,
\[
\sum_{x:\,j\notin V_x}\tilde R_x(j)\le259\,n^3 .
\]
\end{proposition}

\begin{proof}
Fix $x$, $j\notin V=V_x$ and $U,U'$ differing only at $j$, and write
$\eta=\delta(U_j,U'_j)$. Let $\pi$ minimise $\int D\,d\pi$ over couplings of
$\gamma^0_{V,U}$ and $\gamma^0_{V,U'}$, and put $\varrho_v=\int\tilde d_v\,d\pi$.
As in Lemma~\ref{lem:dobS3}, replacing the block of $v$ by an optimal
coupling for $\tilde d_v$ of $\gamma_v(\cdot\mid\sigma,U)$ and
$\gamma_v(\cdot\mid\sigma',U')$ changes only the $v$-term, so $\varrho_v\le\int
\Omega_v\,d\pi$ with $\Omega_v=W^{\tilde d_v}(\gamma_v(\cdot\mid\sigma,U),
\gamma_v(\cdot\mid\sigma',U'))$.

\emph{Pointwise bound.} Changing the two-paths at $v$ one at a time, and then
the outside link $j$ if it is a link at $v$, the triangle inequality for
$W^{\tilde d_v}$ and Lemma~\ref{lem:partial} give $\Omega_v\le\sum_p\Gamma_v(p)
|h_p(\sigma,U)-h_p(\sigma',U')|+r_v\eta$, where $r_v\le1$ if $v$ is lower and
$j$ is an outside link at $v$, $r_v\le c$ if $v$ is upper and $j$ is a link at
$v$, and $r_v=0$ otherwise. For each block vertex $w$ fix $g^*_w$ attaining
$\tilde d_w(\sigma_{B_w},\sigma'_{B_w})$, with $g^*_w=1$ for upper $w$, and put
$\delta_k=\delta(g^*_{w(k)}X_k(\sigma),X_k(\sigma'))$ for $k\in V$, so that
$\sum_{k\in B_w\cap V}\delta_k\le\tilde d_w$. If both links of $p$ at the
opposite corner $v'$ lie in $V$, then $v'$ is lower, an upper vertex carrying
one link of $V$, and by invariance of $h_p$ under $X\mapsto g^*_{v'}X$ at $v'$,
$|h_p(\sigma,U)-h_p(\sigma',U')|$ is at most the sum of their $\delta_k$. If
exactly one lies in $V$, the other is off $V$ and is not an outside link of a
lower vertex (Lemma~\ref{lem:cubestars}(iii)), so $v'$ is upper and the bound
is $\delta_k+[j\in p]\eta$. If none lies in $V$, the bound is $[j\in p]\eta$.
Each $k\in V$ is a factor of a two-path at $v$ for at most one $p$ at $v$, and
these $v$ are the vertices $v_{(\nu,\epsilon)}$ of $k$, so
Lemma~\ref{lem:column} gives, pointwise,
\[
\sum_vu_v\Omega_v\le\lambda\sum_{k\in V}u_{w(k)}\delta_k+B\eta\le\lambda\sum_w
u_w\tilde d_w+B\eta,
\]
with $B=\sum_{p\ni j}u_{v(p)}\Gamma_{v(p)}(p)+\max_vu_vr_v$, where $v(p)$ is
the even corner of $p$ opposite to the even end point of $j$.

\emph{Constants.} By Lemma~\ref{lem:partial}, $u_v\Gamma_v(p)\le\frac{28}{13}c$
in all cases ($\max_zf(z)=f(3)$; with one fixed link $\kappa_0c$ or
$\frac{16}9c$; upper vertices $c$), and only the plaquettes $p\ni j$
containing a link of $V$ contribute. Integrating, $\sum_vu_v\varrho_v\le
\lambda\sum_vu_v\varrho_v+B_x(j)\eta$, with $B_x(j)=\frac{28}{13}c\,P_x(j)+
R_x(j)$, where $P_x(j)$ is the number of plaquettes $p\ni j$ with a link in
$V_x$, and $R_x(j)=1$ if $j$ is an outside link at a lower vertex of the
cube, $c$ if $j$ is a link off $V_x$ at an upper even vertex of the cube, and
$0$ otherwise. All terms are finite, so $\sum_vu_v\varrho_v\le
B_x(j)\eta/(1-\lambda)$ and
\[
W^D(\gamma^0_{V,U},\gamma^0_{V,U'})\le\sum_v\varrho_v\le\frac{B_x(j)\,\eta}
{u_4(1-\lambda)}=\frac{500}9B_x(j)\,\eta .
\]

\emph{Summation over $x$.} A plaquette counted in $P_x(j)$ contains a partner
$k$ of $j$ with $k\in V_x$, and distinct plaquettes give distinct partners.
For $j=(b,\nu)$ the eighteen partners have base points $b+\Delta$ with
$\sum\|\Delta\|_1=18$, $\Delta\in\{0,e_\mu,e_\nu,-e_\mu,-e_\mu,e_\nu-e_\mu\}$
for each $\mu\ne\nu$, and $\#\{x:k\in V_x,\ j\notin V_x\}=|Q_0\setminus(Q_0-
\Delta)|\le n^3\|\Delta\|_1$ with $Q_0=[0,n)^4$, by moving one coordinate at a
time; so $\sum_xP_x(j)\le18n^3$. For $R_x$, let $v$ be the even end point of
$j$. If $v$ is the head of $j=(v-e_a,a)$, it is a lower vertex with $a\in
Z(v)$ for exactly $n^3$ cubes, and an upper vertex of the cube with $j$ off
$V_x$ for at most $3n^3$ cubes, the coordinate equal to $n$ not being $a$; if
$v$ is the base point of $j$, it is never lower with $j\notin V_x$ and is upper
for at most $4n^3$ cubes. So $\sum_xR_x(j)\le(1+3c)n^3$. With $c\le c_0$,
\[
\sum_x\tilde R_x(j)\le\frac{500}9\Bigl(\frac{28}{13}\cdot18c_0+1+3c_0\Bigr)n^3
=\frac{500}9\cdot\frac{4841}{1040}\,n^3<258.61\,n^3 .\qedhere
\]
\end{proof}

\begin{theorem}\label{thm:starsuniform}
Let $G=\SU(2)$ and $|\beta_W|\le\frac7{20}$.
\begin{enumerate}[label=\textup{(\roman*)},leftmargin=*]
\item For every $n\ge260$ and even $L\ge2n+4$, $\BP(n,1-259/n)$ holds; in
particular $\BP(n,\frac12)$ holds for $n\ge518$.
\item For every $n\ge2$ and even $L\ge2n+4$, $\BP(n,\varpi)$ holds with
$\varpi=\frac{106}{1933}\cdot517^{-4}$, independent of $n$ and $L$.
\end{enumerate}
\end{theorem}

\begin{proof}
(i) Let $\mathcal K=L^{-4}\sum_x\mathcal P_x$ and let $h$ be gauge invariant at
every configuration with $\Lip(h)<\infty$, and let $U,U'$ differ only at $i$.
If $i\in V_x$, then $\mathcal P_xh(U)=\mathcal P_xh(U')$. If $i\notin V_x$,
Lemma~\ref{lem:softlip} and any coupling $\pi$ of $\gamma^0_{V_x,U}$ and
$\gamma^0_{V_x,U'}$ give $|\mathcal P_xh(U)-\mathcal P_xh(U')|\le\Lip(h)(\int
D\,d\pi+\delta(U_i,U'_i))$, and the infimum over $\pi$ gives
$\Lip(h)(\tilde R_x(i)+1)\delta(U_i,U'_i)$. As $i$ lies in $n^4$ cubes,
Proposition~\ref{prop:starresponse} gives $\Lip(\mathcal Kh)\le(1-\kappa)
\Lip(h)$ with $\kappa=(1-259/n)n^4/L^4>0$, and $\mathcal Kh$ is again gauge
invariant at every configuration, by gauge covariance of $\gamma^0_{V_x,U}$.
The spectral argument of Proposition~\ref{prop:starcond} gives
$\ip h{(I-\mathcal K)h}\ge\kappa\Var_\mu(h)$, first for such $h$ and then, by
density, for every real $f\in\mathcal A$; and $L^{-4}\sum_x\mu((f-\mathcal
P_xf)^2)=\ip f{(I-\mathcal K)f}$, which is $\BP(n,1-259/n)$.

(ii) For $n\ge518$ use (i). For $2\le n\le517$ use the proof of
Theorem~\ref{thm:stars}, which gives $\BP(n,(1-6c\,s(c))(n-1)^4/(n^4\lceil
(n-1)^4/2\rceil))$ whenever $6c\,s(c)<1$; here $1-6c\,s(c)\ge1-\frac{1827}
{1933}=\frac{106}{1933}$ by Lemma~\ref{lem:svalues}(ii), and $(n-1)^4\ge
\lceil(n-1)^4/2\rceil$, so $\BP(n,\frac{106}{1933}n^{-4})$ holds, and
$\frac{106}{1933}n^{-4}\ge\varpi$.
\end{proof}

Theorem~\ref{thm:starsuniform} extends the uniform range $\frac5{18}$ of
Theorem~\ref{thm:bpweighted} past $\beta_W=\frac13$, where
Proposition~\ref{prop:methodlimit} stops every gauge-fixed single-link
scheme and Proposition~\ref{prop:facebarrier} every scheme that treats face
links as single sites. The faces cost nothing here; the margin is lost at the
cube edges, through the column $h(0,3)$ of Lemma~\ref{lem:column}.\footnote{A
linear programme over weights depending on the block, with the exact
constants for $z(v)=2$, suggests on cubes of side up to $8$ that the method
reaches the star threshold $\beta_W\approx0.3609$ of Theorem~\ref{thm:stars};
with weights depending only on $z(v)$ it stops near $\beta_W\approx0.345$.
These are numerical observations and no statement rests on them.} The value
$\beta_W\le\frac7{20}$ is far below the scaling region of four-dimensional
$\SU(2)$, where $\beta_W$ is of order $2$ and above, and the next proposition
explains why no one-step contraction of this kind reaches it.

\subsection{Single links near flat configurations}

Let $A(\kappa)=\int x_1\,\nu_{\kappa e_1}(dx)$ be the mean resultant length of
$\nu_b$ at $|b|=\kappa$. By symmetry $\int x\,\nu_b(dx)=A(|b|)b/|b|$; $A$ is
nondecreasing, its derivative being the variance of $x_1$, and
$A(\kappa)\to1$ as $\kappa\to\infty$.

\begin{proposition}\label{prop:flat}
Let $\beta>0$, let $e$ be a link whose six plaquettes contain no link fixed at a
value other than $1$, and let $k$ be a partner of $e$. The single-link influence
satisfies
\[
C_{ek}=\sup\frac{W_1(\nu_b,\nu_{b'})}{\delta(U_k,U'_k)}\ge\frac{A(24\beta)}6 .
\]
Hence in any single-link scheme in which every link has at least twelve free
partners, as for the forest of Theorem~\ref{thm:A}, every row sum of the
influence matrix is at least $2A(24\beta)$, which tends to $2$ as
$\beta\to\infty$.
\end{proposition}

\begin{proof}
Take all links equal to $1$, so every staple is $1$ and $b=24\beta\cdot1$.
Change $U_k$ so that the staple of the plaquette containing $e$ and $k$
becomes $\cos\vartheta+t\sin\vartheta$ with a unit imaginary quaternion $t$;
then $\delta(U_k,U'_k)=2\sin\frac\vartheta2$ and $b'=4\beta(5+\cos\vartheta+
t\sin\vartheta)$. The function $x\mapsto\ip xt$ is $1$-Lipschitz, so
$W_1(\nu_b,\nu_{b'})\ge\ip{\int x\,d\nu_{b'}-\int x\,d\nu_b}t=A(|b'|)\,
4\beta\sin\vartheta/|b'|$. Divide by $2\sin\frac\vartheta2$ and let
$\vartheta\to0$. For the limit of $A$: the law of $x_1$ under
$\nu_{\kappa e_1}$ has density proportional to $e^{\kappa t}\sqrt{1-t^2}$ on
$[-1,1]$, and for $0<\epsilon<\frac12$ the mass of $[-1,1-2\epsilon]$ divided
by that of $[1-\epsilon,1]$ is at most $e^{-\kappa\epsilon}\pi/\int_{1-
\epsilon}^1\sqrt{1-t^2}\,dt\to0$.
\end{proof}

At weak coupling the Gibbs measure concentrates near flat configurations, so
even an influence matrix computed only on typical configurations does not
contract there: per link it approaches $\frac16$ per partner, and
$12\cdot\frac16=2$.\footnote{Numerically $A(\kappa)=\frac12$ at
$\kappa\approx2.45$, that is $\beta_W\approx0.41$.} Larger blocks with better
internal estimates push the threshold to some finite $\beta_W$, but the mass at
weak coupling is non-perturbative, $m\sim\Lambda e^{-c\beta}$, and a
contraction can appear only at block sizes of order $e^{c\beta}$ lattice units.
Section~\ref{sec:weak} and Appendix~\ref{app:obstructions} make this precise.

\section{Free massive fields}\label{sec:free}

Here $d\ge1$, $n\ge2$, $L\ge n+1$, $\mathbb T=(\Z/L\Z)^d$, $\mu>0$ is a lattice
mass, $-\Delta$ is the lattice Laplacian and $Q=-\Delta+\mu^2$. For a positive
definite $Q$ on $\R^N$ let $\gamma_Q$ be the centred Gaussian law with
covariance $C=Q^{-1}$, and for $V\subset\{1,\dots,N\}$ let $\mathcal P_V$ be the
conditional expectation given the coordinates outside $V$. For the cubes
$V_x=x+\{0,\dots,n-1\}^d$, $x\in\mathbb T$, put
\begin{equation}\label{eq:ffvarpi}
\varpi(Q)=\inf_f\frac{n^{-d}\sum_x\|f-\mathcal P_{V_x}f\|^2}{\Var_{\gamma_Q}(f)},
\end{equation}
the infimum over nonconstant $f\in L^2(\gamma_Q)$, the best constant in
$\BP(n,\cdot)$. With $s=n\mu$ and $\mu=ma$, $n=\ell/a$, the parameter $s=m\ell$
is the mass times the physical size of the cube.

\begin{lemma}[linear functions suffice]\label{lem:chaos}
For a positive definite $Q$ and any family of coordinate sets $(V_x)$, the
infimum in~\eqref{eq:ffvarpi} equals its infimum over the linear functions
$\ell_h(\phi)=\ip h\phi$, and
\[
\|\ell_h-\mathcal P_V\ell_h\|^2=\ip{h_V}{Q_{VV}^{-1}h_V},\qquad
\Var(\ell_h)=\ip h{Q^{-1}h},
\]
with $h_V$ the restriction of $h$ to $V$ and $Q_{VV}$ the block of $Q$ on $V$.
\end{lemma}

\begin{proof}
The conditional law of $\phi_V$ given $\phi_{V^c}$ is Gaussian with covariance
$Q_{VV}^{-1}$, which gives the formula for $\ell_h$. Write $\phi=Q^{-1/2}\xi$ with
$\xi$ standard Gaussian. Conditioning on $\phi_{V^c}$ is conditioning on the
Gaussian subspace $K_V=\{\ip u\xi:u\in Q^{-1/2}\R^{V^c}\}$, and on the $m$-th
Wiener chaos, identified with the symmetric tensors $\mathrm{Sym}^m(\R^N)$, it
acts as $\pi_V^{\otimes m}$, with $\pi_V$ the orthogonal projection onto $Q^{-1/2}
\R^{V^c}$~\cite[Ch.~4]{Jan97}. The chaoses are orthogonal and invariant under
every $\mathcal P_V$, so it suffices to treat one chaos $m\ge1$. The projection
$I-\pi_V^{\otimes m}$ dominates $(I-\pi_V)\otimes I^{\otimes(m-1)}$, as both are
projections that commute and the range of the second lies in that of the
first. For a symmetric tensor $F$, therefore,
\[
\sum_x\|F-\pi_{V_x}^{\otimes m}F\|^2\ge\Bigl\langle F,\Bigl(\sum_x(I-\pi_{V_x})\Bigr)
\otimes I^{\otimes(m-1)}F\Bigr\rangle\ge\lambda\|F\|^2,
\]
where $\lambda$ is the least eigenvalue of $\sum_x(I-\pi_{V_x})$ on $\R^N$, which is
the infimum of the quotient over linear functions ($m=1$).
\end{proof}

\begin{theorem}[the exact constant]\label{thm:free}
Let $V=\{0,\dots,n-1\}^d$, $-\Delta_D$ its Dirichlet Laplacian and $\mathbf 1$ its
constant function. Then
\[
\varpi(-\Delta+\mu^2)=\frac{\mu^2}{n^d}\ip{\mathbf 1}{(-\Delta_D+\mu^2)^{-1}\mathbf 1},
\]
attained only at multiples of $\ell_{\mathbf 1}$ among linear functions, and
\[
\Bigl(\frac56\Bigr)^d\frac{s^2}{s^2+10d}\le\varpi(-\Delta+\mu^2)\le\frac{s^2}
{s^2+4dn^2\sin^2\frac\pi{2n+2}}\le\frac{s^2}{s^2+16d/9}.
\]
\end{theorem}

\begin{proof}
Since $L\ge n+1$, the block $Q_{V_xV_x}$ is the translate of $Q_{VV}=-\Delta_D+
\mu^2$; put $G_V=Q_{VV}^{-1}$. By Lemma~\ref{lem:chaos}, $n^d\varpi$ is the least
value of $\ip h{Sh}/\ip h{Q^{-1}h}$ with $S=\sum_xR_{V_x}^{\mathsf T}Q_{V_xV_x}^{-1}
R_{V_x}$, that is, the least eigenvalue of $Q^{1/2}SQ^{1/2}$. Translations commute
with $Q$ and $S$, so the plane waves $e_k(y)=e^{ik\cdot y}$ diagonalise it with
eigenvalues $\lambda_k=\hat q(k)\ip{e_k^V}{G_Ve_k^V}$, where $\hat q(k)=\mu^2+
\sum_i(2-2\cos k_i)$ and $e_k^V$ is the restriction to $V$.

Let $B$ map $f$ to the function $u\mapsto\sum_{w\sim u,\,w\notin V}f(w)$ on $V$.
Then $Q_{VV}R_Vf=R_VQf+Bf$, so $Q_{VV}e_k^V=\hat q(k)e_k^V+Be_k$, and pairing
$G_V$ of this with $e_k^V$ gives
\[
\lambda_k=n^d-\sum_{u,u'\in V}\sum_{w\sim u',\,w\notin V}G_V(u,u')\cos\bigl(k\cdot
(w-u)\bigr).
\]
Write $Q_{VV}=(2d+\mu^2)(I-K)$ with $K$ entrywise nonnegative and $\|K\|\le
2d/(2d+\mu^2)<1$; the Neumann series shows that $G_V$ has positive entries.
Hence $\lambda_k-\lambda_0=\sum G_V(u,u')(1-\cos(k\cdot(w-u)))\ge0$, with equality only
at $k=0$ (take $u=u'$ on a face and $w=u+e_i$). Since $Q\mathbf 1=\mu^2\mathbf 1$,
$\lambda_0=\mu^2\ip{\mathbf 1}{G_V\mathbf 1}$, which is the formula.

\emph{Upper bound.} $G_V\le(\lambda_1+\mu^2)^{-1}$ with $\lambda_1=4d\sin^2\frac\pi
{2n+2}$ the least eigenvalue of $-\Delta_D$, and $\|\mathbf 1\|^2=n^d$; then
$\sin t\ge2t/\pi$ gives $n^2\lambda_1\ge4dn^2/(n+1)^2\ge16d/9$.

\emph{Lower bound.} By the variational formula $\ip{\mathbf 1}{G_V\mathbf 1}\ge
\ip{\mathbf 1}w^2/\ip w{Q_{VV}w}$, take $w(u)=\prod_i\tau(u_i)$ with the parabolic
profile $\tau(s)=6(s+1)(n-s)/(n(n+1)(n+2))$ on $\{0,\dots,n-1\}$, zero
outside. Then $\sum\tau=1$, $S_2=\sum\tau^2=\alpha_n/n$ and $D_2=\sum_{s\in\Z}
(\tau(s+1)-\tau(s))^2=\beta_n/n^3$ with
\[
\alpha_n=\frac{6(n^2+2n+2)}{5(n+1)(n+2)}\le\frac65,\qquad
\beta_n=\frac{12n^2}{(n+1)(n+2)}=10\rho_n\alpha_n,\quad
\rho_n=\frac{n^2}{n^2+2n+2}<1,
\]
by the power sums $\sum_{u=1}^nu(n+1-u)=n(n+1)(n+2)/6$, $\sum_{u=1}^nu^2(n+1-u)^2=
n(n+1)(n+2)(n^2+2n+2)/30$ and $\sum_{u=0}^n(n-2u)^2=n(n+1)(n+2)/3$. The Dirichlet
form of the product is $\ip w{-\Delta_Dw}=dS_2^{d-1}D_2$ and $\ip{\mathbf 1}w=1$, so
\[
\varpi\ge\frac{\mu^2}{n^d(\mu^2S_2^d+dS_2^{d-1}D_2)}=\frac1{\alpha_n^d+d\alpha_n^{d-1}
\beta_ns^{-2}}=\alpha_n^{-d}\frac{s^2}{s^2+10d\rho_n}\ge\Bigl(\frac56\Bigr)^d
\frac{s^2}{s^2+10d}.\qedhere
\]
\end{proof}

At $n=2$ the two bounds coincide with $s^2/(s^2+4d)$. For $d=4$ the lower
bound is $0.0118$ at $s=1$ and $0.0438$ at $s=2$; the upper bound vanishes as
$s\to0$, so a massless free field satisfies no cube condition uniform in the
lattice spacing. The same proof, with the Hodge identity of the torus, gives
the cube condition for the free massive vector (Proca) field on gauge-invariant
functions, with the same lower bound and the upper bound $s^2/(1+s^2)$
\cite{BBLRS26}.


\section{Mass generation at one scale implies the cube condition}\label{sec:massgen}

Theorem~\ref{thm:B} reduces the continuum gap to statement (D), a cube
condition uniform along the tuning. This section reduces (D) in turn to a
hypothesis of a different kind, closer to the physics: a uniform decay of
correlations between the gauge-invariant observables of two slabs of a box
separated by a slab of physical width about $\ell/2$, for every boundary
condition. If the maximal correlation is at most $c<\frac12$, then the cube
condition holds with $\varpi=(1-2c)^4/32$ (Theorem~\ref{thm:D}). The value
$\frac12$ is the exact limit of the method (Proposition~\ref{prop:gausschain}
in Appendix~\ref{app:obstructions}); the hypothesis implies exponential decay
of gauge-invariant correlations under every boundary condition
(Proposition~\ref{prop:alldecay}); and it holds in two dimensions along the
weak-coupling tuning (Theorem~\ref{thm:twod}). In four dimensions it is a
form of mass generation at the physical scale and it is not proved here.

\subsection{Boxes, boundary conditions and gauge-invariant observables}

In this section $G\subset\U(N)$ is compact and nontrivial, $\varphi$ is a
positive continuous class function on $G$, $\lambda$ denotes normalised Haar
measure on $G$ or on a product of copies of $G$, and
$\omega_\varphi=\log(\max\varphi/\min\varphi)$. Links of $\Z^4$ are written
$(y,\nu)$ with base point $y$, from $y$ to $y+e_\nu$. A \emph{box} is a set
$Q=\prod_{i=1}^4[\alpha_i,\alpha_i+s_i)\cap\Z$ with integers $s_i\ge1$, its
\emph{widths}, and $V(Q)$ is the set of links with base point in $Q$. For an
exterior configuration $\omega\in G^{E(\Z^4)\setminus V(Q)}$ let
$\nu^\omega_Q$ be the probability measure on $G^{V(Q)}$ with density
proportional to $\prod_{p\cap V(Q)\ne\emptyset}\varphi((\sigma,\omega)_p)$
with respect to $\lambda$, a finite product. It is the conditional law of
$U_{V(Q)}$ given the other links under any Gibbs measure with weight
$\varphi$, and on a torus $\Lambda_L$ with all $s_i\le L-3$ it is the
conditional law of the torus measure.\footnote{The plaquettes meeting $V(Q)$
have their links based in $\prod_i[\alpha_i-1,\alpha_i+s_i]$, which embeds in
$\Lambda_L$ when $s_i\le L-3$.} For $A\subset V(Q)$ write $E_A$ for the
conditional expectation under $\nu^\omega_Q$ given $\sigma_{V(Q)\setminus A}$;
its regular version is integration against the measure on $G^A$ with density
proportional to $\prod_{p\cap A\ne\emptyset}\varphi(U_p)$, which is jointly
continuous in all variables. In particular, if $Q'\subset Q$ is a box, the
conditional law of $\sigma_{V(Q')}$ under $\nu^\omega_Q$ given
$\sigma_{V(Q)\setminus V(Q')}=\eta$ is $\nu^{(\eta,\omega)}_{Q'}$.

A vertex $v$ is \emph{interior} to $Q$ if the eight links incident to $v$ lie
in $V(Q)$. Let $\mathcal G_Q$ be the group of $h\in G^{\Z^4}$ with $h_v=1$ at
every vertex not interior to $Q$, acting by $(h\sigma)_{(y,\nu)}=h_y
\sigma_{(y,\nu)}h_{y+e_\nu}^{-1}$; it fixes every link outside $V(Q)$, whose
end points are not interior. Let $\mathcal A^\omega_Q$ be the set of real
$F\in L^2(\nu^\omega_Q)$ with $F\circ h=F$ almost surely for every $h\in
\mathcal G_Q$: the gauge-invariant observables of the box, the boundary
condition having fixed the gauge on the boundary.

\begin{lemma}\label{lem:boxinv}
Let $Q$ be a box, $\omega$ an exterior configuration and $\nu=\nu^\omega_Q$.
\begin{enumerate}[label=\textup{(\roman*)},leftmargin=*]
\item $\nu$ is invariant under $\mathcal G_Q$.
\item For every $S\subset V(Q)$, each $h\in\mathcal G_Q$ maps
$\sigma(\sigma_S)$ onto itself, and $E_A(F\circ h)=(E_AF)\circ h$ for
$A\subset V(Q)$ and $F\in L^2(\nu)$.
\item $\mathcal A^\omega_Q$ is a closed subspace that contains the constants
and is preserved by every $E_A$. Every $F\in\mathcal A^\omega_Q$ has a
version $\tilde F$ with $\tilde F(h\sigma)=\tilde F(\sigma)$ for all $h$ and
$\sigma$, and if $F$ is $\sigma(\sigma_S)$-measurable, $\tilde F$ can be taken
$\sigma(\sigma_S)$-measurable.
\end{enumerate}
\end{lemma}

\begin{proof}
(i) $h$ multiplies each link on the left and on the right by fixed elements,
so it preserves $\lambda$. Since $h$ fixes the links outside $V(Q)$,
$(h\sigma,\omega)=h\cdot(\sigma,\omega)$, and the holonomy of a plaquette
with base point $y$ is conjugated by $h_y$; $\varphi$ is a class function,
so the density is invariant.

(ii) $(h\sigma)_e$ is a function of $\sigma_e$ alone and $h^{-1}\in\mathcal
G_Q$, so $h$ maps $\sigma(\sigma_S)$ onto itself. For a bounded
$\sigma(\sigma_{V(Q)\setminus A})$-measurable $K$, by (i),
$\nu((F\circ h)K)=\nu(F\,(K\circ h^{-1}))=\nu((E_AF)(K\circ h^{-1}))=
\nu(((E_AF)\circ h)K)$, and $(E_AF)\circ h$ is
$\sigma(\sigma_{V(Q)\setminus A})$-measurable.

(iii) By (i) each map $F\mapsto F\circ h$ is unitary, so its fixed space is
closed, and so is the intersection of these spaces; (ii) gives the
invariance under $E_A$. The group $\mathcal G_Q\cong G^I$, $I$ the finite set
of interior vertices, is compact and acts continuously. For a measurable
version $F_0$ put $\tilde F(\sigma)=\int F_0(h\sigma)\,dh$ where
$\int|F_0(h\sigma)|\,dh<\infty$, and $\tilde F(\sigma)=0$ elsewhere. By the
invariance of Haar measure under right translations, the set where the
integral converges and $\tilde F$ are invariant under every $k\in\mathcal
G_Q$. For each $h$, $F_0(h\sigma)=F_0(\sigma)$ for almost every $\sigma$, so
by Fubini's theorem, for almost every $\sigma$ the equality holds for almost
every $h$, and $\tilde F=F_0$ almost surely. If $F_0$ is a function of
$\sigma_S$, so is $\sigma\mapsto F_0(h\sigma)$ for each $h$, and so is
$\tilde F$.
\end{proof}

\begin{definition}\label{def:rho}
Let $Q$ be a box, $i$ a coordinate, and $Q=Q_-\cup Q_0\cup Q_+$ a splitting
into three nonempty boxes, consecutive in the coordinate $i$ in this order:
the \emph{outer parts} $Q_\pm$ and the \emph{middle} $Q_0$. For an exterior
configuration $\omega$ put $\nu=\nu^\omega_Q$ and
\[
\rho^\omega(Q_-,Q_+)=\sup|\nu(XY)|,
\]
the supremum over $X\in\mathcal A^\omega_Q$ measurable with respect to
$\sigma_{V(Q_-)}$ and $Y\in\mathcal A^\omega_Q$ measurable with respect to
$\sigma_{V(Q_+)}$ with $\nu(X)=\nu(Y)=0$ and $\nu(X^2)=\nu(Y^2)=1$ (the
supremum of the empty set being $0$).
\end{definition}

This is the maximal correlation\footnote{The maximal correlation of two
random variables, the supremum of the correlations of square-integrable
functions of each, was introduced by Hirschfeld and by
Gebelein~\cite{Geb41}; for jointly Gaussian variables it equals the
absolute value of the ordinary correlation.} between the gauge-invariant
observables of the two outer parts, the middle being integrated out, under
the boundary condition $\omega$. By the Cauchy--Schwarz inequality
$\rho^\omega\le1$.

\begin{definition}\label{def:M}
Let $d\ge2$, $m\ge4$ and $w\ge1$ be integers and $0\le c\le1$.
\begin{enumerate}[label=\textup{(\roman*)},leftmargin=*]
\item $(\mathrm M_{d,m}(c))$: $\rho^\omega(Q_-,Q_+)\le c$ for every box $Q$
all of whose widths lie in $[2d-1,md-1]$, every coordinate $i$ in which $Q$
has width $3d-1$, the splitting in which $Q_-$, $Q_0$, $Q_+$ have widths
$d$, $d-1$, $d$ in the coordinate $i$, and every exterior configuration
$\omega$.
\item $(\mathrm M^*_w(c))$: $\rho^\omega(Q_-,Q_+)\le c$ for every box $Q$,
every coordinate $i$, every splitting whose middle has width $w$ in the
coordinate $i$, the outer parts being of any widths, and every $\omega$.
\end{enumerate}
\end{definition}

The boxes of $(\mathrm M_{d,m}(c))$ are fat: all their widths are at least
$2d-1$. The restriction matters. In two dimensions the same condition for
boxes of width one in a transverse direction fails along the weak-coupling
tuning, through a broken global symmetry of a pinned strip, while the fat
form holds there (Theorem~\ref{thm:twod} and Proposition~\ref{prop:strip}
in Appendix~\ref{app:obstructions}). At $\varphi\equiv1$ the links are
independent under $\nu^\omega_Q$, so $\rho^\omega=0$ and both conditions hold
with $c=0$. Along the tuning, with $d\approx\ell/(2a)$, $(\mathrm
M_{d,m}(c))$ is a statement at the physical separation $(d-1)a\approx\ell/2$
in boxes of physical size at most about $m\ell/2$, uniform in the boundary
condition, and stated for all square-integrable gauge-invariant observables
of the two regions rather than for pairs of local ones.

\subsection{Two projections and a chain of projections}

\begin{lemma}\label{lem:friedrichs}
Let $P_1,P_2$ be orthogonal projections on a real Hilbert space $\mathcal K$,
$\mathcal N$ a closed subspace of $\operatorname{ran}P_1\cap
\operatorname{ran}P_2$, and suppose that $|\ip xy|\le c\|x\|\|y\|$ for all
$x\in\operatorname{ran}P_1\ominus\mathcal N$ and $y\in\operatorname{ran}P_2
\ominus\mathcal N$. Put $Q_k=I-P_k$ and $X=Q_1+Q_2$. Then $X^2\ge(1-c)X$,
that is, $Q_1Q_2+Q_2Q_1\ge-c(Q_1+Q_2)$.
\end{lemma}

\begin{proof}
For $n\in\mathcal N$, $P_kn=n$, so $Xn=0$; as $X$ is self-adjoint it also
preserves $\mathcal N^\perp$. Let $g\in\mathcal N^\perp$. Then $P_kg\perp
\mathcal N$, since $\ip{P_kg}n=\ip g{P_kn}=\ip gn=0$, so with $p=\|P_1g\|$
and $q=\|P_2g\|$, $|\ip{P_1g}{P_2g}|\le cpq$. Put $s=\ip g{(P_1+P_2)g}=
p^2+q^2$. Then $\|(P_1+P_2)g\|^2=p^2+q^2+2\ip{P_1g}{P_2g}\le(1+c)(p^2+q^2)$
and $s\le\|g\|\|(P_1+P_2)g\|\le\|g\|((1+c)s)^{1/2}$, so $s\le(1+c)\|g\|^2$
and $\ip g{Xg}=2\|g\|^2-s\ge(1-c)\|g\|^2$. So $X=0$ on $\mathcal N$, and on
$\mathcal N^\perp$ the operator $X-(1-c)$ is nonnegative and commutes with
$X\ge0$; hence $X^2-(1-c)X=X^{1/2}(X-(1-c))X^{1/2}\ge0$ there. Finally
$X^2-X=Q_1Q_2+Q_2Q_1$.
\end{proof}

The angle between $\operatorname{ran}P_1\ominus\mathcal N$ and
$\operatorname{ran}P_2\ominus\mathcal N$ is the Friedrichs angle of the two
ranges, and the lemma is a form of the classical relation between that angle
and the rate of the alternating projections~\cite{Deu01}.

\begin{lemma}\label{lem:chain}
Let $Q_0,\dots,Q_M$ be orthogonal projections on a real Hilbert space
$\mathcal K$ with $Q_tQ_{t'}=Q_{t'}Q_t$ whenever $|t-t'|\ge2$ and
$Q_tQ_{t+1}+Q_{t+1}Q_t\ge-c(Q_t+Q_{t+1})$ for $0\le t<M$, with $0\le c<
\frac12$. Then $H=\sum_tQ_t$ satisfies $H^2\ge(1-2c)H$, and $\ip g{Hg}\ge
(1-2c)\|g\|^2$ for every $g\perp\ker H$.
\end{lemma}

\begin{proof}
$H^2=\sum_tQ_t+\sum_{t<M}(Q_tQ_{t+1}+Q_{t+1}Q_t)+\sum Q_tQ_{t'}$, the last
sum over the ordered pairs with $|t-t'|\ge2$. Each such $Q_tQ_{t'}$ is a
product of commuting projections, hence a projection and nonnegative. The
middle sum is at least $-c\sum_{t<M}(Q_t+Q_{t+1})\ge-2cH$. So $H^2-(1-2c)H
\ge0$, and by the spectral theorem the spectrum of the nonnegative operator
$H$ lies in $\{0\}\cup[1-2c,\infty)$, which gives the second claim.
\end{proof}

\begin{lemma}\label{lem:constant}
Let $Q$ be a box, $\omega$ an exterior configuration, and $R_0,\dots,R_M$
subsets of $V(Q)$ with union $V(Q)$. If $F\in L^2(\nu^\omega_Q)$ is, for every
$t$, almost surely equal to a $\sigma(\sigma_{V(Q)\setminus
R_t})$-measurable function, then $F$ is almost surely constant.
\end{lemma}

\begin{proof}
The density of $\nu=\nu^\omega_Q$ with respect to $\lambda$ is bounded above
and below by positive constants, so $L^2(\nu)=L^2(\lambda)$ with the same
null sets. Let $\Lambda_t$ be the conditional expectation under $\lambda$
given $\sigma_{V(Q)\setminus R_t}$, which integrates out $\sigma_{R_t}$. By
hypothesis $F=F_t$ almost surely with $F_t$ a function of
$\sigma_{V(Q)\setminus R_t}$, so $\Lambda_tF=F_t=F$. By Fubini's theorem
$\Lambda_0\cdots\Lambda_M$ integrates out all of $\bigcup_tR_t=V(Q)$, so
$F=\Lambda_0\cdots\Lambda_MF$ is constant.
\end{proof}

\subsection{One direction}

Fix a box $R=\prod_iI_i$ and a coordinate $i$ with $I_i=[\alpha,\alpha+
(M+2)d-1)$ for an integer $M\ge0$. For $0\le t\le M$ let $R_t$ be the box
obtained from $R$ by replacing $I_i$ with $\alpha+dt+[0,2d-1)$, and for
$t<M$ let $Q_t$ be the box obtained by replacing $I_i$ with $\alpha+dt+
[0,3d-1)$, split into $Q_t^-$, $Q_t^0$, $Q_t^+$ by $\alpha+dt+[0,d)$,
$\alpha+dt+[d,2d-1)$ and $\alpha+dt+[2d-1,3d-1)$. Then
\begin{equation}\label{eq:slabs}
V(Q_t)=V(R_t)\cup V(R_{t+1}),\quad V(Q_t^-)=V(R_t)\setminus V(R_{t+1}),
\quad V(Q_t^+)=V(R_{t+1})\setminus V(R_t),
\end{equation}
and $\bigcup_tV(R_t)=V(R)$, since the intervals $\alpha+dt+[0,2d-1)$ cover
$I_i$.

\begin{proposition}[one direction]\label{prop:slab}
In this setting let $\omega$ be an exterior configuration of $R$,
$\nu=\nu^\omega_R$, $P_t=E_{V(R_t)}$, and suppose that $\rho^{\omega'}(Q^-_t,
Q^+_t)\le c<\frac12$ for every $t<M$ and every exterior configuration
$\omega'$ of $Q_t$. Then every $F\in\mathcal A^\omega_R$ satisfies
\[
(1-2c)\Var_\nu(F)\le\sum_{t=0}^M\nu\bigl((F-P_tF)^2\bigr).
\]
\end{proposition}

\begin{proof}
For $M=0$, $P_0F=\nu(F)$ and the claim is trivial; let $M\ge1$. Put
$\mathcal K=\mathcal A^\omega_R$. By Lemma~\ref{lem:boxinv}(iii) each $P_t$
preserves $\mathcal K$ and restricts to an orthogonal projection of
$\mathcal K$; let $Q_t=I-P_t$ on $\mathcal K$.

\emph{Non-adjacent moves commute.} Let $t'\ge t+2$. The base points of the
links of $V(R_t)$ have $i$-th coordinate at most $\alpha+dt+2d-2$ and those
of $V(R_{t'})$ at least $\alpha+dt+2d$, so no plaquette meets both sets.
Given $\sigma$ off $V(R_t)\cup V(R_{t'})$, the conditional density of
$(\sigma_{V(R_t)},\sigma_{V(R_{t'})})$ is a product of factors each of which
involves the links of only one of the two sets. So they are conditionally
independent, and $P_tP_{t'}F$ and $P_{t'}P_tF$ are both the integral of $F$
against the product of the two conditional laws. Hence $P_tP_{t'}=
P_{t'}P_t$, and the $Q_t$ commute.

\emph{Adjacent moves.} Fix $t<M$ and write $a=V(Q_t^-)$, $b=V(Q_t^0)$,
$c'=V(Q_t^+)$, $e=V(R)\setminus V(Q_t)$. By~\eqref{eq:slabs},
$V(R)\setminus V(R_t)=c'\cup e$ and $V(R)\setminus V(R_{t+1})=a\cup e$, so
the ranges of $P_t$ and of $P_{t+1}$ in $\mathcal K$ are the elements of
$\mathcal K$ measurable with respect to $(\sigma_{c'},\sigma_e)$ and to
$(\sigma_a,\sigma_e)$. Let $\mathcal N$ be the elements of $\mathcal K$
measurable with respect to $\sigma_e$, a closed subspace of both ranges. We
check the hypothesis of Lemma~\ref{lem:friedrichs}. Let $x$ and $y$ be in
the ranges of $P_t$ and $P_{t+1}$ and orthogonal to $\mathcal N$. The
conditional expectation $E_e$ given $\sigma_e$ is $E_{V(Q_t)}$, which
preserves $\mathcal K$, so $E_ex\in\mathcal N$ and $\|E_ex\|^2=\ip x{E_ex}=0$;
likewise $E_ey=0$. By Lemma~\ref{lem:boxinv}(iii) for the box $R$ take
versions $x=\tilde x(\sigma_{c'},\sigma_e)$ and $y=\tilde y(\sigma_a,\sigma_e)$
invariant at every point under $\mathcal G_R$. Every vertex interior to $Q_t$
is interior to $R$, so $\mathcal G_{Q_t}\subset\mathcal G_R$, and
$\mathcal G_{Q_t}$ fixes $\sigma_e$. Hence for each $\eta\in G^e$ the
functions $X_\eta=\tilde x(\cdot,\eta)$ of $\sigma_{V(Q^+_t)}$ and $Y_\eta=
\tilde y(\cdot,\eta)$ of $\sigma_{V(Q^-_t)}$ are invariant under
$\mathcal G_{Q_t}$ at every point. The conditional law of $\sigma_{V(Q_t)}$
given $\sigma_e=\eta$ is $\nu_\eta=\nu^{(\eta,\omega)}_{Q_t}$, so for almost
every $\eta$, $X_\eta,Y_\eta\in L^2(\nu_\eta)$, $\nu_\eta(X_\eta)=(E_ex)(\eta)
=0$, $\nu_\eta(Y_\eta)=0$, and by the hypothesis with $\omega'=(\eta,\omega)$,
\[
|\nu_\eta(X_\eta Y_\eta)|\le c\,\nu_\eta(X_\eta^2)^{1/2}\nu_\eta(Y_\eta^2)^{1/2}.
\]
Integrating over $\eta$ and using the Cauchy--Schwarz inequality,
$|\ip xy|\le c\|x\|\|y\|$, and Lemma~\ref{lem:friedrichs} gives
$Q_tQ_{t+1}+Q_{t+1}Q_t\ge-c(Q_t+Q_{t+1})$ on $\mathcal K$.

\emph{Conclusion.} Lemma~\ref{lem:chain} applies. If $Hg=0$ for $g\in
\mathcal K$, then $\sum_t\|Q_tg\|^2=\ip g{Hg}=0$, so $g=P_tg$ for all $t$ and
$g$ is constant by Lemma~\ref{lem:constant}. For $F\in\mathcal K$ the
function $g=F-\nu(F)$ is orthogonal to the constants, hence to $\ker H$, and
$\ip g{Hg}=\sum_t\nu((F-P_tF)^2)$.
\end{proof}

\subsection{The patch gap and the cube condition}

For $d\ge2$, $m\ge4$ and $r\in\Z^4$ let $\Pi=r+[0,md-1)^4$, the box of
base points of a patch, and for $j\in[0,m-2]^4$ let $B_j=V(r+dj+
[0,2d-1)^4)$, as in Appendix~\ref{app:finite}. For an exterior
configuration $\omega$ of $\Pi$ write $\nu=\nu^\omega_\Pi$ and $P_j=E_{B_j}$.

\begin{theorem}[patch gap from slab decorrelation]\label{thm:patch}
Let $d\ge2$, $m\ge4$, $0\le c<\frac12$, and assume $(\mathrm M_{d,m}(c))$.
Then for every $r$, every $\omega$ and every $F\in\mathcal A^\omega_\Pi$,
\[
(1-2c)^4\Var_\nu(F)\le\sum_{j\in[0,m-2]^4}\nu\bigl((F-P_jF)^2\bigr).
\]
\end{theorem}

\begin{proof}
For $0\le k\le4$ and $\tau\in[0,m-2]^k$ let $R^{(k)}_\tau=r+\prod_{i\le k}
(d\tau_i+[0,2d-1))\times[0,md-1)^{4-k}$ and $P^{(k)}_\tau=
E_{V(R^{(k)}_\tau)}$, under $\nu$. We prove by induction on $k$ that
\[
(1-2c)^k\Var_\nu(F)\le\sum_{\tau\in[0,m-2]^k}\nu\bigl((F-P^{(k)}_\tau F)^2
\bigr)\qquad(F\in\mathcal A^\omega_\Pi).
\]
For $k=0$, $R^{(0)}=\Pi$ and $P^{(0)}F=\nu(F)$. Let $k<4$, fix $\tau$, write
$R=R^{(k)}_\tau$, and split $\sigma=(\sigma',\omega')$ with $\sigma'=
\sigma_{V(R)}$ and $\omega'=\sigma_{V(\Pi)\setminus V(R)}$. Take $F$
invariant at every point under $\mathcal G_\Pi$ (Lemma~\ref{lem:boxinv}(iii)).
Since $\mathcal G_R\subset\mathcal G_\Pi$ fixes $\omega'$, each
$F(\cdot,\omega')$ is invariant under $\mathcal G_R$. The conditional law of
$\sigma'$ given $\omega'$ is $\nu'=\nu^{(\omega',\omega)}_R$, so
\[
\nu\bigl((F-P^{(k)}_\tau F)^2\bigr)=\int\Var_{\nu'}\bigl(F(\cdot,\omega')
\bigr)\,d\nu(\omega').
\]
In the coordinate $k+1$ the box $R$ has width $md-1=((m-2)+2)d-1$. The
boxes $R_u$ and $Q_u$ of Proposition~\ref{prop:slab} built from $R$ with
$i=k+1$ and $M=m-2$ are $R_u=R^{(k+1)}_{(\tau,u)}$ and boxes whose widths are
$2d-1$ in the coordinates $\le k$, $3d-1$ in the coordinate $k+1$ and $md-1$
in the coordinates $>k+1$, split as in $(\mathrm M_{d,m}(c))$; all these
widths lie in $[2d-1,md-1]$. By Proposition~\ref{prop:slab}, for almost every
$\omega'$,
\[
(1-2c)\Var_{\nu'}(F(\cdot,\omega'))\le\sum_{u=0}^{m-2}\nu'\bigl((F(\cdot,
\omega')-P'_uF(\cdot,\omega'))^2\bigr),
\]
with $P'_u$ the conditional expectation under $\nu'$ given $\sigma'$ off
$V(R^{(k+1)}_{(\tau,u)})$. The conditional law of the links of that set given
the other links of $\Pi$ is the same under $\nu$ and under $\nu'$, both
having density proportional to $\prod_p\varphi(U_p)$ over the plaquettes
meeting it; so $P'_uF(\cdot,\omega')=(P^{(k+1)}_{(\tau,u)}F)(\cdot,\omega')$,
and integrating over $\omega'$ gives $\nu((F-P^{(k+1)}_{(\tau,u)}F)^2)$.
Multiplying the induction hypothesis by $1-2c$ and inserting these bounds
proves the case $k+1$. For $k=4$, $V(R^{(4)}_\tau)=B_\tau$.
\end{proof}

\begin{mainthm}[mass generation implies the cube condition]\label{thm:D}
Let $d\ge2$, $0\le c<\frac12$ and $m\ge4$ with $m-1\ge148/(1-2c)^4$, and
assume $(\mathrm M_{d,m}(c))$. Then for every even $L$ with $d\mid L$ and
$L\ge2md$ the measure $\mu_L$ with weight $\varphi$ satisfies
\[
\BP\Bigl(2d-1,\frac{(1-2c)^4}{32}\Bigr).
\]
Consequently, let $(\beta_k,a_k)$ be a tuning with $a_k\to0$ and Wilson
weights, and $d_k\ge2$ integers with $(2d_k-1)a_k\le\ell$. If
$(\mathrm M_{d_k,m}(c))$ holds for every $k$ with $c<\frac12$ and $m$ as
above independent of $k$, then statement (D) holds with $n_k=2d_k-1$ and
$\varpi=(1-2c)^4/32$, and every continuum limit as in Theorem~\ref{thm:B}
along the tori with $d_k\mid L$ satisfies
\[
\sigma(H)\subset\{0\}\cup\Bigl[\frac{(1-2c)^4}{1250\,e^2\ell},\infty\Bigr),
\]
with $0$ a simple eigenvalue; a nontrivial limit has a mass gap at least
this large.
\end{mainthm}

\begin{proof}
Theorem~\ref{thm:patch} says that the gauge-invariant patch gap of
Definition~\ref{def:patchgap} satisfies $\gamma_{d,m}\ge(1-2c)^4$, and
Theorem~\ref{thm:finite} with $\gamma=(1-2c)^4$ gives
$\BP(2d-1,(1-2c)^4/32)$. Along the tuning this holds on an infinite set of
even tori with $L\ge2(2d_k-1)+4$ for each $k$, with $(2d_k-1)a_k\le\ell$, and
Theorem~\ref{thm:B} gives the gap $16\varpi/(625e^2\ell)$ with
$\varpi=(1-2c)^4/32$.
\end{proof}

The patch gap obtained, $(1-2c)^4$, does not decay with $m$, whereas the
finite-size criterion needs only $148/(m-1)$; this slack is what allows a
large but fixed number $m$ of blocks per patch. The threshold $c<\frac12$
cannot be raised for this argument: for the massless Gaussian chain every
adjacent pair of moves has maximal correlation exactly $\frac12$, the
non-adjacent moves commute, and the gap of the slab dynamics tends to zero
(Proposition~\ref{prop:gausschain}).

\subsection{Decay of correlations under every boundary condition}

\begin{lemma}\label{lem:boxmarkov}
Let $Q$ be a box, $\omega$ an exterior configuration, $i$ a coordinate, and
$Y_0,\dots,Y_K$ the sets $V(Q\cap\{x_i\in J_k\})$ for consecutive nonempty
integer intervals $J_0<J_1<\dots<J_K$ whose union is the $i$-th side of $Q$.
Under $\nu^\omega_Q$, for each $k$, the families $(\sigma_{Y_l})_{l<k}$ and
$(\sigma_{Y_l})_{l>k}$ are conditionally independent given $\sigma_{Y_k}$.
\end{lemma}

\begin{proof}
The links in $V(Q)$ of a plaquette are based at points whose $i$-th
coordinates take at most two consecutive values, so they lie in one $Y_l$ or
in two consecutive ones. Grouping the factors of the density by the least
$l$ met, the density of $\nu^\omega_Q$ with respect to $\lambda$ is
$\prod_l\psi_l(\sigma_{Y_l},\sigma_{Y_{l+1}})$, with $\psi_K$ a function of
$\sigma_{Y_K}$, which gives the conditional independence.
\end{proof}

\begin{proposition}\label{prop:alldecay}
Assume $(\mathrm M^*_w(c))$ with $w\ge1$ and $0\le c\le1$. Let $Q$ be a box,
$\omega$ an exterior configuration, $\nu=\nu^\omega_Q$, $i$ a coordinate,
$s\in\Z$, $r\ge1$, and let $f,g\in\mathcal A^\omega_Q$ be measurable with
respect to $\sigma_{V(Q\cap\{x_i\le s\})}$ and $\sigma_{V(Q\cap\{x_i\ge
s+r\})}$. Then
\[
|\Cov_\nu(f,g)|\le c^{\lfloor r/(w+1)\rfloor}\|f-\nu(f)\|\,\|g-\nu(g)\|.
\]
\end{proposition}

\begin{proof}
If one of the two parts of $Q$ is empty, $f$ or $g$ is constant and the
covariance vanishes. Otherwise the $i$-th side of $Q$ contains $[s,s+r]$.
Let $J=\lfloor r/(w+1)\rfloor$; for $J=0$ the claim is the Cauchy--Schwarz
inequality. Let $J\ge1$ and split the $i$-th side into consecutive
intervals: $T_0$, the values $\le s$; for $1\le k\le J$, $M_k=[s+(k-1)(w+1)+
1,\,s+(k-1)(w+1)+w]$; for $1\le k\le J-1$, $T_k=\{s+k(w+1)\}$; and $T_J$, the
values $\ge s+J(w+1)$, which contain the values $\ge s+r$. Use the same
letters for the corresponding sets of links. By Lemma~\ref{lem:boxmarkov}
the sequence $T_0,M_1,T_1,\dots,M_J,T_J$ is Markov, and by
Lemma~\ref{lem:boxinv} the conditional expectations given $\sigma_{T_k}$
preserve $\mathcal A^\omega_Q$.

For $1\le k\le J$ let $u\in\mathcal A^\omega_Q$ be a centred function of
$\sigma_{T_{k-1}}$ and $v=E[u\mid\sigma_{T_k}]$, which is centred and lies
in $\mathcal A^\omega_Q$. The box $Q$ splits in the coordinate $i$ into the
values $\le s+(k-1)(w+1)$, the middle $M_k$ of width $w$, and the values
$\ge s+k(w+1)$; $u$ is a function of the first part and $v$ of the last, so
by $(\mathrm M^*_w(c))$, $\|v\|^2=\nu(uv)\le c\|u\|\|v\|$, that is,
$\|v\|\le c\|u\|$. Let $f_0=f-\nu(f)$ and $u_k=E[f_0\mid\sigma_{T_k}]$. By
the Markov property $u_k=E[u_{k-1}\mid\sigma_{T_k}]$, so $\|u_k\|\le
c\|u_{k-1}\|$ and $\|u_J\|\le c^J\|f_0\|$. As $g$ is
$\sigma_{T_J}$-measurable, $\Cov_\nu(f,g)=\nu(u_J(g-\nu(g)))$, and the
Cauchy--Schwarz inequality gives the claim.
\end{proof}

So $(\mathrm M^*_{d-1}(c))$ with $c<1$, which contains $(\mathrm
M_{d,m}(c))$ for every $m$, gives exponential decay of the covariances of
gauge-invariant observables of a box, uniformly in the box and in the
boundary condition, at rate $\log(1/c)/d$ per lattice spacing, that is, of
order $\log(1/c)/\ell$ in physical units along the tuning. This is close to
the decay hypothesis of Chatterjee~\cite[Def.~2.3]{Cha21}, under which an
area law for Wilson loops follows; there the observables are local
functions of a cube with prescribed boundary, and the geometries differ, so
neither condition is claimed to imply the other. The converse direction is
the weak point of any pairwise form: a bound $K_1e^{-K_2r}$ for pairs of
local functions, summed over the pairs of two slabs of transverse size
$(md)^3$, bounds $\rho^\omega$ only up to a factor of order $(md)^6$, which
is not uniform in the lattice spacing.

\subsection{What is known about the hypothesis}

\emph{Strong coupling at fixed $d$ and $m$.} The hypothesis is not
vacuous: $(\mathrm M_{d,m}(c))$ holds with $c=1-e^{-24\omega_\varphi
(md)^4}$. Let $Q$, the splitting and $\omega$ be as in Definition~\ref{def:M},
$A_1=V(Q_-\cup Q_0)$, $A_2=V(Q_0\cup Q_+)$, $\nu=\nu^\omega_Q$, and let
$\gamma_2$ be the best constant in
\[
\gamma_2\Var_\nu(g)\le\sum_{k=1,2}\nu\bigl((g-E_{A_k}g)^2\bigr),\qquad g\in
L^2(\nu)\text{ real}.
\] Under $\lambda$ the links are independent
and $A_1\cup A_2=V(Q)$, so the Efron--Stein inequality~\cite{ES81} gives the
constant $1$ for $\lambda$. The logarithm of the density of $\nu$ with
respect to $\lambda$ has oscillation $\Omega$ at most $\omega_\varphi$ times
the number of plaquettes meeting $V(Q)$, which is at most $24(md)^4$, and a
change of the reference measure by a density whose logarithm has oscillation
$\Omega$ changes the constant by at most the factor
$e^{-\Omega}$~\cite{HS87}.\footnote{Write $d\nu/d\lambda=e^V/Z$. Then $\Var_\nu(g)\le
Z^{-1}e^{\sup V}\Var_\lambda(g)$, and $\lambda((g-E^\lambda_Ag)^2)\le
Ze^{-\inf V}\nu((g-E_Ag)^2)$, as both conditional expectations project onto
the functions of the links outside $A$.} So $\gamma_2\ge e^{-24\omega_\varphi
(md)^4}$. Now let $X,Y$ be as in Definition~\ref{def:rho} with
$r=\nu(XY)\ge0$, and $g=X+Y$, so $\nu(g)=0$ and $\|g\|^2=2+2r$. Since $X$ is a
function of links off $A_2$, $E_{A_2}X=X$, and $\|X-E_{A_1}X\|^2=1-
\|E_{A_1}X\|^2\le1-r^2$ because $Y$ is a unit vector in the range of $E_{A_1}$
and $|\ip{E_{A_1}X}Y|=|\ip XY|=r$; symmetrically for $Y$. Hence
\[
\gamma_2(2+2r)\le\sum_k\|g-E_{A_k}g\|^2\le2(1-r^2),
\]
so $r\le1-\gamma_2$. The range of
couplings obtained in this way shrinks with $md$, and it is far smaller than
that of Theorem~\ref{thm:C}.

\emph{Two dimensions.} In two dimensions the seam between the outer parts
of a box is one-dimensional, and after an axial gauge the middle slab
integrates out exactly to a convolution power of the weight. This proves the
two-dimensional form of the hypothesis along the weak-coupling tuning, with
$c\to0$ as the physical size of the box grows, uniformly in the boundary
condition (Theorem~\ref{thm:twod} in Appendix~\ref{app:obstructions}). In
four dimensions the seam carries a three-dimensional layer of links, and
$(\mathrm M_{d,m}(c))$ is a bound, uniform in the boundary condition, on the
second singular value of the $(d-1)$-step transfer operator of the box on
gauge-invariant functions. Boundary conditions with values in the centre,
't~Hooft's twisted ones~\cite{tH79}, leave a global conjugation symmetry of
the box, and in the two-dimensional strip such a symmetry is broken at weak
coupling (Proposition~\ref{prop:strip}); in four dimensions the hypothesis
asks that the twisted sectors be disordered at the scale $\ell/2$, which is
the form that confinement takes here. We know no boundary condition that
refutes $(\mathrm M_{d,m}(c))$ along the four-dimensional tuning, and no
proof of it.

\emph{An area law does not suffice.} The next proposition shows that the
Wilson loops of the defining representation carry no information about the
observables that (D) constrains.

\begin{proposition}[central extensions]\label{prop:cover}
Let $\tilde G\subset\U(\tilde N)$ and $G\subset\U(N)$ be compact,
$\pi:\tilde G\to G$ a continuous surjective homomorphism with finite central
kernel $Z$, $\varphi$ a positive continuous class function on $G$ and
$\tilde\varphi=\varphi\circ\pi$. On $\Lambda_L$ let $\tilde\mu$, $\mu$ be the
measures with weights $\tilde\varphi$, $\varphi$, and $\Pi:\tilde G^E\to G^E$
the map $\pi$ applied to every link.
\begin{enumerate}[label=\textup{(\roman*)},leftmargin=*]
\item $\tilde\varphi$ is a positive continuous class function and
$\Pi_*\tilde\mu=\mu$.
\item For every $x$ and $F\in L^2(\mu)$, $\tilde{\mathcal P}_x(F\circ\Pi)=
(\mathcal P_xF)\circ\Pi$ and $\Var_{\tilde\mu}(F\circ\Pi)=\Var_\mu(F)$; if $F$
is gauge invariant, so is $F\circ\Pi$. Hence $\BP(n,\varpi)$ for $\tilde\mu$
implies $\BP(n,\varpi)$ for $\mu$.
\item If some $z\in Z$ acts on $\C^{\tilde N}$ as a scalar $\zeta\ne1$, then
$\tilde\mu(\Tr\tilde U_C)=0$ for every closed lattice path $C$ that traverses
some link exactly once.
\end{enumerate}
\end{proposition}

\begin{proof}
(i) $\pi(hgh^{-1})=\pi(h)\pi(g)\pi(h)^{-1}$, so $\tilde\varphi$ is a class
function. The image under $\pi$ of Haar measure on $\tilde G$ is a
left-invariant probability measure on $G$, hence Haar measure, so
$\Pi_*\tilde\lambda=\lambda$. As $\pi$ is a homomorphism, $\pi(\tilde U_p)=
(\Pi\tilde U)_p$, so $\tilde\mu$ has density $\Psi\circ\Pi$ with respect to
$\tilde\lambda$, with $\Psi(U)=\tilde Z^{-1}\prod_p\varphi(U_p)$, and
$\Pi_*((\Psi\circ\Pi)\tilde\lambda)=\Psi\lambda$; both being probability
measures, $\Psi\lambda=\mu$.

(ii) $\tilde{\mathcal P}_x(F\circ\Pi)(\tilde U)$ is the integral of
$F(\Pi\tilde\sigma,\Pi\tilde U_{\rm out})$ against the density proportional to
$\prod_{p\cap V_x\ne\emptyset}\varphi((\Pi\tilde\sigma,\Pi\tilde U_{\rm
out})_p)$ on $\tilde G^{V_x}$; by the change of variables $\sigma=\Pi\tilde
\sigma$ and (i) on $V_x$ it equals $(\mathcal P_xF)(\Pi\tilde U)$. The variance
identity follows from (i), and $\Pi(\tilde g\tilde U)=\pi(\tilde g)\Pi(\tilde
U)$ for a gauge transformation $\tilde g$ gives the invariance. Then for
gauge-invariant $F$, $\varpi\Var_\mu(F)=\varpi\Var_{\tilde\mu}(F\circ\Pi)\le
n^{-4}\sum_x\tilde\mu((F\circ\Pi-\tilde{\mathcal P}_x(F\circ\Pi))^2)=n^{-4}
\sum_x\mu((F-\mathcal P_xF)^2)$.

(iii) Let $e$ be a link traversed once by $C$ and $T$ the map $\tilde U_e
\mapsto z\tilde U_e$, the other links being fixed. It preserves $\tilde
\lambda$, and as $z$ is central with $\pi(z)=1$, each plaquette variable
through $e$ becomes $z^{\pm1}\tilde U_p$ with $\tilde\varphi(z^{\pm1}\tilde
U_p)=\tilde\varphi(\tilde U_p)$; so $\tilde\mu$ is $T$-invariant. Moving the
central factor out of the product, $\Tr\tilde U_C\circ T=\zeta^{\pm1}\Tr
\tilde U_C$, so $\tilde\mu(\Tr\tilde U_C)=\zeta^{\pm1}\tilde\mu(\Tr\tilde U_C)
=0$.
\end{proof}

For instance, with $\tilde G=G=\U(1)$, $\pi(z)=z^2$ and the Wilson weight
$\varphi$ of $\U(1)$, every Wilson loop of charge one under $\tilde\varphi(z)
=e^{\beta\operatorname{Re}z^2}$ that traverses some link once vanishes for
every $\beta$, so Wilson's area law holds with every string tension, while
by part (ii) the cube condition for $\tilde\varphi$ is at least as hard as for
the $\U(1)$ theory, in which it fails at weak coupling in the Coulomb phase
(Section~\ref{sec:weak}). The vanishing is kinematic, a consequence
of a local symmetry and not of confinement; what the example shows is that
an argument deriving (D) from an area law of the defining representation
cannot be valid for all compact groups and weights. For $\SU(N)$ with the
Wilson action no such example is known, and nothing is claimed about whether
an area law and (D) hold together there.

\section{Weak coupling: what a proof of the mass gap still needs}\label{sec:weak}

The continuum limit of four-dimensional Yang--Mills theory is taken along the
tuning dictated by asymptotic freedom: for $\SU(N)$ the bare coupling
$g^2=2N/\beta_W$ tends to zero like $1/(2b_0\log(1/(a\Lambda)))$ with
$b_0=11N/(48\pi^2)$, the first coefficient of the $\beta$-function
\cite{GW73,Pol73} and a fixed physical scale $\Lambda$, so $\beta_W\to\infty$
logarithmically in $1/a$. Theorems~\ref{thm:A} and~\ref{thm:bpstrong} hold for
$\beta_W<N^2/29$ and $\beta_W<N^2/42$ respectively; at fixed $\beta$ the lattice
gap in physical units, $m_G(\beta)/a$, diverges as $a\to0$, and no continuum
limit exists at fixed $\beta$. The results of this paper therefore say nothing
about the continuum theory without statement (D). Three facts show what (D)
asks for.

\emph{(a) A cube condition is a mass.} By Theorem~\ref{thm:free} the best
constant for a free field is at most $s^2/(s^2+16d/9)$ with $s=m\ell$. At weak
coupling the leading approximation of the Wilson measure, in any gauge, is a
massless Gaussian field, for which the constant vanishes as $a\to0$ at fixed
$\ell$. Hence (D) does not follow from perturbation theory to any finite order:
it holds only if the interaction generates a mass of order $1/\ell$, and it is
a quantitative form of that mass generation.

\emph{(b) The cube must grow.} For $\SU(N)$, $\beta\ge1$ and the torus of side
$2n+4$, $\BP(n,\varpi)$ implies $\varpi\le C_Nn^4(1+\log\beta)/\beta$
\cite{BBLRS26}: on a small torus at weak coupling the Polyakov loops order, and
the cube dynamics can change their centre charge only through the curvature,
which is small. Along the tuning $n_k=\ell/a_k$ grows like a power of $1/a_k$
while $\beta_k$ grows like $\log(1/a_k)$, so this obstruction is consistent with
(D) but excludes every argument at fixed cube size.

\emph{(c) The group must be non-abelian.} For $G=\U(1)$ with the Villain
action at large $\beta$ the theory has a massless photon \cite{FS82}, and the
cube condition fails at every cube size, since by Theorem~\ref{thm:decay} it
would make the field-strength correlations summable \cite{BBLRS26}. Every proof
of (D) must therefore use the non-abelian structure of $G$, exactly as a proof
of the mass gap must.

\emph{The two ends.} Ba{\l}aban's renormalisation group
\cite{Bal84a,Bal85a,Bal89b} controls the effective actions of the lattice
theory from the scale $a$ up to scales at which the running coupling is still
small, uniformly in the lattice spacing, in finite volume. The theorems of
Secs.~\ref{sec:curvature}--\ref{sec:bpstrong} control the theory when the
coupling is large: they need $N\beta$ small, that is, $g$ large. The running
coupling grows towards the infrared, so a proof of (D) would follow the
renormalisation group from the weak-coupling regime, where Ba{\l}aban's
estimates hold, through the regime of coupling of order one, to a scale
$\ell$ at which the effective theory satisfies a strong-coupling criterion of
the kind proved here, and would then transport the resulting cube estimate
back to the fine lattice. Neither the middle step nor the final transport is
available: the effective actions are not of Wilson form, Ba{\l}aban's analysis
stops before the coupling is of order one, and no known criterion accepts the
effective theory at that scale. This gap between the two ends is the Yang--Mills
mass gap problem in the form isolated here.

Statement (E) is the construction problem. Ba{\l}aban's ultraviolet stability
bounds give compactness of suitably smeared observables in finite volume, but
the infinite-volume limit, Euclidean invariance and nontriviality of the limit
are open \cite{JW06}. Given (E), Theorem~\ref{thm:B} reduces the mass gap to
(D) with the explicit constant $\Delta\ge16\varpi/(625e^2\ell)$.


\section{Conclusion}\label{sec:conclusion}

Lattice Yang--Mills theory in four dimensions has a mass gap with the explicit
lower bound $m_G(\beta)$ for every compact simple gauge group in the range
$N\beta<\frac2{29}\kappa_G$, and for $\SU(2)$ up to $\beta_W=\frac13$ with
the bound $-\log(3\beta_W)$ and up to $\beta_W=\frac9{25}$ by resampling stars. A
single
spectral gap $\varpi$ of the dynamics that resamples cubes of a fixed physical
size $\ell$, uniform along the continuum tuning, gives every nontrivial
continuum limit the mass gap $16\varpi/(625e^2\ell)$. That spectral gap holds
at strong coupling, for $\SU(2)$ up to $\beta_W=\frac7{20}$ uniformly in the
cube size and for $\SU(3)$ up to $\beta_W=0.45$, and for massive free fields;
it fails for massless fields, at fixed cube size at weak coupling, and in the
abelian Coulomb phase.

The cube gap follows from a decay of correlations, below the sharp threshold
$\frac12$, between the gauge-invariant observables on the two sides of a slab
of physical width $\ell/2$, uniformly in the boundary condition. This
hypothesis is a form of mass generation at one scale. It holds in two
dimensions along the weak-coupling tuning and fails there for strips one link
wide; in four dimensions it is not proved. Curvature arguments, with or
without gauge fixing, stop at explicit couplings of order one; spectral gaps
for blocks of fixed lattice size force a trivial continuum limit; and a
lattice gap tending to zero at the right rate produces a gap in a continuum
limit that is not identified with Yang--Mills theory. Statements (D) and (E)
remain open, and with them the Clay problem; what this paper contributes is an
explicit strong-coupling theory for every gauge group, a reduction of the
continuum gap to one estimate on one scale and of that estimate to a decay of
correlations across one slab, and a map of what the present methods can reach.


\appendix
\section{The line field}\label{app:table}

Table~\ref{tab:sigma} lists, for each residue class of $x$ modulo two, the
direction $\sigma(x)$ and the values $\pi_x(\mu)=\sigma(x+e_\mu)$. In every row
$\pi_x$ fixes $\sigma(x)$ and is a three-cycle on the other directions.

\begin{table}[h]
\caption{The line field: $x$ modulo two (digits $x_0x_1x_2x_3$), $\sigma(x)$, and
$\sigma(x+e_\mu)$ for $\mu=0,1,2,3$.}
\label{tab:sigma}
\begin{tabular}{cccccc|cccccc}
\toprule
$x$&$\sigma$&$0$&$1$&$2$&$3$&$x$&$\sigma$&$0$&$1$&$2$&$3$\\
\midrule
0000 & 0 & 0 & 2 & 3 & 1 & 0001 & 1 & 2 & 1 & 3 & 0 \\
1000 & 0 & 0 & 3 & 1 & 2 & 1001 & 2 & 1 & 3 & 2 & 0 \\
0100 & 2 & 3 & 0 & 2 & 1 & 0101 & 1 & 3 & 1 & 0 & 2 \\
1100 & 3 & 2 & 0 & 1 & 3 & 1101 & 3 & 1 & 2 & 0 & 3 \\
0010 & 3 & 1 & 2 & 0 & 3 & 0011 & 3 & 2 & 0 & 1 & 3 \\
1010 & 1 & 3 & 1 & 0 & 2 & 1011 & 2 & 3 & 0 & 2 & 1 \\
0110 & 2 & 1 & 3 & 2 & 0 & 0111 & 0 & 0 & 3 & 1 & 2 \\
1110 & 1 & 2 & 1 & 3 & 0 & 1111 & 0 & 0 & 2 & 3 & 1 \\
\bottomrule
\end{tabular}
\end{table}


\section{The blocks of Lemma~\ref{lem:local}}\label{app:blocks}

\emph{Diagonal blocks.} Let $w_a=e^{i\phi}$. On the diagonal part
$f_p=-\cos\phi\,(y_1+y_2-y_3-y_4)^2$ and $c_p=2(y_1y_4-y_1y_2+y_2y_3-y_3y_4)$, so
\[
\mathcal F=\sum_i\alpha_iy_i^2+\cos\phi\,(y_1+y_2-y_3-y_4)^2-\tfrac95(y_1y_4-y_1y_2
+y_2y_3-y_3y_4),
\]
with $y_k=0$, $\alpha_{k\pm1}=\frac{21}{10}$, $\alpha_{k+2}=\frac54$. It is affine in
$\cos\phi$, so positivity at $\cos\phi=\pm1$ suffices. For $k=1$ and $\cos\phi=1$
the matrix in $(y_2,y_3,y_4)$ is
\[
\begin{pmatrix}\frac{31}{10}&-\frac{19}{10}&-1\\[2pt]-\frac{19}{10}&\frac94&
\frac{19}{10}\\[2pt]-1&\frac{19}{10}&\frac{31}{10}\end{pmatrix},
\]
with leading principal minors $\frac{31}{10}$, $\frac{673}{200}$,
$\frac{8421}{2000}$. In general the leading principal minors are
$\frac{31}{10},\frac{673}{200},\frac{8421}{2000}$ (fixed link $e_1$ or $e_4$),
$\frac{31}{10},\frac{861}{100},\frac{8421}{2000}$ ($e_2$) and
$\frac94,\frac{673}{200},\frac{8421}{2000}$ ($e_3$) at $\cos\phi=1$, and
$\frac{11}{10},\frac{53}{200},\frac{21}{2000}$ ($e_1$, $e_4$),
$\frac{11}{10},\frac{21}{100},\frac{21}{2000}$ ($e_2$) and
$\frac14,\frac{53}{200},\frac{21}{2000}$ ($e_3$) at $\cos\phi=-1$; all are
positive.

\emph{Pair blocks.} Let $w_a=e^{i(\chi+\vartheta)}$ and $w_b=e^{i(\chi-\vartheta)}$. The
terms of $f_p$ carry $w_a+\bar w_b=2\cos\chi\,e^{i\vartheta}$ or its conjugate, and
$c_p$ carries $w_a\bar w_b=e^{2i\vartheta}$, so the Hermitian matrix $M(m,\vartheta)$ of
$\mathcal F$, $m=\cos\chi$, is affine in $m$ with $M(-1,\vartheta)=M(1,\vartheta+\pi)$,
and $m=1$ suffices. For the fixed link $e_2$, in the variables $(z_1,z_3,z_4)$
and with $c=\cos\vartheta$,
\[
M=\begin{pmatrix}\frac{21}5+2c&-2e^{i\vartheta}&-\frac95-2e^{i\vartheta}\\[2pt]
-2e^{-i\vartheta}&\frac{21}5+2c&\frac95+2e^{-i\vartheta}\\[2pt]
-\frac95-2e^{-i\vartheta}&\frac95+2e^{i\vartheta}&\frac52+2c\end{pmatrix}.
\]
With $u=\cos^2(\vartheta/2)\in[0,1]$ its leading principal minors are
\[
\tfrac{20u+11}5,\qquad\tfrac{(20u+1)(20u+21)}{25},\qquad
\tfrac{400u(40u^2-23u+4)+21}{250},
\]
and $40u^2-23u+4$ has negative discriminant $529-640$. The other positions give
the same determinant and the leading minors $\frac{20u+11}5$ or $\frac{8u+1}2$
and $\frac{800u^2-180u+53}{50}$, whose discriminant $32400-169600$ is negative.
Every block is positive definite.


\section{The single-link law of SU(2)}\label{app:sharp}

This appendix proves the sharp quarter, $W_1(\nu_b,\nu_{b'})\le\frac14|b-b'|$
for the one-link law of $\SU(2)$ in the chordal metric, which
the proof of Theorem~\ref{thm:A} for $\SU(2)$, Lemma~\ref{lem:onelink}(i) and the whole of
Sec.~\ref{sec:stars} use, together with the exact value $\frac14$ at $b=0$,
which shows that the single-link argument of Sec.~\ref{sec:state} stops at
$\beta_W=\frac13$. The argument is that of~\cite{BB26a}; it is given in full
so that the paper is self-contained. The computations it calls on are in
Appendix~\ref{app:vmfcomp}, and the facts about Bessel functions in
Appendix~\ref{app:bessel}.

\subsection{The von Mises--Fisher law}

Identify $\SU(2)$ with the unit quaternions, that is with the unit sphere
$S^3\subset\R^4$, so that $\ReTr U=2x_0(U)$ and $\ReTr(UV)=2\ip u{\bar v}$
for the Euclidean inner product of $\R^4$ and quaternion conjugation
$v\mapsto\bar v$. Fix a link $e$ and the field off $e$. Every plaquette
through $e$ can be rotated cyclically so that its trace reads $\ReTr(Q_eA_p)$
with $A_p\in\SU(2)$, the staple of $e$ in $p$. The action therefore depends on
$Q_e$ only through
\[
N\beta\sum_{p\ni e}\ReTr(Q_eA_p)=\ip{x}{b},\qquad
b=4\beta\sum_{p\ni e}\bar A_p\in\R^4,
\]
with $x=Q_e\in S^3$. Since there are six staples, each of norm one,
\begin{equation}\label{q-eq:kappamax}
\kappa:=|b|\le y:=24|\beta| .
\end{equation}
The conditional law of $Q_e$ is the von Mises--Fisher law\footnote{On the
sphere $S^{p-1}\subset\R^p$ the von Mises--Fisher law with mean direction
$\mu$ and concentration $\kappa$ has density proportional to
$e^{\kappa\ip x\mu}$ with respect to the surface measure, and its normalising
constant is $(2\pi)^{p/2}I_{p/2-1}(\kappa)/\kappa^{p/2-1}$; it is the maximum
entropy law with prescribed mean resultant, and it plays on the sphere the
role of the Gaussian in $\R^p$. For $p=4$ the Bessel index is $1$. The mean
resultant length $\E\ip x\mu=I_{p/2}(\kappa)/I_{p/2-1}(\kappa)$ is the
function $A$ of Lemma~\ref{q-lem:moments}. See Mardia and Jupp~\cite{MJ00},
Ch.~9, for these facts and for the tangent-normal decomposition used in
Appendix~\ref{app:vmfcomp}.}
\begin{equation}\label{q-eq:vmf}
\nu_b(dx)=Z(\kappa)^{-1}e^{\ip xb}\,\sigma(dx)
\end{equation}
on the unit sphere, with $\sigma$ the normalised surface measure. We use
coordinates in which $b=\kappa e_1$. The density of $t=x_1$ is proportional
to $(1-t^2)^{1/2}e^{\kappa t}$ on $[-1,1]$, so $Z(\kappa)=2I_1(\kappa)/\kappa$
with the modified Bessel function $I_1$, normalised so that $Z(0)=1$.

\begin{lemma}\label{q-lem:vmf}
For every $\kappa\ge0$ the law $\nu_b$ on the unit sphere $S^3$ satisfies
the Poincar\'e inequality $\Var_{\nu_b}(g)\le\frac13\int|\nabla g|^2d\nu_b$,
and the constant $\frac13$ is attained at $\kappa=0$.
\end{lemma}

\begin{proof}
In the polar angle $\theta\in(0,\pi)$ from $e_1$ the generator of the
reversible diffusion of $\nu_b$ is $L=\partial_\theta^2+b(\theta)
\partial_\theta+\sin^{-2}\theta\,\Delta_{S^2}$ with $b(\theta)=2\cot\theta-
\kappa\sin\theta$. On spherical harmonics of degree $\ell$ on $S^2$ it acts as
a Sturm--Liouville operator $L_\ell$ with potential
$\ell(\ell+1)\sin^{-2}\theta$. A ground state comparison of Barta's
type~\cite{Bar37} with the trial function $\psi=\sin\theta\,
e^{a\kappa\cos\theta}$ gives $-L_1\psi\ge3\psi$ for $a=-\frac15$, and the same
trial function with $a=-\frac25$ gives the bound $3$ for the derivatives of
zonal eigenfunctions; Appendix~\ref{app:vmfcomp} carries out both computations.
Hence every nonzero eigenvalue of $-L$ is at least $3$. At $\kappa=0$ the
first eigenvalue of the Laplacian of the unit $S^3$ is $3$.
\end{proof}

\subsection{Moments}

Let $A(\kappa)=\E_{\nu_b}x_1=Z'(\kappa)/Z(\kappa)=I_2(\kappa)/I_1(\kappa)$.

\begin{lemma}\label{q-lem:moments}
For every $\kappa>0$, in coordinates with $b=\kappa e_1$,
\begin{enumerate}[label=\textup{(\roman*)},leftmargin=*]
\item $\E x_2^2=\E x_3^2=\E x_4^2=A/\kappa$, $\E(1-x_1^2)=3A/\kappa$, the
covariance matrix of $x$ is $\operatorname{diag}(A',A/\kappa,A/\kappa,
A/\kappa)$, and $A'=1-A^2-3A/\kappa$;
\item $\displaystyle\frac{\kappa}{2+\sqrt{\kappa^2+4}}\le
A(\kappa)\le\frac{\kappa}{1+\sqrt{\kappa^2+9}}$;
\item $A'\le A/\kappa\le\bigl(1+\sqrt{\kappa^2+9}\bigr)^{-1}\le\frac14$;
\item $A(\kappa)\ge\kappa\big/\bigl(\frac32+\sqrt{\kappa^2+25/4}\bigr)$.
\end{enumerate}
\end{lemma}

\begin{proof}
(i) Integrating by parts against $\frac{d}{dt}(1-t^2)^{3/2}=-3t(1-t^2)^{1/2}$
gives $\kappa\,\E(1-x_1^2)=3\E x_1$, and the three coordinates
$x_2,x_3,x_4$ are exchangeable under $\nu_b$, with $\sum_ix_i^2=1$. The
off-diagonal covariances vanish by the reflections $x_i\mapsto-x_i$,
$i\ge2$, and $A'=\Var x_1$ because $A=(\log Z)'$.
(ii) By (i), $A$ solves the Riccati equation $A'=1-A^2-3A/\kappa$ with
$A(\kappa)=\kappa/4+O(\kappa^3)$. Appendix~\ref{app:vmfcomp} shows that the left
side of (ii) is a subsolution and the right side a supersolution of this
equation, with the same behaviour at $0$, and concludes by comparison.
(iii) The middle inequality is (ii). For the first, put $s=\sqrt{\kappa^2+4}$.
By (ii), $4A/\kappa\ge4/(2+s)$ and $1-A^2\le1-\kappa^2/(2+s)^2$, and
\[
\begin{aligned}
1-\frac{\kappa^2}{(2+s)^2}\le\frac4{2+s}
&\iff(2+s)^2-\kappa^2\le4(2+s)\\
&\iff s^2-\kappa^2\le4,
\end{aligned}
\]
which holds with equality. So $A'=1-A^2-3A/\kappa\le A/\kappa$.
(iv) is the comparison of (ii) with a third barrier, which improves the
lower bound of (ii); it is carried out in the same appendix.
\end{proof}

\subsection{The Lipschitz rate of the family}

Let $\delta$ be the chordal metric $\delta(x,x')=|x-x'|$ on $S^3$ and $W_1$
the Kantorovich distance for $\delta$. For $c\in\R^4$ and a unit vector
$h\in\R^4$ put
\begin{equation}\label{q-eq:rate}
L(c,h)=\sup_g\Cov_{\nu_c}\bigl(g,\ip hx\bigr),
\end{equation}
the supremum over $1$-Lipschitz functions $g$ for $\delta$. Since
$\frac{d}{dt}\nu_{c+th}(g)=\Cov_{\nu_{c+th}}(g,\ip hx)$, the
Kantorovich--Rubinstein duality gives, for all $b,b'$,
\begin{equation}\label{q-eq:w1segment}
W_1(\nu_b,\nu_{b'})\le|b-b'|\sup_{c\in[b,b']}L\Bigl(c,\frac{b'-b}{|b'-b|}
\Bigr).
\end{equation}
By rotation invariance $L(c,h)$ depends only on $\kappa=|c|$ and on the angle
$\varphi$ between $h$ and $c^\perp$, and we write $L(\kappa,\varphi)$.

\begin{lemma}\label{q-lem:rate}
For all $c$ and $h$, $L(c,h)\le1/(2\sqrt3)$. Moreover:
\begin{enumerate}[label=\textup{(\roman*)},leftmargin=*]
\item at $c=0$, $L(0,h)=\frac14$ for every $h$;
\item if $h\perp c$, then $L(c,h)=\E_{\nu_c}x_2^2=A(\kappa)/\kappa\le\frac14$;
\item if $\kappa\ge\sqrt{88}/3\approx3.127$, then $L(c,h)\le\frac14$ for
every $h$.
\end{enumerate}
\end{lemma}

\begin{proof}
A $1$-Lipschitz function for $\delta$ is $1$-Lipschitz for the geodesic
distance, which is larger, so $|\nabla g|\le1$ almost everywhere, and
Lemma~\ref{q-lem:vmf} gives $\Var g\le\frac13$. By Lemma~\ref{q-lem:moments}(i),
(iii) the largest eigenvalue of the covariance of $x$ is $A/\kappa\le\frac14$,
so $\Var\ip hx\le A/\kappa$. The Cauchy--Schwarz inequality gives
\begin{equation}\label{q-eq:csbound}
L(c,h)\le\bigl(\tfrac13A/\kappa\bigr)^{1/2}\le\Bigl(3\bigl(1+
\sqrt{\kappa^2+9}\bigr)\Bigr)^{-1/2}\le\frac1{2\sqrt3}.
\end{equation}
(i) and (ii) are reflection arguments. Let $R$ be the reflection in the
hyperplane $h^\perp$, and suppose $\nu_c$ is $R$-invariant, which is the case
when $c=0$ or $h\perp c$. Then $\E\ip hx=0$ and
\[
\begin{aligned}
\Cov(g,\ip hx)&=\E\bigl[(g(x)-g(Rx))\ip hx;\ \ip hx>0\bigr]\\
&\le\E\bigl[2\ip hx^2;\ \ip hx>0\bigr]=\E\ip hx^2,
\end{aligned}
\]
since $\delta(x,Rx)=2|\ip hx|$. Equality holds for $g=\ip hx$, which is
$1$-Lipschitz. At $c=0$, $\E\ip hx^2=\frac14$; for $h\perp c$, $\E\ip
hx^2=A/\kappa$. (iii) By~\eqref{q-eq:csbound}, $L\le\frac14$ as soon as
$1+\sqrt{\kappa^2+9}\ge\frac{16}3$, that is $\kappa^2\ge\frac{88}9$.
\end{proof}

The reflection settles the directions
$h\perp c$ and the point $c=0$; for an oblique $h$ the reflection no longer
preserves $\nu_c$, and the Cauchy--Schwarz bound~\eqref{q-eq:csbound} is too
weak for $\kappa<\sqrt{88}/3$. The next theorem closes this case.

\begin{theorem}[The sharp quarter]\label{q-thm:quarter}
For every $c\in\R^4$ and every unit vector $h\in\R^4$,
\[
L(c,h)\le\tfrac14 ,
\]
with equality at $c=0$. Equivalently, $W_1(\nu_b,\nu_{b'})\le\frac14|b-b'|$
for all $b,b'\in\R^4$, and the constant $\frac14$ is the best possible.
\end{theorem}

The equivalence is~\eqref{q-eq:w1segment} in one direction and the derivative
formula above it in the other; the sharpness is Lemma~\ref{q-lem:rate}(i). The
metric matters.\footnote{In the geodesic metric of $S^3$ the statement is
false already at $c=0$: the function $\arcsin x_2$ is geodesically
$1$-Lipschitz and its covariance with $x_2$ under the uniform law is
$\frac8{9\pi}\approx0.283>\frac14$. A proof therefore has to use that the
chordal distance between two points is shorter than their geodesic distance,
that is, it has to see the sphere from the ball it bounds. The proof below
does exactly this: it extends a chordally Lipschitz function to the ball and
integrates by parts there.} The proof has four steps. A Lipschitz function
on $S^3$ extends to the unit ball $B\subset\R^4$, and Green's formula in $B$
bounds the covariance by a negative Sobolev norm of the density
(Lemma~\ref{q-lem:harmonic}). This norm splits into a zonal and a transverse
part, both explicit series in modified Bessel functions
(Lemma~\ref{q-lem:spectral}). The transverse part is at most $\frac14$ because a
certain power series has nonnegative coefficients
(Proposition~\ref{q-prop:transverse}), and the zonal part is at most $\frac14$
by a Riccati comparison (Proposition~\ref{q-prop:zonal}).

\subsection{Harmonic extension}

Let $\sigma$ be the normalised surface measure of $S^3$ and $H_k$ the space
of spherical harmonics of degree $k$, the restrictions to $S^3$ of harmonic
homogeneous polynomials of degree $k$ on $\R^4$. The spaces $H_k$ are
orthogonal in $L^2(\sigma)$ and span it; we write $u_k$ for the component of
$u$ in $H_k$, and for $u$ with $\int u\,d\sigma=0$ we put
\begin{equation}\label{q-eq:Qnorm}
\mathcal Q(u)=\sum_{k\ge1}\frac1k\,\|u_k\|_{L^2(\sigma)}^2 .
\end{equation}
This is $\ip u{\Lambda^{-1}u}$ for the Dirichlet-to-Neumann operator
$\Lambda$ of $B$, which acts on $H_k$ as multiplication by $k$.

\begin{lemma}\label{q-lem:harmonic}
Let $u$ be a real-analytic function on $S^3$ with $\int u\,d\sigma=0$ and let
$g$ be $1$-Lipschitz for the chordal metric. Then
\[
\int_{S^3}g\,u\,d\sigma\le\tfrac12\,\mathcal Q(u)^{1/2}.
\]
\end{lemma}

\begin{proof}
The chordal metric is the Euclidean metric of $\R^4$ restricted to $S^3$,
so McShane's extension $G(y)=\inf_{x\in S^3}\bigl(g(x)+|x-y|\bigr)$
is $1$-Lipschitz on $\R^4$ and equals $g$ on $S^3$~\cite{McS34}; by
Rademacher's theorem $|\nabla G|\le1$ almost everywhere. Let
$\psi=\sum_{k\ge1}k^{-1}r^ku_k(\omega)$ in polar coordinates $y=r\omega$. Each
term is a harmonic polynomial, the series converges with all derivatives on
the closed ball because $u$ is real-analytic,\footnote{For real-analytic
$u$ the norms $\|u_k\|$ decay geometrically in $k$, and the supremum of a
function in $H_k$ is at most $(\dim H_k)^{1/2}=k+1$ times its $L^2(\sigma)$
norm, so the series for $\psi$ and for each of its derivatives converge
uniformly on $\bar B$.} and $\partial_r\psi=u$ on $S^3$. Green's formula for
the Lipschitz function $G$ and the smooth harmonic function $\psi$ reads
$\int_{S^3}G\,\partial_r\psi\,dS=\int_B\nabla G\cdot\nabla\psi\,dy$. With
$|S^3|=2\pi^2$ and $|B|=\pi^2/2$ the Cauchy--Schwarz inequality gives
\[
\int g\,u\,d\sigma=\frac1{2\pi^2}\int_B\nabla G\cdot\nabla\psi\,dy\le
\frac{|B|^{1/2}}{2\pi^2}\Bigl(\int_B|\nabla\psi|^2dy\Bigr)^{1/2}.
\]
Green's formula once more gives $\int_B|\nabla\psi|^2=\int_{S^3}\psi\,
\partial_r\psi\,dS=2\pi^2\sum_kk^{-1}\|u_k\|^2=2\pi^2\mathcal Q(u)$, and the
right side becomes $\frac12\mathcal Q(u)^{1/2}$.
\end{proof}

The lemma is a form of Thomson's principle. By Kantorovich--Rubinstein
duality, $\sup_g\int gu\,d\sigma$ is the least cost of a flow in $\R^4$ that
carries the positive part of $u\,d\sigma$ to its negative part, the cost
being the total length travelled by the mass.\footnote{This is the
Beckmann formulation of the transport problem with cost $|x-y|$: the
optimal cost equals the least $\int|V|$ over vector measures $V$ on $\R^4$
with divergence $u\,d\sigma$. Restricting the flow to the closed ball does
not change the infimum, since projecting onto $\bar B$ shortens paths. The
gradient $\nabla\psi$ on $B$ is one admissible flow. It minimises the energy
$\int|V|^2$ rather than the cost $\int|V|$, and the Cauchy--Schwarz inequality
converts one into the other at the price of the factor $|B|^{1/2}$.} The
harmonic field $\nabla\psi$ is one such flow, through the interior of the
ball, and the lemma computes its cost. At $c=0$ the flow is optimal: for
$u=\ip hx$ one has $\psi=\ip hy$, the field $\nabla\psi=h$ is constant,
the Cauchy--Schwarz inequality is an equality, and $\mathcal Q(u)=\|u\|^2=
\frac14$.

\subsection{Spectral decomposition}

Fix $c=\kappa e_1$ with $\kappa>0$ and write $h=\cos\varphi\,e_1+\sin\varphi
\,e_2$, which is no loss of generality by rotation invariance. Let
$\rho=d\nu_c/d\sigma=e^{\kappa x_1}/Z(\kappa)$, $Z(\kappa)=\int e^{\kappa x_1}
d\sigma=2I_1(\kappa)/\kappa$, and
\[
u=\bigl(\ip hx-\E_{\nu_c}\ip hx\bigr)\rho ,
\]
so that $\Cov_{\nu_c}(g,\ip hx)=\int gu\,d\sigma$ and $\int u\,d\sigma=0$.
Let $U_n$ be the Chebyshev polynomials of the second kind and put
\begin{equation}\label{q-eq:dn}
d_n(\kappa)=\frac{(n+1)\,I_{n+1}(\kappa)}{I_1(\kappa)},\qquad n\ge0 .
\end{equation}

\begin{lemma}\label{q-lem:spectral}
With this notation
\[
\mathcal Q(u)=\cos^2\varphi\,\mathcal Q_z(\kappa)+\sin^2\varphi\,
\mathcal Q_x(\kappa),
\]
where
\[
\mathcal Q_z=\sum_{n\ge1}\frac{d_n'(\kappa)^2}n,\qquad
\mathcal Q_x=\frac1{3\kappa^2}\sum_{n\ge1}(n+2)\,d_n(\kappa)^2 .
\]
\end{lemma}

\begin{proof}
The functions $U_n(x_1)$, $n\ge0$, are the zonal harmonics: $U_n(x_1)\in
H_n$, and they are orthonormal in $L^2(\sigma)$ (Appendix~\ref{app:bessel}).
The expansion $e^{\kappa t}=\frac2\kappa\sum_n(n+1)I_{n+1}(\kappa)U_n(t)$ of
the same appendix gives $\rho=\sum_nd_n(\kappa)U_n(x_1)$, with $d_0=1$.

Write $u=\cos\varphi\,u_z+\sin\varphi\,u_x$ with $u_z=(x_1-A)\rho$ and
$u_x=x_2\rho$; here $\E_{\nu_c}x_1=A$ and $\E_{\nu_c}x_2=0$. The reflection
$x_2\mapsto-x_2$ maps each $H_k$ to itself, fixes $u_z$ and reverses $u_x$,
so the components of $u_z$ and $u_x$ in each $H_k$ are orthogonal and the
cross term in $\mathcal Q(u)$ vanishes.

Since $A=Z'/Z$, we have $u_z=\partial_\kappa\rho=\sum_nd_n'(\kappa)U_n(x_1)$,
and the component of $u_z$ in $H_n$ is $d_n'U_n(x_1)$, of norm $|d_n'|$. This
gives $\mathcal Q_z$.

For $u_x$, let $y(t)=\cos t\,e_1+\sin t\,e_2$. The law $\nu_{\kappa y(t)}$ has
density $\rho_t=e^{\kappa\ip{y(t)}x}/Z(\kappa)=\sum_nd_n(\kappa)U_n(\ip{y(t)}x)$,
and $\partial_t\rho_t|_{t=0}=\kappa x_2\rho$. The function
$U_n(\ip{y(t)}x)$ lies in $H_n$ for every $t$, hence so does its derivative
$x_2U_n'(x_1)$ at $t=0$, and the component of $u_x$ in $H_n$ is
$\kappa^{-1}d_nx_2U_n'(x_1)$. Its norm follows from the eigenvalue
$n(n+2)$ of the Laplacian of $S^3$ on $H_n$. For $f=U_n(x_1)$,
\[
\begin{aligned}
n(n+2)&=n(n+2)\|f\|^2=\int|\nabla_{S^3}f|^2d\sigma\\
&=\int U_n'(x_1)^2(1-x_1^2)\,d\sigma ,
\end{aligned}
\]
since the tangential gradient of $x_1$ has squared length $1-x_1^2$. Given
$x_1$, the point $(x_2,x_3,x_4)$ is uniform on the sphere of radius
$(1-x_1^2)^{1/2}$, so $\E[x_2^2\mid x_1]=\frac13(1-x_1^2)$ and
$\|x_2U_n'(x_1)\|^2=\frac13n(n+2)$. Summing $n^{-1}\kappa^{-2}d_n^2\cdot\frac13
n(n+2)$ over $n\ge1$ gives $\mathcal Q_x$.
\end{proof}

By Lemmas~\ref{q-lem:harmonic} and~\ref{q-lem:spectral},
\begin{equation}\label{q-eq:Lbyq}
L(\kappa,\varphi)\le\tfrac12\max\bigl(\mathcal Q_z(\kappa),\mathcal Q_x(
\kappa)\bigr)^{1/2},
\end{equation}
and Theorem~\ref{q-thm:quarter} follows once both parts are at most
$\frac14$. Both are $\frac14$ at $\kappa=0$, by continuity from the values
$d_1\sim\kappa/2$, $d_1'\to\frac12$ and $d_n=O(\kappa^n)$.

\subsection{Three Bessel sums}

The series of Lemma~\ref{q-lem:spectral} reduce to closed forms. Write $I_j$
for $I_j(\kappa)$.

\begin{lemma}\label{q-lem:besselsums}
For every $\kappa>0$,
\begin{align*}
\sum_{j\ge1}j^2I_j^2&=\frac\kappa4\,I_1(2\kappa),\qquad
\sum_{j\ge1}j^3I_j^2=\frac{\kappa^2}4\bigl(I_0^2+I_1^2\bigr),\\
I_1(2\kappa)&=2\sum_{j\ge1}I_{j-1}I_j .
\end{align*}
\end{lemma}

The first and third identities follow from Parseval's formula for
$e^{\kappa\cos\theta}=\sum_{j\in\Z}I_je^{ij\theta}$, and the second from the
recurrences of $I_j$ by telescoping; Appendix~\ref{app:bessel} gives the
proofs.

\subsection{The transverse part}

\begin{proposition}\label{q-prop:transverse}
$\mathcal Q_x(\kappa)\le\frac14$ for every $\kappa\ge0$.
\end{proposition}

\begin{proof}
By~\eqref{q-eq:dn}, $\mathcal Q_x=(3\kappa^2I_1^2)^{-1}\sum_{j\ge2}j^2(j+1)
I_j^2$, and the sum equals $\sum_{j\ge1}(j^3+j^2)I_j^2-2I_1^2$. By
Lemma~\ref{q-lem:besselsums}, $\mathcal Q_x\le\frac14$ is equivalent to
\begin{equation}\label{q-eq:Dfunction}
\mathcal D(\kappa):=(8+2\kappa^2)I_1^2-\kappa^2I_0^2-\kappa I_1(2\kappa)\ge0 .
\end{equation}
Each term is an entire function of $\kappa$ that is even after the
factor $\kappa^2$ is removed, so $\mathcal D(\kappa)=\sum_{n\ge0}\mathcal D_n
\kappa^{2n}$. Lemma~\ref{q-lem:Dcoef} of Appendix~\ref{app:bessel} computes the
coefficients: $\mathcal D_0=\mathcal D_1=0$, and for $n=m+1\ge2$
\[
\mathcal D_n=\frac{N_m}{4^m\,m!\,(m+1)!\,(m+1)(m+2)},
\]
with
\[
\begin{aligned}
N_m={}&\binom{2m}m\bigl(4(2m+1)+(m-1)(m+1)(m+2)\bigr)\\
&-4^m(m+1)(m+2).
\end{aligned}
\]
The numerator $N_m$ vanishes for $m=0,1,2$ and is positive for $m\ge3$
(Lemma~\ref{q-lem:Dcoef}). So every coefficient of $\mathcal D$ is nonnegative
and $\mathcal D\ge0$.
\end{proof}

The first nonzero coefficient is $\mathcal D_4=\frac1{2304}$, so
$\frac14-\mathcal Q_x=\mathcal D/(12\kappa^2I_1^2)$ is of order $\kappa^4$ at $0$: the transverse part
leaves the value $\frac14$ very slowly, which is why no estimate that loses a
constant can prove the proposition.

\subsection{The zonal part}

Let $A=I_2/I_1$ as before, so that $d_1=2A$ and, by
Lemma~\ref{q-lem:moments}(i), $A'=1-A^2-3A/\kappa$.

\begin{proposition}\label{q-prop:zonal}
For $0<\kappa\le\sqrt{88}/3$,
\[
\mathcal Q_z(\kappa)\le\frac14+4A'^2+\frac{A^2}4-\frac A{2\kappa}-
\frac{2A^2}{\kappa^2}<\frac14 .
\]
\end{proposition}

\begin{proof}
\emph{The higher modes.} With $r_\nu=I_{\nu+1}/I_\nu$ the recurrence
$I_\nu'=I_{\nu+1}+\frac\nu\kappa I_\nu$ gives
\[
\frac{d_n'}{d_n}=\frac{I_{n+1}'}{I_{n+1}}-\frac{I_1'}{I_1}=\frac n\kappa+
r_{n+1}-r_1 .
\]
By Lemma~\ref{q-lem:ratio} of Appendix~\ref{app:bessel}, $r_{n+1}\le r_1$, and
$r_1=A\le\kappa/4$ by Lemma~\ref{q-lem:moments}(iii). Hence $-n/\kappa\le
n/\kappa-\kappa/4\le d_n'/d_n\le n/\kappa$ for $n\ge2$ and $\kappa^2\le16$, so
$d_n'^2\le n^2d_n^2/\kappa^2$. With $d_1'=2A'$,
\[
\mathcal Q_z\le4A'^2+\frac1{\kappa^2}\sum_{n\ge2}n\,d_n^2=4A'^2+\frac{S_3-
S_2-4I_2^2}{\kappa^2I_1^2},
\]
where $S_2=\sum_{j\ge1}j^2I_j^2$ and $S_3=\sum_{j\ge1}j^3I_j^2$; the
second form uses $n\,d_n^2I_1^2=(j-1)j^2I_j^2$ with $j=n+1$.

\emph{The closed form.} The recurrences $I_0-I_2=\frac2\kappa I_1$ and
$I_1-I_3=\frac4\kappa I_2$ give $I_0/I_1=2/\kappa+A$ and $I_2I_3/I_1^2=A-
4A^2/\kappa$. By Lemma~\ref{q-lem:besselsums},
$S_3/(\kappa^2I_1^2)=\frac14(2/\kappa+A)^2+\frac14$, and, keeping three
terms of the third identity,
\[
\frac{S_2}{\kappa^2I_1^2}\ge\frac{I_0I_1+I_1I_2+I_2I_3}{2\kappa I_1^2}=
\frac1{2\kappa}\Bigl(\frac2\kappa+3A-\frac{4A^2}\kappa\Bigr).
\]
Collecting terms, $\mathcal Q_z\le\frac14+4A'^2+T(A)$ with $T(a)=\frac{a^2}4-
\frac a{2\kappa}-\frac{2a^2}{\kappa^2}$, which is the first inequality.

\emph{The barrier.} Let $s=\sqrt{\kappa^2+25/4}$ and $A_\flat=\kappa/(\frac32
+s)$. By Lemma~\ref{q-lem:moments}(iv), $A_\flat\le A\le\kappa/4$. The function
$a\mapsto1-a^2-3a/\kappa$ decreases on $a\ge0$, and $A'\ge0$ because $A'=
\Var x_1$, so $0\le A'\le\alpha:=1-A_\flat^2-3A_\flat/\kappa$. For $\kappa^2\le
12$ the quadratic $T$ decreases on $[0,\kappa/4]$, since
\[
T(a)-T(a')=(a'-a)\,\frac{2\kappa-(a+a')(\kappa^2-8)}{4\kappa^2}
\]
and $(a+a')(\kappa^2-8)\le\frac\kappa2(\kappa^2-8)\le2\kappa$ there. So
$\mathcal Q_z\le\frac14+4\alpha^2+T(A_\flat)$. With $\kappa^2=s^2-\frac{25}4$
and $A_\flat/\kappa=2/(2s+3)$ both terms are rational in $s$:
$\alpha=16/(2s+3)^2$, and
\begin{equation}\label{q-eq:zonalfactor}
\begin{aligned}
4\alpha^2+T(A_\flat)&=\frac{(2s-5)\,p(s)}{4(2s+3)^4},\\
p(s)&=8s^3+28s^2-98s-695 .
\end{aligned}
\end{equation}
For $0<\kappa\le\sqrt{88}/3$ we have $\frac52<s\le\sqrt{577}/6<4.004$. On
$s\ge\frac52$ the cubic $p$ increases, since $p'(s)=24s^2+56s-98>0$, and
$p(4.004)<-120$. So the right side of~\eqref{q-eq:zonalfactor} is negative.
\end{proof}

\subsection{Proof of Theorem~\ref{q-thm:quarter}}

At $c=0$ the theorem is Lemma~\ref{q-lem:rate}(i), and for $|c|\ge\sqrt{88}/3$
it is Lemma~\ref{q-lem:rate}(iii). For $0<|c|<\sqrt{88}/3$,
Propositions~\ref{q-prop:transverse} and~\ref{q-prop:zonal} give $\mathcal Q_x,
\mathcal Q_z\le\frac14$, and~\eqref{q-eq:Lbyq} gives $L\le\frac14$. \qed


\section{Computations for the single-link law}\label{app:vmfcomp}

\subsection{Barta's inequality}

Let $\rho(\theta)=\sin^2\theta\,e^{\kappa\cos\theta}$ on $(0,\pi)$, let $V$
be a continuous potential, and let
\[
T=-\rho^{-1}\partial_\theta\,\rho\,\partial_\theta+V .
\] If $\psi>0$ is smooth on $(0,\pi)$ with
$\rho\psi^2=O(\sin^4\theta)$ at the end points and $T\psi\ge\lambda_0\psi$,
then
\[
\int(|h'|^2+Vh^2)\rho\,d\theta\ge\lambda_0\int h^2\rho\,d\theta
\]
for every smooth $h$ vanishing to first order at $0$ and $\pi$. Indeed, writing
$h=\psi u$ and integrating by parts, the boundary terms $\rho\psi\psi'u^2$
vanish and
\[
\int(|h'|^2+Vh^2)\rho=\int\rho\psi^2|u'|^2+\int\frac{T\psi}\psi\,h^2\rho
\ \ge\ \lambda_0\int h^2\rho .
\]
This is the inequality of Barta~\cite{Bar37} in the form of a ground state
transformation.

\subsection{The computations of Lemma~\ref{q-lem:vmf}}

The generator of the reversible diffusion of $\nu_b$ on $S^3$ is
$L=\Delta_{S^3}+\kappa\nabla x_1\cdot\nabla$. In the polar angle $\theta$
from $e_1$ the round metric is $d\theta^2+\sin^2\theta\,g_{S^2}$, so
$\Delta_{S^3}=\partial_\theta^2+2\cot\theta\,\partial_\theta+\sin^{-2}\theta\,
\Delta_{S^2}$, and $\nabla x_1\cdot\nabla=-\sin\theta\,\partial_\theta$. On
$h(\theta)Y_\ell$ with a spherical harmonic $Y_\ell$ of degree $\ell$ on $S^2$
the operator acts as $L_\ell=\partial_\theta^2+b\,\partial_\theta-\ell(\ell+1)
\sin^{-2}\theta$, $b=2\cot\theta-\kappa\sin\theta=\rho'/\rho$, so
$-L_\ell=T$ with $V=\ell(\ell+1)\sin^{-2}\theta$.

\emph{Degree $\ell\ge1$.} Take $\psi=\sin\theta\,e^{a\kappa\cos\theta}$. Then
$\psi'/\psi=\cot\theta-a\kappa\sin\theta$ and
$\psi''/\psi=(\psi'/\psi)^2+(\psi'/\psi)'$, and a short computation gives
\[
\frac{-L_1\psi}\psi=3+(1+5a)\kappa\cos\theta-a(1+a)\kappa^2\sin^2\theta .
\]
For $a=-\frac15$ this is $3+\frac4{25}\kappa^2\sin^2\theta\ge3$. Since
$V=\ell(\ell+1)\sin^{-2}\theta\ge2\sin^{-2}\theta$ for $\ell\ge1$, Barta's
inequality gives $-L_\ell\ge3$ on every sector with $\ell\ge1$.

\emph{Degree $\ell=0$.} Let $g(\theta)$ be a zonal eigenfunction,
$-L_0g=\lambda g$, $g$ not constant, so that $g'(0)=g'(\pi)=0$. Differentiating
$g''+bg'=-\lambda g$ gives, for $h=g'$,
\[
-\bigl(h''+bh'\bigr)-b'h=\lambda h,\qquad -b'=\frac2{\sin^2\theta}+\kappa
\cos\theta,
\]
so $h$ is an eigenfunction of $T$ with $V=2\sin^{-2}\theta+\kappa\cos\theta$,
with eigenvalue $\lambda$, and $h\not\equiv0$ vanishes at both ends. With the
same $\psi$,
\[
\frac{T\psi}\psi=3+(2+5a)\kappa\cos\theta-a(1+a)\kappa^2\sin^2\theta,
\]
which for $a=-\frac25$ is $3+\frac6{25}\kappa^2\sin^2\theta\ge3$. Barta's
inequality applied to $h$ gives $\lambda\ge3$. Together with the case
$\ell\ge1$ every nonzero eigenvalue of $-L$ is at least $3$, which is
Lemma~\ref{q-lem:vmf}; the spectrum is discrete because $S^3$ is compact and
the drift is smooth.

\subsection{Moments}

The density of $t=x_1$ under $\nu_b$ is proportional to $(1-t^2)^{1/2}e^{
\kappa t}$. For every smooth $\phi$ on $[-1,1]$,
\[
\frac d{dt}\Bigl[(1-t^2)^{3/2}e^{\kappa t}\Bigr]=(1-t^2)^{1/2}e^{\kappa t}
\bigl(\kappa(1-t^2)-3t\bigr),
\]
and integration by parts gives $\E[\phi'(t)(1-t^2)]=\E\bigl[\phi(t)\bigl(3t-
\kappa(1-t^2)\bigr)\bigr]$. With $\phi=1$ this is $\kappa\E(1-x_1^2)=3\E x_1$,
the identity of Lemma~\ref{q-lem:moments}(i). With $\phi=t^k$ it gives the
recursion
\[
m_{k+2}=m_k-\frac{k+3}\kappa\,m_{k+1}+\frac k\kappa\,m_{k-1},\qquad m_k=\E
x_1^k,
\]
from $m_0=1$ and $m_1=A$; all polynomial moments of $\nu_b$ are rational in
$\kappa$ and $A(\kappa)$.\footnote{For instance $m_2=1-3A/\kappa$,
$m_3=A-3/\kappa+12A/\kappa^2$ and $m_4=m_2-5m_3/\kappa+2A/\kappa$. The joint moments
with $x_2,x_3,x_4$ follow from exchangeability:
$\E[\phi(x_1)x_2^2]=\frac13\E[\phi(x_1)(1-x_1^2)]$.}

\subsection{The Riccati barriers}

By Lemma~\ref{q-lem:moments}(i), $A$ solves $A'=\mathcal F(\kappa,A):=1-A^2-
3A/\kappa$ on $(0,\infty)$, and $A(\kappa)=\frac\kappa4-\frac{\kappa^3}{96}+
O(\kappa^5)$ from the series of $I_1$ and $I_2$. Let $s=\sqrt{\kappa^2+4}$,
$u=\sqrt{\kappa^2+9}$, $\phi_-=\kappa/(2+s)$ and $\phi_+=\kappa/(1+u)$. Direct
differentiation gives
\[
\begin{aligned}
\phi_-'-\mathcal F(\kappa,\phi_-)&=-\frac{\kappa^2}{s(s+2)^2}\le0,\\
\phi_+'-\mathcal F(\kappa,\phi_+)&=\frac{(u-3)^2}{u(1+u)^2}\ge0,
\end{aligned}
\]
and $\phi_\pm=\frac\kappa4+O(\kappa^3)$.\footnote{For the first identity write
$\phi_-=(s-2)/\kappa$, which follows from $s^2-4=\kappa^2$; then
$\phi_-'=(\kappa^2/s-s+2)/\kappa^2$ and $1-\phi_-^2-3\phi_-/\kappa=
(\kappa^2-(s-2)^2-3(s-2))/\kappa^2$, and the difference simplifies with
$s^2=\kappa^2+4$. The second is the same computation with
$\phi_+=\kappa(u-1)/(u^2-1)=\kappa(u-1)/(\kappa^2+8)$ in place of $\phi_-$
and $u^2=\kappa^2+9$ in place of $s^2=\kappa^2+4$. After clearing the denominator $\kappa^2(1+u)^2$ the
numerator becomes $\kappa^2(u-3)^2/u$, which is the stated residual.} The comparison is singular at $0$
because of the term $3A/\kappa$, and we carry it out with an integrating
factor. Let $D=A-\phi_-$. Then
\[
\begin{aligned}
D'&=\mathcal F(\kappa,A)-\mathcal F(\kappa,\phi_-)-\bigl(\phi_-'-\mathcal F(
\kappa,\phi_-)\bigr)\\
&=-\Bigl(A+\phi_-+\frac3\kappa\Bigr)D+R,
\end{aligned}
\]
with $R=\kappa^2/(s(s+2)^2)\ge0$. Put $\mu(\kappa)=\kappa^3\exp\int_0^\kappa
(A+\phi_-)$. Then $(\mu D)'=\mu R\ge0$, and $\mu D=O(\kappa^6)\to0$ as
$\kappa\to0$, so $\mu D\ge0$ and $A\ge\phi_-$. The same argument with
$D=\phi_+-A$ and $R=(u-3)^2/(u(1+u)^2)$ gives $A\le\phi_+$. This is
Lemma~\ref{q-lem:moments}(ii).

For (iv) let $s=\sqrt{\kappa^2+25/4}$ and $\phi_\flat=\kappa/(\frac32+s)$.
Since $s^2-\frac94=\kappa^2+4$, the same differentiation gives
\[
\phi_\flat'-\mathcal F(\kappa,\phi_\flat)=-\frac{5\,(s-\frac52)}{2\,s\,
(\frac32+s)^2}\le0,
\]
so $\phi_\flat$ is a subsolution, and $\phi_\flat=\frac\kappa4-\frac{\kappa^3}
{80}+O(\kappa^5)$.\footnote{With $\phi_\flat=\kappa(s-\frac32)/(\kappa^2+4)$ one
has $\phi_\flat'=(\frac32 s+\frac{25}4)/(s(\frac32+s)^2)$, using $ds/d\kappa=
\kappa/s$; and $\mathcal F(\kappa,\phi_\flat)=1-\phi_\flat^2-3/(\frac32+s)$.
Clearing the denominator $s(\frac32+s)^2$ and replacing $\kappa^2$ by
$s^2-\frac{25}4$ leaves $-\frac52 s+\frac{25}4$ in the numerator, which is the
stated residual. The expansion follows from $s=\frac52+\frac{\kappa^2}5+
O(\kappa^4)$.} Hence $A-\phi_\flat=\frac{\kappa^3}{480}+O(\kappa^5)$, and the
integrating-factor argument above, with $D=A-\phi_\flat$, $R=\frac{5(s-5/2)}
{2s(3/2+s)^2}\ge0$ and $\mu=\kappa^3\exp\int_0^\kappa(A+\phi_\flat)$, gives
$\mu D=O(\kappa^6)$ at $0$, $(\mu D)'=\mu R\ge0$, and so $A\ge\phi_\flat$. Since
$\frac32+\sqrt{\kappa^2+25/4}<2+\sqrt{\kappa^2+4}$ for $\kappa>0$, this
improves the lower bound of (ii).


\section{Bessel functions}\label{app:bessel}

This appendix proves the auxiliary statements used in Appendix~\ref{app:sharp}
for Theorem~\ref{q-thm:quarter}. Throughout, $I_j=I_j(\kappa)$ with $\kappa>0$,
and $I_{-j}=I_j$ for integers $j$. We use the series
\[
I_\nu(\kappa)=\sum_{i\ge0}\frac{(\kappa/2)^{2i+\nu}}{i!\,(i+\nu)!},
\]
the integral $I_j(\kappa)=\pi^{-1}\int_0^\pi e^{\kappa\cos\theta}\cos
j\theta\,d\theta$, the Fourier series $e^{\kappa\cos\theta}=\sum_{j\in\Z}
I_je^{ij\theta}$ and the recurrences
\begin{equation}\label{q-eq:besselrec}
I_{j-1}-I_{j+1}=\frac{2j}\kappa\,I_j,\quad I_j'=I_{j+1}+\frac j\kappa\,I_j ,
\end{equation}
all of which are in Watson's treatise~\cite{Wat44}, \S\S3.7, 3.71,
2.22. All series below converge absolutely, because $0<I_j\le(\kappa/2)^j
e^\kappa/j!$.

\subsection{Zonal harmonics}

Under $\sigma$ the coordinate $x_1$ has the density $\frac2\pi(1-t^2)^{1/2}$
on $[-1,1]$, the area of the two-sphere $\{x_1=t\}$ being proportional to
$1-t^2$ and its normal speed to $(1-t^2)^{-1/2}$. The Chebyshev polynomials
$U_n(\cos\theta)=\sin((n+1)\theta)/\sin\theta$ satisfy
\[
\int U_nU_m\,d\sigma=\frac2\pi\int_0^\pi\sin((n+1)\theta)\sin((m+1)\theta)
\,d\theta=\delta_{nm},
\]
and $U_n=C_n^{(1)}$ is the Gegenbauer polynomial of index $\frac{4-2}2=1$, so
$U_n(\ip yx)$ is the zonal harmonic of degree $n$ on $S^3$ with pole $y$:
$|x|^nU_n(\ip yx/|x|)$ is a harmonic polynomial on $\R^4$~\cite{SW71},
Ch.~IV, \S2.

\begin{lemma}\label{q-lem:chebexp}
$e^{\kappa t}=\frac2\kappa\sum_{n\ge0}(n+1)I_{n+1}(\kappa)\,U_n(t)$ in
$L^2\bigl(\frac2\pi(1-t^2)^{1/2}dt\bigr)$, and uniformly on $[-1,1]$.
\end{lemma}

\begin{proof}
With $t=\cos\theta$ and $2\sin\theta\sin((n+1)\theta)=\cos n\theta-\cos(
n+2)\theta$, the $n$-th coefficient is
\[
\begin{aligned}
\frac2\pi\int_0^\pi e^{\kappa\cos\theta}\sin((n+1)\theta)\sin\theta\,d\theta
&=I_n-I_{n+2}\\
&=\frac{2(n+1)}\kappa I_{n+1},
\end{aligned}
\]
by~\eqref{q-eq:besselrec}. Since $|U_n|\le n+1$ and $I_{n+1}$ decays
factorially, the series converges uniformly.
\end{proof}

\subsection{Proof of Lemma~\ref{q-lem:besselsums}}

Let $F(\theta)=e^{\kappa\cos\theta}=\sum_jI_je^{ij\theta}$. Then
\[
F'(\theta)=-\kappa\sin\theta\,F(\theta)=\sum_jijI_je^{ij\theta},
\]
and Parseval's formula gives
\[
\begin{aligned}
\sum_{j\in\Z}j^2I_j^2&=\frac{\kappa^2}{2\pi}\int_0^{2\pi}\sin^2\theta\,
e^{2\kappa\cos\theta}d\theta\\
&=\frac{\kappa^2}2\bigl(I_0(2\kappa)-I_2(2\kappa)\bigr)=\frac\kappa2I_1(2\kappa),
\end{aligned}
\]
the last step by~\eqref{q-eq:besselrec} at the argument $2\kappa$. Halving gives
the first identity. The coefficient of $e^{i\theta}$ in $F^2=e^{2\kappa\cos
\theta}$ is $I_1(2\kappa)=\sum_{j\in\Z}I_jI_{1-j}$, and the terms with $j\le0$
repeat those with $j\ge1$; this is the third identity.

For the second, $jI_j=\frac\kappa2(I_{j-1}-I_{j+1})$ by~\eqref{q-eq:besselrec}, so
\[
\sum_{j\ge1}j^3I_j^2=\frac{\kappa^2}4\sum_{j\ge1}j\bigl(I_{j-1}^2+I_{j+1}^2-
2I_{j-1}I_{j+1}\bigr).
\]
Shifting indices gives
\[
\sum_{j\ge1}j\bigl(I_{j-1}^2+I_{j+1}^2\bigr)=I_0^2+\sum_{m\ge1}2mI_m^2,
\qquad
\sum_{m\ge1}2mI_m^2=\kappa\sum_{m\ge1}(I_{m-1}I_m-I_mI_{m+1})=\kappa I_0I_1,
\]
the last by telescoping. For the mixed term write $j=\frac12((j-1)+(j+1))$
and apply~\eqref{q-eq:besselrec} to $(j-1)I_{j-1}$ and to $(j+1)I_{j+1}$,
using $I_{-1}=I_1$ at $j=1$:
\[
\begin{aligned}
jI_{j-1}I_{j+1}=\frac\kappa4\bigl[&(I_{j-1}I_j-I_jI_{j+1})\\
&+(I_{j-2}I_{j+1}-I_{j-1}I_{j+2})\bigr].
\end{aligned}
\]
Both brackets telescope, to $I_0I_1$ and to $I_{-1}I_2=I_1I_2$, so
$\sum_{j\ge1}jI_{j-1}I_{j+1}=\frac\kappa4(I_0I_1+I_1I_2)$. Altogether
\[
\sum_{j\ge1}j^3I_j^2=\frac{\kappa^2}4\Bigl(I_0^2+\frac\kappa2I_1(I_0-I_2)
\Bigr)=\frac{\kappa^2}4\bigl(I_0^2+I_1^2\bigr). \qquad\qed
\]

\subsection{\texorpdfstring{The coefficients of $\mathcal D$}{The coefficients of D}}

\begin{lemma}\label{q-lem:Dcoef}
Write $\mathcal D(\kappa)=\sum_{n\ge0}\mathcal D_n\kappa^{2n}$ for the function
of~\eqref{q-eq:Dfunction}. Then $\mathcal D_0=\mathcal D_1=0$, and for $m\ge1$
\[
\mathcal D_{m+1}=\frac{N_m}{4^m\,m!\,(m+1)!\,(m+1)(m+2)},
\]
with $N_m=\binom{2m}m\bigl(4(2m+1)+(m-1)(m+1)(m+2)\bigr)-4^m(m+1)(m+2)$.
Moreover $N_m=0$ for $m\le2$ and $N_m>0$ for $m\ge3$.
\end{lemma}

\begin{proof}
\emph{The coefficients.} Multiplying the series of $I_0$ and $I_1$ and
collecting powers,
\[
\begin{gathered}
I_0^2=\sum_{m\ge0}\alpha_m\kappa^{2m},\qquad I_1^2=\sum_{m\ge0}\beta_m
\kappa^{2m+2},\\
\kappa I_1(2\kappa)=\sum_{m\ge0}\frac{\kappa^{2m+2}}{m!\,(m+1)!},
\end{gathered}
\]
where
\[
\begin{aligned}
\alpha_m&=\frac1{4^m}\sum_{i=0}^m\frac1{i!^2(m-i)!^2},\\
\beta_m&=\frac1{4^{m+1}}\sum_{i=0}^m\frac1{i!\,(i+1)!\,(m-i)!\,(m-i+1)!}.
\end{aligned}
\]
By Vandermonde's identity, $\sum_i\binom mi^2=\binom{2m}m$ and
$\sum_i\binom mi\binom{m+2}{i+1}=\binom{2m+2}{m+1}$, so
\[
\alpha_m=\frac{\binom{2m}m}{4^m\,m!^2},\qquad\beta_m=\frac{\binom{2m+2}{m+1}}
{4^{m+1}m!\,(m+2)!}.
\]
The coefficient of $\kappa^{2m+2}$ in $\mathcal D$ is $\mathcal D_{m+1}=8\beta_m
+2\beta_{m-1}-\alpha_m-1/(m!(m+1)!)$, with $\beta_{-1}=0$; for $m=0$ it is
$\frac84-1-1=0$, and $\mathcal D_0=0$ because every term of $\mathcal D$ is
$O(\kappa^2)$. For $m\ge1$ put $\Delta=4^m\,m!\,(m+1)!\,(m+1)(m+2)$ and
$b=\binom{2m}m$. With $\binom{2m+2}{m+1}=\frac{2(2m+1)}{m+1}b$,
\[
\begin{aligned}
8\beta_m&=\frac{4(2m+1)\,b}\Delta,\qquad\alpha_m=\frac{(m+1)^2(m+2)\,b}\Delta,\\
2\beta_{m-1}&=\frac{2m(m+1)(m+2)\,b}\Delta,\\
\frac1{m!(m+1)!}&=\frac{4^m(m+1)(m+2)}\Delta,
\end{aligned}
\]
and $2m-(m+1)=m-1$ gives $\mathcal D_{m+1}=N_m/\Delta$.

\emph{The sign.} Directly, $N_0=N_1=N_2=0$ ($2-2$, $24-24$, $192-192$),
$N_3=80$, $N_4=1140$ and $N_5=10416$. For $m\ge6$ we first show
$16^m\le4m\binom{2m}m^2$ for $m\ge1$ by induction: it is an equality at
$m=1$, and $(m+1)\binom{2m+2}{m+1}=2(2m+1)\binom{2m}m$ gives
\[
\begin{aligned}
4(m+1)\binom{2m+2}{m+1}^2&=\frac{16(2m+1)^2}{m+1}\binom{2m}m^2\\
&\ge\frac{16(2m+1)^2}{4m(m+1)}16^m\ge16^{m+1}.
\end{aligned}
\]
Since $(m-1)^2\ge4m$ for $m\ge6$, it follows that $\binom{2m}m^2(m-1)^2\ge
16^m$, that is $\binom{2m}m(m-1)\ge4^m$, and
\[
\begin{aligned}
N_m={}&\Bigl(\binom{2m}m(m-1)-4^m\Bigr)(m+1)(m+2)\\
&+4(2m+1)\binom{2m}m>0 .
\end{aligned}
\]
\end{proof}

\subsection{Monotonicity of Bessel ratios}

\begin{lemma}\label{q-lem:ratio}
For $n\ge1$ and $\kappa>0$, $I_{n+2}(\kappa)/I_{n+1}(\kappa)\le I_2(\kappa)/
I_1(\kappa)$.
\end{lemma}

\begin{proof}
We use a lemma of Biernacki and Krzy\.z~\cite{BK55}: if $f=\sum a_jx^j$ and
$g=\sum b_jx^j$ have positive coefficients and $a_j/b_j$ is nonincreasing in
$j$, then $f/g$ is nonincreasing on $x>0$. Indeed
\[
\begin{aligned}
f'g-fg'&=\sum_{i,j}(i-j)a_ib_jx^{i+j-1}\\
&=\frac12\sum_{i,j}(i-j)b_ib_j\Bigl(\frac{a_i}{b_i}-\frac{a_j}{b_j}\Bigr)
x^{i+j-1}\le0 .
\end{aligned}
\]
Take $x=\kappa^2$, $f=I_{n+1}(\kappa)/\kappa^{n+1}$ and $g=I_1(\kappa)/\kappa$.
Then $a_j/b_j=2^{-n}(j+1)!/(j+n+1)!$ decreases in $j$, so $I_{n+1}/(\kappa^n
I_1)$ is nonincreasing in $\kappa$. Its logarithmic derivative is, by the
second recurrence of~\eqref{q-eq:besselrec},
\[
\begin{aligned}
&\frac{n+1}\kappa+\frac{I_{n+2}}{I_{n+1}}-\frac n\kappa-\frac1\kappa-\frac{I_2}
{I_1}\\
&\qquad=\frac{I_{n+2}}{I_{n+1}}-\frac{I_2}{I_1}\le0 .
\end{aligned}\qedhere
\]
\end{proof}


\section{A finite-size criterion for the cube condition}\label{app:finite}

This appendix proves the finite-size criterion used in
Sec.~\ref{sec:massgen}: a gap for the dynamics that resamples the blocks of
one patch of $m^4$ blocks, uniform in the configuration outside the patch,
gives the cube condition on every large torus. Blocks that share no
plaquette give commuting projections, and this replaces the disjointness
used in the finite-size criteria of Knabe~\cite{Kna88} and of Gosset and
Mozgunov~\cite{GM16} for quantum spin systems.\footnote{In those criteria
the local terms of a Hamiltonian act on overlapping sets of sites, and the
gap of the whole system is bounded below by the gap of a finite patch minus
a correction that decays with the size of the patch. Here the local terms
are the complements $I-\mathcal P_j$ of conditional expectations, and two
of them commute as soon as the blocks are at distance at least two.}

\emph{Setting.} The notation is that of Sec.~\ref{sec:cube} and of the
setting of Sec.~\ref{sec:massgen}. Fix integers $d\ge2$ and $m\ge4$ with
$d\mid L$, put $K=L/d$, $n=2d-1$, and fix a shift $a\in[0,d)^4$. For $j$ in
the coarse torus $T_K=(\Z/K\Z)^4$ let $r_j=a+dj$, let $B_j=V_{r_j}$ be the
set of links with base point in $r_j+[0,n)^4$, and write $\mathcal P_j=
\mathcal P_{r_j}$. Displacements $v=j'-j$ in $T_K$ are represented in
$(-K/2,K/2]^4$. Two sets $A,A'$ of links are \emph{non-interacting} if no
plaquette contains a link of each; for a set $A$ of links, $P_A$ is the
conditional expectation under $\mu$ given the links outside $A$.

\begin{lemma}\label{lem:fcommute}
\begin{enumerate}[label=\textup{(\roman*)},leftmargin=*]
\item For non-interacting sets $A,A'$ of links, $P_AP_{A'}=P_{A'}P_A=
P_{A\cup A'}$, and $(I-P_A)(I-P_{A'})$ is an orthogonal projection.
\item If $K\ge4$ and $|j'-j|_\infty\ge2$, then $B_j$ and $B_{j'}$ are
non-interacting.
\item The blocks $B_j$, $j\in T_K$, cover $E$.
\item If $f\in L^2(\mu)$ and $\mathcal P_jf=f$ for every $j\in T_K$, then
$f$ is almost surely constant.
\end{enumerate}
\end{lemma}

\begin{proof}
(i) Every link lies on a plaquette, so $A\cap A'=\emptyset$. Given the
links outside $A\cup A'$, the law of $(U_A,U_{A'})$ has density
proportional to $\prod_{p\cap(A\cup A')\ne\emptyset}\varphi(U_p)$, and each
factor involves the links of exactly one of $A$, $A'$. The density is
therefore a product of a function of $U_A$ and a function of $U_{A'}$,
with coefficients depending on the links outside $A\cup A'$. Hence $U_A$
and $U_{A'}$ are conditionally independent, the conditional law
$\gamma_A$ of $U_A$ given the links outside $A$ does not depend on
$U_{A'}$, and likewise for $A'$. So $P_AP_{A'}f$ and $P_{A'}P_Af$ are both
the integral of $f$ against $\gamma_A\otimes\gamma_{A'}$, which is
$P_{A\cup A'}f$, and $(I-P_A)(I-P_{A'})=I-P_A-P_{A'}+P_{A\cup A'}$ is
self-adjoint and idempotent.

(ii) Let $|v_i|\ge2$ for $v=j'-j$. In the coordinate $i$ the base points of
the links of $B_j$ and of $B_{j'}$ lie in $dj_i+a_i+[0,2d-2]$ and in
$dj_i+a_i+dv_i+[0,2d-2]$ modulo $L$. If $v_i\ge2$, their distance in
$\Z/L\Z$ is at least $\min\{(v_i-2)d+2,(K-v_i-2)d+2\}\ge2$, since $v_i\le
K/2$ and $K\ge4$; the case $v_i\le-2$ is symmetric. A plaquette meeting both
blocks would contain two links whose base points differ by at least two in
the coordinate $i$, while the base points of the links of a plaquette
differ by at most one in each coordinate.

(iii) The sets $r_j+[0,d)^4$ partition $\Lambda_L$, and $[0,d)\subset
[0,n)$.

(iv) The density of $\mu$ with respect to the Haar measure $h$ of $G^E$ is
bounded above and below by positive constants, so $\mu$ and $h$ have the
same null sets and $L^2(\mu)=L^2(h)$. If $\mathcal P_jf=f$, then $f$ is
almost surely a function of $U_{E\setminus B_j}$, and it is fixed by the
conditional expectation $Q_j$ of $h$ given $U_{E\setminus B_j}$, which
integrates out $U_{B_j}$. By Fubini's theorem the $Q_j$ commute and their
product integrates out the links of $\bigcup_jB_j=E$, so $f=\int f\,dh$
almost surely.
\end{proof}

In four dimensions the displacements $v$ with $|v|_\infty=1$ are $80$ in
number, and those with $|v|_\infty\ge2$ give commuting projections. For
$s\in T_K$ let $S_s=s+[0,m-1]^4$ be the \emph{patch} at $s$ and $S^\circ_s=
s+[0,m-2]^4$ the set of its $(m-1)^4$ block indices, and put
\[
H=\sum_{j\in T_K}(I-\mathcal P_j),\qquad H_s=\sum_{j\in S^\circ_s}(I-
\mathcal P_j).
\]
These are bounded nonnegative operators on $L^2(\mu)$ and preserve the
subspace $\mathcal A$ of gauge-invariant functions.

\begin{lemma}\label{lem:knabe}
Let $m\ge4$, $p=m-1$, $K\ge2m$, let $\mathcal K$ be $L^2(\mu)$ or
$\mathcal A$, and suppose that $H_s^2\ge\gamma H_s$ on $\mathcal K$ for
every $s\in T_K$, with $\gamma>0$. Then $H^2\ge\varepsilon H$ on
$\mathcal K$ with
\[
\varepsilon=1-\frac{(1-\gamma)p^4+72p^3-168p^2+96p-16}{p^3(p-2)}\ge
\frac{\gamma p-74}{p-2}.
\]
In particular $\varepsilon\ge\gamma/2$ if $m-1\ge148/\gamma$.
\end{lemma}

\begin{proof}
All operators below preserve $\mathcal K$, and a restriction of an
orthogonal projection to an invariant closed subspace is an orthogonal
projection; so it suffices to argue on $L^2(\mu)$ with inequalities read on
$\mathcal K$. Write $\mathbf P_j=I-\mathcal P_j$ and $A_v=\sum_j\mathbf P_j
\mathbf P_{j+v}$. Since $K\ge2m$, two indices $j$ and $j+v$ lie in a common
$S^\circ_s$ only if $|v_i|\le m-2$ for all $i$, and then for exactly
$c_v=\prod_{i=1}^4(p-|v_i|)$ points $s$. Let $O$ be the set of the $80$
displacements with $|v|_\infty=1$, $D$ the set of those with
$|v|_\infty\ge2$, and $D'$ the set of $v\in D$ with all $|v_i|\le m-2$. For
$v\in D$, $\mathbf P_j\mathbf P_{j+v}$ is an orthogonal projection by
Lemma~\ref{lem:fcommute}(i),(ii), so $A_v\ge0$. Summing the hypothesis over
$s$,
\[
\sum_sH_s^2=c_0H+\sum_{v\in O}c_vA_v+\sum_{v\in D'}c_vA_v\ge\gamma\sum_sH_s
=\gamma c_0H .
\]
Put $c_*=p^3(p-2)$; then $c_v\le c_*$ for $v\in D'$, since one factor of
$c_v$ is at most $p-2$ and the others are at most $p$. For orthogonal
projections $\mathbf P,\mathbf P'$, $(\mathbf P\pm\mathbf P')^2\ge0$ gives
$-(\mathbf P+\mathbf P')\le\mathbf P\mathbf P'+\mathbf P'\mathbf P\le
\mathbf P+\mathbf P'$, so $-2H\le A_v+A_{-v}\le2H$. As $O$ is symmetric and
$c_{-v}=c_v$,
\[
\sum_{v\in O}(c_*-c_v)A_v\ge-\Sigma H,\qquad\Sigma=\sum_{v\in O}|c_*-c_v|,
\]
and therefore
\[
c_*\sum_{v\in O}A_v=\sum_{v\in O}c_vA_v+\sum_{v\in O}(c_*-c_v)A_v\ge
-\bigl((1-\gamma)c_0+\Sigma\bigr)H-\sum_{v\in D'}c_vA_v .
\]
Since $H^2=\sum_vA_v=H+\sum_{v\in O}A_v+\sum_{v\in D}A_v$,
\[
H^2\ge\Bigl(1-\frac{(1-\gamma)c_0+\Sigma}{c_*}\Bigr)H+\sum_{v\in D'}\Bigl(1-
\frac{c_v}{c_*}\Bigr)A_v+\sum_{v\in D\setminus D'}A_v ,
\]
and the last two sums are nonnegative. It remains to compute $\Sigma$. A
displacement $v\in O$ with $k$ nonzero coordinates has $c_v=p^{4-k}(p-1)^k$,
and there are $8$, $24$, $32$ and $16$ of them for $k=1,2,3,4$. The
differences $c_v-c_*$ are
\[
p^3,\qquad p^2,\qquad-p(p^2-3p+1),\qquad-2p(p-1)(p-2)+1,
\]
with the signs shown since $p\ge3$. Hence
$\Sigma=8p^3+24p^2+32p(p^2-3p+1)+16(2p(p-1)(p-2)-1)=72p^3-168p^2+96p-16$,
and with $c_0=p^4$ this is the formula for $\varepsilon$. As $-168p^2+96p
-16<0$, $\varepsilon\ge1-((1-\gamma)p+72)/(p-2)=(\gamma p-74)/(p-2)$, and if
$\gamma p\ge148$ this is at least $(\gamma p-74)/p\ge\gamma/2$.
\end{proof}

For a patch $S_s$ let $\Pi_s=\bigcup_{j\in S^\circ_s}B_j$, the set of links
with base point in $r_s+[0,md-1)^4$. For $L\ge md+2$ the patch and the
plaquettes meeting it embed in $\Z^4$, and the conditional law of
$U_{\Pi_s}$ under $\mu$ given the links outside $\Pi_s$ is the measure
$\nu^\omega_{Q}$ of Sec.~\ref{sec:massgen} for the box $Q=r_s+[0,md-1)^4$,
with $\omega$ the configuration outside. For $j\in S^\circ_s$ write
$\mathcal P^\omega_j$ for the conditional expectation under $\nu^\omega_Q$
given the links of $\Pi_s\setminus B_j$.

\begin{definition}\label{def:patchgap}
The \emph{gauge-invariant patch gap} $\gamma_{d,m}$ is the largest
$\gamma\ge0$ such that
\[
\gamma\Var_{\nu^\omega_Q}(F)\le\sum_{j\in S^\circ_s}\nu^\omega_Q\bigl((F-
\mathcal P^\omega_jF)^2\bigr)
\]
for every patch box $Q=r+[0,md-1)^4$ of $\Z^4$, every exterior
configuration $\omega$ and every $F\in\mathcal A^\omega_Q$.
\end{definition}

By translation invariance $\gamma_{d,m}$ depends only on $G$, $\varphi$,
$d$ and $m$. It is the gap, on gauge-invariant functions, of the dynamics
that resamples the $(m-1)^4$ blocks of one patch, uniform in the
configuration outside the patch.

\begin{theorem}[finite-size criterion]\label{thm:finite}
Let $d\ge2$, $m\ge4$, $d\mid L$, $L$ even, $L\ge2md$, and suppose that
$\gamma_{d,m}\ge\gamma>0$ and that the number $\varepsilon$ of
Lemma~\ref{lem:knabe} is positive. Then
\[
\frac{\varepsilon}{16}\Var_\mu(f)\le\frac1{n^4}\sum_{x\in\Lambda_L}\mu\bigl(
(f-\mathcal P_xf)^2\bigr)\qquad\text{for every real }f\in\mathcal A,
\]
that is, $\BP(2d-1,\varepsilon/16)$. If $m-1\ge148/\gamma$, then
$\varepsilon/16\ge\gamma/32$.
\end{theorem}

\begin{proof}
Fix $a$. Since $L\ge2md$, $K\ge2m\ge8$.

\emph{The patch inequality on the torus.} Let $F\in\mathcal A$; replacing $F$
by $\int F(gU)\,dg$, the integral over all gauge transformations of
$\Lambda_L$ with Haar measure, we may take $F$ invariant at every point.
Write $U=(\sigma,\omega)$ with $\sigma$ the links of $\Pi_s$ and $\omega$ the
others. Under $\mu$ the conditional law of $\sigma$ given $\omega$ is
$\nu^\omega_Q$, and for $j\in S^\circ_s$ the conditional law of $U_{B_j}$
given the other links is the same under $\mu$ and under $\nu^\omega_Q$, so
$(\mathcal P_jF)(\cdot,\omega)=\mathcal P^\omega_j(F(\cdot,\omega))$ for almost
every $\omega$ and $H_s$ acts on each fibre as $H^\omega_s=\sum_{j\in
S^\circ_s}(I-\mathcal P^\omega_j)$. A gauge transformation of $\Lambda_L$
that is trivial at the vertices not interior to $Q$ fixes $\omega$, so each
fibre $F(\cdot,\omega)$ lies in $\mathcal A^\omega_Q$. On $\mathcal
A^\omega_Q$ the operator $H^\omega_s$ is self-adjoint and nonnegative,
preserves the constants and their orthogonal complement
(Lemma~\ref{lem:boxinv}(iii)), and by Definition~\ref{def:patchgap} is at
least $\gamma$ on that complement. So its spectrum there lies in
$\{0\}\cup[\gamma,\infty)$ and $(H^\omega_s)^2\ge\gamma H^\omega_s$.
Integrating over $\omega$, $\|H_sF\|^2\ge\gamma\ip F{H_sF}$ for $F\in
\mathcal A$, that is, $H_s^2\ge\gamma H_s$ on $\mathcal A$.

\emph{Conclusion.} By Lemma~\ref{lem:knabe} with $\mathcal K=\mathcal A$,
$H^2\ge\varepsilon H$ on $\mathcal A$, so the spectrum of $H$ on $\mathcal
A$ lies in $\{0\}\cup[\varepsilon,\infty)$. If $F\in\mathcal A$ and $HF=0$,
then $\sum_j\|F-\mathcal P_jF\|^2=\ip F{HF}=0$, so $F$ is constant by
Lemma~\ref{lem:fcommute}(iv). Hence for real $f\in\mathcal A$ the function
$g=f-\mu(f)\in\mathcal A$ is orthogonal to the kernel of $H$ on $\mathcal
A$, and
\[
\varepsilon\Var_\mu(f)\le\ip g{Hg}=\sum_{j\in T_K}\mu\bigl((f-\mathcal
P_{r_j}f)^2\bigr).
\]
As $a$ runs over $[0,d)^4$ and $j$ over $T_K$, $r_j=a+dj$ runs over
$\Lambda_L$ once. Averaging over $a$ gives $\varepsilon d^4\Var_\mu(f)\le
\sum_x\mu((f-\mathcal P_xf)^2)$, and $d^4\ge n^4/16$. Since $L\ge8d\ge2n+4$
this is $\BP(2d-1,\varepsilon/16)$, and the last claim is
Lemma~\ref{lem:knabe}.
\end{proof}

The passage from the patch to the torus costs a factor: at $\varphi\equiv1$
the links are independent, $\gamma_{d,m}\ge1$ by the Efron--Stein
inequality~\cite{ES81}, and the theorem gives $\BP(2d-1,\varepsilon/16)$
with $\varepsilon<1$, against $\BP(n,1)$ from Proposition~\ref{prop:bpsup}.

\section{The threshold one half, two-dimensional models and limits of methods}
\label{app:obstructions}

This appendix proves the statements about the hypothesis of
Section~\ref{sec:massgen} and about other approaches to (D) that are used in
the text. Section~\ref{ssec:gausschain} shows that the threshold $c<\frac12$ of
Theorem~\ref{thm:D} is the exact limit of the slab mechanism, the boundary value
being attained by a massless Gaussian chain. Sections~\ref{ssec:twodfat}
and~\ref{ssec:strips} treat two dimensions: the slab hypothesis for fat boxes
holds along the weak-coupling tuning, uniformly in the boundary condition, while
the same condition for strips of width one fails. Sections~\ref{ssec:curvlimits}
to~\ref{ssec:transfer} prove that four other routes cannot give (D) without an
input of the same strength: curvature of the Wilson measure, entropy
factorisation, tensorisation over blocks of fixed lattice size, and the
transfer matrix taken alone.

\subsection{The Gaussian chain}\label{ssec:gausschain}

For a probability measure and $\sigma$-algebras $\mathcal F_1,\mathcal F_2$ let
$\rho(\mathcal F_1,\mathcal F_2)=\sup|E[XY]|$, the supremum over centred $X\in
L^2(\mathcal F_1)$ and $Y\in L^2(\mathcal F_2)$ with $E[X^2]=E[Y^2]=1$; this is
the maximal correlation of Definition~\ref{def:rho}, without the constraint of
gauge invariance.

\begin{lemma}\label{lem:gmarkov}
Let $(U,V,M,W,X)$ be random variables with values in standard Borel spaces whose
joint law has the form $\psi_1(u,v)\psi_2(v,m,w)\psi_3(w,x)\,d\lambda$ for a
product measure $\lambda$ and nonnegative functions $\psi_i$. Then
$\rho(\sigma(U,V),\sigma(W,X))\le\rho(\sigma(V),\sigma(W))$.
\end{lemma}

\begin{proof}
By the factorisation, the conditional law of $U$ given $(V,M,W,X)$ has density
proportional to $\psi_1(\cdot,V)$ and that of $X$ given $(U,V,M,W)$ has density
proportional to $\psi_3(W,\cdot)$, and given $(V,M,W)$ the variables $U$ and
$X$ are conditionally independent. Let $f=f(U,V)$ and $g=g(W,X)$ be centred and
square integrable. Then $E[f\mid V,M,W,X]=\tilde f(V)$ with $\tilde f=E[f\mid
V]$, and $E[g\mid U,V,M,W]=\tilde g(W)$ with $\tilde g=E[g\mid W]$. Hence
$E[fg]=E[E[fg\mid V,M,W]]=E[\tilde f(V)\tilde g(W)]$, where $\tilde f$ and
$\tilde g$ are centred with $\|\tilde f\|\le\|f\|$ and $\|\tilde g\|\le\|g\|$.
\end{proof}

\begin{lemma}[Gebelein]\label{lem:gebelein}
If $(V,W)$ is a nondegenerate Gaussian vector in $\R^2$ with correlation
coefficient $r$, then $\rho(\sigma(V),\sigma(W))=|r|$.
\end{lemma}

\begin{proof}
The choice $X=V$, $Y=\pm W$, standardised, gives $\rho\ge|r|$. For the converse
standardise $V$ and $W$. Let $H_n$ be the Hermite polynomials with generating
function $\sum_nH_n(x)s^n/n!=e^{sx-s^2/2}$; then $h_n=H_n/\sqrt{n!}$ is a
complete orthonormal system of $L^2(\R,\gamma)$, $\gamma$ the standard Gaussian
law~\cite[Sec.~1.1.1]{Nua06}. From $E[e^{sV-s^2/2}e^{uW-u^2/2}]=e^{rsu}$ and a
comparison of coefficients one gets $E[h_j(V)h_k(W)]=\delta_{jk}r^j$; the double
series converges absolutely, since $H_n(x)=E(x+iZ)^n$ with $Z\sim\gamma$ gives
$|H_n(x)|\le E(|x|+|Z|)^n$ and $Ee^{t(|V|+|W|+2|Z|)}<\infty$ for every $t$.
For centred $f=\sum_{j\ge1}a_jh_j$ and $g=\sum_{j\ge1}b_jh_j$ in $L^2(\gamma)$
the form $(f,g)\mapsto E[f(V)g(W)]$ is bounded by $\|f\|\|g\|$, so passing to
the limit in the partial sums, $E[f(V)g(W)]=\sum_{j\ge1}a_jb_jr^j$, whose
modulus is at most $|r|\sum_j|a_jb_j|\le|r|\|f\|\|g\|$.
\end{proof}

The inequality $\rho\le|r|$ is Gebelein's theorem~\cite{Geb41}.

\emph{The chain.} Let $\mu\ge0$ and define $\kappa\ge0$ by $2\cosh\kappa=2+
\mu^2$. For integers $p<q$ with $q-p\ge2$ and real boundary values $z_p,z_q$ let
$\nu$ be the Gaussian law of $(z_{p+1},\dots,z_{q-1})$ with density proportional
to
\[
\exp\Bigl(-\tfrac12\sum_{j=p}^{q-1}(z_{j+1}-z_j)^2-\tfrac12\mu^2\sum_{j=p+1}^{q-1}
z_j^2\Bigr),
\]
a nearest-neighbour Gaussian chain, massless for $\mu=0$. Its precision matrix
on the $q-p-1$ interior sites is $\mathsf L=\operatorname{tridiag}(-1,2+\mu^2,
-1)$, independent of $z_p$ and $z_q$, which only shift the mean, and its
conditional laws on subintervals are chains of the same form.

\begin{lemma}\label{lem:green}
Put $T=q-p$ and, for $0\le i\le T$, $u(i)=\sinh(\kappa i)$ and $v(i)=\sinh(
\kappa(T-i))$ if $\mu>0$, $u(i)=i$ and $v(i)=T-i$ if $\mu=0$. For $0<i\le j<T$,
$\Cov_\nu(z_{p+i},z_{p+j})=A\,u(i)v(j)$ with $A>0$ independent of $i$ and $j$;
for $\mu=0$, $A=1/T$. Hence the correlation of $z_{p+i}$ and $z_{p+j}$ is
$(u(i)v(j)/(v(i)u(j)))^{1/2}$.
\end{lemma}

\begin{proof}
Both $u$ and $v$ satisfy $(2+\mu^2)f(l)=f(l-1)+f(l+1)$ (for $\mu>0$ because
$\sinh(x-\kappa)+\sinh(x+\kappa)=2\cosh\kappa\sinh x$), $u(0)=v(T)=0$, and both
are positive on $(0,T)$. The covariance matrix is $G=\mathsf L^{-1}$, which is
symmetric. For fixed $j$ the column $G(\cdot,j)$, extended by $0$ at $0$ and
$T$, satisfies the recursion at every $l\ne j$, so $G(i,j)=\alpha_ju(i)$ for
$i\le j$ and $G(i,j)=\beta_jv(i)$ for $i\ge j$: a solution of the second-order
recursion on $[0,j]$ that vanishes at $0$ is a multiple of $u$, and similarly
for $v$. For $i<j$, $\alpha_ju(i)=G(i,j)=G(j,i)=\beta_iv(j)$, so $\alpha_j/v(j)
=\beta_i/u(i)$ for all $i<j$, a constant $A$; and $\alpha_ju(j)=\beta_jv(j)$
gives $G(j,j)=Au(j)v(j)$, so $A>0$. For $\mu=0$ the $j$-th entry of $\mathsf L
G(\cdot,j)$ equals one, that is $A[2j(T-j)-(j-1)(T-j)-j(T-j-1)]=AT=1$.
\end{proof}

\emph{Slab geometry.} Let $s\ge1$, $s\le w\le2s-1$ and $M\ge1$. Take the sites
$I=[0,sM+w)$ with the two exterior sites $-1$ and $sM+w$, the blocks
$R_t=[st,st+w)$ for $0\le t\le M$, and, for $t<M$, $Q_t=R_t\cup R_{t+1}$ split
into $Q_t^-=R_t\setminus R_{t+1}=[st,st+s)$, $Q_t^0=R_t\cap R_{t+1}=[st+s,
st+w)$ and $Q_t^+=R_{t+1}\setminus R_t=[st+w,st+s+w)$. The condition $w\le2s-1$
says that $R_t$ and $R_{t+2}$ contain no adjacent sites, so that, exactly as in
Proposition~\ref{prop:slab}, the resampling moves of non-adjacent blocks
commute; Section~\ref{sec:massgen} uses $s=d$ and $w=2d-1$. Let $\nu$ be the
chain on $I$ with any exterior values, $P_t$ the conditional expectation given
the sites outside $R_t$, and $\rho_t$ the supremum, over all values outside
$Q_t$, of $\rho(\sigma(z_{Q_t^-}),\sigma(z_{Q_t^+}))$ under the conditional law
of $z_{Q_t}$.

\begin{proposition}[the Gaussian chain]\label{prop:gausschain}
\leavevmode
\begin{enumerate}[label=\textup{(\roman*)},leftmargin=*]
\item For every $t<M$, $\rho_t=\sinh(\kappa s)/\sinh(\kappa(w+1))$ if $\mu>0$,
and $\rho_t=s/(w+1)$ if $\mu=0$. For $w=2s-1$ this is $1/(2\cosh(\kappa s))$,
and $\frac12$ for $\mu=0$. For $\mu=0$, $\rho_t\ge\frac12$ in every admissible
geometry, with equality only for $w=2s-1$.
\item For $\mu=0$ the best constant $\gamma_M$ in the inequality $\gamma_M
\Var_\nu(F)\le\sum_{t=0}^M\nu((F-P_tF)^2)$ for all $F\in L^2(\nu)$ satisfies
$\gamma_M\le16/(3M)$.
\item For a subinterval split into consecutive nonempty intervals of $A'$, $r$
and $C'$ sites, with the sites outside it conditioned, the maximal correlation
of the two outer intervals is
\[
\Bigl(\frac{\sinh(\kappa A')\sinh(\kappa C')}{\sinh(\kappa(A'+r+1))\sinh(\kappa(
C'+r+1))}\Bigr)^{1/2}\le e^{-\kappa(r+1)}\quad\text{if }\mu>0,
\]
with supremum $e^{-\kappa(r+1)}$ over $A'$ and $C'$, and $(A'C'/((A'+r+1)(C'+r
+1)))^{1/2}$ if $\mu=0$, with supremum $1$.
\end{enumerate}
\end{proposition}

\begin{proof}
(i) Given the values outside $Q_t$, the law of $z_{Q_t}$ is the chain with
$p=st-1$, $q=st+s+w$, $T=s+w+1$. Its density factorises as in
Lemma~\ref{lem:gmarkov} with $U=z_{[st,st+s-1)}$, $V=z_{st+s-1}$,
$M=z_{Q^0_t}$, $W=z_{st+w}$ and $X=z_{(st+w,st+s+w)}$, the interactions being
between nearest neighbours; so $\rho_t\le\rho(\sigma(V),\sigma(W))$, with
equality since $V$ and $W$ are functions of the two outer parts. By
Lemma~\ref{lem:gebelein} this is the modulus of the correlation of $V$ and $W$,
which by Lemma~\ref{lem:green} with $i=s$, $j=w+1$ and $T=s+w+1$ equals
$(u(s)v(w+1)/(v(s)u(w+1)))^{1/2}$. Here $v(w+1)=u(s)$ and $v(s)=u(w+1)$, which
gives the formulas; they do not depend on the conditioning values. For $w=2s-1$,
$\sinh(\kappa s)/\sinh(2\kappa s)=1/(2\cosh\kappa s)$. For $\mu=0$, $w+1\le2s$
gives $s/(w+1)\ge\frac12$, with equality only when $w=2s-1$.

(ii) Let $T=sM+w+1\ge sM\ge2$ and take $F=z_j$ with $j=\lfloor T/2\rfloor-1\in
I$. By Lemma~\ref{lem:green} with $p=-1$, $\Var_\nu(F)=i(T-i)/T$ with $i=
\lfloor T/2\rfloor$, so $\Var_\nu(F)\ge(T^2-1)/(4T)\ge3T/16$. If $j\notin R_t$,
then $P_tF=F$. If $j\in R_t$, then $\nu((F-P_tF)^2)$ is the mean conditional
variance of $z_j$ given $z_{st-1}$ and $z_{st+w}$, the variance of a chain of
$w+1$ steps at an interior site, which is at most $(w+1)/4\le s/2$ by
Lemma~\ref{lem:green}. The $t$ with $j\in R_t$ satisfy $(j-w)/s<t\le j/s$, an
interval of length $w/s<2$, so there are at most two. Hence $\sum_t\nu((F-P_t
F)^2)\le s$, and $\gamma_M\le16s/(3T)\le16/(3M)$.

(iii) Put the left conditioned site at $0$. The outer parts are $[1,A']$ and
$[A'+r+1,A'+r+C']$, and the right conditioned site is $T=A'+r+C'+1$. As in (i),
by Lemmas~\ref{lem:gmarkov} and~\ref{lem:gebelein}, the maximal correlation is
the modulus of the correlation of the inner ends $z_{A'}$ and $z_{A'+r+1}$,
which by Lemma~\ref{lem:green} is $(u(A')v(A'+r+1)/(v(A')u(A'+r+1)))^{1/2}$.
With $u_0(x)=\sinh\kappa x$, respectively $u_0(x)=x$, one has $u(i)=u_0(i)$ and
$v(i)=u_0(T-i)$, so $v(A'+r+1)=u_0(C')$ and $v(A')=u_0(C'+r+1)$, which gives the
formulas. For $\mu>0$ the inequality $\sinh x/\sinh(x+y)\le e^{-y}$ for $x,y\ge
0$, which is equivalent to $e^{-x-y}\le e^{-x+y}$, bounds both factors by
$e^{-\kappa(r+1)}$, and both tend to $e^{-\kappa(r+1)}$ as $A',C'\to\infty$.
For $\mu=0$ the expression tends to $1$.
\end{proof}

Three consequences explain the role of the chain in Section~\ref{sec:massgen}.
First, for $w=2s-1$ the massless chain satisfies the hypothesis of
Proposition~\ref{prop:slab} with $c=\frac12$ for every exterior condition, its
non-adjacent moves commute, and yet by (ii) no inequality $\gamma\Var_\nu(F)\le
\sum_t\nu((F-P_tF)^2)$ holds with $\gamma>0$ independent of $M$; in any geometry
with commuting non-adjacent moves the same holds for every $c\ge s/(w+1)\ge
\frac12$. So the threshold $\frac12$ of Lemma~\ref{lem:chain} cannot be raised
by an argument that uses only these two properties, and no slab geometry of
this kind improves it. Second, for $w=2s-1$ the slab constant of the chain is
$1/(2\cosh\kappa s)$, which is below $\frac12$ if and only if $\mu>0$: along a
tuning $s=\ell/(2a)$ with lattice mass $\mu=\mu_{\rm phys}a$ one has $\kappa/
\mu\to1$, so the constant tends to $1/(2\cosh(\mu_{\rm phys}\ell/2))<\frac12$.
In the chain the hypothesis is exactly a mass at the scale $\ell/2$. Third, the
condition $(\mathrm M^*_r(c))$ with a middle of width $r$ has in the chain
the best constant $e^{-\kappa(r+1)}$ by (iii), and for $\mu=0$ it fails for
every $c<1$ and every $r$; the iteration of Proposition~\ref{prop:alldecay} reproduces the
true decay up to the integer part of the ratio of widths, so it loses nothing,
but it only transports a mass present at the starting scale.

\subsection{Two dimensions: fat boxes}\label{ssec:twodfat}

In this section and the next, $G$ is a compact group with normalised Haar
measure $\lambda$, also on products of copies of $G$, and $\varphi$ is a
positive continuous class function on $G$ with $\varphi(g^{-1})=\varphi(g)$; for
$G=\SU(2)$ the last condition holds for every class function, since every
element of $\SU(2)$ is conjugate to its inverse. Multiplying $\varphi$ by a
constant changes no measure, so we normalise $\int\varphi\,d\lambda=1$. For
class functions put $(\alpha*\beta)(g)=\int\alpha(gy^{-1})\beta(y)\,dy$ and let
$\alpha^{*k}$ be the $k$-fold convolution power.

Links of $\Z^2$ are $(y,\nu)$, $\nu\in\{1,2\}$, from $y$ to $y+e_\nu$, and the
plaquette with base point $z$ has holonomy $U_p=U_{(z,1)}U_{(z+e_1,2)}
U_{(z+e_2,1)}^{-1}U_{(z,2)}^{-1}$. Boxes $Q=([\alpha_1,\alpha_1+s_1)\times[
\alpha_2,\alpha_2+s_2))\cap\Z^2$, the sets $V(Q)$, the measures $\nu^\omega_Q$,
the groups $\mathcal G_Q$, the spaces $\mathcal A^\omega_Q$, splittings and the
maximal correlation $\rho^\omega(Q_-,Q_+)$ are defined as in
Section~\ref{sec:massgen} and Definition~\ref{def:rho}, with $\Z^2$ in place of
$\Z^4$; a vertex $v$ is interior to $Q$ when its four incident links lie in
$V(Q)$, that is, when $v$, $v-e_1$ and $v-e_2$ lie in $Q$. The
\emph{two-dimensional form} of $(\mathrm M_{d,m}(c))$ is the condition
$\rho^\omega(Q_-,Q_+)\le c$ for every box of $\Z^2$ whose two widths lie in
$[2d-1,md-1]$, every coordinate in which it has width $3d-1$, the splitting into
widths $d$, $d-1$, $d$ in that coordinate, and every exterior configuration
$\omega$.

\begin{lemma}[merging]\label{lem:merge}
Let $\alpha,\beta$ be continuous class functions on $G$.
\begin{enumerate}[label=\textup{(\roman*)},leftmargin=*]
\item $\alpha*\beta$ is a continuous class function, positive if $\alpha$ and
$\beta$ are, and $\int\alpha(Ax)\beta(x^{-1}B)\,dx=(\alpha*\beta)(AB)$ for all
$A,B\in G$.
\item Let $\alpha_0,\dots,\alpha_M$ be continuous class functions, $M\ge1$, and
$g_j,c_j,V_0,V_{M+1}\in G$. Then
\[
\int\prod_{j=0}^M\alpha_j\bigl(V_{j+1}g_jV_j^{-1}c_j\bigr)\,dV_1\cdots dV_M=
(\alpha_0*\cdots*\alpha_M)\bigl(V_{M+1}\,g_M\cdots g_0\,V_0^{-1}\,c_0\cdots c_M
\bigr).
\]
\end{enumerate}
\end{lemma}

\begin{proof}
(i) The substitution $y=x^{-1}B$ preserves Haar measure and turns the integral
into $\int\alpha(ABy^{-1})\beta(y)\,dy$. For $h\in G$ the substitution $y=hy'
h^{-1}$ gives $(\alpha*\beta)(hgh^{-1})=\int\alpha(hgy'^{-1}h^{-1})\beta(hy'
h^{-1})\,dy'=(\alpha*\beta)(g)$. Continuity and positivity are clear.

(ii) For $M=1$: since $\alpha_0$ and $\alpha_1$ are class functions,
$\alpha_0(V_1g_0V_0^{-1}c_0)=\alpha_0(g_0V_0^{-1}c_0V_1)$ and $\alpha_1(V_2g_1
V_1^{-1}c_1)=\alpha_1(V_1^{-1}c_1V_2g_1)$, so by (i) the integral over $V_1$ is
$(\alpha_0*\alpha_1)(g_0V_0^{-1}c_0c_1V_2g_1)=(\alpha_0*\alpha_1)(V_2g_1g_0
V_0^{-1}c_0c_1)$. For $M\ge2$ integrate $V_1,\dots,V_{M-1}$ first; by
induction the result is $\beta(V_Mg_{M-1}\cdots g_0V_0^{-1}c_0\cdots c_{M-1})$
with $\beta=\alpha_0*\cdots*\alpha_{M-1}$, and the case $M=1$ with
$g_{M-1}\cdots g_0$ and $c_0\cdots c_{M-1}$ in place of $g_0$ and $c_0$ gives
the claim.
\end{proof}

Fix a box, after translation $Q=[0,s_1)\times[0,s)$ with $s_1=3d-1$, $d\ge2$,
$s\ge1$, split in the first coordinate into $Q_-=[0,d)\times[0,s)$, $Q_0=
[d,2d-1)\times[0,s)$ and $Q_+=[2d-1,3d-1)\times[0,s)$. Write $h(x,y)=
\sigma_{((x,y),1)}$ and $v(x,y)=\sigma_{((x,y),2)}$. The interior vertices are
$I=[1,s_1-1]\times[1,s-1]$. Let
\[
F'=\{((x,y),2):1\le x\le s_1-1,\ 0\le y\le s-2\},
\]
the vertical links entering the interior vertices from below ($F'=\emptyset$
if $s=1$), and $E'=V(Q)\setminus F'$: all $h(x,y)$, all $v(0,y)$, and
$v(x,s-1)$ for $x\ge1$.

\begin{lemma}[axial gauge]\label{lem:axial}
For $\sigma\in G^{V(Q)}$ define $H(\sigma)\in\mathcal G_Q$ recursively in the
second coordinate by $H_u=\sigma_{(u-e_2,2)}^{-1}H_{u-e_2}$ for $u\in I$, with
$H_w=1$ for $w\notin I$, and put $\tau(\sigma)=(H(\sigma)^{-1}\sigma)|_{E'}$.
\begin{enumerate}[label=\textup{(\roman*)},leftmargin=*]
\item $(H(\sigma)^{-1}\sigma)_e=1$ for $e\in F'$, and $H(\sigma)$ depends only
on $\sigma_{F'}$.
\item Under $\lambda$ on $G^{V(Q)}$, $\tau(\sigma)$ has law $\lambda$ on $G^{E'}$.
\item Under $\nu=\nu^\omega_Q$, $\tau(\sigma)$ has density proportional to
$\mathcal W(\tau)=\prod_{p\cap V(Q)\ne\emptyset}\varphi((\tau,1_{F'},\omega)_p)$.
\item If $X\in\mathcal A^\omega_Q$ is measurable with respect to
$\sigma_{V(Q_-)}$, there is a Borel function $\bar X$ with $X=\bar X(\tau(
\sigma)_{E'\cap V(Q_-)})$ $\nu$-almost surely; likewise for $Q_+$.
\end{enumerate}
\end{lemma}

\begin{proof}
(i) For $u\in I$, $(H^{-1}\sigma)_{(u-e_2,2)}=H_{u-e_2}^{-1}\sigma_{(u-e_2,2)}
H_u=1$ by the definition of $H_u$. The recursion terminates, since $u-e_2$ has
a smaller second coordinate, and it uses only the links $(u-je_2,2)$, $1\le j
\le u_2$, which lie in $F'$.

(ii) Under $\lambda$ the families $\sigma_{F'}$ and $\sigma_{E'}$ are
independent with laws $\lambda$. Given $\sigma_{F'}$, $H$ is fixed, and for a
link $e$ of $E'$ from $x$ to $y$, $\tau_e=H_x^{-1}\sigma_eH_y$ is a two-sided
translate of $\sigma_e$. Translations preserve $\lambda$, so the conditional
law of $\tau$ given $\sigma_{F'}$ is $\lambda$, and so is its law.

(iii) Let $w(\sigma)=\prod_{p\cap V(Q)\ne\emptyset}\varphi((\sigma,\omega)_p)$.
Every $h\in\mathcal G_Q$ fixes the links outside $V(Q)$, whose end points are
not interior, so $(h\sigma,\omega)=h\cdot(\sigma,\omega)$ and each plaquette
holonomy is conjugated; as $\varphi$ is a class function, $w(h\sigma)=
w(\sigma)$. Hence $w(\sigma)=w(H(\sigma)^{-1}\sigma)=\mathcal W(\tau(\sigma))$
by (i), and for bounded Borel $\Psi$, by (ii), $\int\Psi(\tau(\sigma))
w(\sigma)\,d\lambda(\sigma)=\int\Psi(\tau)\mathcal W(\tau)\,d\lambda(\tau)$.

(iv) The group $\mathcal G_Q\cong G^I$ acts on $G^{V(Q_-)}$ link by link. Let
$X=f(\sigma_{V(Q_-)})$ almost surely with $f$ Borel, and put $\bar f(\eta)=
\int f(h\eta)\,dh$ where the integral converges absolutely and $\bar f=0$
elsewhere; by invariance of Haar measure both the set of convergence and $\bar f$
are invariant under every $k\in\mathcal G_Q$. Since $\nu$ and $\lambda$ have the
same null sets, the density being bounded above and below, and $f(h\eta)=
f(\eta)$ for $\nu$-almost every configuration for each $h$, Fubini's theorem
gives $\bar f=f$ almost surely. Then $X=\bar f(\sigma_{V(Q_-)})=\bar f((H(
\sigma)^{-1}\sigma)_{V(Q_-)})$ almost surely, and by (i) the right side is a
function of $\tau(\sigma)_{E'\cap V(Q_-)}$, the coordinates in $F'$ being $1$.
\end{proof}

\begin{lemma}[near-product densities]\label{lem:nearproduct}
Let $(S_1,\mu_1)$ and $(S_2,\mu_2)$ be probability spaces and $\pi$ a
probability measure on $S_1\times S_2$ with density $r$ with respect to
$\mu_1\otimes\mu_2$, $|r-1|\le\delta<1$. If $X\in L^2(\pi)$ depends only on
$s_1$, $Y\in L^2(\pi)$ only on $s_2$, and $\pi(X)=\pi(Y)=0$, then $|\pi(XY)|\le
\frac{\delta+\delta^2}{1-\delta}\|X\|_{L^2(\pi)}\|Y\|_{L^2(\pi)}$. If $r=k/Z'$
with $1-\varepsilon\le k\le1+\varepsilon$ and $Z'=\int k\,d(\mu_1\otimes\mu_2)$,
then $\delta=2\varepsilon/(1-\varepsilon)$ is admissible, and for $\varepsilon\le
\frac1{10}$ the bound is at most $4\varepsilon$.
\end{lemma}

\begin{proof}
The marginal density $m_1(s_1)=\int r(s_1,\cdot)\,d\mu_2$ lies in $[1-\delta,1+
\delta]$, so $\int X^2d\mu_1\le(1-\delta)^{-1}\|X\|^2_{L^2(\pi)}$; and $\int X\,
d\mu_1=\int X(1-m_1)\,d\mu_1$ since $\int Xm_1\,d\mu_1=\pi(X)=0$, so $|\int X\,
d\mu_1|\le\delta\int|X|\,d\mu_1$. The same holds for $Y$. Now $\pi(XY)=\iint XY
(r-1)\,d\mu_1d\mu_2+\int X\,d\mu_1\int Y\,d\mu_2$, and by the Cauchy--Schwarz
inequality the two terms together are at most $(\delta+\delta^2)(\int X^2d\mu_1
\int Y^2d\mu_2)^{1/2}$, which gives the first claim. If $r=k/Z'$, then $Z'\in
[1-\varepsilon,1+\varepsilon]$ and $r\in[\frac{1-\varepsilon}{1+\varepsilon},
\frac{1+\varepsilon}{1-\varepsilon}]$, so $|r-1|\le2\varepsilon/(1-\varepsilon)$.
For $\varepsilon\le\frac1{10}$, $\delta\le\frac{20}9\varepsilon\le\frac29$ and
$\frac{\delta(1+\delta)}{1-\delta}\le\frac{11}7\delta\le\frac{220}{63}
\varepsilon<4\varepsilon$.
\end{proof}

For $G=\SU(2)$ write $\chi_j$ for the character of spin $j\in\frac12\Z_{\ge0}$,
$d_j=2j+1$ and $c_j=j(j+1)$, and let
\[
p_t=\sum_jd_je^{-tc_j}\chi_j\qquad(t>0)
\]
be the heat kernel. It is a uniformly convergent series of real class
functions, $\int p_t=1$, $p_t*p_{t'}=p_{t+t'}$ because $\chi_i*\chi_j=
\delta_{ij}\chi_j/d_j$, and $p_t>0$, as the heat kernel of a compact connected
Riemannian manifold~\cite[Chs.~7--8]{Gri09}. It is the weight of the
\emph{heat-kernel action}, for which the lattice theory in two dimensions
reproduces the continuum Yang--Mills measure on every lattice
\cite{GKS89,Dri89,Lev03}, and the two-dimensional tuning is $t=g^2a^2$ with a
coupling $g$ of the dimension of a mass. With $n=2j+1$, $c_j=(n^2-1)/4$ and
$|\chi_j|\le d_j$,
\begin{equation}\label{eq:hk}
\|p_\theta-1\|_\infty\le\sum_{n\ge2}n^2e^{-\theta(n^2-1)/4}\le\sum_{k\ge1}
(k+1)^2q^k\le\frac{4q}{(1-q)^3},\qquad q=e^{-3\theta/4},
\end{equation}
using $n^2-1\ge3(n-1)$ and $\sum_{k\ge1}(k+1)^2q^k=(4q-3q^2+q^3)/(1-q)^3$. For
$\theta\ge6$, $q\le e^{-4.5}<\frac1{20}$ and the right side is at most $5q\le
0.06$.

\begin{theorem}[two dimensions, fat boxes]\label{thm:twod}
\leavevmode
\begin{enumerate}[label=\textup{(\roman*)},leftmargin=*]
\item Let $d\ge2$, $s\ge1$, let $Q$ be a box of $\Z^2$ of widths $3d-1$ and $s$,
in either order, split in the coordinate of width $3d-1$ into widths $d$,
$d-1$, $d$, and put $N=1+(d-1)(s+1)$ and $\varepsilon_N=\|\varphi^{*N}-1
\|_\infty$. If $\varepsilon_N\le\frac1{10}$, then $\rho^\omega(Q_-,Q_+)\le4
\varepsilon_N$ for every exterior configuration $\omega$.
\item Let $G=\SU(2)$, $\varphi=p_t$ with $t=g^2a^2$, $d\ge2$ and $gda\ge\sqrt6$.
For every box of widths $3d-1$ and $s\ge2d-1$ and every $\omega$,
$\rho^\omega(Q_-,Q_+)\le20e^{-3Nt/4}$ with $N=1+(d-1)(s+1)$. In particular the
two-dimensional form of $(\mathrm M_{d,m}(c))$ holds for every $m\ge4$ with
\[
c=20\,e^{-3g^2(da)^2/4}.
\]
\end{enumerate}
\end{theorem}

\begin{proof}
(i) \emph{Reduction to the split in the first coordinate.} The exchange of the
two coordinates maps links to links, boxes to boxes, interior vertices to
interior vertices and $\mathcal G_Q$ to $\mathcal G_{Q'}$, and it maps the
holonomy of the plaquette at $z$ to the inverse of the holonomy of the
plaquette at the image of $z$; since $\varphi(g^{-1})=\varphi(g)$, it maps
$\nu^\omega_Q$ to $\nu^{\omega'}_{Q'}$. So we may take $Q$ as above.

\emph{Reduction to the axial model.} By Lemma~\ref{lem:axial}(iii) and (iv),
$\rho^\omega(Q_-,Q_+)$ is at most the supremum of $|\pi(\bar X\bar Y)|$ over
centred unit vectors $\bar X\in L^2(\pi)$ depending on $\tau_-=\tau_{E'\cap
V(Q_-)}$ and $\bar Y$ depending on $\tau_+=\tau_{E'\cap V(Q_+)}$, where $\pi$ is
the law on $G^{E'}$ with density proportional to $\mathcal W$.

\emph{Which plaquettes couple the three parts.} Every plaquette meeting $V(Q)$
contains a link of $E'$: the plaquette at $z\in Q$ contains $h(z)$, the
plaquette at $(-1,y)$ contains $v(0,y)$, and the plaquette at $(x,-1)$ contains
$h(x,0)$. The plaquette at $(x,y)$ contains $h(x,y)$, $h(x,y+1)$ and $v(x,y)$,
which lie in the part containing the column $x$, and $v(x+1,y)$, which lies in
the part containing the column $x+1$ and belongs to $F'$, so is fixed to $1$,
unless $y=s-1$. Hence the only plaquettes that contain free links of two
different parts are the two seam plaquettes at $(d-1,s-1)$, joining $Q_-$ and
$Q_0$, and at $(2d-2,s-1)$, joining $Q_0$ and $Q_+$, and no plaquette contains
free links of both $Q_-$ and $Q_+$. So $\mathcal W=\mathcal W_-(\tau_-)\mathcal
W_+(\tau_+)\widehat{\mathcal W}$, where $\widehat{\mathcal W}$ is the product of
the factors of the plaquettes that contain a free link of $Q_0$.

\emph{Integration of the middle.} Write $H_x=h(x,s-1)$, $V_x=v(x,s-1)$ and
$e_x=\sigma_{((x,s),1)}$, an exterior link. For a column $x\in[d,2d-2]$ the
factors of $\widehat{\mathcal W}$ are: the bottom plaquette at $(x,-1)$, which
equals $\varphi(h(x,0)^{-1}c_x)$ with $c_x$ a product of exterior links, by the
cyclic invariance of $\varphi$; the plaquettes at $(x,y)$, $0\le y\le s-2$,
which equal $\varphi(h(x,y)h(x,y+1)^{-1})$, both vertical links being in $F'$;
and the top plaquette $\varphi(H_xV_{x+1}e_x^{-1}V_x^{-1})$. Integrating
$h(x,0),\dots,h(x,s-2)$ successively with Lemma~\ref{lem:merge}(i), each step
having the form $\int\varphi(h_{y+1}^{-1}h_y)\psi(h_y^{-1}c_x)\,dh_y$, gives
$\varphi^{*s}(H_x^{-1}c_x)$, and then $\int\varphi(V_{x+1}e_x^{-1}V_x^{-1}H_x)
\varphi^{*s}(H_x^{-1}c_x)\,dH_x=\varphi^{*(s+1)}(V_{x+1}e_x^{-1}V_x^{-1}c_x)$,
again by Lemma~\ref{lem:merge}(i). The seam plaquette at $(d-1,s-1)$ is
$\varphi(H_{d-1}V_de_{d-1}^{-1}V_{d-1}^{-1})=\varphi(V_de_{d-1}^{-1}V_{d-1}^{-1}
H_{d-1})$, with $H_{d-1}$ and $V_{d-1}$ free links of $Q_-$, and the seam
plaquette at $(2d-2,s-1)$ is the top plaquette of the column $2d-2$, which
contains $V_{2d-1}$, a free link of $Q_+$. Lemma~\ref{lem:merge}(ii) with
$M=d-1$, $\alpha_0=\varphi$, $\alpha_j=\varphi^{*(s+1)}$ for $j\ge1$,
$V_0=V_{d-1}$, $c_0=H_{d-1}$, $g_j=e^{-1}_{d-1+j}$ and $c_j=c_{d-1+j}$ gives,
after a cyclic rearrangement,
\[
\int\widehat{\mathcal W}\,d\lambda(\tau_0)=\varphi^{*N}(\Phi_-\Phi_+),\qquad
\Phi_-=V_{d-1}^{-1}H_{d-1}C,\quad\Phi_+=V_{2d-1}\Gamma,
\]
with $C=c_d\cdots c_{2d-2}$ and $\Gamma=e_{2d-2}^{-1}\cdots e_{d-1}^{-1}$
exterior, and $N=1+(d-1)(s+1)$.

\emph{Conclusion.} Under $\pi$ the law of the pair $(\tau_-,\tau_+)$ has a
density proportional to
\[
\mathcal W_-(\tau_-)\,\mathcal W_+(\tau_+)\,\varphi^{*N}(\Phi_-\Phi_+).
\]
Let $\mu_1$ and $\mu_2$ be the probability measures proportional to
$\mathcal W_-\lambda$ and to $\mathcal W_+\lambda$. Then the law has density
$k/Z'$ with respect to $\mu_1\otimes\mu_2$, where
$k=\varphi^{*N}(\Phi_-\Phi_+)$ takes values in
$[1-\varepsilon_N,1+\varepsilon_N]$, and Lemma~\ref{lem:nearproduct} gives the
bound $4\varepsilon_N$.

(ii) For $s\ge2d-1$, $N\ge1+2d(d-1)\ge d^2$, so $Nt\ge g^2d^2a^2\ge6$. Since
$\varphi^{*N}=p_{Nt}$, \eqref{eq:hk} gives $\varepsilon_N\le5e^{-3Nt/4}\le0.06$,
and (i) gives $\rho^\omega\le20e^{-3Nt/4}\le20e^{-3g^2d^2a^2/4}$.
\end{proof}

The heat-kernel weight has string tension $\frac34g^2$, and $Nt$ is $g^2$ times
the area of the middle slab up to boundary terms, so the bound reads
$\rho\le20e^{-\frac34g^2\cdot\mathrm{area}}$: in two dimensions the slab
decorrelation of Section~\ref{sec:massgen} is the confinement of the middle
slab, and it holds uniformly in the boundary condition with $c\to0$ as the
physical size $da$ of the box grows. The proof rests on one fact. After the
axial gauge the plaquettes straddling each seam contain a single free link, so
the middle slab integrates out exactly to a function $\varphi^{*N}(\Phi_-\Phi_+)$
of one group element, uniformly close to $1$ once the middle slab is confined.
In four dimensions the same gauge leaves on each seam the free links of a
three-dimensional layer, and integrating out the middle produces the
$(d-1)$-step transfer operator of the box, a kernel on of order $(md)^3$ copies
of $G$; the hypothesis $(\mathrm M_{d,m}(c))$ is a bound on its second singular
value on gauge-invariant functions, uniform in the boundary condition, and the
argument above has no four-dimensional counterpart.

\subsection{Two dimensions: thin strips}\label{ssec:strips}

The restriction of $(\mathrm M_{d,m}(c))$ to boxes all of whose widths are at
least $2d-1$ is needed. We show that for strips of width one the same bound
fails along the two-dimensional tuning for every $c<1$.

Identify $\SU(2)$ with the unit quaternions $g=g_0+g_1\mathbf i+g_2\mathbf j+g_3
\mathbf k$; then $\chi:=\chi_{1/2}$ satisfies $\chi(g)=\Tr g=2g_0$, $\chi(g^{-1})
=\chi(g)$, and $\ip gh_{\R^4}=\frac12\chi(g^{-1}h)$. For $\epsilon\in\{1,-1\}$
and $a,b,c>0$ with $T=a+b+c$ let $\mu^\epsilon_{a,b,c}$ be the probability
measure on $\SU(2)^2$ with density
\[
p_a(x)\,p_b(x^{-1}y)\,p_c(\epsilon y)/p_T(\epsilon)
\]
with respect to $\lambda\otimes\lambda$: a Brownian bridge on $\SU(2)$ from $1$
to $\epsilon$ observed at the times $a$ and $a+b$. The normalisation follows
from $\int p_a(x)p_b(x^{-1}y)\,dx=p_{a+b}(y)$ and $\int p_{a+b}(y)p_c(\epsilon y)
\,dy=p_T(\epsilon)$, both instances of Lemma~\ref{lem:merge}(i).

\begin{lemma}[the strip is a bridge]\label{lem:striplaw}
Let $Q=[0,3d-1)\times\{0\}$, $d\ge2$, split in the first coordinate, let
$\varphi=p_t$, and let $\omega_\epsilon$ be the exterior configuration with all
links equal to $1$ except $\sigma_{((3d-1,0),2)}=\epsilon$. Then $\mathcal G_Q$
is trivial, so $\mathcal A^{\omega_\epsilon}_Q=L^2(\nu^{\omega_\epsilon}_Q)$,
and with $h_x=h(x,0)$ and $v_x=v(x,0)$ the pair $(x,y)=(v_{d-1}^{-1}h_{d-1},
v_{2d-1}^{-1})$, formed by a function of $\sigma_{V(Q_-)}$ and a function of
$\sigma_{V(Q_+)}$, has law $\mu^\epsilon_{a,b,c}$ with $a=c=2dt$ and
$b=(2d-1)t$.
\end{lemma}

\begin{proof}
No vertex $v$ of $Q$ has $v-e_2\in Q$, so there are no interior vertices. The
plaquettes meeting $V(Q)$ are those at $(-1,0)$, with holonomy $v_0$; at
$(x,-1)$, with holonomy $h_x^{-1}$; and at $(x,0)$, $0\le x\le3d-2$, with
holonomy $h_xv_{x+1}v_x^{-1}$, where $v_{3d-1}:=\epsilon$. So the density of
$\nu$ is proportional to $p_t(v_0)\prod_xp_t(h_x)\prod_xp_t(h_xv_{x+1}v_x^{-1})$.
For $x\ne d-1$, $\int p_t(h)p_t(hg)\,dh=(p_t*p_t)(g)=p_{2t}(g)$ by
Lemma~\ref{lem:merge}(i) and $p_t(h^{-1})=p_t(h)$. Integrating then
$v_0,\dots,v_{d-2}$, $v_{2d},\dots,v_{3d-2}$ and $v_d,\dots,v_{2d-2}$ with
Lemma~\ref{lem:merge}(i), in the form $\int\alpha(gy^{-1})\beta(yh)\,dy=(\alpha*
\beta)(gh)$, leaves the density of $(v_{d-1},h_{d-1},v_{2d-1})$ proportional to
\[
p_{(2d-1)t}(v_{d-1})\,p_t(h_{d-1})\,p_{(2d-1)t}(v_{2d-1}v_{d-1}^{-1}h_{d-1})\,
p_{2dt}(\epsilon v_{2d-1}^{-1}).
\]
Substituting $x=v_{d-1}^{-1}h_{d-1}$ for $h_{d-1}$, which preserves Haar measure
for fixed $v_{d-1}$, and integrating $v_{d-1}$ gives $(p_{(2d-1)t}*p_t)(x)=
p_{2dt}(x)$. With $y=v_{2d-1}^{-1}$ the third factor is $p_{(2d-1)t}(y^{-1}x)=
p_b(x^{-1}y)$, and the last is $p_{2dt}(\epsilon y)$.
\end{proof}

\begin{lemma}[moments of the bridge]\label{lem:moments}
Under $\mu^\epsilon_{a,b,c}$, $E[x_1]=E[y_1]=0$ and
\[
E[x_1y_1]=\tfrac13\bigl(\tfrac12D-\tfrac14C\bigr),\quad E[x_1^2]=\tfrac13
\bigl(1-\tfrac14A\bigr),\quad E[y_1^2]=\tfrac13\bigl(1-\tfrac14B\bigr),
\]
where $D=E\chi(x^{-1}y)$, $C=E[\chi(x)\chi(y)]$, $A=E\chi(x)^2$ and $B=E\chi(
y)^2$ are given, with $q_\theta(n)=e^{-\theta(n^2-1)/4}$ and $Z=p_T(\epsilon)$, by
\begin{align*}
Z&=\tfrac12\textstyle\sum_{n\in\Z}\epsilon^{n-1}n^2q_T(n),\qquad
DZ=\sum_{n\in\Z}\epsilon^{n-1}n(n+1)q_{a+c}(n)q_b(n+1),\\
CZ&=\textstyle\sum_{n\in\Z}\epsilon^n\bigl[(n+1)^2q_a(n+1)q_c(n+1)+(n^2-1)
q_a(n+1)q_c(n-1)\bigr]q_b(n),\\
AZ&=2Z+\textstyle\sum_{n\in\Z}\epsilon^{n-1}n(n+2)q_a(n+2)q_{b+c}(n),\\
BZ&=2Z+\textstyle\sum_{n\in\Z}\epsilon^{n-1}n(n+2)q_c(n+2)q_{a+b}(n).
\end{align*}
\end{lemma}

\begin{proof}
\emph{Symmetry.} Conjugation $(x,y)\mapsto(gxg^{-1},gyg^{-1})$ preserves
$\mu^\epsilon$, since $\epsilon$ is central and the $p_\theta$ are class
functions. It fixes $x_0$ and $y_0$ and rotates $(x_1,x_2,x_3)$ and
$(y_1,y_2,y_3)$ by the same element of $\SO(3)$. Conjugation by $\mathbf j$
maps $\mathbf i$ to $-\mathbf i$, so $E[x_1]=E[y_1]=0$; conjugation by
$\frac12(1+\mathbf i+\mathbf j+\mathbf k)$ permutes $\mathbf i,\mathbf j,
\mathbf k$ cyclically, so $E[x_1y_1]=\frac13E\sum_kx_ky_k$ and $E[x_1^2]=\frac13
E\sum_kx_k^2$. Finally $\sum_kx_ky_k=\ip xy-x_0y_0=\frac12\chi(x^{-1}y)-\frac14
\chi(x)\chi(y)$ and $\sum_kx_k^2=1-\frac14\chi(x)^2$.

\emph{Reduction to character sums.} Write class functions as $\sum_{n\ge1}F_n
\chi_{(n-1)/2}$. Then $p_\theta$ has coefficients $nq_\theta(n)$, and
$p_\theta(\epsilon\,\cdot)$ has coefficients $\epsilon^{n-1}nq_\theta(n)$, since
$\chi_j(\epsilon g)=\epsilon^{2j}\chi_j(g)$; moreover $\int(\sum F_n\chi)(\sum
G_n\chi)=\sum F_nG_n$, and $(\alpha*p_\theta*\gamma)(1)=\sum_n\alpha_n\gamma_n
q_\theta(n)$ because $\chi_i*\chi_j=\delta_{ij}\chi_j/d_j$ and $\chi_j(1)=d_j$.
By the Clebsch--Gordan rule $\chi\chi_j=\chi_{j-1/2}+\chi_{j+1/2}$ for $j\ge
\frac12$ and $\chi\chi_0=\chi$, multiplication by $\chi$ acts on coefficients
by $F_n\mapsto F_{n+1}+F_{n-1}$ with $F_0:=0$. The marginal of $x$ is $p_a(x)
p_{b+c}(\epsilon x)/Z$ and that of $y$ is $p_{a+b}(y)p_c(\epsilon y)/Z$;
substituting $z=x^{-1}y$, the law of $z$ is $p_b(z)p_{a+c}(\epsilon z)/Z$; and
$C=((\chi p_a)*p_b*\gamma)(1)/Z$ with $\gamma(w)=\chi(w)p_c(\epsilon w)$. All
series converge absolutely. Hence, with $P^{(\theta)}_n=nq_\theta(n)$ and
$\alpha^{(\theta)}_n=P^{(\theta)}_{n+1}+P^{(\theta)}_{n-1}$,
\begin{gather*}
Z=\sum_{n\ge1}\epsilon^{n-1}n^2q_T(n),\quad DZ=\sum_{n\ge1}\epsilon^{n-1}
\alpha^{(b)}_n\,nq_{a+c}(n),\quad CZ=\sum_{n\ge1}\epsilon^n\alpha^{(a)}_n
\alpha^{(c)}_nq_b(n),\\
AZ=\sum_{n\ge1}\epsilon^{n-1}(\alpha^{(a)}_{n+1}+\alpha^{(a)}_{n-1})\,nq_{b+c}
(n),\quad BZ=\sum_{n\ge1}\epsilon^{n-1}(\alpha^{(c)}_{n+1}+\alpha^{(c)}_{n-1})\,
nq_{a+b}(n),
\end{gather*}
where for $CZ$ we used that the coefficients of $\gamma$ are $\sum_\pm
\epsilon^{n\pm1-1}P^{(c)}_{n\pm1}=\epsilon^n\alpha^{(c)}_n$, and $\alpha_0:=0$.

\emph{Symmetrisation.} Extend $P^{(\theta)}_n=nq_\theta(n)$ to $n\in\Z$; it is
odd, so $\alpha^{(\theta)}$, now defined for all $n$, is odd with
$\alpha^{(\theta)}_0=0$ and agrees with the above for $n\ge1$, and so is
$\alpha_{n+1}+\alpha_{n-1}$. Each summand above is then an even function of $n$
vanishing at $0$, as $\epsilon^{n\pm1}$ is even in $n$, so $\sum_{n\ge1}=\frac12
\sum_{n\in\Z}$. Expanding the products and using the substitution $n\mapsto-n$,
under which $P_{n-1}\mapsto-P_{n+1}$ and $P_{n-2}\mapsto-P_{n+2}$, to identify
the terms in pairs gives the stated formulas; for $AZ$ the middle term
$2P^{(a)}_n\cdot nq_{b+c}(n)$ gives $2Z$, since $q_aq_{b+c}=q_T$.
\end{proof}

\begin{lemma}[Poisson summation]\label{lem:poisson}
Let $0<T\le1$, $\lambda=T/4$, $\gamma=\sqrt{4\pi/T}$, $\mu\in\R$, $\widetilde
Q(m)=q_2m^2+q_1m+q_0$ real, and $f(x)=\widetilde Q(x+\mu)e^{-\lambda(x+\mu)^2}$.
\begin{enumerate}[label=\textup{(\roman*)},leftmargin=*]
\item $|\sum_{n\in\Z}f(n)-\gamma(q_0+2q_2/T)|\le30\gamma(|q_0|+4|q_1|/T+16|q_2|/
T^2)e^{-4\pi^2/T}$.
\item If $q_2=1$, $|q_1|\le2$ and $|q_0|\le1$, then $\sum_{n\in\Z}(-1)^{n-1}f(n)
=\Lambda_0(\cos\pi\mu-\mathcal E)$ with $\Lambda_0=\frac{8\pi^2}{T^2}\gamma
e^{-\pi^2/T}$ and $|\mathcal E|\le0.22\,T$.
\end{enumerate}
\end{lemma}

\begin{proof}
$f$ is a Schwartz function. With $\widehat f(\xi)=\int f(x)e^{-2\pi i\xi x}dx$,
the substitution $y=x+\mu$ and the formulas $\int e^{-\lambda y^2-i\omega y}dy=
\gamma e^{-\omega^2/(4\lambda)}=:G(\omega)$, $\int ye^{-\lambda y^2-i\omega y}dy
=iG'(\omega)$, $\int y^2e^{-\lambda y^2-i\omega y}dy=-G''(\omega)$, $G'=-\frac
\omega{2\lambda}G$ and $G''=(\frac{\omega^2}{4\lambda^2}-\frac1{2\lambda})G$ give
\[
\widehat f(\xi)=e^{2\pi i\xi\mu}\gamma e^{-\pi^2\xi^2/\lambda}\Bigl[q_0-iq_1
\frac{\pi\xi}\lambda+q_2\Bigl(\frac1{2\lambda}-\frac{\pi^2\xi^2}{\lambda^2}
\Bigr)\Bigr].
\]
(i) By the Poisson summation formula~\cite[Ch.~5]{SS03}, $\sum_nf(n)=\sum_k
\widehat f(k)$, and $\widehat f(0)=\gamma(q_0+2q_2/T)$. For $k\ne0$, since
$\lambda\le\frac14$, each of $1$, $\pi|k|/\lambda$ and $\frac1{2\lambda}+
\frac{\pi^2k^2}{\lambda^2}$ is at most $M_k\lambda^{-e}$ with $e=0,1,2$
respectively and $M_k=1+\pi|k|+\pi^2k^2$; so $|\widehat f(k)|\le\gamma e^{-\pi^2
k^2/\lambda}M_k(|q_0|+|q_1|/\lambda+|q_2|/\lambda^2)$, and $\sum_{k\ne0}e^{-\pi^2
k^2/\lambda}M_k\le2e^{-\pi^2/\lambda}(M_1+\sum_{k\ge2}M_ke^{-12\pi^2(k-1)})\le
30e^{-\pi^2/\lambda}$, using $k^2-1\ge3(k-1)$ and $M_1<14.1$. Put $\lambda=T/4$.

(ii) Poisson summation for $f(x)e^{i\pi x}$ gives $\sum_n(-1)^nf(n)=\sum_{\xi\in
\frac12+\Z}\widehat f(\xi)$. The terms $\xi=\pm\frac12$ are complex conjugate
and add up to $2\gamma e^{-\pi^2/(4\lambda)}[\cos(\pi\mu)(q_0+\frac1{2\lambda}-
\frac{\pi^2}{4\lambda^2})+\sin(\pi\mu)\frac{q_1\pi}{2\lambda}]$. Hence
$\sum(-1)^{n-1}f(n)=\Lambda_0(\cos\pi\mu-\mathcal E)$ with $\Lambda_0=\frac{\pi^2}
{2\lambda^2}\gamma e^{-\pi^2/(4\lambda)}$, the stated value, and $\mathcal E=
\mathcal E_1+\mathcal E_2$, where $\mathcal E_1=\frac{2\lambda^2}{\pi^2}[\cos(\pi
\mu)(q_0+\frac1{2\lambda})+\sin(\pi\mu)\frac{q_1\pi}{2\lambda}]$ and $\mathcal
E_2=\Lambda_0^{-1}\sum_{|\xi|\ge3/2}\widehat f(\xi)$. First, $|\mathcal E_1|\le
\frac2{\pi^2}(\lambda^2+\frac\lambda2+\pi\lambda)\le\frac2{\pi^2}(\frac14+\frac12+
\pi)\lambda<0.79\lambda$. Second, for $|\xi|\ge\frac32$, $|\widehat f(\xi)|/
\Lambda_0\le e^{-\pi^2(\xi^2-1/4)/\lambda}(\frac{2\lambda^2}{\pi^2}+\frac{4\lambda
|\xi|}\pi+\frac\lambda{\pi^2}+2\xi^2)\le3\xi^2e^{-\pi^2(\xi^2-1/4)/\lambda}$, and
with $\xi=\pm(k+\frac12)$, $\xi^2-\frac14=k^2+k\ge2k$, so $|\mathcal E_2|\le6
\sum_{k\ge1}(k+\frac12)^2e^{-2\pi^2k/\lambda}\le14e^{-2\pi^2/\lambda}\le56
e^{-8\pi^2}\lambda$, as $e^{-2\pi^2/\lambda}/\lambda$ increases on $(0,\frac14]$.
So $|\mathcal E|\le0.8\lambda\le0.22\,T$.
\end{proof}

\begin{lemma}[the bridge from $1$ to $1$]\label{lem:bridge1}
Let $0<T\le1$ and suppose that $\kappa:=ac/T$, $\tau_x:=a(b+c)/T$ and $\tau_y:=
c(a+b)/T$ are all at least $T/10$. Under $\mu^1_{a,b,c}$,
\[
\operatorname{Corr}(x_1,y_1)\ge r_G(1-T),\qquad r_G=\Bigl(\frac{ac}{(a+b)(b+c)}
\Bigr)^{1/2}.
\]
\end{lemma}

Here $r_G$ is the correlation of a Euclidean Brownian bridge on $[0,T]$ at the
times $a$ and $a+b$.

\begin{proof}
\emph{Main terms.} Complete the squares in the exponents of
Lemma~\ref{lem:moments} with $\epsilon=1$: $q_T(n)=e^{T/4}e^{-Tn^2/4}$;
$q_{a+c}(n)q_b(n+1)=e^{E_D}e^{-T(n+\mu_D)^2/4}$ with $\mu_D=b/T$ and $E_D=
\frac{b^2}{4T}+\frac{a+c}4$; $q_a(n+1)q_c(n+1)q_b(n)=e^{E_1}e^{-T(n+\mu_1)^2/4}$
with $\mu_1=\frac{a+c}T$ and $E_1=\frac{(a+c)^2}{4T}+\frac b4$; $q_a(n+1)q_c(n-1)
q_b(n)=e^{E_2}e^{-T(n+\mu_2)^2/4}$ with $\mu_2=\frac{a-c}T$ and $E_2=\frac{(a-c)^2}
{4T}+\frac b4$; and $q_a(n+2)q_{b+c}(n)=e^{E_A}e^{-T(n+\mu_A)^2/4}$ with $\mu_A=
\frac{2a}T$ and $E_A=\frac{a^2}T-\frac{3a-b-c}4$. In the variable $m=n+\mu$ the
polynomial factors become $m^2$ for $Z$, $m^2+(1-2\mu_D)m+\mu_D^2-\mu_D$ for
$D$, $(m+1-\mu_1)^2$ with $1-\mu_1=b/T$ and $(m-\mu_2)^2-1$ for $C$, and
$(m-\mu_A)(m-\mu_A+2)$ for $A$. The main terms $\gamma(q_0+2q_2/T)$ of
Lemma~\ref{lem:poisson}(i), divided by the main term $\gamma e^{T/4}/T$ of $Z$,
give, after the simplifications $T-b=a+c$, $(a-c)^2-T(a+c)=-4ac-b(a+c)$ and
$a^2/T-a=-\tau_x$,
\begin{gather*}
D_0=e^{-u}(2-4u),\qquad C_0=e^{-u}\Bigl(2+\frac{b^2}T\Bigr)+e^{-u-\kappa}
\Bigl(2+\frac{(a-c)^2}T-T\Bigr),\\
A_0=2+e^{-\tau_x}(2-4\tau_x),\qquad B_0=2+e^{-\tau_y}(2-4\tau_y),\qquad
u=\frac{b(a+c)}{4T}.
\end{gather*}
Put $P=\frac{b(T+a+c)}{4T}$. Using $\frac{(a-c)^2-T^2}{4T}=-\frac{(2a+b)(2c+b)}
{4T}$ and $\frac{(2a+b)(2c+b)}{4T}=P+\kappa$, one finds
\[
X_0:=\tfrac12D_0-\tfrac14C_0=e^{-u}\bigl[(\tfrac12-P)(1-e^{-\kappa})+\kappa
e^{-\kappa}\bigr],\qquad V^x_0:=1-\tfrac14A_0=\tfrac12(1-e^{-\tau_x})+\tau_x
e^{-\tau_x}.
\]
Since $ac\le T^2/4$ and $b(a+c)\le T^2/4$, we have $\kappa\le T/4$, $u\le T/16$
and $P\le b/2\le T/2\le\frac12$. With $1-e^{-\kappa}\ge\kappa-\kappa^2/2$,
$e^{-\kappa}\ge1-\kappa$ and $e^{-u}\ge1-u$,
\[
X_0\ge(1-u)\kappa\bigl(\tfrac32-P-\tfrac54\kappa\bigr)\ge\tfrac32\kappa\bigl(1-
\tfrac T{16}\bigr)\bigl(1-\tfrac{13}{24}T\bigr)\ge\tfrac32\kappa\bigl(1-\tfrac{29}
{48}T\bigr).
\]
Also $V^x_0\le\frac32\tau_x$ and $V^y_0:=1-\frac14B_0\le\frac32\tau_y$.

\emph{Errors.} Each of the five sums is its main term times $1+\theta$ with
$|\theta|\le\eta:=750e^{-4\pi^2/T}/T$. This follows from
Lemma~\ref{lem:poisson}(i), with $T\le1$, the coefficient bounds $q_2=1$,
$|q_1|\le2$ and $|q_0|\le1$ read off above, and the lower bounds $2/T-1/4$,
$2/T$, $2/T-1$ and $2/T-1\ge1/T$ for the brackets $q_0+2q_2/T$ of the main
terms: the relative errors are at most $240e^{-4\pi^2/T}/T$ for $Z$ and $750
e^{-4\pi^2/T}/T$ for the others. As $e^{-4\pi^2/T}/T^3$ increases on $(0,1]$,
$\eta\le750e^{-4\pi^2}T^2<10^{-14}T^2$. Since $|D_0|\le2$, both parts of $C_0$
are at most $3$ in absolute value, and $0\le A_0-2\le2$, we get $|D-D_0|\le6
\eta$, $|C-C_0|\le18\eta$ and $|A-A_0|,|B-B_0|\le6\eta$, hence $|X-X_0|\le7.5
\eta$ and $|V^x-V^x_0|,|V^y-V^y_0|\le1.5\eta$, where $X=\frac12D-\frac14C$,
$V^x=1-\frac14A$ and $V^y=1-\frac14B$.

\emph{Conclusion.} By Lemma~\ref{lem:moments}, $\operatorname{Corr}(x_1,y_1)=X/
(V^xV^y)^{1/2}$. Using $\kappa,\tau_x,\tau_y\ge T/10$,
\[
\operatorname{Corr}(x_1,y_1)\ge\frac{\frac32\kappa(1-\frac{29}{48}T-50\eta/T)}
{\frac32(\tau_x\tau_y)^{1/2}(1+10\eta/T)}\ge r_G\Bigl(1-\frac{29}{48}T-60\frac
\eta T\Bigr)\ge r_G(1-T),
\]
since $\kappa/(\tau_x\tau_y)^{1/2}=r_G$ and $60\eta/T<10^{-12}T$.
\end{proof}

\begin{lemma}[the bridge from $1$ to $-1$]\label{lem:bridgem1}
Let $s_0>0$ and $0<T\le\min(1,s_0^2/4)$, and suppose $\sin(\pi a/T)\ge s_0$ and
$\sin(\pi c/T)\ge s_0$. Under $\mu^{-1}_{a,b,c}$,
\[
\operatorname{Corr}(x_1,y_1)\ge1-\frac{5T}{2s_0^2}.
\]
\end{lemma}

\begin{proof}
Write $r_a=a/T$, $r_b=b/T$, $r_c=c/T$, $s_a=\sin\pi r_a$ and $s_c=\sin\pi r_c$.
The sums of Lemma~\ref{lem:moments} with $\epsilon=-1$ have the exponents and
polynomial factors of the proof of Lemma~\ref{lem:bridge1}, all with $q_2=1$,
$|q_1|\le2$ and $|q_0|\le1$, and carry the sign $(-1)^{n-1}$, except $CZ$, which
carries $(-1)^n$. Lemma~\ref{lem:poisson}(ii) gives
\begin{gather*}
Z=\tfrac12e^{T/4}\Lambda_0(1-\mathcal E_Z),\qquad D=2e^{-u}\frac{\cos\pi r_b-
\mathcal E_D}{1-\mathcal E_Z},\\
C=-2\,\frac{e^{-u}(\cos\pi(r_a+r_c)-\mathcal E_1)+e^{-u-\kappa}(\cos\pi(r_a-r_c)
-\mathcal E_2)}{1-\mathcal E_Z},\\
A=2+2e^{-\tau_x}\frac{\cos2\pi r_a-\mathcal E_A}{1-\mathcal E_Z},\qquad
B=2+2e^{-\tau_y}\frac{\cos2\pi r_c-\mathcal E_B}{1-\mathcal E_Z},
\end{gather*}
with all $|\mathcal E_\bullet|\le0.22T$ and $u,\kappa,\tau_x,\tau_y$ as before;
the exponential prefactors are those of the main terms in the proof of
Lemma~\ref{lem:bridge1}. Each quotient has the form $w(\cos\theta-\mathcal E)/(1-
\mathcal E_Z)$ with $w\in[1-T/4,1]$, as $u+\kappa,\tau_x,\tau_y\le T/4$; since
$|1/(1-\mathcal E_Z)|\le1.29$ and $|1/(1-\mathcal E_Z)-1|\le0.29T$, it differs
from $\cos\theta$ by at most $0.25T\cdot1.29+0.29T+0.22T\cdot1.29<T$. Hence
\[
\bigl|\tfrac D2-\cos\pi r_b\bigr|\le T,\quad\bigl|\tfrac C4+\cos\pi r_a\cos\pi
r_c\bigr|\le T,\quad\bigl|\tfrac A4-\cos^2\pi r_a\bigr|\le\tfrac T2,\quad\bigl|
\tfrac B4-\cos^2\pi r_c\bigr|\le\tfrac T2,
\]
using $\frac12(\cos\pi(r_a+r_c)+\cos\pi(r_a-r_c))=\cos\pi r_a\cos\pi r_c$. Since
$r_b=1-r_a-r_c$, $\cos\pi r_b=-\cos\pi(r_a+r_c)$, so $\cos\pi r_b+\cos\pi r_a
\cos\pi r_c=s_as_c$. Thus $X=\frac12D-\frac14C\ge s_as_c-2T>0$, and $0<V^x=1-
\frac14A\le s_a^2+\frac T2$ and $0<V^y\le s_c^2+\frac T2$, positivity coming
from Lemma~\ref{lem:moments} and $T\le s_0^2/4$. By Lemma~\ref{lem:moments} and
$\sqrt{(1+x)(1+y)}\le1+\frac{x+y}2$,
\[
\operatorname{Corr}(x_1,y_1)\ge\frac{s_as_c-2T}{s_as_c(1+\frac T{4s_a^2}+\frac T
{4s_c^2})}\ge\frac{1-2T/s_0^2}{1+T/(2s_0^2)}\ge1-\frac{5T}{2s_0^2}.
\qedhere
\]
\end{proof}

\begin{proposition}[thin strips]\label{prop:strip}
Let $G=\SU(2)$, $\varphi=p_t$, $d\ge2$, let $Q=[0,3d-1)\times\{0\}$ be split in
the first coordinate into widths $d$, $d-1$, $d$, put $T=(6d-1)t$, and let
$\omega_{\pm1}$ be as in Lemma~\ref{lem:striplaw}.
\begin{enumerate}[label=\textup{(\roman*)},leftmargin=*]
\item If $T\le1$, then $\rho^{\omega_1}(Q_-,Q_+)\ge\frac{2d}{4d-1}(1-T)$.
\item If $T\le\frac3{16}$, then $\rho^{\omega_{-1}}(Q_-,Q_+)\ge1-\frac{10}3T$.
\item Along the tuning $t_k=g^2a_k^2$, $a_k\to0$, with integers $d_k\ge2$ and
$d_ka_k\le\ell$, for every $m\ge4$ and every $c<1$ the bound $\rho^\omega(Q_-,
Q_+)\le c$ fails for all large $k$ for the strip of widths $3d_k-1$ and $1$,
both at most $md_k-1$; already the trivial boundary condition gives $\rho\ge
\frac12(1-6g^2\ell a_k)$.
\end{enumerate}
\end{proposition}

\begin{proof}
By Lemma~\ref{lem:striplaw} the functions $X=x_1$ and $Y=y_1$ of
$\sigma_{V(Q_-)}$ and $\sigma_{V(Q_+)}$ lie in $\mathcal A^{\omega_\epsilon}_Q$,
and their joint law is that of $\mu^\epsilon_{a,b,c}$ with $a=c=2dt$ and $b=(2d-
1)t$; so $\rho^{\omega_\epsilon}\ge|\operatorname{Corr}(x_1,y_1)|$. (i) Here
$\kappa/T=4d^2/(6d-1)^2\ge\frac19$, $\tau_x/T=\tau_y/T=2d(4d-1)/(6d-1)^2\ge\frac15$
for $d\ge2$, and $r_G=2d/(4d-1)$; apply Lemma~\ref{lem:bridge1}. (ii) $a/T=c/T=
2d/(6d-1)\in(\frac13,\frac12)$, so $\sin(\pi a/T)\ge\sin\frac\pi3=\frac{\sqrt3}2=
:s_0$ with $s_0^2/4=\frac3{16}$; apply Lemma~\ref{lem:bridgem1}. (iii)
$T_k=(6d_k-1)t_k\le6g^2(d_ka_k)a_k\le6g^2\ell a_k\to0$, so for large $k$, by
(ii), $\rho^{\omega_{-1}}\ge1-\frac{10}3T_k>c$, and by (i), $\rho^{\omega_1}\ge
\frac{2d_k}{4d_k-1}(1-T_k)\ge\frac12(1-6g^2\ell a_k)$.
\end{proof}

The mechanism is a bridge. The fixed boundary pins the holonomy of the loop
surrounding the strip, and both outer parts read off the same collective
variable. For the condition $\omega_{-1}$ the measure is invariant under the
global conjugation $\sigma\mapsto g\sigma g^{-1}$ of all links of $V(Q)$, the
exterior links being central, and part (ii) shows that this global $\SU(2)$
symmetry is broken at small $T$: the bridge from $1$ to $-1$ concentrates near
the two-sphere of minimising geodesics, whose point is read off by both outer
parts, and the slab correlation tends to one. Theorem~\ref{thm:twod} and
Proposition~\ref{prop:strip} together show that in two dimensions, along the
weak-coupling tuning of a non-abelian theory, the slab hypothesis holds for fat
boxes uniformly in the boundary condition with $c\to0$, and fails for strips
of width one for every $c<1$. In four dimensions the boundary conditions with
values in the centre are 't~Hooft's twisted ones~\cite{tH79}, and for fat boxes
they are the natural test of $(\mathrm M_{d,m}(c))$; we have no estimate for
them along the four-dimensional tuning.

\subsection{Curvature}\label{ssec:curvlimits}

Theorem~\ref{thm:A} is a curvature argument on the gauge-fixed measure, and its
range is limited by the Hessian bound of Theorem~\ref{thm:hess}. We show that
no curvature argument applied to the Wilson measure itself, with or without the
gauge, can be uniform along the weak-coupling tuning. We use the notation of
Section~\ref{sec:curvature}: $G=\SU(N)$ with the metric $\ip XY=-\Tr(XY)$,
tangent vectors $v=(X_eQ_e)$, the transported vectors $Y_1,\dots,Y_4$ of a
plaquette, and $S_0=\sum_p\ReTr Q_p$. For a probability measure with density
proportional to $e^{-V}$ on a compact Riemannian manifold, the
\emph{Bakry--\'Emery constant} is $\inf(\Ric+\Hess V)(v,v)/|v|^2$ over all
points and nonzero tangent vectors; by Bochner's formula it is the largest $K$
with $\Gamma_2(f)\ge K\Gamma(f)$ pointwise for all smooth $f$, and it controls
the Poincar\'e and logarithmic Sobolev constants of the Langevin
dynamics~\cite{BGL14}. For a plaquette put $Z_p=Y_1+Y_2-Y_3-Y_4$.

\begin{lemma}[central plaquettes]\label{lem:central}
Let $Q$ be a field and $p$ a plaquette with $Q_p=wI$ central. Along the
geodesics $Q_e(t)=e^{tX_e}Q_e$,
\[
\frac{d^2}{dt^2}\Big|_{t=0}\ReTr Q_p(t)=-\operatorname{Re}(w)\,|Z_p|^2 .
\]
If $\varphi$ is a smooth class function on $G$, then $\frac{d^2}{dt^2}|_0
\varphi(Q_p(t))=H_w(Z_p,Z_p)$ for an $\operatorname{Ad}$-invariant quadratic
form $H_w$ on $\g$; for $G=\SU(2)$, $H_{\pm1}=h_\pm|\cdot|^2$ with real constants
$h_\pm$.
\end{lemma}

\begin{proof}
As in Section~\ref{sec:curvature}, $Q_p(t)=e^{tY_1}e^{tY_2}We^{-tY_3}e^{-tY_4}$
with $W=Q_p$. Since $W=wI$ is central, $Q_p(t)=w\,c(t)$ with $c(t)=e^{tY_1}
e^{tY_2}e^{-tY_3}e^{-tY_4}$. With $a=(1,1,-1,-1)$, $c''(0)=\sum_ia_i^2Y_i^2+2
\sum_{i<j}a_ia_jY_iY_j$, and as $\Tr(Y_iY_j)=\Tr(Y_jY_i)$, $\Tr c''(0)=\Tr(Z_p^2)
=-|Z_p|^2$. Hence $\frac{d^2}{dt^2}\ReTr Q_p=\operatorname{Re}(w\Tr Z_p^2)=
-\operatorname{Re}(w)|Z_p|^2$, the trace of a product of two skew-Hermitian
matrices being real. For $\varphi$, the curve $\gamma(t)=Q_p(t)$ has
$\gamma(0)=wI$ and $\gamma'(0)=Z_pwI$, and $(\varphi\circ\gamma)''(0)=\Hess
\varphi(\gamma'(0),\gamma'(0))+d\varphi(\nabla_t\gamma'(0))$. The gradient of
$\varphi$ at $wI$ is $YwI$ with $Y$ fixed by every $\operatorname{Ad}_g$, because
conjugation by $g$ is an isometry that fixes $wI$ and preserves $\varphi$; the
centre of $\mathfrak{su}(N)$ is trivial, so $Y=0$ and the second term vanishes.
The same isometries show that $H_w(Z,Z)=\Hess\varphi_{wI}(ZwI,ZwI)$ is
$\operatorname{Ad}$-invariant, and $\operatorname{Ad}$ is irreducible on
$\mathfrak{su}(2)$, so $H_{\pm1}$ is a multiple of the metric.
\end{proof}

The configuration
\begin{equation}\label{eq:piflux}
U^\pi_{x,\mu}=(-1)^{\sum_{\nu<\mu}x_\nu}I
\end{equation}
has every plaquette equal to $-I$: for $\mu<\nu$ the link $(x+e_\mu,\nu)$ sees
the coordinate $x_\mu$ raised by one and the link $(x+e_\nu,\mu)$ does not, so
the four signs of the plaquette at $x$ multiply to $-1$. It is defined on every
box and on $\Lambda_L$ for even $L$, and lies in $\SU(N)$ for even $N$. Put
$s(x)=(-1)^{|x|}$ with $|x|=\sum_\mu x_\mu$.

\begin{theorem}[full space]\label{thm:fullspace}
Let $N$ be even, $L\ge4$ even, and $\mu_L$ the $\SU(N)$ Wilson measure on
$\Lambda_L$ with density proportional to $e^{N\beta S_0}$. Its Bakry--\'Emery
constant satisfies
\[
K_{\rm full}\le\frac N2-16N|\beta|=\frac N2-\frac{16|\beta_W|}N .
\]
For $\SU(2)$, $K_{\rm full}\le1-8|\beta_W|$.
\end{theorem}

\begin{proof}
By Lemma~\ref{lem:ricci} and Table~\ref{tab:groups}, $\Ric=\frac N2\ip\cdot\cdot$.
Let $c\in\R^4$ with $\sum_\mu c_\mu=0$ and $c\ne0$, let $E\in\mathfrak{su}(N)$
with $|E|=1$, and put $X_{x,\mu}=c_\mu s(x)E$. For $\beta\ge0$ take the field
$U^\pi$, at which all transports are $\operatorname{Ad}_{\pm I}=\mathrm{id}$.
For the plaquette at $x$ in the plane $\mu<\nu$, $Y_1=c_\mu s(x)E$,
$Y_2=-c_\nu s(x)E$, $Y_3=-c_\mu s(x)E$ and $Y_4=c_\nu s(x)E$, so $Z_p=2s(x)(c_\mu-
c_\nu)E$. By Lemma~\ref{lem:central} with $w=-1$,
\[
\Hess S_0(v,v)=\sum_p|Z_p|^2=4L^4\sum_{\mu<\nu}(c_\mu-c_\nu)^2=4L^4\Bigl(4\sum_\mu
c_\mu^2-\Bigl(\sum_\mu c_\mu\Bigr)^2\Bigr)=16|v|^2 .
\]
With $V=-N\beta S_0$ this gives $(\Ric+\Hess V)(v,v)=(\frac N2-16N\beta)|v|^2$.
For $\beta<0$ take the field $Q\equiv I$, where Lemma~\ref{lem:central} with
$w=1$ gives $\Hess S_0(v,v)=-16|v|^2$ for the same $v$, hence $(\Ric+\Hess V)(v,
v)=(\frac N2-16N|\beta|)|v|^2$. Finally $\beta_W=N^2\beta$, and $\beta_W=4\beta$
for $\SU(2)$.
\end{proof}

Without the gauge the bound $|\Hess S_0(v,v)|\le16|v|^2$ holds for every field
and vector~\cite{BB26a}, so Theorem~\ref{thm:fullspace} is an equality,
$K_{\rm full}=\frac N2-16N|\beta|$. The Ricci curvature grows with $N$, but so
does the Hessian, and the criterion $K_{\rm full}>0$ holds exactly for
$|\beta_W|<N^2/32$; at fixed $N$ along the tuning $\beta_W\to\infty$ the
constant tends to $-\infty$ linearly. The Bakry--\'Emery constant is also the
best rate of contraction of the Langevin dynamics in the quadratic Kantorovich
distance~\cite{vRS05}, so no synchronous coupling of the dynamics does better.
The gauge of Section~\ref{sec:linefield} improves the constant $16$ to
$\frac{29}2$, and the next theorem shows that it cannot improve it below
$\frac{38}3$.

\begin{theorem}[the gauge slice]\label{thm:slice}
Let $G=\SU(2)$, $m\ge1$, $\Lambda_m=[0,2m]^4\cap\Z^4$ with free boundary, and
let $K_F(\Lambda_m)$ be the Bakry--\'Emery constant of the gauge-fixed measure
$\mu^F_{\Lambda_m}$ of Lemma~\ref{lem:gaugefix}. Then
\[
K_F(\Lambda_m)\le1-2|\beta|\rho_m,\qquad\rho_m=\frac{304\,m^4}{24\,m^4+8(2m+1)^3}
\xrightarrow[m\to\infty]{}\frac{38}3 .
\]
Hence the constant $\frac{29}2$ of Theorem~\ref{thm:hess} cannot be replaced by
any number below $\frac{38}3$, and asymptotically in $m$
\[
1-\tfrac{29}4|\beta_W|\le K_F\le1-\tfrac{19}3|\beta_W| ,
\]
so no Bakry--\'Emery argument for the gauge-fixed measure reaches $|\beta_W|=
\frac3{19}$.
\end{theorem}

\begin{proof}
\emph{The field.} For $\beta\ge0$ start from $U^\pi$ on $\Lambda_m$ and apply
the gauge transformation of the proof of Lemma~\ref{lem:gaugefix} along the
forest $F\cap E(\Lambda_m)$; the recursion $g_{x_{i+1}}=g_{x_i}U_{(x_i,x_{i+1})}$
produces $g_x\in\{\pm I\}$, so the result $Q$ takes values in $\{\pm I\}$,
equals $I$ on $F$, and has all plaquettes conjugate to those of $U^\pi$, that is
$-I$. For $\beta<0$ take $Q\equiv I$. For $v$ vanishing on $F$ the Hessian of
the gauge-fixed action is $\Hess S_{\Lambda_m}(v,v)$, which by
Lemma~\ref{lem:central} equals $\pm\sum_{p\in P(\Lambda_m)}|Z_p|^2$, with $+$ for
$\beta\ge0$ and $-$ for $\beta<0$; all transports are trivial.

\emph{The vector.} Let $c=(1,-1,0,0)$, $|E|=1$, and put $X_{x,\mu}=c_\mu s(x)E$
on the active links of $E(\Lambda_m)$ and $X=0$ on $F$. Let $p$ be a plaquette at
$x$ in the plane $\mu<\nu$ with all its links in $E(\Lambda_m)$; it contains
exactly one link of $F$. Without that link the computation of the proof of
Theorem~\ref{thm:fullspace} gives $Z_p=2s(x)(c_\mu-c_\nu)E$, the links $e_1,e_3$
of direction $\mu$ contributing $s(x)c_\mu E$ each and the links $e_2,e_4$ of
direction $\nu$ contributing $-s(x)c_\nu E$ each. Removing the link of $F$
gives $Z_p=s(x)(c_\mu-2c_\nu)E$ if that link has direction $\mu$, and $Z_p=
s(x)(2c_\mu-c_\nu)E$ if it has direction $\nu$.

\emph{Counting.} Let $B=[0,2m-1]^4$; every plaquette with base point in $B$ has
its links in $E(\Lambda_m)$. Fix $\mu<\nu$. The link of $F$ of the plaquette at
$x$ has direction $\mu$ exactly when $\sigma(x)=\mu$ or $\sigma(x+e_\nu)=\mu$,
and the two cases exclude each other. By Table~\ref{tab:sigma}, $\sigma$ has
period two and takes each value at four of the sixteen residues. $B$ and
$B+e_\nu$ are disjoint unions of $m^4$ translates of $\{0,1\}^4$, each
containing every residue once, so each case occurs for $4m^4$ base points in
$B$; of the $16m^4$ plaquettes of the plane with base point in $B$, exactly
$8m^4$ have their link of $F$ in direction $\mu$ and $8m^4$ in direction $\nu$.
Hence, with $|c|^2=2$ and $\sum_\mu c_\mu=0$,
\[
\sum_{p\in P(\Lambda_m)}|Z_p|^2\ge8m^4\sum_{\mu<\nu}\bigl[(c_\mu-2c_\nu)^2+
(2c_\mu-c_\nu)^2\bigr]=8m^4\Bigl(19|c|^2-4\Bigl(\sum_\mu c_\mu\Bigr)^2\Bigr)=
304m^4,
\]
using $\sum_{\mu<\nu}(5c_\mu^2+5c_\nu^2-8c_\mu c_\nu)=15|c|^2-4((\sum_\mu c_\mu
)^2-|c|^2)$. On the other side, for $\mu\in\{0,1\}$ the links $(x,\mu)$ with
$x\in B$ lie in $E(\Lambda_m)$, and exactly $16m^4-4m^4=12m^4$ of them are
active; they contribute $24m^4$ to $|v|^2$. The remaining links of
$E(\Lambda_m)$ with $X\ne0$ have direction $0$ or $1$ and base point in
$[0,2m]^4\setminus B$, at most $(2m+1)^4-(2m)^4\le4(2m+1)^3$ in each direction,
and each contributes at most $1$. So $|v|^2\le24m^4+8(2m+1)^3$, and
$|\Hess S_{\Lambda_m}(v,v)|\ge\rho_m|v|^2$ with the sign that gives $\Hess V(v,
v)\le-2|\beta|\rho_m|v|^2$, $V=-2\beta S_{\Lambda_m}$.

\emph{Conclusion.} $(\Ric+\Hess V)(v,v)\le(1-2|\beta|\rho_m)|v|^2$. If
Theorem~\ref{thm:hess} held with a constant $c'<\frac{38}3$, then $\rho_m\le c'$
for all $m$, contradicting $\rho_m\to\frac{38}3$. The lower bound for $K_F$ is
Lemma~\ref{lem:weighted} with $\lambda=0$, and $\beta_W=4\beta$.
\end{proof}

\begin{proposition}[plaquette actions at a coarse scale]\label{prop:coarse}
Let $L\ge4$ be even and $\nu(dU)\propto\exp(\sum_p\phi(U_p))\prod_edU_e$ on
$\SU(2)^{E(\Lambda_L)}$ with $\phi(U)=g(\frac12\ReTr U)$, $g\in C^2([-1,1])$. The
Bakry--\'Emery constant $K$ of $\nu$ satisfies
\[
K\le1-8\max\bigl(g'(-1),-g'(1)\bigr).
\]
\end{proposition}

\begin{proof}
$\frac12\ReTr(\pm e^{tZ})=\pm(1-\frac{t^2}4|Z|^2)+O(t^3)$, since $\Tr Z^2=
-|Z|^2$; so in Lemma~\ref{lem:central} $h_+=-\frac12g'(1)$ and $h_-=\frac12
g'(-1)$. The potential is $-\sum_p\phi(U_p)$. At $U^\pi$, with the vector of
the proof of Theorem~\ref{thm:fullspace}, $\Hess\sum_p\phi(U_p)(v,v)=h_-\sum_p
|Z_p|^2=16h_-|v|^2$, and at $Q\equiv I$ it is $16h_+|v|^2$. Hence $K\le1-16
h_\pm$.
\end{proof}

For the Wilson weight $g(s)=4\beta s$ this is Theorem~\ref{thm:fullspace}. For
an effective plaquette action at a coarse scale the Bakry--\'Emery route needs
$|g'(\pm1)|<\frac18$, an effective theory deep in strong coupling, and producing
such an effective measure from the weak-coupling lattice theory, and
transporting the resulting estimate back, are the two missing steps of
Section~\ref{sec:weak}.

\subsection{Entropy, link covariances and blocks of fixed size}
\label{ssec:entropy}

In the setting of Section~\ref{sec:cube}, say that $\mu$ satisfies \emph{block
entropy factorisation} with constant $C$ if $\operatorname{Ent}_\mu(F)\le C
n^{-4}\sum_x\mu(\operatorname{Ent}_{\gamma^0_{V_x}}F)$ for all bounded
gauge-invariant $F\ge0$, where $\operatorname{Ent}_{\gamma^0_{V_x}}F$ is the
entropy of $F$ under $\gamma^0_{V_x}(\cdot\mid U)$.

\begin{proposition}\label{prop:entbp}
Block entropy factorisation with constant $C$ implies $\BP(n,1/C)$.
\end{proposition}

\begin{proof}
Let $\Phi(s)=s\log s$. For $|u|\le\frac12$, $\Phi(1+u)=u+\frac{u^2}2+r(u)$ with
$r(u)=\sum_{k\ge3}(-u)^k/(k(k-1))$ and $|r(u)|\le|u|^3\sum_{k\ge3}2^{3-k}/(k(k-1))
\le|u|^3$. Let $g$ be real, bounded and gauge invariant, and $0<\varepsilon\le1/
(2\|g\|_\infty)$. Then $\operatorname{Ent}_\mu(1+\varepsilon g)=\mu\Phi(1+
\varepsilon g)-\Phi(1+\varepsilon\mu g)=\frac{\varepsilon^2}2\Var_\mu g+R$ with
$|R|\le2\varepsilon^3\|g\|_\infty^3$, and in the same way
$\operatorname{Ent}_{\gamma^0_{V_x}}(1+\varepsilon g)=\frac{\varepsilon^2}2
\Var_{\gamma^0_{V_x}}g+R_x$ with $|R_x|\le2\varepsilon^3\|g\|_\infty^3$. Inserting
$F=1+\varepsilon g$, dividing by $\varepsilon^2/2$ and letting $\varepsilon\to0$
gives $\Var_\mu g\le Cn^{-4}\sum_x\mu(\Var_{\gamma^0_{V_x}}g)=Cn^{-4}\sum_x\mu((g-
\mathcal P_xg)^2)$. Truncations of $f\in\mathcal A$ are bounded and gauge
invariant and converge in $L^2(\mu)$, and both sides are continuous on
$L^2(\mu)$.
\end{proof}

So entropy factorisation at the physical scale, uniform along the tuning,
implies (D) with $\varpi=1/C$, and every obstruction to the cube condition
applies to it: it fails for $\U(1)$ at weak coupling and at fixed cube size
(Section~\ref{sec:weak}). The methods that derive block factorisation from
spatial mixing uniformly in the boundary condition~\cite{CP21} need, at the
physical scale, an input of the strength of (D) itself.

\begin{proposition}[link covariances]\label{prop:linkcov}
Let $\mu$ be a gauge-invariant probability measure on $\SU(2)^E$, $E$ the links
of a box or of $\Lambda_L$ with $L\ge3$, and identify $U_e$ with a unit vector
$\sigma_e\in\R^4$. Then $E\sigma_e=0$, $E\sigma_e\sigma_e^T=\frac14I_4$ and
$E\sigma_e\sigma_{e'}^T=0$ for $e\ne e'$.
\end{proposition}

\begin{proof}
Distinct links have distinct sets of end points, so $e$ has an end point $x$
that is not an end point of $e'$. The gauge transformation $g$ at $x$ alone
maps $\sigma_e$ to $\ell_g\sigma_e$, where $\ell_g$ is left multiplication by
$g$ or right multiplication by $g^{-1}$, an orthogonal map, and fixes
$\sigma_{e'}$. By invariance $E\sigma_e\sigma_{e'}^T=\int\ell_g\,dg\;E\sigma_e
\sigma_{e'}^T=0$, since $\ell_g$ is linear in $g$ and $\int g\,dg=0$. The same
argument gives $E\sigma_e=0$ and shows that the law of $\sigma_e$ is invariant
under multiplication, hence Haar measure, so $E\sigma_e\sigma_e^T=\frac14I_4$.
\end{proof}

The influence and covariance matrices that spectral independence and
localisation schemes ask to be bounded are therefore exactly $\frac14I$ for the
Wilson measure at every $\beta$, weak coupling included, and carry no
information. Those methods need the bound for all tilted measures, and the
tilts break gauge invariance. Their output is a tensorisation inequality over
blocks of bounded lattice size, uniform in the volume, and the next
proposition shows that such an inequality cannot hold uniformly along the
tuning when the continuum limit is nontrivial.

\begin{proposition}[fixed blocks force triviality]\label{prop:fixedblocks}
Let $a_k=a_02^{-k}$ and $\beta_k$ be given, and suppose that for every $k$
there is $C_k<\infty$ such that the Wilson measure $\mu_{\beta_k,L}$ satisfies
\[
\Var_\mu(f)\le C_k\sum_{e\in E}\mu(\Var_ef)\qquad(f\in\mathcal A)
\]
for all $L$ in an infinite set $\mathcal L_k$ of even integers $L\ge8$, where
$\Var_ef$ is the conditional variance given the links other than $e$. Then
$\BP(2,1/(64C_k))$ holds on these tori, every torus limit along $\mathcal L_k$
has transfer-matrix gap at least $1/(5000e^2C_k)$, and every continuum limit
of such torus limits in the sense of Definition~\ref{def:limit} has
$\sigma(H)\subset\{0\}\cup[E_*,\infty)$ with $E_*=\limsup_k(5000e^2C_ka_k)^{-1}$.
In particular, if $C_ka_k\to0$, every continuum limit is trivial.
\end{proposition}

\begin{proof}
For $e\in V_x$ the $\sigma$-algebra of the links off $V_x$ is contained in that
of the links other than $e$, so $\mu(\Var_ef)=\|f-E[f\mid e^c]\|^2\le\|f-
\mathcal P_xf\|^2$. With $n=2$ each link lies in $16$ cubes $V_x$ and $|V_x|=4
\cdot2^4=64$, so
\[
\sum_e\mu(\Var_ef)\le\frac1{16}\sum_x\sum_{e\in V_x}\mu((f-\mathcal P_xf)^2)=4
\sum_x\mu((f-\mathcal P_xf)^2),
\]
and $\Var_\mu f\le64C_k\cdot2^{-4}\sum_x\mu((f-\mathcal P_xf)^2)$, which is
$\BP(2,1/(64C_k))$ on tori with $L\ge8=2\cdot2+4$. Corollary~\ref{cor:latticegap}
gives the gap $m_k=\varpi/(2e^2q_2)$ with $q_2=\frac{625}{16}$, that is $m_k=1/
(5000e^2C_k)$. For a continuum limit, as in the proof of Theorem~\ref{thm:B},
for $f\in\mathcal T$, $\varphi=\psi(f)-\ip\Omega{\psi(f)}\Omega$ and dyadic
$\tau>0$, $\ip\varphi{e^{-\tau H}\varphi}=\lim_k\ip{\varphi_k}{\mathbb T_k^{\tau/
a_k}\varphi_k}\le\liminf_ke^{-m_k\tau/a_k}\|\varphi_k\|^2=e^{-E_*\tau}\|\varphi
\|^2$, using $\|\varphi_k\|^2\to\|\varphi\|^2$, the case $\tau=0$ of the
definition. These $\varphi$ are dense in $\Omega^\perp$, and the spectral
argument of the proof of Theorem~\ref{thm:B} gives $\sigma(H|_{\Omega^\perp})
\subset[E_*,\infty)$. If $E_*=\infty$, then $\ip\varphi{e^{-\tau H}\varphi}=0$
for all dyadic $\tau>0$, and strong continuity of $e^{-\tau H}$ at $\tau=0$
gives $\varphi=0$, so $\Omega^\perp=\{0\}$.
\end{proof}

The same holds for blocks of any fixed lattice size $b$ in place of single
links, with $\BP(b,\cdot)$. So the tensorisation constants of single links or of
bounded blocks must grow at least like $1/a_k$ along the tuning whenever the
continuum limit is nontrivial; a usable version must work with blocks of
physical size, and is then a form of (D).

\subsection{The transfer matrix alone}\label{ssec:transfer}

The next results show what the continuum gap needs and that a lattice gap
tending to zero at the right rate is, by itself, a kinematic statement. In this
section Definition~\ref{def:limit} is used with $a_k=a_02^{-j_k}$ for a strictly
increasing sequence of integers $j_k$, with $\tau$ in $\mathbb D=\{a_0p2^{-q}:
p,q\in\Z_{\ge0}\}$, and with an arbitrary complex vector space $\mathcal T$ and
linear maps $O_k$ into the gauge-invariant local observables at nonnegative
times; the proof of Theorem~\ref{thm:B} uses nothing more. The states
$\omega_k$ are translation invariant and reflection positive, with transfer
matrices $0\le\mathbb T_k\le1$ on $\mathcal H_k$, vacua $\Omega_k$ and maps
$\Psi_k$ from observables at nonnegative times onto dense subspaces, as in
Section~\ref{sec:lattice}. Put $\gamma_k=\|\mathbb T_k|_{\Omega_k^\perp}\|$,
$m_k=-\log\gamma_k\in[0,\infty]$, and $\mathrm C_k(f;s)=\omega_k(\vartheta(
\overline{O_kf})\tau_{s/a_k}O_kf)-|\omega_k(O_kf)|^2$.

\begin{proposition}[the spectral input]\label{prop:spectralinput}
Let $(\mathcal H,\Omega,H,\psi)$ be a continuum limit of $(\omega_k,O_k)$ and
$\Delta>0$. The following are equivalent.
\begin{enumerate}[label=\textup{(\roman*)},leftmargin=*]
\item $\sigma(H|_{\Omega^\perp})\subset[\Delta,\infty)$.
\item For every $f\in\mathcal T$ there is $C_f<\infty$ with $\lim_k\mathrm C_k(
f;s)\le C_fe^{-\Delta s}$ for every $s\in\mathbb D$.
\end{enumerate}
\end{proposition}

\begin{proof}
For $f\in\mathcal T$ let $\phi_f=\psi f-\ip\Omega{\psi f}\Omega$. Since $e^{-sH}
\Omega=\Omega$ and $e^{-sH}$ is self-adjoint, $\ip{\phi_f}{e^{-sH}\phi_f}=\ip{
\psi f}{e^{-sH}\psi f}-|\ip\Omega{\psi f}|^2$, so $\lim_k\mathrm C_k(f;s)$ exists
and equals $\ip{\phi_f}{e^{-sH}\phi_f}$. If (i) holds, the spectral theorem
gives (ii) with $C_f=\|\phi_f\|^2$. Conversely, let $\mu_f$ be the spectral
measure of $H$ at $\phi_f$. If $\mu_f([0,\Delta-\eta])=c>0$ for some $\eta>0$,
then $\ip{\phi_f}{e^{-sH}\phi_f}\ge ce^{-(\Delta-\eta)s}>C_fe^{-\Delta s}$ for
large $s\in\mathbb D$, which contradicts (ii); so the spectral projection $P$ of
$H$ for $[0,\Delta)$ satisfies $P\phi_f=0$. The $\phi_f$ are dense in
$\Omega^\perp$, so $P=0$ on $\Omega^\perp$.
\end{proof}

So the continuum gap needs exponential decay, at physical times and only
asymptotically in $k$, of the autocorrelations of the observables that generate
the limit; it does not need a lattice gap, since lattice states of physical
energy below $\Delta$ are allowed if they decouple from every $O_k(f)$.

\begin{lemma}\label{lem:limsupgap}
Let $(\mathcal H,\Omega,H,\psi)$ be a continuum limit of $(\omega_k,O_k)$ along
a subsequence $\Lambda$ and $\bar E=\limsup_{k\in\Lambda}m_k/a_k\in[0,\infty]$.
Then $\ip\phi{e^{-sH}\phi}\le e^{-\bar Es}\|\phi\|^2$ for all $\phi\in\Omega^\perp$
and $s\ge0$; so $\sigma(H|_{\Omega^\perp})\subset[\bar E,\infty)$, and $\Omega^\perp
=\{0\}$ if $\bar E=\infty$.
\end{lemma}

\begin{proof}
For $f\in\mathcal T$ and $k\in\Lambda$ put $O=O_kf$ and $\phi_k=\Psi_k(O)-
\omega_k(O)\Omega_k\perp\Omega_k$. Since $\mathbb T_k\ge0$ is self-adjoint and
fixes $\Omega_k$, it preserves $\Omega_k^\perp$, and for $s\in\mathbb D$ and
$t=s/a_k\in\Z$, $\mathrm C_k(f;s)=\ip{\phi_k}{\mathbb T_k^t\phi_k}\le\gamma_k^t
\|\phi_k\|^2=e^{-(m_k/a_k)s}\|\phi_k\|^2$. Along a further subsequence $m_k/a_k
\to\bar E$, $\mathrm C_k(f;s)\to\ip{\phi_f}{e^{-sH}\phi_f}$ and $\|\phi_k\|^2=
\mathrm C_k(f;0)\to\|\phi_f\|^2$, with $e^{-\infty s}=0$ for $s>0$. The $\phi_f$
are dense in $\Omega^\perp$, both sides are continuous in $\phi$, and $s\mapsto
\ip\phi{e^{-sH}\phi}$ is continuous, so the inequality holds on $\Omega^\perp$
for all $s\ge0$. The spectral statement follows as in
Proposition~\ref{prop:spectralinput}; if $\bar E=\infty$, then $\ip\phi{e^{-sH}
\phi}=0$ for $s>0$, so $\phi=0$ by strong continuity.
\end{proof}

\begin{theorem}[gap tuning]\label{thm:gaptuning}
Suppose $0<m_k<\infty$ for all $k$ and $m_k\to0$, and let $E>0$ and $a_0>0$.
For each $k$ let $j_k$ be the integer with $E\le m_k2^{j_k}/a_0<2E$, pass to a
subsequence on which $j_k$ is strictly increasing, and put $a_k=a_02^{-j_k}$.
\begin{enumerate}[label=\textup{(\alph*)},leftmargin=*]
\item Every continuum limit of $(\omega_k,O_k)$ along any subsequence, for any
$\mathcal T$ and $O_k$, has $\sigma(H|_{\Omega^\perp})\subset[E,\infty)$.
\item There are a subsequence, a space $\mathcal T$ with maps $O_k$, and a
nontrivial continuum limit along it with $\inf\sigma(H|_{\Omega^\perp})=\lim_k
m_k/a_k\in[E,2E]$.
\end{enumerate}
\end{theorem}

\begin{proof}
The integer $j_k$ exists and is unique, because the condition reads
\[
2^{j_k}\in[Ea_0/m_k,\,2Ea_0/m_k),
\]
and $j_k\to\infty$ since $m_k\to0$. Part (a) is
Lemma~\ref{lem:limsupgap}, since $m_k/a_k\ge E$.

(b) Pass to a subsequence with $m_k/a_k\to\bar E\in[E,2E]$ and fix $k$. As
$\gamma_k>0$, $\Omega_k^\perp\ne\{0\}$, and $\gamma_k$, the norm of the positive
operator $\mathbb T_k|_{\Omega_k^\perp}$, lies in its spectrum; so the spectral
projection $P_k$ of $\mathbb T_k|_{\Omega_k^\perp}$ for $[\gamma_ke^{-a_k/k},
\gamma_k]$ is nonzero. Choose a unit vector $v_k$ in its range. The vectors
$\Psi_k(F-\omega_k(F))$ are the projections of $\Psi_k(F)$ onto $\Omega_k^\perp$,
hence dense in it; choose $F'_k$ with $\|u'_k-v_k\|\le\frac1{10}$ for $u'_k=
\Psi_k(F'_k-\omega_k(F'_k))$, and put $F_k=(F'_k-\omega_k(F'_k))/\|u'_k\|$ and
$u_k=\Psi_k(F_k)$. Then $\|u_k\|=1$, $u_k\perp\Omega_k$ and $\|P_ku_k\|\ge(1-
\frac1{10})/\frac{11}{10}=\frac9{11}$. Let $c_k(t)=\ip{u_k}{\mathbb T_k^tu_k}$.
Since $\mathbb T_k^t\ge0$ commutes with $P_k$ and has norm $\gamma_k^t$ on
$\Omega_k^\perp$,
\begin{equation}\label{eq:twosided}
\tfrac{81}{121}e^{-(m_k+a_k/k)t}\le\ip{P_ku_k}{\mathbb T_k^tP_ku_k}\le c_k(t)\le
e^{-m_kt}.
\end{equation}
Let $\sigma_k$ be the spectral measure of $\mathbb T_k$ at $u_k$, a probability
measure on $[0,1]$, and $\rho_k$ its image under $x\mapsto-a_k^{-1}\log x$, a
probability measure on the compact space $[0,\infty]$, with $0\mapsto\infty$.
Then $c_k(s/a_k)=\int e^{-sx}\,d\rho_k(x)$ for $s\in\mathbb D$, with $e^{-s\infty}
=0$ for $s>0$. Pass to a subsequence with $\rho_k\to\rho$ weakly on $[0,\infty]$.
For $s>0$ the function $x\mapsto e^{-sx}$ is continuous on $[0,\infty]$, so with
$\rho'=\rho|_{[0,\infty)}$,
\[
B(s):=\lim_kc_k(s/a_k)=\int_{[0,\infty)}e^{-sx}\,d\rho'(x)\qquad(s\in\mathbb D,
\ s>0).
\]
Let $\mathcal T$ be the free complex vector space on the symbols $\mathbf1$ and
$j_s$, $s\in\mathbb D$, $s>0$; put $O_k(\mathbf1)=1$ and $O_k(j_s)=\tau_{s/a_k}
F_k$ for the $k$ with $s/a_k\in\Z$, extended linearly. Let $\mathcal H=\C\Omega
\oplus L^2([0,\infty),\rho')$ with $\Omega=(1,0)$, $H=0\oplus x$, $\psi(\mathbf1)
=\Omega$ and $\psi(j_s)=(0,e^{-sx})$. Since $\Psi_k(O_kj_s)=\mathbb T_k^{s/a_k}
u_k$, for $s,s'>0$ and $r$ in $\mathbb D$,
\[
\omega_k\bigl(\vartheta(\overline{O_kj_s})\tau_{r/a_k}O_kj_{s'}\bigr)=c_k\bigl((s
+s'+r)/a_k\bigr)\to B(s+s'+r)=\ip{\psi j_s}{e^{-rH}\psi j_{s'}};
\]
moreover $\omega_k(\tau_{r/a_k}O_kj_s)=\ip{\Omega_k}{\mathbb T_k^{(r+s)/a_k}u_k}=
\ip{\Omega_k}{u_k}=0=\ip\Omega{e^{-rH}\psi j_s}$, the same holds with the roles
exchanged, and $\omega_k(1)=1=\ip\Omega{e^{-rH}\Omega}$. By sesquilinearity the
identities of Definition~\ref{def:limit} hold on $\mathcal T$. The span of
$\{e^{-sx}:s\in\mathbb D,s>0\}$ is an algebra of functions in $C_0([0,\infty))$
that separates points and vanishes nowhere, hence dense in $C_0([0,\infty))$ by
the Stone--Weierstrass theorem and dense in $L^2(\rho')$; so $\psi$ has dense
range, and $H\ge0$ is self-adjoint with $H\Omega=0$.

By~\eqref{eq:twosided} with $t=s/a_k$, $\frac{81}{121}e^{-\bar Es}\le B(s)\le e^{
-\bar Es}$ for $s\in\mathbb D$, $s>0$. The lower bound gives $\rho'\ne0$, so
$\mathcal H\ne\C\Omega$. The upper bound gives $\rho'([0,\bar E-\eta])=0$ for
every $\eta>0$, and the lower bound gives $\rho'([\bar E,\bar E+\eta])>0$ for
every $\eta>0$. Since $\sigma(H|_{\Omega^\perp})$ is the support of $\rho'$, its
infimum is $\bar E$.
\end{proof}

Theorem~\ref{thm:gaptuning} uses no property of the gauge theory: it applies
equally to the free scalar field with lattice masses tending to zero. A
nontrivial continuum limit with a gap, in the sense of
Definition~\ref{def:limit}, therefore exists as soon as the lattice gaps
$m(\beta_k)$ are positive and tend to zero along a sequence $\beta_k\to\infty$,
and neither property is known for four-dimensional $\SU(N)$. The limit so
obtained is generated by one sequence of observables chosen through spectral
projections; nothing identifies it with Yang--Mills theory, since it carries no
Euclidean invariance, no locality and no Wilson loops. The physical content of
(E) is the convergence of fixed physical observables, smeared fields and Wilson
loops of fixed physical size, and by Lemma~\ref{lem:limsupgap} a gap coming
from the transfer matrices requires $m_k/a_k$ to stay bounded below along a
subsequence: along the asymptotic-freedom tuning, the lattice gap must decay at
exactly the asymptotic-freedom rate.


\section*{Acknowledgements}

The von Mises--Fisher modulus and the Hessian bound without gauge fixing are those
of~\cite{BB26a}; the line field, the curvature bound in the gauge, the cube
condition and its consequences follow the longer account~\cite{BBLR26,BBLRS26}.
The constructive framework is that of Osterwalder, Schrader and Seiler, the
curvature criterion that of Bakry and \'Emery, and the uniqueness criterion
that of Dobrushin. The maximal correlation of Gaussian variables is
Gebelein's, the finite-size criterion follows Knabe and Gosset--Mozgunov, the
two-dimensional analysis rests on the heat kernel of the continuum measure of
Driver, Gross, King, Sengupta and L\'evy, and the twisted boundary conditions
are 't~Hooft's. The Lean~4 formalisation of the strong-coupling analysis
described in the introduction is the work of I.~J.~Sessford~\cite{Ses26}.

\section*{Declarations}

\emph{Data availability.} No data were generated. The source of this paper is
archived at
\href{https://doi.org/10.5281/zenodo.23444338}{doi:10.5281/zenodo.23444338}, and
the longer account of this work, with further results and finite checks, at
\href{https://doi.org/10.5281/zenodo.24085959}{doi:10.5281/zenodo.24085959}.

\emph{Funding.} This research received no external funding.

\emph{Conflict of interest.} The authors declare no conflict of interest.

\emph{Ethics.} Not applicable: the work involves no human participants,
animals or personal data.



\begin{thebibliography}{99}

\bibitem{AC25}
A.~Adhikari and S.~Cao, Correlation decay for finite lattice gauge theories at weak coupling, Ann. Probab. \textbf{53} (2025), 140--174, \href{https://doi.org/10.1214/24-AOP1702}{doi:10.1214/24-AOP1702}.

\bibitem{Bal84a}
T.~Ba{\l}aban, Propagators and renormalization transformations for lattice
gauge theories.~I, Comm. Math. Phys. \textbf{95} (1984), 17--40,
\href{https://doi.org/10.1007/BF01215753}{doi:10.1007/BF01215753}.

\bibitem{Bal85a}
T.~Ba{\l}aban, Averaging operations for lattice gauge theories, Comm. Math.
Phys. \textbf{98} (1985), 17--51,
\href{https://doi.org/10.1007/BF01211042}{doi:10.1007/BF01211042}.

\bibitem{Bal89b}
T.~Ba{\l}aban, Large field renormalization. II. Localization,
exponentiation, and bounds for the $\mathbf R$ operation, Comm. Math. Phys.
\textbf{122} (1989), 355--392,
\href{https://doi.org/10.1007/BF01238433}{doi:10.1007/BF01238433}.

\bibitem{Bar37}
J.~Barta, Sur la vibration fondamentale d'une membrane, C. R. Acad. Sci. Paris \textbf{204} (1937), 472--473.

\bibitem{BB26a}
D.~Bhattacharjee and U.~Bhattacharya, Mass gap of lattice Yang--Mills theory
up to the Dobrushin ceiling, in \emph{Lattice Yang--Mills theory: spacing-uniform
obstructions, the mass gap up to the Dobrushin ceiling, and Lean certificate},
Zenodo, release v2.0.2, 2026,
\href{https://doi.org/10.5281/zenodo.23449466}{doi:10.5281/zenodo.23449466}.

\bibitem{BBLR26}
D.~Bhattacharjee, U.~Bhattacharya, J.~M.~Lockwood and T.~Romanov, One cube
estimate implies the Yang--Mills mass gap, Hexagon preprint 2610.00341v7,
2026, \url{https://hexagonmath.org/2610.00341v7}.

\bibitem{BBLRS26}
D.~Bhattacharjee, U.~Bhattacharya, J.~M.~Lockwood, T.~Romanov and S.~Singha
Roy, Yang--Mills theory and the mass gap, research archive, Zenodo, release
v4.44.0, 2026,
\href{https://doi.org/10.5281/zenodo.24085959}{doi:10.5281/zenodo.24085959}.

\bibitem{BE85}
D.~Bakry and M.~\'Emery, Diffusions hypercontractives, in \emph{S\'eminaire
de Probabilit\'es XIX}, Lecture Notes in Math. 1123, Springer, Berlin, 1985,
177--206, \href{https://doi.org/10.1007/BFb0075847}{doi:10.1007/BFb0075847}.

\bibitem{BGL14}
D.~Bakry, I.~Gentil and M.~Ledoux, \emph{Analysis and Geometry of Markov
Diffusion Operators}, Grundlehren Math. Wiss. 348, Springer, Cham, 2014,
\href{https://doi.org/10.1007/978-3-319-00227-9}{doi:10.1007/978-3-319-00227-9}.

\bibitem{BK55}
M.~Biernacki and J.~Krzy\.z, On the monotonity of certain functionals in the theory of analytic functions, Ann. Univ. Mariae Curie-Sk{\l}odowska Sect. A \textbf{9} (1955), 135--147.

\bibitem{Cha21}
S.~Chatterjee, A probabilistic mechanism for quark confinement, Comm. Math. Phys. \textbf{385} (2021), 1007--1039, \href{https://doi.org/10.1007/s00220-021-04086-y}{doi:10.1007/s00220-021-04086-y}.

\bibitem{Chat19}
S.~Chatterjee, Yang--Mills for probabilists, in \emph{Probability and Analysis in Interacting Physical Systems}, Springer Proc. Math. Stat. 283, Springer, Cham, 2019, 1--16, \href{https://doi.org/10.1007/978-3-030-15338-0_1}{doi:10.1007/978-3-030-15338-0\_1}.

\bibitem{CNS25}
S.~Cao, R.~Nissim and S.~Sheffield, Expanded regimes of area law for lattice Yang--Mills theories, preprint, 2025, \href{https://arxiv.org/abs/2505.16585}{arXiv:2505.16585}.

\bibitem{CP21}
P.~Caputo and D.~Parisi, Block factorization of the relative entropy via spatial mixing, Comm. Math. Phys. \textbf{388} (2021), 793--818, \href{https://doi.org/10.1007/s00220-021-04237-1}{doi:10.1007/s00220-021-04237-1}.

\bibitem{Cre77}
M.~Creutz, Gauge fixing, the transfer matrix, and confinement on a lattice,
Phys. Rev. D \textbf{15} (1977), 1128--1136,
\href{https://doi.org/10.1103/PhysRevD.15.1128}{doi:10.1103/PhysRevD.15.1128}.

\bibitem{Deu01}
F.~Deutsch, \emph{Best Approximation in Inner Product Spaces}, CMS Books in Mathematics, Springer, New York, 2001, \href{https://doi.org/10.1007/978-1-4684-9298-9}{doi:10.1007/978-1-4684-9298-9}.

\bibitem{Dim13}
J.~Dimock, The renormalization group according to Balaban, I. Small fields, Rev. Math. Phys. \textbf{25} (2013), 1330010, \href{https://doi.org/10.1142/S0129055X13300100}{doi:10.1142/S0129055X13300100}.

\bibitem{Dob68}
R.~L.~Dobrushin, The description of a random field by means of conditional
probabilities and conditions of its regularity, Theory Probab. Appl.
\textbf{13} (1968), 197--224,
\href{https://doi.org/10.1137/1113026}{doi:10.1137/1113026}.

\bibitem{Dob70}
R.~L.~Dobrushin, Prescribing a system of random variables by conditional
distributions, Theory Probab. Appl. \textbf{15} (1970), 458--486,
\href{https://doi.org/10.1137/1115049}{doi:10.1137/1115049}.

\bibitem{Dri89}
B.~K.~Driver, YM$_2$: continuum expectations, lattice convergence, and lassos, Comm. Math. Phys. \textbf{123} (1989), 575--616, \href{https://doi.org/10.1007/BF01218586}{doi:10.1007/BF01218586}.

\bibitem{ES81}
B.~Efron and C.~Stein, The jackknife estimate of variance, Ann. Statist. \textbf{9} (1981), 586--596, \href{https://doi.org/10.1214/aos/1176345462}{doi:10.1214/aos/1176345462}.

\bibitem{FS82}
J.~Fr\"ohlich and T.~Spencer, Massless phases and symmetry restoration in abelian gauge theories and spin systems, Comm. Math. Phys. \textbf{83} (1982), 411--454, \href{https://doi.org/10.1007/BF01213610}{doi:10.1007/BF01213610}.

\bibitem{Geb41}
H.~Gebelein, Das statistische Problem der Korrelation als Variations- und Eigenwertproblem und sein Zusammenhang mit der Ausgleichsrechnung, Z. Angew. Math. Mech. \textbf{21} (1941), 364--379, \href{https://doi.org/10.1002/zamm.19410210604}{doi:10.1002/zamm.19410210604}.

\bibitem{GKS89}
L.~Gross, C.~King and A.~Sengupta, Two dimensional Yang--Mills theory via stochastic differential equations, Ann. Physics \textbf{194} (1989), 65--112, \href{https://doi.org/10.1016/0003-4916(89)90032-8}{doi:10.1016/0003-4916(89)90032-8}.

\bibitem{GM16}
D.~Gosset and E.~Mozgunov, Local gap threshold for frustration-free spin systems, J. Math. Phys. \textbf{57} (2016), 091901, \href{https://doi.org/10.1063/1.4962337}{doi:10.1063/1.4962337}.

\bibitem{Gri09}
A.~Grigor'yan, \emph{Heat Kernel and Analysis on Manifolds}, AMS/IP Studies in Advanced Mathematics 47, Amer. Math. Soc., Providence, RI, 2009, \href{https://doi.org/10.1090/amsip/047}{doi:10.1090/amsip/047}.

\bibitem{GW73}
D.~J.~Gross and F.~Wilczek, Ultraviolet behavior of non-Abelian gauge theories, Phys. Rev. Lett. \textbf{30} (1973), 1343--1346, \href{https://doi.org/10.1103/PhysRevLett.30.1343}{doi:10.1103/PhysRevLett.30.1343}.

\bibitem{HS87}
R.~Holley and D.~Stroock, Logarithmic Sobolev inequalities and stochastic Ising models, J. Stat. Phys. \textbf{46} (1987), 1159--1194, \href{https://doi.org/10.1007/BF01011161}{doi:10.1007/BF01011161}.

\bibitem{Jan97}
S.~Janson, \emph{Gaussian Hilbert Spaces}, Cambridge Tracts in Mathematics
\textbf{129}, Cambridge University Press, Cambridge, 1997,
\href{https://doi.org/10.1017/CBO9780511526169}{doi:10.1017/CBO9780511526169}.

\bibitem{JW06}
A.~Jaffe and E.~Witten, Quantum Yang--Mills theory, in \emph{The Millennium
Prize Problems}, Clay Mathematics Institute and American Mathematical
Society, Cambridge, MA and Providence, RI, 2006, 129--152,
\url{https://www.claymath.org/wp-content/uploads/2022/06/yangmills.pdf}.

\bibitem{Kac90}
V.~G.~Kac, \emph{Infinite-Dimensional Lie Algebras}, 3rd ed., Cambridge University Press, Cambridge, 1990, \href{https://doi.org/10.1017/CBO9780511626234}{doi:10.1017/CBO9780511626234}.

\bibitem{Kal21}
O.~Kallenberg, \emph{Foundations of Modern Probability}, 3rd ed., Springer,
Cham, 2021, \href{https://doi.org/10.1007/978-3-030-61871-1}{doi:10.1007/978-3-030-61871-1}.

\bibitem{Kna88}
S.~Knabe, Energy gaps and elementary excitations for certain VBS-quantum antiferromagnets, J. Stat. Phys. \textbf{52} (1988), 627--638, \href{https://doi.org/10.1007/BF01019721}{doi:10.1007/BF01019721}.

\bibitem{Lev03}
T.~L\'evy, Yang--Mills measure on compact surfaces, Mem. Amer. Math. Soc. \textbf{166} (2003), no.~790, \href{https://doi.org/10.1090/memo/0790}{doi:10.1090/memo/0790}.

\bibitem{McS34}
E.~J.~McShane, Extension of range of functions, Bull. Amer. Math. Soc. \textbf{40} (1934), 837--842, \href{https://doi.org/10.1090/S0002-9904-1934-05978-0}{doi:10.1090/S0002-9904-1934-05978-0}.

\bibitem{Mil76}
J.~Milnor, Curvatures of left invariant metrics on Lie groups, Adv.
Math. \textbf{21} (1976), 293--329,
\href{https://doi.org/10.1016/S0001-8708(76)80002-3}{doi:10.1016/S0001-8708(76)80002-3}.

\bibitem{MJ00}
K.~V.~Mardia and P.~E.~Jupp, \emph{Directional Statistics}, Wiley, Chichester, 2000, \href{https://doi.org/10.1002/9780470316979}{doi:10.1002/9780470316979}.

\bibitem{MRS93}
J.~Magnen, V.~Rivasseau and R.~S\'en\'eor, Construction of YM$_4$ with an infrared cutoff, Comm. Math. Phys. \textbf{155} (1993), 325--383, \href{https://doi.org/10.1007/BF02097397}{doi:10.1007/BF02097397}.

\bibitem{Nua06}
D.~Nualart, \emph{The Malliavin Calculus and Related Topics}, 2nd ed., Probability and its Applications, Springer, Berlin, 2006, \href{https://doi.org/10.1007/3-540-28329-3}{doi:10.1007/3-540-28329-3}.

\bibitem{OS73}
K.~Osterwalder and R.~Schrader, Axioms for Euclidean Green's functions,
Comm. Math. Phys. \textbf{31} (1973), 83--112,
\href{https://doi.org/10.1007/BF01645738}{doi:10.1007/BF01645738}.

\bibitem{OS75}
K.~Osterwalder and R.~Schrader, Axioms for Euclidean Green's functions.~II,
Comm. Math. Phys. \textbf{42} (1975), 281--305,
\href{https://doi.org/10.1007/BF01608978}{doi:10.1007/BF01608978}.

\bibitem{OSei78}
K.~Osterwalder and E.~Seiler, Gauge field theories on a lattice, Ann. Physics
\textbf{110} (1978), 440--471,
\href{https://doi.org/10.1016/0003-4916(78)90039-8}{doi:10.1016/0003-4916(78)90039-8}.

\bibitem{Pol73}
H.~D.~Politzer, Reliable perturbative results for strong interactions?, Phys. Rev. Lett. \textbf{30} (1973), 1346--1349, \href{https://doi.org/10.1103/PhysRevLett.30.1346}{doi:10.1103/PhysRevLett.30.1346}.

\bibitem{Sei82}
E.~Seiler, \emph{Gauge Theories as a Problem of Constructive Quantum Field
Theory and Statistical Mechanics}, Lecture Notes in Phys. 159, Springer,
Berlin, 1982, \href{https://doi.org/10.1007/3-540-11559-5}{doi:10.1007/3-540-11559-5}.

\bibitem{Ses26}
I.~J.~Sessford, The mass gap of $\SU(N)$ gauge theory seen by an observer of
fixed physical resolution, Lean~4 development and preprint, Zenodo, release
v0.3.1, 2026,
\href{https://doi.org/10.5281/zenodo.24073388}{doi:10.5281/zenodo.24073388};
\url{https://github.com/Agience/entroptics-mass-gap}.

\bibitem{SS03}
E.~M.~Stein and R.~Shakarchi, \emph{Fourier Analysis: An Introduction}, Princeton Lectures in Analysis I, Princeton University Press, Princeton, NJ, 2003.

\bibitem{SW71}
E.~M.~Stein and G.~Weiss, \emph{Introduction to Fourier Analysis on Euclidean Spaces}, Princeton University Press, Princeton, NJ, 1971, \href{https://doi.org/10.1515/9781400883899}{doi:10.1515/9781400883899}.

\bibitem{SZZ23}
H.~Shen, R.~Zhu and X.~Zhu, A stochastic analysis approach to lattice
Yang--Mills at strong coupling, Comm. Math. Phys. \textbf{400} (2023),
805--851, \href{https://doi.org/10.1007/s00220-022-04609-1}{doi:10.1007/s00220-022-04609-1}.

\bibitem{tH79}
G.~'t~Hooft, A property of electric and magnetic flux in non-Abelian gauge theories, Nuclear Phys. B \textbf{153} (1979), 141--160, \href{https://doi.org/10.1016/0550-3213(79)90595-9}{doi:10.1016/0550-3213(79)90595-9}.

\bibitem{Vil09}
C.~Villani, \emph{Optimal Transport: Old and New}, Grundlehren Math. Wiss.
338, Springer, Berlin, 2009,
\href{https://doi.org/10.1007/978-3-540-71050-9}{doi:10.1007/978-3-540-71050-9}.

\bibitem{vRS05}
M.-K.~von Renesse and K.-T.~Sturm, Transport inequalities, gradient estimates, entropy, and Ricci curvature, Comm. Pure Appl. Math. \textbf{58} (2005), 923--940, \href{https://doi.org/10.1002/cpa.20060}{doi:10.1002/cpa.20060}.

\bibitem{Wat44}
G.~N.~Watson, \emph{A Treatise on the Theory of Bessel Functions}, 2nd ed., Cambridge University Press, Cambridge, 1944.

\bibitem{Wil74}
K.~G.~Wilson, Confinement of quarks, Phys.\ Rev.\ D \textbf{10} (1974), no.~8, 2445--2459,
\href{https://doi.org/10.1103/PhysRevD.10.2445}{doi:10.1103/\allowbreak PhysRevD.10.2445}.

\end{thebibliography}

\end{document}
